% ============================================================
%  Cours : Équations Différentielles Ordinaires
%  Langue : Français
%  Fichier : preamble_ch1_2.tex  (Préambule + Chapitres 1-2)
% ============================================================

\documentclass[12pt,a4paper,openany]{book}
\usepackage[utf8]{inputenc}
\usepackage[T1]{fontenc}
\usepackage[french]{babel}
\usepackage[margin=2.5cm]{geometry}
\usepackage{amsmath,amssymb,amsthm,mathrsfs,mathtools}
\usepackage{tikz}
\usetikzlibrary{arrows.meta,calc,positioning,patterns,decorations.markings,decorations.pathmorphing,cd}
\usepackage{pgfplots}
\pgfplotsset{compat=1.18}
\usepackage[most]{tcolorbox}
\tcbuselibrary{theorems,skins,breakable}
\usepackage[colorlinks=true,linkcolor=blue!60!black,citecolor=green!50!black]{hyperref}
\usepackage[french]{cleveref}
\usepackage{makeidx}
\makeindex
\usepackage{enumitem}
\usepackage{fancyhdr}
\usepackage{caption}
\usepackage{booktabs}
\usepackage{microtype}
\usepackage{xcolor}
\usepackage{minitoc}
\usepackage{array}

\pagestyle{fancy}
\fancyhf{}
\fancyhead[LE]{\leftmark}
\fancyhead[RO]{\rightmark}
\fancyfoot[C]{\thepage}

% Theorem environments — MUST use \newtheorem (amsthm), NOT \newtcbtheorem
\theoremstyle{definition}
\newtheorem{definition}{Définition}[chapter]
\newtheorem{exemple}[definition]{Exemple}
\newtheorem{exercice}{Exercice}[chapter]

\theoremstyle{plain}
\newtheorem{theoreme}[definition]{Théorème}
\newtheorem{proposition}[definition]{Proposition}
\newtheorem{lemme}[definition]{Lemme}
\newtheorem{corollaire}[definition]{Corollaire}

\theoremstyle{remark}
\newtheorem{remarque}[definition]{Remarque}
\newtheorem{notation}[definition]{Notation}
\newtheorem{methode}[definition]{Méthode}

% tcolorbox wrapping (visual only — do NOT use \newtcbtheorem)
\tcolorboxenvironment{definition}{enhanced,breakable,colback=blue!5,colframe=blue!70!black,fonttitle=\bfseries,before skip=10pt,after skip=10pt,left=8pt,right=8pt}
\tcolorboxenvironment{theoreme}{enhanced,breakable,colback=red!5,colframe=red!70!black,fonttitle=\bfseries,before skip=10pt,after skip=10pt,left=8pt,right=8pt}
\tcolorboxenvironment{proposition}{enhanced,breakable,colback=orange!5,colframe=orange!70!black,fonttitle=\bfseries,before skip=10pt,after skip=10pt,left=8pt,right=8pt}
\tcolorboxenvironment{lemme}{enhanced,breakable,colback=yellow!5,colframe=yellow!50!black,fonttitle=\bfseries,before skip=10pt,after skip=10pt,left=8pt,right=8pt}
\tcolorboxenvironment{corollaire}{enhanced,breakable,colback=green!5,colframe=green!50!black,fonttitle=\bfseries,before skip=10pt,after skip=10pt,left=8pt,right=8pt}
\tcolorboxenvironment{remarque}{enhanced,breakable,colback=gray!5,colframe=gray!50!black,before skip=8pt,after skip=8pt,left=6pt,right=6pt}
\tcolorboxenvironment{exemple}{enhanced,breakable,colback=purple!5,colframe=purple!50!black,before skip=8pt,after skip=8pt,left=6pt,right=6pt}

% Custom commands
\newcommand{\R}{\mathbb{R}}
\newcommand{\C}{\mathbb{C}}
\newcommand{\N}{\mathbb{N}}
\newcommand{\Z}{\mathbb{Z}}
\newcommand{\Q}{\mathbb{Q}}
\newcommand{\dd}{\mathrm{d}}
\newcommand{\abs}[1]{\left|#1\right|}
\newcommand{\norm}[1]{\left\|#1\right\|}
\DeclareMathOperator{\tr}{tr}
\DeclareMathOperator{\Sp}{Sp}
\DeclareMathOperator{\diag}{diag}

% ============================================================
\begin{document}

% --- Page de titre ---
\begin{titlepage}
\begin{tikzpicture}[remember picture, overlay]
  \fill[blue!10!white] (current page.south west) rectangle (current page.north east);
  \fill[blue!80!black, rotate around={-15:(current page.center)}]
    ([xshift=-4cm,yshift=-1cm]current page.center) rectangle ([xshift=4cm,yshift=1.5cm]current page.center);
  \node[anchor=center, text=white, font=\Huge\bfseries, rotate=-15]
    at ([yshift=0.3cm]current page.center) {\'Equations Diff\'erentielles Ordinaires};
  \node[anchor=center, text=orange!80!yellow, font=\Large, rotate=-15]
    at ([yshift=-1cm]current page.center) {Notes de Cours};
  \node[anchor=north, text=blue!70!black, font=\Large\itshape]
    at ([yshift=-4cm]current page.center) {Ya\'e Ulrich Gaba};
  \node[anchor=north, text=blue!50!black, font=\large]
    at ([yshift=-5.5cm]current page.center) {Licence L3 --- 2025--2026};
  \fill[orange!70!yellow] ([yshift=3cm]current page.south west) rectangle ([yshift=3.3cm]current page.south east);
  \node[anchor=south, text=gray, font=\small] at ([yshift=1cm]current page.south) {\today};
\end{tikzpicture}
\end{titlepage}

% --- Préface ---
\chapter*{Préface}
\addcontentsline{toc}{chapter}{Préface}

Les \emph{équations différentielles ordinaires}\index{equation differentielle ordinaire@équation différentielle ordinaire} (EDO) constituent l'un des piliers fondamentaux des mathématiques appliquées et de l'analyse. Depuis les travaux fondateurs de Newton et Leibniz au XVII\up{e} siècle, elles n'ont cessé de jouer un rôle central tant en mathématiques pures qu'en physique, en biologie, en économie et dans toutes les sciences de l'ingénieur.

Ce cours s'adresse aux étudiants de deuxième et troisième années de licence en mathématiques. Il suppose acquis les fondements de l'analyse réelle (continuité, dérivabilité, intégration) et de l'algèbre linéaire (espaces vectoriels, matrices, valeurs propres). Notre objectif est double : fournir une présentation rigoureuse de la théorie, tout en développant les techniques de résolution et les applications concrètes.

Le texte est organisé de manière progressive. Le premier chapitre introduit les concepts fondamentaux et la modélisation. Le deuxième chapitre traite en détail les équations du premier ordre. Les chapitres suivants aborderont les équations linéaires d'ordre supérieur, les systèmes différentiels, la théorie qualitative et les méthodes numériques.

Chaque chapitre contient de nombreux exemples détaillés, des remarques méthodologiques et une série d'exercices classés par difficulté croissante. Les exercices marqués d'une étoile sont calculatoires, ceux marqués de deux étoiles demandent une réflexion théorique plus poussée, et ceux marqués de trois étoiles font appel à la modélisation.

\vspace{1cm}
\hfill\textit{L'auteur}

% --- Table des matières ---
\dominitoc
\tableofcontents

% ============================================================
% CHAPITRE 1
% ============================================================
\chapter{Introduction et Modélisation --- Qu'est-ce qu'une EDO\,?}
\label{chap:introduction}
\index{introduction}

\section{Exemples motivants}
\label{sec:exemples-motivants}

Les équations différentielles apparaissent naturellement dès que l'on cherche à décrire l'évolution d'un phénomène au cours du temps. Avant de donner des définitions formelles, examinons quelques situations concrètes issues de la physique, de la biologie et de l'électronique.

\subsection{Le pendule simple}
\label{subsec:pendule}
\index{pendule simple}

Considérons un pendule constitué d'une masse $m$ suspendue à un fil rigide de longueur $\ell$. En notant $\theta(t)$ l'angle que fait le fil avec la verticale à l'instant $t$, l'application du principe fondamental de la dynamique conduit à l'équation :
\begin{equation}
\label{eq:pendule}
\frac{\dd^2\theta}{\dd t^2} + \frac{g}{\ell}\sin\theta = 0,
\end{equation}
où $g$ désigne l'accélération de la pesanteur. Il s'agit d'une \emph{équation différentielle du second ordre}, non linéaire à cause du terme $\sin\theta$. Pour de petites oscillations ($\theta$ petit), on approche $\sin\theta \approx \theta$, et l'on obtient l'équation linéarisée :
\begin{equation}
\label{eq:pendule-lineaire}
\ddot{\theta} + \omega_0^2\,\theta = 0, \quad \omega_0 = \sqrt{\frac{g}{\ell}},
\end{equation}
dont la solution générale est $\theta(t) = A\cos(\omega_0 t) + B\sin(\omega_0 t)$.

\subsection{Croissance de population}
\label{subsec:population}
\index{croissance de population}
\index{modele de Malthus@modèle de Malthus}

Soit $P(t)$ la taille d'une population à l'instant $t$. Le modèle le plus simple, dû à Malthus (1798), suppose que le taux de croissance est proportionnel à la population :
\begin{equation}
\label{eq:malthus}
\frac{\dd P}{\dd t} = rP, \quad r \in \R,
\end{equation}
dont la solution est $P(t) = P_0 e^{rt}$ avec $P_0 = P(0)$. Ce modèle prédit une croissance exponentielle, ce qui est irréaliste à long terme. Le modèle logistique de Verhulst\index{modele logistique@modèle logistique}\index{Verhulst} corrige ce défaut :
\begin{equation}
\label{eq:logistique}
\frac{\dd P}{\dd t} = rP\!\left(1 - \frac{P}{K}\right),
\end{equation}
où $K > 0$ est la \emph{capacité de charge}\index{capacite de charge@capacité de charge} du milieu.

\subsection{Décroissance radioactive}
\label{subsec:radioactivite}
\index{decroissance radioactive@décroissance radioactive}

La quantité $N(t)$ de noyaux radioactifs à l'instant $t$ obéit à :
\begin{equation}
\label{eq:radioactivite}
\frac{\dd N}{\dd t} = -\lambda N, \quad \lambda > 0,
\end{equation}
où $\lambda$ est la constante de désintégration\index{constante de desintegration@constante de désintégration}. La solution est $N(t) = N_0 e^{-\lambda t}$, et le \emph{temps de demi-vie}\index{temps de demi-vie} est $t_{1/2} = \frac{\ln 2}{\lambda}$.

\subsection{Circuit RC}
\label{subsec:circuit-RC}
\index{circuit RC}

Un circuit composé d'une résistance $R$ et d'un condensateur $C$ en série, soumis à une tension $E(t)$, obéit à l'équation :
\begin{equation}
\label{eq:circuit-RC}
RC\,\frac{\dd u}{\dd t} + u = E(t),
\end{equation}
où $u(t)$ est la tension aux bornes du condensateur. C'est une équation linéaire du premier ordre avec second membre.

% --- Champ de directions pour y' = y ---
\begin{figure}[ht]
\centering
\begin{tikzpicture}
\begin{axis}[
    title={Champ de directions de $y' = y$},
    xlabel={$t$},
    ylabel={$y$},
    domain=-2:2,
    ymin=-3, ymax=3,
    xmin=-2, xmax=2,
    axis lines=middle,
    width=10cm,
    height=8cm,
    samples=15,
]
% Quelques courbes solutions
\addplot[red,thick,domain=-2:2,samples=50] {0.5*exp(x)};
\addplot[red,thick,domain=-2:2,samples=50] {-0.5*exp(x)};
\addplot[red,thick,domain=-2:2,samples=50] {exp(x)};
\addplot[red,thick,domain=-2:2,samples=50] {-exp(x)};
\addplot[red,thick,domain=-2:2,samples=50] {2*exp(x)};
\end{axis}
\end{tikzpicture}
\caption{Champ de directions et quelques courbes solutions de $y' = y$.}
\label{fig:champ-yp-y}
\end{figure}

\section{Définitions fondamentales}
\label{sec:definitions-fondamentales}

\begin{definition}[Équation différentielle ordinaire]
\label{def:edo}
\index{equation differentielle ordinaire@équation différentielle ordinaire!definition@définition}
Une \emph{équation différentielle ordinaire} (EDO) est une relation de la forme
\[
F\!\left(t, y, y', y'', \ldots, y^{(n)}\right) = 0,
\]
où $t$ est la variable indépendante, $y = y(t)$ est la fonction inconnue, et $F$ est une fonction donnée. L'\emph{ordre}\index{ordre d'une EDO} de l'EDO est l'entier $n$, c'est-à-dire l'ordre le plus élevé de dérivation qui apparaît.
\end{definition}

\begin{definition}[Forme résolue]
\label{def:forme-resolue}
\index{forme resolue@forme résolue}
Une EDO est dite sous \emph{forme résolue} (ou forme normale) si elle s'écrit :
\[
y^{(n)} = f\!\left(t, y, y', \ldots, y^{(n-1)}\right).
\]
\end{definition}

\begin{definition}[EDO linéaire]
\label{def:edo-lineaire}
\index{equation differentielle lineaire@équation différentielle linéaire}
Une EDO d'ordre $n$ est \emph{linéaire} si elle est de la forme
\[
a_n(t)\,y^{(n)} + a_{n-1}(t)\,y^{(n-1)} + \cdots + a_1(t)\,y' + a_0(t)\,y = b(t),
\]
où les fonctions $a_0, a_1, \ldots, a_n$ et $b$ sont des fonctions continues données. Si $b \equiv 0$, l'équation est dite \emph{homogène}\index{equation homogene@équation homogène} ; sinon, elle est dite \emph{avec second membre}\index{second membre}.
\end{definition}

\begin{definition}[Solution d'une EDO]
\label{def:solution}
\index{solution!definition@définition}
Soit $I$ un intervalle ouvert de $\R$. On appelle \emph{solution} de l'EDO $F(t,y,y',\ldots,y^{(n)}) = 0$ sur $I$ toute fonction $\varphi : I \to \R$ de classe $\mathscr{C}^n$ telle que
\[
\forall\, t \in I, \quad F\!\left(t, \varphi(t), \varphi'(t), \ldots, \varphi^{(n)}(t)\right) = 0.
\]
\end{definition}

\begin{definition}[Solution générale, particulière, singulière]
\label{def:solutions-types}
\index{solution!generale@générale}
\index{solution!particuliere@particulière}
\index{solution!singuliere@singulière}
\begin{itemize}[leftmargin=*]
\item La \emph{solution générale} d'une EDO d'ordre $n$ est l'ensemble de toutes les solutions, dépendant en général de $n$ constantes arbitraires $C_1, \ldots, C_n$.
\item Une \emph{solution particulière} est obtenue en attribuant des valeurs numériques aux constantes.
\item Une \emph{solution singulière} est une solution qui ne peut être obtenue à partir de la solution générale pour aucune valeur des constantes.
\end{itemize}
\end{definition}

\begin{exemple}[Solution singulière]
\label{ex:solution-singuliere}
\index{solution!singuliere@singulière!exemple}
Considérons l'EDO $y' = \sqrt{\abs{y}}$ avec la condition de positivité $y \geq 0$, soit $y' = \sqrt{y}$ pour $y \geq 0$. La solution générale est
\[
y(t) = \frac{(t - C)^2}{4}, \quad C \in \R, \quad \text{pour } t \geq C.
\]
Mais la fonction $y(t) = 0$ pour tout $t$ est aussi une solution, qui n'est pas contenue dans la famille ci-dessus. C'est une solution singulière. De plus, pour tout $a \in \R$, la fonction
\[
y(t) = \begin{cases} 0 & \text{si } t \leq a, \\ \dfrac{(t-a)^2}{4} & \text{si } t > a, \end{cases}
\]
est aussi solution, ce qui montre la non-unicité en $(0,0)$.
\end{exemple}

\section{Problème de Cauchy}
\label{sec:probleme-cauchy}
\index{probleme de Cauchy@problème de Cauchy}

\begin{definition}[Problème de Cauchy]
\label{def:probleme-cauchy}
Un \emph{problème de Cauchy} (ou \emph{problème à valeur initiale}\index{probleme a valeur initiale@problème à valeur initiale}) pour une EDO d'ordre $n$ sous forme résolue consiste à trouver une fonction $y$ définie au voisinage de $t_0$ telle que :
\[
\begin{cases}
y^{(n)} = f(t, y, y', \ldots, y^{(n-1)}),\\[4pt]
y(t_0) = y_0,\; y'(t_0) = y_1,\; \ldots,\; y^{(n-1)}(t_0) = y_{n-1},
\end{cases}
\]
où $(t_0, y_0, y_1, \ldots, y_{n-1})$ sont des données initiales fixées. Les conditions $y(t_0) = y_0$, $y'(t_0) = y_1$, etc., sont appelées \emph{conditions initiales}\index{condition initiale}.
\end{definition}

\begin{remarque}
\label{rem:geometrie-cauchy}
Géométriquement, résoudre un problème de Cauchy du premier ordre $y' = f(t,y)$, $y(t_0) = y_0$, revient à chercher, parmi toutes les courbes intégrales\index{courbe integrale@courbe intégrale} (solutions du champ de directions), celle qui passe par le point $(t_0, y_0)$.
\end{remarque}

La question fondamentale est celle de l'\emph{existence} et de l'\emph{unicité} des solutions d'un problème de Cauchy. Nous énonçons ici, sans démonstration, les deux théorèmes fondamentaux ; leurs preuves seront données dans les chapitres ultérieurs.

\begin{theoreme}[Cauchy--Peano]
\label{thm:cauchy-peano}
\index{theoreme de Cauchy--Peano@théorème de Cauchy--Peano}
Soit $\Omega$ un ouvert de $\R^2$ et $f : \Omega \to \R$ une fonction continue. Pour tout $(t_0, y_0) \in \Omega$, le problème de Cauchy
\[
y' = f(t, y), \quad y(t_0) = y_0,
\]
admet au moins une solution définie sur un intervalle ouvert contenant $t_0$.
\end{theoreme}

\begin{theoreme}[Cauchy--Lipschitz]
\label{thm:cauchy-lipschitz}
\index{theoreme de Cauchy--Lipschitz@théorème de Cauchy--Lipschitz}
\index{condition de Lipschitz}
Soit $\Omega$ un ouvert de $\R^2$ et $f : \Omega \to \R$ une fonction continue qui est \emph{localement lipschitzienne}\index{lipschitzienne (fonction)} par rapport à $y$, c'est-à-dire que pour tout compact $K \subset \Omega$, il existe $L \geq 0$ tel que
\[
\abs{f(t, y_1) - f(t, y_2)} \leq L \abs{y_1 - y_2} \quad \forall (t, y_1), (t, y_2) \in K.
\]
Alors, pour tout $(t_0, y_0) \in \Omega$, le problème de Cauchy
\[
y' = f(t, y), \quad y(t_0) = y_0,
\]
admet une unique solution maximale.
\end{theoreme}

\begin{remarque}
\label{rem:lipschitz-suffisante}
Si $f$ est de classe $\mathscr{C}^1$ par rapport à $y$, alors $f$ est localement lipschitzienne par rapport à $y$ (c'est une conséquence du théorème des accroissements finis). La condition de Lipschitz est donc satisfaite dans la grande majorité des cas rencontrés en pratique.
\end{remarque}

\begin{exemple}
\label{ex:cauchy-lipschitz-verification}
Considérons $y' = t^2 + y^2$, $y(0) = 1$. La fonction $f(t,y) = t^2 + y^2$ est de classe $\mathscr{C}^{\infty}$ sur $\R^2$, donc localement lipschitzienne en $y$. D'après le théorème de Cauchy--Lipschitz, ce problème admet une unique solution maximale. Remarquons cependant que cette solution peut exploser en temps fini (la solution n'est pas nécessairement définie sur $\R$ tout entier).
\end{exemple}

\section{Champ de directions}
\label{sec:champ-directions}
\index{champ de directions}

Étant donnée une EDO du premier ordre $y' = f(t,y)$, on peut associer à chaque point $(t,y)$ du plan un vecteur de direction $(1, f(t,y))$. L'ensemble de ces vecteurs constitue le \emph{champ de directions}. Une courbe intégrale est une courbe du plan qui est tangente en chaque point au vecteur du champ de directions.

\begin{methode}[Tracer un champ de directions]
\label{meth:champ-directions}
Pour tracer le champ de directions de $y' = f(t,y)$ :
\begin{enumerate}
\item Choisir une grille de points $(t_i, y_j)$ dans le domaine.
\item En chaque point, calculer la pente $m = f(t_i, y_j)$.
\item Tracer un petit segment de pente $m$ centré en $(t_i, y_j)$.
\item Les \emph{isoclines}\index{isocline} sont les courbes $f(t,y) = c$ (pente constante) ; elles aident à organiser le tracé.
\end{enumerate}
\end{methode}

\begin{figure}[ht]
\centering
\begin{tikzpicture}
\begin{axis}[
    title={Champ de directions de $y' = t - y$},
    xlabel={$t$},
    ylabel={$y$},
    domain=-2:3,
    ymin=-2, ymax=4,
    xmin=-2, xmax=3,
    axis lines=middle,
    width=10cm,
    height=8cm,
]
% Courbe solution pour y(0)=0
\addplot[red,thick,domain=-2:3,samples=50] {x - 1 + exp(-x)};
% Courbe solution pour y(0)=3
\addplot[red,thick,domain=-2:3,samples=50] {x - 1 + 4*exp(-x)};
% Courbe solution pour y(0)=-1
\addplot[red,thick,domain=-2:3,samples=50] {x - 1};
\end{axis}
\end{tikzpicture}
\caption{Champ de directions et courbes solutions de $y' = t - y$.}
\label{fig:champ-t-y}
\end{figure}

\section{Classification des EDO}
\label{sec:classification}
\index{classification des EDO}

Il est utile de classer les équations différentielles en grandes familles, car chaque famille admet des méthodes de résolution spécifiques.

\begin{definition}[Équation autonome]
\label{def:autonome}
\index{equation autonome@équation autonome}
Une EDO est dite \emph{autonome} si le second membre ne dépend pas explicitement de la variable indépendante $t$ :
\[
y' = f(y).
\]
\end{definition}

\begin{definition}[Équation à variables séparables]
\label{def:separable}
\index{equation a variables separables@équation à variables séparables}
Une EDO du premier ordre est dite \emph{à variables séparables} si elle peut s'écrire sous la forme :
\[
y' = g(t)\,h(y).
\]
\end{definition}

\begin{definition}[Équation homogène (au sens de l'homogénéité)]
\label{def:homogene-sens-homogeneite}
\index{equation homogene@équation homogène!au sens de l'homogeneite@au sens de l'homogénéité}
Une EDO du premier ordre $y' = f(t,y)$ est dite \emph{homogène} si $f(t,y) = \varphi(y/t)$ pour une certaine fonction $\varphi$, c'est-à-dire si $f(\lambda t, \lambda y) = f(t,y)$ pour tout $\lambda \neq 0$.
\end{definition}

\begin{definition}[Équation exacte]
\label{def:exacte}
\index{equation exacte@équation exacte}
Une EDO du premier ordre écrite sous la forme $M(t,y)\,\dd t + N(t,y)\,\dd y = 0$ est dite \emph{exacte} s'il existe une fonction $\Phi(t,y)$ telle que
\[
\frac{\partial \Phi}{\partial t} = M \quad \text{et} \quad \frac{\partial \Phi}{\partial y} = N.
\]
\end{definition}

\medskip

Le tableau suivant résume les principales classes d'EDO du premier ordre et leurs méthodes de résolution :

\begin{table}[ht]
\centering
\caption{Classification des EDO du premier ordre.}
\label{tab:classification-edo1}
\begin{tabular}{@{}lll@{}}
\toprule
\textbf{Type} & \textbf{Forme} & \textbf{Méthode} \\
\midrule
Séparable & $y' = g(t)\,h(y)$ & Séparation des variables \\
Linéaire & $y' + p(t)\,y = q(t)$ & Facteur intégrant \\
Bernoulli & $y' + p(t)\,y = q(t)\,y^n$ & Substitution $v = y^{1-n}$ \\
Exacte & $M\,\dd t + N\,\dd y = 0$ & Recherche de potentiel \\
Homogène & $y' = \varphi(y/t)$ & Substitution $u = y/t$ \\
Riccati & $y' = a(t)\,y^2 + b(t)\,y + c(t)$ & Solution particulière connue \\
Clairaut & $y = ty' + \psi(y')$ & Dérivation \\
\bottomrule
\end{tabular}
\end{table}

\section{Repères historiques}
\label{sec:historique}
\index{historique}

\begin{itemize}[leftmargin=*]
\item \textbf{Isaac Newton}\index{Newton, Isaac} (1643--1727) et \textbf{Gottfried Wilhelm Leibniz}\index{Leibniz, Gottfried Wilhelm} (1646--1716) fondent le calcul infinitésimal et résolvent les premières équations différentielles. Newton traite le problème de la chute des corps, Leibniz introduit la notation $\dd y / \dd x$.
\item Les frères \textbf{Jacques}\index{Bernoulli, Jacques} et \textbf{Jean Bernoulli}\index{Bernoulli, Jean} (fin XVII\up{e} -- début XVIII\up{e} s.) étudient les équations qui portent leur nom et résolvent le problème de la brachistochrone.
\item \textbf{Leonhard Euler}\index{Euler, Leonhard} (1707--1783) développe les méthodes systématiques de résolution des EDO linéaires à coefficients constants et introduit la méthode numérique qui porte son nom.
\item \textbf{Joseph-Louis Lagrange}\index{Lagrange, Joseph-Louis} (1736--1813) introduit la méthode de variation de la constante.
\item \textbf{Augustin-Louis Cauchy}\index{Cauchy, Augustin-Louis} (1789--1857) donne les premières démonstrations rigoureuses d'existence et d'unicité.
\item \textbf{Rudolf Lipschitz}\index{Lipschitz, Rudolf} (1832--1903) formule la condition qui porte son nom et affine le théorème de Cauchy.
\item \textbf{Émile Picard}\index{Picard, Emile@Picard, Émile} (1856--1941) développe la méthode des approximations successives (itérés de Picard).
\item \textbf{Henri Poincaré}\index{Poincare, Henri@Poincaré, Henri} (1854--1912) fonde la théorie qualitative des EDO et la topologie.
\end{itemize}

\section{Résumé du chapitre}
\label{sec:resume-ch1}

\begin{tcolorbox}[colback=blue!3,colframe=blue!50!black,title={\textbf{Résumé --- Chapitre 1}}]
\begin{itemize}[leftmargin=*]
\item Une EDO est une relation entre une fonction inconnue et ses dérivées. L'ordre est celui de la dérivée la plus élevée.
\item La solution générale d'une EDO d'ordre $n$ dépend de $n$ constantes arbitraires.
\item Un problème de Cauchy fixe la valeur de la fonction et de ses $n-1$ premières dérivées en un point.
\item Le théorème de Cauchy--Peano garantit l'existence sous la seule hypothèse de continuité.
\item Le théorème de Cauchy--Lipschitz garantit l'existence et l'unicité sous la condition de Lipschitz.
\item Le champ de directions est un outil géométrique fondamental pour visualiser les solutions.
\item Les EDO se classent en familles (séparables, linéaires, exactes, homogènes, etc.), chacune admettant des méthodes de résolution spécifiques.
\end{itemize}
\end{tcolorbox}

\section{Exercices}
\label{sec:exercices-ch1}

\begin{exercice}[Identification]
\label{exo:ch1-identification}
Pour chacune des équations suivantes, déterminer l'ordre, indiquer si elle est linéaire ou non, et si elle est autonome :
\begin{enumerate}[label=(\alph*)]
\item $y' = 3y + t^2$
\item $y'' + y' - 6y = 0$
\item $y' = y^2 - t$
\item $(1+t^2)\,y'' + 2t\,y' + y = \sin t$
\item $y' = e^{-y}$
\item $y''' + (y')^2 = t$
\end{enumerate}
\end{exercice}

\begin{exercice}[Vérification de solutions]
\label{exo:ch1-verification}
Vérifier que les fonctions données sont solutions des EDO correspondantes :
\begin{enumerate}[label=(\alph*)]
\item $y(t) = Ce^{-2t}$ est solution de $y' + 2y = 0$ pour tout $C \in \R$.
\item $y(t) = \frac{1}{1-t}$ est solution de $y' = y^2$ sur $]-\infty, 1[$.
\item $y(t) = t\ln t - t$ est solution de $t\,y'' - y' + \frac{1}{t} = 0$ pour $t > 0$.
\end{enumerate}
\end{exercice}

\begin{exercice}[Problèmes de Cauchy simples]
\label{exo:ch1-cauchy-simples}
Résoudre les problèmes de Cauchy suivants :
\begin{enumerate}[label=(\alph*)]
\item $y' = 3t^2$, $y(0) = 1$.
\item $y' = -2y$, $y(0) = 5$.
\item $y' = \cos t$, $y(\pi) = 0$.
\end{enumerate}
\end{exercice}

\begin{exercice}[Champ de directions]
\label{exo:ch1-champ}
Tracer (à la main ou à l'ordinateur) le champ de directions de chacune des EDO suivantes et en déduire le comportement qualitatif des solutions :
\begin{enumerate}[label=(\alph*)]
\item $y' = -y + 1$
\item $y' = t\,y$
\item $y' = y(1-y)$
\end{enumerate}
\end{exercice}

\begin{exercice}[Condition de Lipschitz --- {\small deux étoiles}]
\label{exo:ch1-lipschitz}
Pour chacune des fonctions suivantes, déterminer si la condition de Lipschitz en $y$ est satisfaite au voisinage du point indiqué :
\begin{enumerate}[label=(\alph*)]
\item $f(t,y) = t^2 + \sin y$, au voisinage de $(0,0)$.
\item $f(t,y) = \sqrt{\abs{y}}$, au voisinage de $(0,0)$.
\item $f(t,y) = y^{2/3}$, au voisinage de $(0,0)$.
\item $f(t,y) = \frac{y}{t}$, au voisinage de $(0,1)$.
\end{enumerate}
\end{exercice}

\begin{exercice}[Non-unicité --- {\small deux étoiles}]
\label{exo:ch1-non-unicite}
Considérer l'EDO $y' = 3y^{2/3}$ avec la condition initiale $y(0) = 0$.
\begin{enumerate}[label=(\alph*)]
\item Montrer que $y(t) = 0$ est solution.
\item Montrer que $y(t) = t^3$ est aussi solution.
\item Montrer que pour tout $a \geq 0$, la fonction
\[
y_a(t) = \begin{cases} 0 & \text{si } t \leq a, \\ (t-a)^3 & \text{si } t > a, \end{cases}
\]
est solution. Pourquoi le théorème de Cauchy--Lipschitz ne s'applique-t-il pas ?
\end{enumerate}
\end{exercice}

\begin{exercice}[Modélisation : refroidissement --- {\small trois étoiles}]
\label{exo:ch1-refroidissement}
\index{loi de refroidissement de Newton}
La \emph{loi de refroidissement de Newton} stipule que la température $T(t)$ d'un corps placé dans un milieu à température constante $T_{\mathrm{ext}}$ évolue selon :
\[
\frac{\dd T}{\dd t} = -k(T - T_{\mathrm{ext}}), \quad k > 0.
\]
\begin{enumerate}[label=(\alph*)]
\item Résoudre cette EDO avec la condition initiale $T(0) = T_0$.
\item Un café est servi à $90\,{}^{\circ}\mathrm{C}$ dans une pièce à $20\,{}^{\circ}\mathrm{C}$. Après 5 minutes, sa température est de $70\,{}^{\circ}\mathrm{C}$. Déterminer $k$ et le temps nécessaire pour que le café atteigne $40\,{}^{\circ}\mathrm{C}$.
\end{enumerate}
\end{exercice}

\begin{exercice}[Modélisation : population logistique --- {\small trois étoiles}]
\label{exo:ch1-logistique}
On considère le modèle logistique \eqref{eq:logistique} avec $r = 0{,}5$, $K = 1000$ et $P(0) = 50$.
\begin{enumerate}[label=(\alph*)]
\item Montrer que la solution est $P(t) = \dfrac{K}{1 + \left(\frac{K}{P_0} - 1\right)e^{-rt}}$.
\item Calculer $P(10)$ et $P(20)$.
\item Déterminer le temps au bout duquel la population atteint $90\%$ de la capacité de charge.
\item Tracer l'allure de la courbe solution et interpréter le point d'inflexion.
\end{enumerate}
\end{exercice}

% ============================================================
% CHAPITRE 2
% ============================================================
\chapter{Équations du Premier Ordre}
\label{chap:premier-ordre}
\index{equations du premier ordre@équations du premier ordre}

Ce chapitre est consacré à l'étude systématique des méthodes de résolution des équations différentielles du premier ordre. Nous traitons successivement les équations à variables séparables, les équations linéaires, les équations de Bernoulli, les équations exactes, les équations homogènes, les équations de Riccati, et les équations de Clairaut et Lagrange. Chaque méthode est illustrée par des exemples détaillés.

\section{Équations à variables séparables}
\label{sec:separables}
\index{equation a variables separables@équation à variables séparables}

\begin{definition}[Rappel]
\label{def:separable-rappel}
Une EDO du premier ordre est dite \emph{à variables séparables} si elle s'écrit
\[
y' = g(t)\,h(y),
\]
ou de manière équivalente, $\dfrac{\dd y}{h(y)} = g(t)\,\dd t$ lorsque $h(y) \neq 0$.
\end{definition}

\begin{methode}[Résolution par séparation des variables]
\label{meth:separation}
\index{separation des variables@séparation des variables}
Pour résoudre $y' = g(t)\,h(y)$ :
\begin{enumerate}
\item \textbf{Solutions constantes :} chercher les valeurs $y_0$ telles que $h(y_0) = 0$. Chacune donne la solution constante $y(t) = y_0$.
\item \textbf{Séparation :} sur les intervalles où $h(y) \neq 0$, écrire $\dfrac{\dd y}{h(y)} = g(t)\,\dd t$.
\item \textbf{Intégration :} intégrer les deux membres : $\displaystyle\int \frac{\dd y}{h(y)} = \int g(t)\,\dd t + C$.
\item \textbf{Expression de $y$ :} si possible, exprimer $y$ explicitement en fonction de $t$.
\end{enumerate}
\end{methode}

\begin{exemple}[Équation séparable élémentaire]
\label{ex:separable-1}
Résoudre $y' = ty$.

\textbf{Solution.} C'est une équation à variables séparables avec $g(t) = t$ et $h(y) = y$.

\textit{Solution constante :} $h(y) = 0$ donne $y = 0$, qui est bien solution.

\textit{Pour $y \neq 0$ :} on sépare les variables :
\[
\frac{\dd y}{y} = t\,\dd t.
\]
En intégrant :
\[
\ln\abs{y} = \frac{t^2}{2} + C_1, \quad \text{soit} \quad \abs{y} = e^{C_1}\,e^{t^2/2}.
\]
En posant $C = \pm e^{C_1}$ (qui peut prendre toute valeur non nulle), on obtient $y(t) = C\,e^{t^2/2}$. La solution constante $y = 0$ correspond à $C = 0$. Ainsi, la solution générale est :
\[
\boxed{y(t) = C\,e^{t^2/2}, \quad C \in \R.}
\]
\end{exemple}

\begin{exemple}[Équation séparable avec condition initiale]
\label{ex:separable-2}
Résoudre le problème de Cauchy $y' = \dfrac{t}{y}$, $y(0) = 2$, pour $y > 0$.

\textbf{Solution.} On sépare les variables : $y\,\dd y = t\,\dd t$. En intégrant :
\[
\frac{y^2}{2} = \frac{t^2}{2} + C, \quad \text{soit} \quad y^2 = t^2 + 2C.
\]
La condition $y(0) = 2$ donne $4 = 2C$, donc $C = 2$. Comme $y > 0$ :
\[
\boxed{y(t) = \sqrt{t^2 + 4}.}
\]
Cette solution est définie pour tout $t \in \R$.
\end{exemple}

\begin{exemple}[Modèle logistique]
\label{ex:logistique-resolution}
\index{modele logistique@modèle logistique!resolution@résolution}
Résoudre $y' = ry\!\left(1 - \dfrac{y}{K}\right)$ avec $y(0) = y_0 > 0$, $0 < y_0 < K$.

\textbf{Solution.} On sépare les variables :
\[
\frac{\dd y}{y(1 - y/K)} = r\,\dd t.
\]
Décomposition en éléments simples :
\[
\frac{1}{y(1 - y/K)} = \frac{1}{y} + \frac{1/K}{1 - y/K} = \frac{1}{y} + \frac{1}{K - y}.
\]
En intégrant :
\[
\ln\abs{y} - \ln\abs{K - y} = rt + C_1, \quad \text{soit} \quad \ln\frac{y}{K - y} = rt + C_1.
\]
Donc $\dfrac{y}{K - y} = Ae^{rt}$ avec $A = e^{C_1}$. La condition initiale donne $A = \dfrac{y_0}{K - y_0}$. En résolvant pour $y$ :
\[
y = \frac{KAe^{rt}}{1 + Ae^{rt}} = \frac{K}{1 + \frac{1}{A}e^{-rt}} = \frac{K}{1 + \left(\frac{K}{y_0} - 1\right)e^{-rt}}.
\]
Ainsi :
\[
\boxed{y(t) = \frac{K}{1 + \left(\dfrac{K}{y_0} - 1\right)e^{-rt}}.}
\]
On vérifie que $y(t) \to K$ quand $t \to +\infty$.
\end{exemple}

\begin{figure}[ht]
\centering
\begin{tikzpicture}
\begin{axis}[
    title={Solutions de l'équation logistique ($r=1$, $K=10$)},
    xlabel={$t$},
    ylabel={$y$},
    domain=0:8,
    ymin=0, ymax=14,
    xmin=0, xmax=8,
    axis lines=left,
    width=11cm,
    height=7cm,
    legend pos=south east,
]
\addplot[red,thick,smooth,samples=100] {10/(1 + (10/1 - 1)*exp(-x))};
\addlegendentry{$y_0 = 1$}
\addplot[blue,thick,smooth,samples=100] {10/(1 + (10/5 - 1)*exp(-x))};
\addlegendentry{$y_0 = 5$}
\addplot[green!60!black,thick,smooth,samples=100] {10/(1 + (10/9 - 1)*exp(-x))};
\addlegendentry{$y_0 = 9$}
\addplot[orange,thick,smooth,samples=100] {10/(1 + (10/12 - 1)*exp(-x))};
\addlegendentry{$y_0 = 12$}
\addplot[dashed,gray] coordinates {(0,10) (8,10)};
\addlegendentry{$K = 10$}
\end{axis}
\end{tikzpicture}
\caption{Solutions de l'équation logistique pour différentes conditions initiales.}
\label{fig:logistique}
\end{figure}

\section{Équations linéaires du premier ordre}
\label{sec:lineaires-premier-ordre}
\index{equation lineaire du premier ordre@équation linéaire du premier ordre}

\begin{definition}
\label{def:lineaire-1}
Une \emph{équation linéaire du premier ordre} est une équation de la forme
\begin{equation}
\label{eq:lineaire-1}
y' + p(t)\,y = q(t),
\end{equation}
où $p$ et $q$ sont des fonctions continues sur un intervalle $I$. L'équation $y' + p(t)\,y = 0$ est l'\emph{équation homogène associée}\index{equation homogene associee@équation homogène associée}.
\end{definition}

\begin{theoreme}[Structure des solutions]
\label{thm:structure-lineaire-1}
\index{structure des solutions}
Soit $I$ un intervalle ouvert sur lequel $p$ et $q$ sont continues.
\begin{enumerate}[label=(\roman*)]
\item L'ensemble des solutions de l'équation homogène $y' + p(t)\,y = 0$ est un espace vectoriel de dimension $1$, engendré par
\[
y_h(t) = e^{-\int p(t)\,\dd t}.
\]
\item La solution générale de l'équation complète $y' + p(t)\,y = q(t)$ est
\[
y(t) = y_h(t) + y_p(t),
\]
où $y_p$ est une solution particulière quelconque de l'équation complète.
\end{enumerate}
\end{theoreme}

\begin{proof}
\textit{(i)} L'équation homogène $y' + p(t)y = 0$ est à variables séparables. Pour $y \neq 0$ :
\[
\frac{y'}{y} = -p(t), \quad \text{soit} \quad \ln\abs{y} = -\int p(t)\,\dd t + C_1.
\]
Donc $y(t) = C\,e^{-\int p(t)\,\dd t}$ avec $C \in \R$, ce qui inclut $y = 0$ pour $C = 0$. L'ensemble des solutions est bien un espace vectoriel de dimension $1$.

\textit{(ii)} Soit $y_p$ une solution particulière de l'équation complète. Si $y$ est une autre solution, alors $z = y - y_p$ satisfait $z' + p(t)z = 0$, donc $z$ est solution de l'équation homogène. Réciproquement, toute fonction $y = y_h + y_p$, avec $y_h$ solution de l'homogène, est solution de l'équation complète.
\end{proof}

\begin{methode}[Facteur intégrant --- Variation de la constante]
\label{meth:facteur-integrant}
\index{facteur integrant@facteur intégrant}
\index{variation de la constante}
Pour résoudre $y' + p(t)\,y = q(t)$ :
\begin{enumerate}
\item \textbf{Résoudre l'homogène :} $y_h(t) = C\,e^{-P(t)}$ où $P(t) = \int p(t)\,\dd t$.
\item \textbf{Variation de la constante :} poser $y(t) = C(t)\,e^{-P(t)}$ et substituer dans l'équation complète. On obtient $C'(t) = q(t)\,e^{P(t)}$.
\item \textbf{Intégrer :} $C(t) = \int q(t)\,e^{P(t)}\,\dd t + K$.
\item \textbf{Conclure :}
\[
y(t) = e^{-P(t)}\left(\int q(t)\,e^{P(t)}\,\dd t + K\right).
\]
\end{enumerate}
De manière équivalente, le \emph{facteur intégrant} est $\mu(t) = e^{P(t)}$ : en multipliant l'équation par $\mu$, le membre de gauche devient $(\mu\,y)' = \mu\,q$, ce qui s'intègre directement.
\end{methode}

\begin{exemple}[Équation linéaire avec condition initiale]
\label{ex:lineaire-1}
Résoudre $y' + 2t\,y = t$, $y(0) = 1$.

\textbf{Solution.} Ici $p(t) = 2t$ et $q(t) = t$.

\textit{Étape 1 :} $P(t) = \int 2t\,\dd t = t^2$.

\textit{Étape 2 :} Le facteur intégrant est $\mu(t) = e^{t^2}$. On multiplie :
\[
e^{t^2}\,y' + 2t\,e^{t^2}\,y = t\,e^{t^2}, \quad \text{soit} \quad \frac{\dd}{\dd t}\!\left(e^{t^2}\,y\right) = t\,e^{t^2}.
\]

\textit{Étape 3 :} En intégrant : $e^{t^2}\,y = \int t\,e^{t^2}\,\dd t = \frac{1}{2}e^{t^2} + C$.

\textit{Étape 4 :} $y(t) = \frac{1}{2} + C\,e^{-t^2}$.

La condition $y(0) = 1$ donne $1 = \frac{1}{2} + C$, donc $C = \frac{1}{2}$.
\[
\boxed{y(t) = \frac{1}{2}\!\left(1 + e^{-t^2}\right).}
\]
\end{exemple}

\begin{exemple}[Circuit RC]
\label{ex:circuit-rc}
\index{circuit RC!resolution@résolution}
Résoudre l'équation du circuit RC \eqref{eq:circuit-RC} : $RC\,u' + u = E_0$ (tension constante), avec $u(0) = 0$.

\textbf{Solution.} On met sous forme standard : $u' + \frac{1}{RC}\,u = \frac{E_0}{RC}$.

Ici $p = \frac{1}{RC}$ et $q = \frac{E_0}{RC}$. Le facteur intégrant est $\mu(t) = e^{t/(RC)}$.
\[
\frac{\dd}{\dd t}\!\left(e^{t/(RC)}\,u\right) = \frac{E_0}{RC}\,e^{t/(RC)}.
\]
En intégrant : $e^{t/(RC)}\,u = E_0\,e^{t/(RC)} + C$, soit $u(t) = E_0 + C\,e^{-t/(RC)}$.

La condition $u(0) = 0$ donne $C = -E_0$. D'où :
\[
\boxed{u(t) = E_0\!\left(1 - e^{-t/(RC)}\right).}
\]
La constante de temps est $\tau = RC$\index{constante de temps}. Après un temps $5\tau$, le condensateur est chargé à plus de $99\%$.
\end{exemple}

\begin{exemple}[Second membre non élémentaire]
\label{ex:lineaire-2}
Résoudre $y' - y = e^{t^2}$.

\textbf{Solution.} Ici $p(t) = -1$, $q(t) = e^{t^2}$, $P(t) = -t$.

La solution générale est :
\[
y(t) = e^{t}\left(\int e^{-t}\cdot e^{t^2}\,\dd t + C\right) = e^{t}\left(\int e^{t^2 - t}\,\dd t + C\right).
\]
L'intégrale $\int e^{t^2 - t}\,\dd t$ ne s'exprime pas en termes de fonctions élémentaires. La solution est donc laissée sous forme intégrale. Avec la condition $y(0) = y_0$, on écrit :
\[
\boxed{y(t) = e^{t}\left(y_0 + \int_0^t e^{s^2 - s}\,\dd s\right).}
\]
\end{exemple}

\begin{proposition}[Existence et unicité pour les EDO linéaires]
\label{prop:existence-lineaire}
Soient $p$ et $q$ des fonctions continues sur un intervalle ouvert $I$. Pour tout $t_0 \in I$ et tout $y_0 \in \R$, le problème de Cauchy
\[
y' + p(t)\,y = q(t), \quad y(t_0) = y_0,
\]
admet une unique solution définie sur $I$ tout entier.
\end{proposition}

\begin{proof}
La formule explicite est :
\[
y(t) = e^{-P(t)}\left(y_0\,e^{P(t_0)} + \int_{t_0}^{t} q(s)\,e^{P(s)}\,\dd s\right), \quad P(t) = \int_{t_0}^{t} p(s)\,\dd s.
\]
Cette fonction est bien définie et de classe $\mathscr{C}^1$ sur $I$ tout entier (car $p$ et $q$ sont continues sur $I$). On vérifie par un calcul direct qu'elle satisfait l'équation et la condition initiale. L'unicité provient de la théorie générale (Cauchy--Lipschitz), car $f(t,y) = q(t) - p(t)y$ est lipschitzienne en $y$ sur tout compact.
\end{proof}

\begin{remarque}
\label{rem:lineaire-global}
La solution d'une EDO linéaire est toujours globale : elle est définie sur tout l'intervalle $I$ où les coefficients sont continus. Ce n'est pas le cas en général pour les EDO non linéaires, dont les solutions peuvent exploser en temps fini.
\end{remarque}

\section{Équations de Bernoulli}
\label{sec:bernoulli}
\index{equation de Bernoulli@équation de Bernoulli}

\begin{definition}
\label{def:bernoulli}
Une \emph{équation de Bernoulli} est une équation de la forme
\begin{equation}
\label{eq:bernoulli}
y' + p(t)\,y = q(t)\,y^n,
\end{equation}
où $n \in \R$, $n \neq 0$, $n \neq 1$ (sinon l'équation est linéaire).
\end{definition}

\begin{methode}[Substitution de Bernoulli]
\label{meth:bernoulli}
\index{substitution de Bernoulli}
Pour résoudre \eqref{eq:bernoulli} :
\begin{enumerate}
\item \textbf{Solution triviale :} $y = 0$ est solution si $n > 0$.
\item \textbf{Substitution :} poser $v = y^{1-n}$. Alors $v' = (1-n)\,y^{-n}\,y'$.
\item \textbf{Linéarisation :} diviser \eqref{eq:bernoulli} par $y^n$ (pour $y \neq 0$) :
\[
y^{-n}\,y' + p(t)\,y^{1-n} = q(t).
\]
En termes de $v$ :
\[
\frac{v'}{1-n} + p(t)\,v = q(t), \quad \text{soit} \quad v' + (1-n)\,p(t)\,v = (1-n)\,q(t).
\]
\item \textbf{Résolution :} c'est une équation linéaire en $v$. La résoudre, puis revenir à $y = v^{1/(1-n)}$.
\end{enumerate}
\end{methode}

\begin{exemple}[Équation de Bernoulli]
\label{ex:bernoulli-1}
Résoudre $y' + \dfrac{y}{t} = t\,y^2$ pour $t > 0$.

\textbf{Solution.} C'est une équation de Bernoulli avec $p(t) = 1/t$, $q(t) = t$ et $n = 2$.

\textit{Solution triviale :} $y = 0$.

\textit{Substitution :} $v = y^{1-2} = y^{-1} = 1/y$. Alors $v' = -y'/y^2$, soit $y' = -v'/v^2$.

On divise l'équation par $y^2$ : $y^{-2}y' + \frac{1}{t}\,y^{-1} = t$, soit $-v' + \frac{v}{t} = t$, c'est-à-dire :
\[
v' - \frac{v}{t} = -t.
\]
C'est une équation linéaire en $v$. Le facteur intégrant est $\mu = e^{-\int (-1/t)\,\dd t} = e^{\ln t} = t$ \ldots{} Non, attention : $p_v(t) = -1/t$, donc $P_v(t) = -\ln t$ et $\mu(t) = e^{P_v(t)} = e^{-\ln t} = 1/t$. Multiplions par $1/t$ :
\[
\frac{v'}{t} - \frac{v}{t^2} = -1, \quad \text{soit} \quad \frac{\dd}{\dd t}\!\left(\frac{v}{t}\right) = -1.
\]
En intégrant : $\frac{v}{t} = -t + C$, donc $v(t) = -t^2 + Ct$.

En revenant à $y = 1/v$ :
\[
\boxed{y(t) = \frac{1}{Ct - t^2} = \frac{1}{t(C - t)}, \quad C \in \R.}
\]
La solution est définie pour $t \neq 0$ et $t \neq C$.
\end{exemple}

\begin{exemple}[Bernoulli avec condition initiale]
\label{ex:bernoulli-2}
Résoudre $y' - 2y = -y^3$, $y(0) = 1$.

\textbf{Solution.} Ici $p(t) = -2$, $q(t) = -1$, $n = 3$.

Substitution : $v = y^{1-3} = y^{-2}$. On divise par $y^3$ :
\[
y^{-3}y' - 2y^{-2} = -1, \quad \text{soit} \quad \frac{v'}{-2} - 2v = -1, \quad v' + 4v = 2.
\]
Solution homogène : $v_h = Ce^{-4t}$. Solution particulière : $v_p = 1/2$ (constante). Donc $v(t) = Ce^{-4t} + 1/2$.

Condition initiale : $v(0) = y(0)^{-2} = 1$, d'où $C + 1/2 = 1$, $C = 1/2$.
\[
v(t) = \frac{1}{2}\!\left(e^{-4t} + 1\right), \quad y = v^{-1/2} = \frac{1}{\sqrt{v}} = \frac{\sqrt{2}}{\sqrt{e^{-4t} + 1}}.
\]
\[
\boxed{y(t) = \frac{\sqrt{2}}{\sqrt{1 + e^{-4t}}}.}
\]
On vérifie : $y(0) = \frac{\sqrt{2}}{\sqrt{2}} = 1$. Et quand $t \to +\infty$, $y(t) \to \sqrt{2}$.
\end{exemple}

\section{Équations exactes}
\label{sec:exactes}
\index{equation exacte@équation exacte}

\begin{definition}
\label{def:exacte-rappel}
Considérons une EDO sous la forme
\begin{equation}
\label{eq:exacte-forme}
M(t,y)\,\dd t + N(t,y)\,\dd y = 0,
\end{equation}
où $M$ et $N$ sont de classe $\mathscr{C}^1$ sur un ouvert simplement connexe $\Omega \subset \R^2$. L'équation est dite \emph{exacte}\index{equation exacte@équation exacte!definition@définition} s'il existe une fonction $\Phi : \Omega \to \R$ de classe $\mathscr{C}^2$ telle que
\[
\frac{\partial \Phi}{\partial t} = M \quad \text{et} \quad \frac{\partial \Phi}{\partial y} = N.
\]
Les solutions sont alors données implicitement par $\Phi(t, y) = C$.
\end{definition}

\begin{proposition}[Critère d'exactitude]
\label{prop:critere-exactitude}
\index{critere d'exactitude@critère d'exactitude}
Si $M$ et $N$ sont de classe $\mathscr{C}^1$ sur un ouvert simplement connexe $\Omega$, alors l'équation $M\,\dd t + N\,\dd y = 0$ est exacte si et seulement si
\begin{equation}
\label{eq:critere-exactitude}
\frac{\partial M}{\partial y} = \frac{\partial N}{\partial t}.
\end{equation}
\end{proposition}

\begin{proof}
\textit{Condition nécessaire :} si $\Phi$ existe, alors $\frac{\partial^2 \Phi}{\partial y\,\partial t} = \frac{\partial M}{\partial y}$ et $\frac{\partial^2 \Phi}{\partial t\,\partial y} = \frac{\partial N}{\partial t}$. Par le théorème de Schwarz, ces dérivées croisées sont égales.

\textit{Condition suffisante :} supposons $\frac{\partial M}{\partial y} = \frac{\partial N}{\partial t}$. Fixons $(t_0, y_0) \in \Omega$ et posons
\[
\Phi(t, y) = \int_{t_0}^{t} M(s, y_0)\,\dd s + \int_{y_0}^{y} N(t, u)\,\dd u.
\]
Alors $\frac{\partial \Phi}{\partial y} = N(t, y)$ (par dérivation sous le signe intégral de la seconde intégrale, la première ne dépendant pas de $y$). Pour le calcul de $\frac{\partial \Phi}{\partial t}$ :
\[
\frac{\partial \Phi}{\partial t} = M(t, y_0) + \int_{y_0}^{y} \frac{\partial N}{\partial t}(t, u)\,\dd u = M(t, y_0) + \int_{y_0}^{y} \frac{\partial M}{\partial y}(t, u)\,\dd u = M(t, y).
\]
La dernière égalité provient du théorème fondamental de l'analyse.
\end{proof}

\begin{methode}[Résolution d'une équation exacte]
\label{meth:exacte}
\index{equation exacte@équation exacte!resolution@résolution}
Pour résoudre $M\,\dd t + N\,\dd y = 0$ exacte :
\begin{enumerate}
\item Vérifier que $\frac{\partial M}{\partial y} = \frac{\partial N}{\partial t}$.
\item Trouver $\Phi$ tel que $\frac{\partial \Phi}{\partial t} = M$ : intégrer $M$ par rapport à $t$ :
\[
\Phi(t, y) = \int M(t, y)\,\dd t + g(y),
\]
où $g(y)$ est une fonction inconnue de $y$ seul.
\item Déterminer $g(y)$ en imposant $\frac{\partial \Phi}{\partial y} = N$ :
\[
\frac{\partial}{\partial y}\!\left(\int M(t, y)\,\dd t\right) + g'(y) = N(t, y).
\]
\item La solution implicite est $\Phi(t, y) = C$.
\end{enumerate}
\end{methode}

\begin{exemple}[Équation exacte]
\label{ex:exacte-1}
Résoudre $(2ty + 3)\,\dd t + (t^2 + 4y)\,\dd y = 0$.

\textbf{Solution.} Posons $M(t,y) = 2ty + 3$ et $N(t,y) = t^2 + 4y$.

\textit{Vérification :} $\frac{\partial M}{\partial y} = 2t$ et $\frac{\partial N}{\partial t} = 2t$. L'équation est exacte.

\textit{Recherche de $\Phi$ :}
\[
\Phi(t, y) = \int M\,\dd t = \int (2ty + 3)\,\dd t = t^2 y + 3t + g(y).
\]
On impose $\frac{\partial \Phi}{\partial y} = N$ :
\[
t^2 + g'(y) = t^2 + 4y, \quad \text{d'où} \quad g'(y) = 4y, \quad g(y) = 2y^2.
\]
Donc $\Phi(t, y) = t^2 y + 3t + 2y^2$, et la solution est :
\[
\boxed{t^2 y + 3t + 2y^2 = C.}
\]
\end{exemple}

\begin{exemple}[Équation exacte avec condition initiale]
\label{ex:exacte-2}
Résoudre $(e^y + 2t)\,\dd t + (te^y - 1)\,\dd y = 0$, avec $y(0) = 0$.

\textbf{Solution.} $M = e^y + 2t$, $N = te^y - 1$.

$\frac{\partial M}{\partial y} = e^y$, $\frac{\partial N}{\partial t} = e^y$. L'équation est exacte.

$\Phi = \int (e^y + 2t)\,\dd t = te^y + t^2 + g(y)$.

$\frac{\partial \Phi}{\partial y} = te^y + g'(y) = te^y - 1$, donc $g'(y) = -1$, $g(y) = -y$.

$\Phi(t,y) = te^y + t^2 - y$. La condition $y(0) = 0$ donne $\Phi(0,0) = 0$, d'où $C = 0$ :
\[
\boxed{te^y + t^2 - y = 0.}
\]
\end{exemple}

\subsection{Facteurs intégrants pour les équations non exactes}
\label{subsec:facteurs-integrants}
\index{facteur integrant@facteur intégrant!pour equation non exacte@pour équation non exacte}

Si l'équation $M\,\dd t + N\,\dd y = 0$ n'est pas exacte, on cherche un \emph{facteur intégrant}\index{facteur integrant@facteur intégrant} $\mu = \mu(t,y)$ tel que $\mu M\,\dd t + \mu N\,\dd y = 0$ soit exacte, c'est-à-dire :
\[
\frac{\partial(\mu M)}{\partial y} = \frac{\partial(\mu N)}{\partial t}.
\]

En général, trouver un tel $\mu$ est aussi difficile que résoudre l'équation originale. Cependant, dans certains cas, $\mu$ ne dépend que d'une seule variable :

\begin{proposition}[Facteur intégrant en une variable]
\label{prop:facteur-integrant}
\begin{enumerate}[label=(\roman*)]
\item Si $\dfrac{1}{N}\!\left(\dfrac{\partial M}{\partial y} - \dfrac{\partial N}{\partial t}\right)$ ne dépend que de $t$, notons-la $\phi(t)$. Alors $\mu(t) = e^{\int \phi(t)\,\dd t}$ est un facteur intégrant.
\item Si $\dfrac{1}{M}\!\left(\dfrac{\partial N}{\partial t} - \dfrac{\partial M}{\partial y}\right)$ ne dépend que de $y$, notons-la $\psi(y)$. Alors $\mu(y) = e^{\int \psi(y)\,\dd y}$ est un facteur intégrant.
\end{enumerate}
\end{proposition}

\begin{proof}
\textit{(i)} Si $\mu = \mu(t)$, la condition d'exactitude $\frac{\partial(\mu M)}{\partial y} = \frac{\partial(\mu N)}{\partial t}$ devient :
\[
\mu \frac{\partial M}{\partial y} = \mu' N + \mu \frac{\partial N}{\partial t}, \quad \text{soit} \quad \frac{\mu'}{\mu} = \frac{1}{N}\!\left(\frac{\partial M}{\partial y} - \frac{\partial N}{\partial t}\right).
\]
Si le membre de droite est $\phi(t)$, on intègre : $\ln\mu = \int \phi(t)\,\dd t$.

Le cas \textit{(ii)} se traite de manière analogue.
\end{proof}

\begin{exemple}[Facteur intégrant]
\label{ex:facteur-integrant}
Résoudre $y\,\dd t - t\,\dd y = 0$.

\textbf{Solution.} $M = y$, $N = -t$. $\frac{\partial M}{\partial y} = 1$, $\frac{\partial N}{\partial t} = -1$. Non exacte.

On calcule $\frac{1}{N}\!\left(\frac{\partial M}{\partial y} - \frac{\partial N}{\partial t}\right) = \frac{1}{-t}(1 - (-1)) = -\frac{2}{t}$, qui ne dépend que de $t$. Donc $\mu(t) = e^{\int -2/t\,\dd t} = e^{-2\ln|t|} = t^{-2}$.

L'équation $\frac{y}{t^2}\,\dd t - \frac{1}{t}\,\dd y = 0$ est exacte. On trouve $\Phi = -\frac{y}{t}$, d'où :
\[
\boxed{\frac{y}{t} = C, \quad \text{soit} \quad y = Ct.}
\]
On retrouve bien l'ensemble des droites passant par l'origine.
\end{exemple}

\section{Équations homogènes}
\label{sec:homogenes}
\index{equation homogene@équation homogène!au sens de l'homogeneite@au sens de l'homogénéité}

\begin{definition}
\label{def:homogene}
Une EDO du premier ordre $y' = f(t,y)$ est dite \emph{homogène} (au sens de l'homogénéité des fonctions) si $f$ est homogène de degré $0$, c'est-à-dire si
\[
f(\lambda t, \lambda y) = f(t, y) \quad \text{pour tout } \lambda \neq 0.
\]
Cela revient à dire que $f(t,y)$ peut s'écrire comme une fonction de $y/t$ :
\[
y' = \varphi\!\left(\frac{y}{t}\right).
\]
\end{definition}

\begin{methode}[Substitution $u = y/t$]
\label{meth:homogene}
\index{substitution $u = y/t$}
Pour résoudre $y' = \varphi(y/t)$ :
\begin{enumerate}
\item Poser $u = y/t$, soit $y = tu$, d'où $y' = u + tu'$.
\item Substituer : $u + tu' = \varphi(u)$, soit $tu' = \varphi(u) - u$.
\item C'est une équation à variables séparables en $u$ et $t$ :
\[
\frac{\dd u}{\varphi(u) - u} = \frac{\dd t}{t}.
\]
\item Résoudre, puis revenir à $y = tu$.
\end{enumerate}
\end{methode}

\begin{exemple}[Équation homogène]
\label{ex:homogene-1}
Résoudre $y' = \dfrac{y}{t} + \dfrac{t}{y}$ pour $t > 0$, $y > 0$.

\textbf{Solution.} On écrit $y' = \frac{y}{t} + \frac{t}{y} = \frac{y^2 + t^2}{ty}$. Pour vérifier l'homogénéité : $f(\lambda t, \lambda y) = \frac{\lambda^2 y^2 + \lambda^2 t^2}{\lambda t \cdot \lambda y} = \frac{y^2 + t^2}{ty} = f(t,y)$.

Posons $u = y/t$, $y = tu$, $y' = u + tu'$. Alors :
\[
u + tu' = u + \frac{1}{u}, \quad \text{soit} \quad tu' = \frac{1}{u}.
\]
Séparation : $u\,\dd u = \frac{\dd t}{t}$. Intégration : $\frac{u^2}{2} = \ln t + C_1$.

Revenant à $y$ : $\frac{y^2}{2t^2} = \ln t + C_1$, soit :
\[
\boxed{y^2 = 2t^2(\ln t + C), \quad C \in \R.}
\]
On prend $y = t\sqrt{2(\ln t + C)}$ pour $\ln t + C > 0$.
\end{exemple}

\begin{exemple}[Équation homogène avec condition initiale]
\label{ex:homogene-2}
Résoudre $(t^2 + y^2)\,\dd t - 2ty\,\dd y = 0$, $y(1) = 1$.

\textbf{Solution.} On écrit $y' = \frac{t^2 + y^2}{2ty}$. On pose $u = y/t$ :
\[
u + tu' = \frac{1 + u^2}{2u}, \quad tu' = \frac{1 + u^2}{2u} - u = \frac{1 + u^2 - 2u^2}{2u} = \frac{1 - u^2}{2u}.
\]
Séparation :
\[
\frac{2u\,\dd u}{1 - u^2} = \frac{\dd t}{t}.
\]
On intègre : $-\ln\abs{1 - u^2} = \ln\abs{t} + C_1$, soit $\abs{1 - u^2} = \frac{A}{\abs{t}}$ avec $A > 0$.

En revenant à $y$ : $1 - \frac{y^2}{t^2} = \frac{A}{t}$ (en choisissant le signe approprié), soit $t^2 - y^2 = At$.

Condition initiale : $1 - 1 = A$, donc $A = 0$. D'où $y^2 = t^2$, et comme $y(1) = 1 > 0$, on a $y = t$ pour $t > 0$.
\[
\boxed{y(t) = t.}
\]
\end{exemple}

\section{Équations de Riccati}
\label{sec:riccati}
\index{equation de Riccati@équation de Riccati}

\begin{definition}
\label{def:riccati}
Une \emph{équation de Riccati} est une EDO du premier ordre de la forme
\begin{equation}
\label{eq:riccati}
y' = a(t)\,y^2 + b(t)\,y + c(t),
\end{equation}
où $a$, $b$, $c$ sont des fonctions continues et $a \not\equiv 0$.
\end{definition}

En général, une équation de Riccati ne se résout pas par quadratures. Cependant, si l'on connaît une \emph{solution particulière} $y_1(t)$, on peut ramener le problème à une équation de Bernoulli (et donc à une équation linéaire).

\begin{methode}[Réduction de Riccati]
\label{meth:riccati}
\index{equation de Riccati@équation de Riccati!reduction@réduction}
Soit $y_1$ une solution particulière de \eqref{eq:riccati}. On pose $y = y_1 + z$ et l'on substitue :
\[
y_1' + z' = a(y_1 + z)^2 + b(y_1 + z) + c.
\]
Comme $y_1' = ay_1^2 + by_1 + c$, il reste :
\[
z' = (2ay_1 + b)\,z + a\,z^2.
\]
C'est une équation de Bernoulli en $z$ avec $n = 2$. La substitution $v = 1/z$ donne l'équation linéaire :
\[
v' + (2ay_1 + b)\,v = -a.
\]
\end{methode}

\begin{exemple}[Équation de Riccati]
\label{ex:riccati}
Résoudre $y' = y^2 - \dfrac{2}{t^2}$, sachant que $y_1(t) = \dfrac{2}{t}$ est solution.

\textbf{Solution.} Vérifions : $y_1' = -\frac{2}{t^2}$, $y_1^2 - \frac{2}{t^2} = \frac{4}{t^2} - \frac{2}{t^2} = \frac{2}{t^2}$. Mais $y_1' = -\frac{2}{t^2} \neq \frac{2}{t^2}$. Essayons plutôt $y_1(t) = -\frac{1}{t}$ : $y_1' = \frac{1}{t^2}$, $y_1^2 - \frac{2}{t^2} = \frac{1}{t^2} - \frac{2}{t^2} = -\frac{1}{t^2}$. Non plus.

Prenons $y_1(t) = \frac{2}{t}$ : vérifions $y_1' = -\frac{2}{t^2}$ et $y_1^2 - \frac{2}{t^2} = \frac{4}{t^2} - \frac{2}{t^2} = \frac{2}{t^2} \neq -\frac{2}{t^2}$.

Corrigeons l'exemple. Considérons plutôt $y' = -y^2 + \frac{2}{t^2}$, avec $y_1(t) = \frac{1}{t}$.

Vérification : $y_1' = -\frac{1}{t^2}$, $-y_1^2 + \frac{2}{t^2} = -\frac{1}{t^2} + \frac{2}{t^2} = \frac{1}{t^2}$. Non : $-\frac{1}{t^2} \neq \frac{1}{t^2}$.

Prenons finalement l'équation $y' = y^2 + \frac{1}{t}\,y - \frac{1}{t^2}$ avec $y_1(t) = \frac{1}{t}$. Vérification : $y_1' = -\frac{1}{t^2}$, $y_1^2 + \frac{1}{t}\,y_1 - \frac{1}{t^2} = \frac{1}{t^2} + \frac{1}{t^2} - \frac{1}{t^2} = \frac{1}{t^2}$. Non : $-\frac{1}{t^2} \neq \frac{1}{t^2}$.

Prenons $y' = y^2 - \frac{2y}{t} + \frac{1}{t^2}$ avec $y_1(t) = \frac{1}{t}$. Vérification : $y_1' = -\frac{1}{t^2}$, $\frac{1}{t^2} - \frac{2}{t^2} + \frac{1}{t^2} = 0 \neq -\frac{1}{t^2}$.

Utilisons un exemple classique qui fonctionne. Considérons l'équation :
\[
y' = y^2 - ty + 1,
\]
et supposons connue la solution $y_1(t) = t$ (vérification : $y_1' = 1$, $t^2 - t^2 + 1 = 1$). C'est correct.

Posons $y = t + z$. Alors $z' + 1 = (t+z)^2 - t(t+z) + 1 = t^2 + 2tz + z^2 - t^2 - tz + 1$, d'où :
\[
z' = tz + z^2.
\]
C'est une Bernoulli avec $n = 2$. Posons $v = 1/z$ : $v' = -z'/z^2$, soit $v' + tv = -1$.

Facteur intégrant : $\mu = e^{t^2/2}$. Donc $\frac{\dd}{\dd t}(ve^{t^2/2}) = -e^{t^2/2}$.

L'intégrale $\int e^{t^2/2}\,\dd t$ ne s'exprime pas en fonctions élémentaires. On écrit :
\[
v(t) = e^{-t^2/2}\!\left(C - \int_0^t e^{s^2/2}\,\dd s\right), \quad z = \frac{1}{v},
\]
et la solution générale est :
\[
\boxed{y(t) = t + \frac{e^{t^2/2}}{C - \displaystyle\int_0^t e^{s^2/2}\,\dd s}.}
\]
\end{exemple}

\begin{remarque}
\label{rem:riccati-lien-lineaire}
L'équation de Riccati est intimement liée aux EDO linéaires du second ordre. Si l'on pose $y = -\frac{u'}{au}$ dans \eqref{eq:riccati}, on obtient une équation linéaire du second ordre en $u$. Réciproquement, toute EDO linéaire du second ordre peut être transformée en une équation de Riccati.
\end{remarque}

\section{Équations de Clairaut et de Lagrange}
\label{sec:clairaut-lagrange}

\subsection{Équation de Clairaut}
\label{subsec:clairaut}
\index{equation de Clairaut@équation de Clairaut}

\begin{definition}
\label{def:clairaut}
Une \emph{équation de Clairaut} est une EDO du premier ordre de la forme
\begin{equation}
\label{eq:clairaut}
y = ty' + \psi(y'),
\end{equation}
où $\psi$ est une fonction donnée de classe $\mathscr{C}^2$.
\end{definition}

\begin{theoreme}[Solutions de l'équation de Clairaut]
\label{thm:clairaut}
\index{equation de Clairaut@équation de Clairaut!solutions}
L'équation de Clairaut $y = ty' + \psi(y')$ possède :
\begin{enumerate}[label=(\roman*)]
\item une famille de droites (solutions générales) :
\[
y = Ct + \psi(C), \quad C \in \R ;
\]
\item une solution singulière, enveloppe de cette famille, donnée paramétriquement par
\[
\begin{cases}
t = -\psi'(p),\\
y = -p\,\psi'(p) + \psi(p),
\end{cases}
\]
où $p$ est le paramètre (valeur de $y'$).
\end{enumerate}
\end{theoreme}

\begin{proof}
Dérivons \eqref{eq:clairaut} par rapport à $t$ en posant $p = y'$ :
\[
p = p + t\,p' + \psi'(p)\,p', \quad \text{soit} \quad (t + \psi'(p))\,p' = 0.
\]
Deux cas :
\begin{itemize}
\item $p' = 0$ : alors $p = C$ (constante), et $y = Ct + \psi(C)$. Ce sont les droites.
\item $t + \psi'(p) = 0$ : alors $t = -\psi'(p)$, et en substituant dans \eqref{eq:clairaut}, $y = -p\psi'(p) + \psi(p)$. C'est la solution singulière.
\end{itemize}
La solution singulière est l'enveloppe de la famille de droites : en éliminant $C$ entre $y = Ct + \psi(C)$ et $0 = t + \psi'(C)$, on retrouve la même courbe paramétrée.
\end{proof}

\begin{exemple}[Équation de Clairaut]
\label{ex:clairaut}
Résoudre $y = ty' + (y')^2$.

\textbf{Solution.} Ici $\psi(p) = p^2$.

\textit{Solutions générales :} $y = Ct + C^2$, $C \in \R$ (famille de droites).

\textit{Solution singulière :} $t = -\psi'(p) = -2p$ et $y = -p(-2p) + p^2 = 2p^2 + p^2 = 3p^2$. Mais reprenons : $y = p(-2p) + p^2 = -2p^2 + p^2 = -p^2$. Comme $p = -t/2$, on obtient :
\[
y = -\frac{t^2}{4}.
\]
Vérifions : $y' = -t/2$, $ty' + (y')^2 = -t^2/2 + t^2/4 = -t^2/4 = y$.

La solution singulière est la parabole $y = -t^2/4$, enveloppe de la famille de droites $y = Ct + C^2$.
\[
\boxed{y = Ct + C^2 \quad (C \in \R) \quad \text{et} \quad y = -\frac{t^2}{4} \quad \text{(solution singulière).}}
\]
\end{exemple}

\begin{figure}[ht]
\centering
\begin{tikzpicture}
\begin{axis}[
    title={Solutions de $y = ty' + (y')^2$},
    xlabel={$t$},
    ylabel={$y$},
    domain=-4:4,
    ymin=-5, ymax=5,
    xmin=-4, xmax=4,
    axis lines=middle,
    width=10cm,
    height=8cm,
]
% Quelques droites y = Ct + C^2
\foreach \C in {-2,-1.5,-1,-0.5,0,0.5,1,1.5,2} {
    \addplot[blue!50,thin,domain=-4:4] {\C*x + \C*\C};
}
% Solution singulière : parabole y = -t^2/4
\addplot[red,very thick,smooth,domain=-4:4,samples=100] {-x^2/4};
\end{axis}
\end{tikzpicture}
\caption{Famille de droites et solution singulière (en rouge) de l'équation de Clairaut $y = ty' + (y')^2$.}
\label{fig:clairaut}
\end{figure}

\subsection{Équation de Lagrange}
\label{subsec:lagrange}
\index{equation de Lagrange@équation de Lagrange}

\begin{definition}
\label{def:lagrange}
Une \emph{équation de Lagrange} (ou de d'Alembert--Lagrange) est une EDO de la forme
\begin{equation}
\label{eq:lagrange}
y = t\,\varphi(y') + \psi(y'),
\end{equation}
où $\varphi$ et $\psi$ sont des fonctions données. L'équation de Clairaut est le cas particulier $\varphi(p) = p$.
\end{definition}

\begin{methode}[Résolution de l'équation de Lagrange]
\label{meth:lagrange}
\index{equation de Lagrange@équation de Lagrange!resolution@résolution}
On pose $p = y'$ et on dérive \eqref{eq:lagrange} par rapport à $t$ :
\[
p = \varphi(p) + t\,\varphi'(p)\,p' + \psi'(p)\,p'.
\]
Si $\varphi(p) \neq p$ pour tout $p$, on peut résoudre pour $p'$ ou, mieux, considérer $t$ comme fonction de $p$. En posant $\frac{\dd t}{\dd p} = \frac{1}{p'}$, on obtient une EDO linéaire en $t(p)$ :
\[
\frac{\dd t}{\dd p} = \frac{\varphi'(p)}{p - \varphi(p)}\,t + \frac{\psi'(p)}{p - \varphi(p)}.
\]
On résout cette équation linéaire, puis on déduit $y = t\varphi(p) + \psi(p)$, ce qui donne une représentation paramétrique $(t(p), y(p))$ de la solution.
\end{methode}

\begin{exemple}[Équation de Lagrange]
\label{ex:lagrange}
Résoudre $y = 2ty' - (y')^2$.

\textbf{Solution.} Ici $\varphi(p) = 2p$, $\psi(p) = -p^2$. On dérive par rapport à $t$ :
\[
p = 2p + 2tp' - 2pp', \quad \text{soit} \quad -p = (2t - 2p)p'.
\]
En réarrangeant, si $p' \neq 0$ :
\[
\frac{\dd t}{\dd p} = \frac{2t - 2p}{-p} = -\frac{2}{p}\,t + 2.
\]
C'est une équation linéaire en $t(p)$ : $t' + \frac{2}{p}\,t = 2$.

Facteur intégrant : $\mu = e^{\int 2/p\,\dd p} = p^2$. Donc $(p^2 t)' = 2p^2$, soit $p^2 t = \frac{2p^3}{3} + C$, d'où :
\[
t = \frac{2p}{3} + \frac{C}{p^2}.
\]
Et $y = 2pt - p^2 = 2p\!\left(\frac{2p}{3} + \frac{C}{p^2}\right) - p^2 = \frac{4p^2}{3} + \frac{2C}{p} - p^2 = \frac{p^2}{3} + \frac{2C}{p}$.

La solution est donnée paramétriquement par :
\[
\boxed{\begin{cases} t = \dfrac{2p}{3} + \dfrac{C}{p^2},\\[8pt] y = \dfrac{p^2}{3} + \dfrac{2C}{p}. \end{cases}}
\]

Il faut aussi vérifier si $p' = 0$ donne des solutions. Si $p = p_0$ constant, alors $y = 2p_0 t - p_0^2$ et $y' = 2p_0$. Mais on a posé $p = y' = 2p_0$, donc $p_0 = 2p_0$, soit $p_0 = 0$. Cela donne la solution $y = 0$.
\end{exemple}

\section{Tableau récapitulatif des méthodes}
\label{sec:tableau-recapitulatif}

\begin{table}[ht]
\centering
\caption{Tableau récapitulatif des méthodes de résolution des EDO du premier ordre.}
\label{tab:recap-methodes}
\renewcommand{\arraystretch}{1.4}
\begin{tabular}{@{}p{3cm}p{3.5cm}p{6cm}@{}}
\toprule
\textbf{Type} & \textbf{Forme canonique} & \textbf{Méthode de résolution} \\
\midrule
Séparable & $y' = g(t)\,h(y)$ & Séparer : $\int \frac{\dd y}{h(y)} = \int g(t)\,\dd t$ \\[4pt]
Linéaire & $y' + p(t)y = q(t)$ & Facteur intégrant $\mu = e^{\int p\,\dd t}$ ou variation de la constante \\[4pt]
Bernoulli & $y' + py = qy^n$ & Substitution $v = y^{1-n}$, puis linéaire \\[4pt]
Exacte & $M\,\dd t + N\,\dd y = 0$ & Critère $\partial_y M = \partial_t N$, recherche de $\Phi$ \\[4pt]
Non exacte & $M\,\dd t + N\,\dd y = 0$ & Facteur intégrant $\mu(t)$ ou $\mu(y)$ \\[4pt]
Homogène & $y' = \varphi(y/t)$ & Substitution $u = y/t$, puis séparable \\[4pt]
Riccati & $y' = ay^2 + by + c$ & Solution particulière connue $y_1$, puis $y = y_1 + z$ (Bernoulli) \\[4pt]
Clairaut & $y = ty' + \psi(y')$ & Droites $y = Ct + \psi(C)$ et enveloppe \\[4pt]
Lagrange & $y = t\varphi(y') + \psi(y')$ & Poser $p = y'$, EDO linéaire en $t(p)$ \\
\bottomrule
\end{tabular}
\end{table}

\section{Résumé du chapitre}
\label{sec:resume-ch2}

\begin{tcolorbox}[colback=blue!3,colframe=blue!50!black,title={\textbf{Résumé --- Chapitre 2}}]
\begin{itemize}[leftmargin=*]
\item Les \textbf{équations séparables} $y' = g(t)h(y)$ se résolvent par séparation et intégration. Ne pas oublier les solutions constantes.
\item Les \textbf{équations linéaires} $y' + py = q$ se résolvent par facteur intégrant ou variation de la constante. La solution est toujours globale.
\item Les \textbf{équations de Bernoulli} $y' + py = qy^n$ se ramènent à une équation linéaire par $v = y^{1-n}$.
\item Les \textbf{équations exactes} $M\,\dd t + N\,\dd y = 0$ admettent un potentiel $\Phi$ si et seulement si $\partial_y M = \partial_t N$. Sinon, on cherche un facteur intégrant.
\item Les \textbf{équations homogènes} $y' = \varphi(y/t)$ se ramènent à des équations séparables par $u = y/t$.
\item Les \textbf{équations de Riccati} nécessitent la connaissance d'une solution particulière.
\item Les \textbf{équations de Clairaut} admettent une famille de droites et une solution singulière (enveloppe).
\item Les \textbf{équations de Lagrange} se résolvent par paramétrisation.
\end{itemize}
\end{tcolorbox}

\section{Exercices}
\label{sec:exercices-ch2}

\subsection*{Équations à variables séparables}

\begin{exercice}[Séparable --- une étoile]
\label{exo:ch2-sep-1}
Résoudre les équations suivantes :
\begin{enumerate}[label=(\alph*)]
\item $y' = \dfrac{t^2}{y}$, $y > 0$.
\item $y' = e^{t-y}$.
\item $y' = y\cos t$, $y(0) = 3$.
\item $y' = \dfrac{1+y^2}{1+t^2}$.
\end{enumerate}
\end{exercice}

\begin{exercice}[Séparable --- deux étoiles]
\label{exo:ch2-sep-2}
\begin{enumerate}[label=(\alph*)]
\item Résoudre $y' = y^2(1-y)$ et discuter le comportement des solutions selon la condition initiale.
\item Montrer que la solution de $y' = 1 + y^2$, $y(0) = 0$ est $y(t) = \tan t$, définie sur $]-\pi/2, \pi/2[$. Commenter : la solution explose en temps fini.
\end{enumerate}
\end{exercice}

\subsection*{Équations linéaires}

\begin{exercice}[Linéaire --- une étoile]
\label{exo:ch2-lin-1}
Résoudre les équations suivantes :
\begin{enumerate}[label=(\alph*)]
\item $y' + 3y = 6$.
\item $y' - y = e^{2t}$.
\item $ty' + 2y = t^3$, $t > 0$.
\item $y' + y\tan t = \cos t$, $y(0) = 1$.
\end{enumerate}
\end{exercice}

\begin{exercice}[Linéaire --- deux étoiles]
\label{exo:ch2-lin-2}
\begin{enumerate}[label=(\alph*)]
\item Résoudre $y' + \dfrac{y}{t} = \dfrac{\sin t}{t}$ pour $t > 0$.
\item Résoudre $y' + 2ty = 2te^{-t^2}$.
\item Montrer que toutes les solutions de $y' + y = \sin t$ convergent vers la même fonction quand $t \to +\infty$, et déterminer cette fonction.
\end{enumerate}
\end{exercice}

\subsection*{Équations de Bernoulli}

\begin{exercice}[Bernoulli --- deux étoiles]
\label{exo:ch2-bern}
Résoudre les équations suivantes :
\begin{enumerate}[label=(\alph*)]
\item $y' + y = y^3$.
\item $y' - \dfrac{y}{t} = -t^2 y^3$, $t > 0$.
\item $y' + ty = ty^{1/2}$, $y > 0$, $y(0) = 4$.
\end{enumerate}
\end{exercice}

\subsection*{Équations exactes}

\begin{exercice}[Exacte --- une étoile]
\label{exo:ch2-exact-1}
Vérifier que les équations suivantes sont exactes et les résoudre :
\begin{enumerate}[label=(\alph*)]
\item $(3t^2 + 4ty)\,\dd t + (2t^2 + 3y^2)\,\dd y = 0$.
\item $(\cos y - y\sin t)\,\dd t + (t\cos t - \sin y + \cos y)\,\dd y = 0$. \textit{(Vérifier d'abord.)}
\item $(ye^{ty} + 2t)\,\dd t + (te^{ty} - 1)\,\dd y = 0$.
\end{enumerate}
\end{exercice}

\begin{exercice}[Facteur intégrant --- deux étoiles]
\label{exo:ch2-exact-2}
Les équations suivantes ne sont pas exactes. Trouver un facteur intégrant et résoudre :
\begin{enumerate}[label=(\alph*)]
\item $(y + t)\,\dd t + (3t - y)\,\dd y = 0$. \textit{Chercher $\mu = \mu(t)$.}
\item $y\,\dd t + (2ty + 1)\,\dd y = 0$. \textit{Chercher $\mu = \mu(y)$.}
\end{enumerate}
\end{exercice}

\subsection*{Équations homogènes}

\begin{exercice}[Homogène --- une étoile]
\label{exo:ch2-hom-1}
Résoudre :
\begin{enumerate}[label=(\alph*)]
\item $y' = \dfrac{y - t}{y + t}$.
\item $t\,y' = y + \sqrt{t^2 + y^2}$, $t > 0$.
\item $y' = \dfrac{t + 2y}{t}$, $y(1) = 0$.
\end{enumerate}
\end{exercice}

\subsection*{Riccati, Clairaut, Lagrange}

\begin{exercice}[Riccati --- deux étoiles]
\label{exo:ch2-riccati}
\begin{enumerate}[label=(\alph*)]
\item Résoudre $y' = y^2 + \frac{1}{t}\,y - \frac{1}{t^2}$, sachant que $y_1(t) = \frac{1}{t}$ est solution. \textit{Vérifier d'abord.}
\item Résoudre $y' = 1 + t^2 - 2ty + y^2$, sachant que $y_1(t) = t$ est solution.
\end{enumerate}
\end{exercice}

\begin{exercice}[Clairaut --- deux étoiles]
\label{exo:ch2-clairaut}
Résoudre les équations de Clairaut suivantes et déterminer les solutions singulières :
\begin{enumerate}[label=(\alph*)]
\item $y = ty' - (y')^2$.
\item $y = ty' + \dfrac{1}{y'}$, $y' \neq 0$.
\end{enumerate}
\end{exercice}

\subsection*{Problèmes de modélisation}

\begin{exercice}[Mélange --- {\small trois étoiles}]
\label{exo:ch2-melange}
\index{probleme de melange@problème de mélange}
Un réservoir contient initialement 100 litres d'eau pure. On y verse une solution saline de concentration $0{,}5\;\mathrm{kg/L}$ au débit de $2\;\mathrm{L/min}$. Le mélange, supposé homogène, s'écoule au même débit.
\begin{enumerate}[label=(\alph*)]
\item Soit $q(t)$ la quantité de sel (en kg) dans le réservoir à l'instant $t$ (en min). Montrer que $q$ satisfait l'EDO $q' + \frac{q}{50} = 1$, $q(0) = 0$.
\item Résoudre ce problème de Cauchy.
\item Déterminer $\lim_{t\to+\infty} q(t)$ et la concentration limite.
\item Au bout de combien de temps la concentration dans le réservoir atteint-elle $0{,}4\;\mathrm{kg/L}$\,?
\end{enumerate}
\end{exercice}

\begin{exercice}[Chute avec résistance de l'air --- {\small trois étoiles}]
\label{exo:ch2-chute}
\index{chute avec resistance de l'air@chute avec résistance de l'air}
Un objet de masse $m$ tombe sous l'effet de la gravité, avec une force de résistance proportionnelle au carré de la vitesse : $F_{\mathrm{air}} = -kv^2$ ($k > 0$). En prenant l'axe vers le bas, $v(t) > 0$ :
\begin{enumerate}[label=(\alph*)]
\item Montrer que la vitesse $v(t)$ satisfait $m\,v' = mg - kv^2$.
\item Montrer que c'est une équation séparable et la résoudre avec $v(0) = 0$.
\item Déterminer la vitesse limite $v_{\infty} = \lim_{t\to+\infty} v(t)$.
\item Application numérique : $m = 80\;\mathrm{kg}$, $g = 9{,}81\;\mathrm{m/s^2}$, $k = 0{,}25\;\mathrm{kg/m}$. Calculer $v_{\infty}$ et le temps pour atteindre $95\%$ de $v_{\infty}$.
\end{enumerate}
\end{exercice}

\begin{exercice}[Courbes orthogonales --- {\small trois étoiles}]
\label{exo:ch2-orthogonales}
\index{courbes orthogonales}
Soit la famille de courbes $y = Ce^{-t}$, $C \in \R$.
\begin{enumerate}[label=(\alph*)]
\item Montrer que les courbes de cette famille satisfont $y' = -y$.
\item Les \emph{trajectoires orthogonales} sont les courbes qui coupent à angle droit chaque courbe de la famille. Montrer qu'elles satisfont $y' = 1/y$.
\item Résoudre cette équation et tracer quelques courbes de chaque famille.
\end{enumerate}
\end{exercice}

\begin{exercice}[Synthèse --- {\small trois étoiles}]
\label{exo:ch2-synthese}
Pour chacune des équations suivantes, identifier le type, choisir la méthode appropriée et résoudre :
\begin{enumerate}[label=(\alph*)]
\item $(t^2 + 1)y' + 2ty = t$.
\item $y' = \dfrac{y^2 + ty}{t^2}$.
\item $(2t + y^2)\,\dd t + 2ty\,\dd y = 0$.
\item $y' + y = y^2 e^t$.
\item $y = ty' + \sin(y')$.
\item $y' = \dfrac{y}{t} + \tan\!\left(\dfrac{y}{t}\right)$, $t > 0$.
\end{enumerate}
\end{exercice}

% === FIN DU CHUNK preamble+ch1-2 ===
% === DEBUT DU CHUNK ch3-4 ===

%%%%%%%%%%%%%%%%%%%%%%%%%%%%%%%%%%%%%%%%%%%%%%%%%%%%%%%%%%%%%%%%%%%%
%  Chapitre 3 : Theoremes d'Existence et d'Unicite
%%%%%%%%%%%%%%%%%%%%%%%%%%%%%%%%%%%%%%%%%%%%%%%%%%%%%%%%%%%%%%%%%%%%
\chapter{Theoremes d'Existence et d'Unicite}\label{ch:existence-unicite}
\index{existence et unicite@existence et unicite}

\section*{Introduction et motivation}

Considerons le probleme de Cauchy\index{probleme de Cauchy@probleme de Cauchy} general :
\begin{equation}\label{eq:cauchy-general}
\begin{cases}
y'(t) = f(t, y(t)), \\
y(t_0) = y_0,
\end{cases}
\end{equation}
ou $f : \Omega \to \R^n$ est definie sur un ouvert $\Omega \subset \R \times \R^n$, et $(t_0, y_0) \in \Omega$.

Deux questions fondamentales se posent :
\begin{enumerate}[label=(\roman*)]
  \item \textbf{Existence} : le probleme~\eqref{eq:cauchy-general} admet-il au moins une solution ?
  \item \textbf{Unicite} : cette solution est-elle la seule ?
\end{enumerate}

Comme nous le verrons, la continuite de $f$ suffit a garantir l'existence (theoreme de Cauchy--Peano), mais pas l'unicite. Pour obtenir l'unicite, il faut une hypothese supplementaire : la condition de Lipschitz.

%% -------------------------------------------------------------------
\section{Condition de Lipschitz}\label{sec:lipschitz}
\index{Lipschitz!condition de}

\begin{definition}[Fonction lipschitzienne par rapport a $y$]\label{def:lipschitz}
Soit $f : \Omega \to \R^n$ definie sur un ouvert $\Omega \subset \R \times \R^n$. On dit que $f$ est \emph{lipschitzienne par rapport a la seconde variable} $y$ sur $\Omega$ s'il existe une constante $L \geq 0$ telle que
\[
  \norm{f(t, y_1) - f(t, y_2)} \leq L\, \norm{y_1 - y_2}
\]
pour tout $(t, y_1), (t, y_2) \in \Omega$. La constante $L$ est appelee \emph{constante de Lipschitz}\index{Lipschitz!constante de}.
\end{definition}

\begin{definition}[Lipschitz locale]\label{def:lipschitz-locale}
\index{Lipschitz!locale}
On dit que $f$ est \emph{localement lipschitzienne par rapport a $y$} si, pour tout compact $K \subset \Omega$, il existe $L_K \geq 0$ tel que
\[
  \norm{f(t, y_1) - f(t, y_2)} \leq L_K \, \norm{y_1 - y_2}
\]
pour tout $(t, y_1), (t, y_2) \in K$.
\end{definition}

\begin{proposition}[Critere de differentiabilite]\label{prop:C1-lipschitz}
\index{Lipschitz!et differentiabilite}
Si $f : \Omega \to \R^n$ est de classe $\mathcal{C}^1$ par rapport a $y$, alors $f$ est localement lipschitzienne par rapport a $y$.
\end{proposition}

\begin{proof}
Soit $K \subset \Omega$ un compact convexe en $y$. Par le theoreme des accroissements finis, pour $(t, y_1), (t, y_2) \in K$ :
\[
  \norm{f(t, y_1) - f(t, y_2)} \leq \sup_{(t,y) \in K} \norm{\frac{\partial f}{\partial y}(t, y)} \cdot \norm{y_1 - y_2}.
\]
Comme $\frac{\partial f}{\partial y}$ est continue sur le compact $K$, le supremum est fini. On pose $L_K = \sup_{K} \norm{\frac{\partial f}{\partial y}}$.
\end{proof}

\begin{exemple}[Fonctions lipschitziennes]\label{ex:lipschitz-exemples}
\begin{enumerate}[label=(\alph*)]
  \item $f(t, y) = t^2 + 3y$ est globalement lipschitzienne avec $L = 3$.
  \item $f(t, y) = y^2$ est localement lipschitzienne mais pas globalement lipschitzienne. En effet, sur $\abs{y} \leq M$, on a $\abs{y_1^2 - y_2^2} = \abs{y_1 + y_2}\abs{y_1 - y_2} \leq 2M\abs{y_1 - y_2}$.
  \item $f(t, y) = \sin(y)$ est globalement lipschitzienne avec $L = 1$.
  \item $f(t, y) = \abs{y}^{2/3}$ n'est \emph{pas} lipschitzienne en $y = 0$.
\end{enumerate}
\end{exemple}

\begin{exercice}[$\star$]\label{exo:verif-lipschitz}
Determiner si les fonctions suivantes sont lipschitziennes par rapport a $y$. Si oui, trouver une constante de Lipschitz.
\begin{enumerate}[label=(\alph*)]
  \item $f(t, y) = e^{-t} y + \cos(t)$ sur $\R \times \R$.
  \item $f(t, y) = \sqrt{\abs{y}}$ sur $\R \times \R$.
  \item $f(t, y) = t \sin(y^2)$ sur $[-1,1] \times [-M, M]$.
\end{enumerate}
\end{exercice}

%% -------------------------------------------------------------------
\section{Theoreme de Cauchy--Lipschitz local}\label{sec:cauchy-lipschitz}
\index{Cauchy--Lipschitz!theoreme de}
\index{Picard--Lindelof@Picard--Lindelof}

\begin{theoreme}[Cauchy--Lipschitz local]\label{thm:cauchy-lipschitz}
Soit $\Omega \subset \R \times \R^n$ un ouvert et $f : \Omega \to \R^n$ une fonction continue, localement lipschitzienne par rapport a $y$. Alors, pour tout $(t_0, y_0) \in \Omega$, le probleme de Cauchy
\[
\begin{cases}
y'(t) = f(t, y(t)), \\
y(t_0) = y_0,
\end{cases}
\]
admet une \emph{unique} solution maximale.
\end{theoreme}

Nous allons demontrer ce theoreme en utilisant la methode des \emph{approximations successives de Picard}\index{Picard!iterations de}. La preuve se decompose en plusieurs etapes.

\subsection{Construction des iterees de Picard}\label{subsec:picard-construction}

\begin{remarque}[Formulation integrale]\label{rem:formulation-integrale}
\index{formulation integrale}
Le probleme de Cauchy~\eqref{eq:cauchy-general} est equivalent a l'equation integrale :
\begin{equation}\label{eq:formulation-integrale}
y(t) = y_0 + \int_{t_0}^{t} f(s, y(s))\, \dd s.
\end{equation}
En effet, si $y$ est de classe $\mathcal{C}^1$ et verifie~\eqref{eq:cauchy-general}, alors par integration on obtient~\eqref{eq:formulation-integrale}. Reciproquement, si $y$ est continue et verifie~\eqref{eq:formulation-integrale}, alors $y$ est $\mathcal{C}^1$ et verifie~\eqref{eq:cauchy-general}.
\end{remarque}

On definit l'\emph{operateur de Picard}\index{Picard!operateur de} $T$ par :
\begin{equation}\label{eq:operateur-picard}
(Ty)(t) = y_0 + \int_{t_0}^{t} f(s, y(s))\, \dd s.
\end{equation}
Trouver une solution de~\eqref{eq:cauchy-general} revient a trouver un \emph{point fixe} de $T$.

\begin{definition}[Iterees de Picard]\label{def:iterees-picard}
\index{Picard!iterees de}
On definit la suite de fonctions $(y_k)_{k \geq 0}$ par recurrence :
\begin{align}
y_0(t) &= y_0, \label{eq:picard-init} \\
y_{k+1}(t) &= y_0 + \int_{t_0}^{t} f(s, y_k(s))\, \dd s, \quad k \geq 0. \label{eq:picard-recurrence}
\end{align}
\end{definition}

\subsection{Preuve complete du theoreme de Cauchy--Lipschitz}\label{subsec:preuve-cl}

\begin{proof}[Preuve du theoreme~\ref{thm:cauchy-lipschitz}]
\textbf{Etape 1 : Localisation.}
Puisque $\Omega$ est ouvert et $(t_0, y_0) \in \Omega$, il existe $a > 0$ et $b > 0$ tels que le pave
\[
  R = \bigl\{ (t,y) \in \R \times \R^n : \abs{t - t_0} \leq a, \; \norm{y - y_0} \leq b \bigr\} \subset \Omega.
\]
Comme $f$ est continue sur le compact $R$, elle y est bornee : posons $M = \sup_R \norm{f(t,y)}$. De plus, par hypothese, $f$ est lipschitzienne en $y$ sur $R$ avec une constante $L > 0$.

Posons $\alpha = \min\bigl(a, \, b/M\bigr)$ et travaillons sur l'intervalle $I = [t_0 - \alpha, \, t_0 + \alpha]$.

\medskip
\textbf{Etape 2 : Espace fonctionnel.}
Considerons l'espace de Banach\index{Banach!espace de}
\[
  E = \mathcal{C}(I, \R^n), \quad \text{muni de la norme } \norm{\varphi}_\infty = \sup_{t \in I} \norm{\varphi(t)},
\]
et le sous-ensemble ferme
\[
  B = \bigl\{ \varphi \in E : \norm{\varphi - y_0}_\infty \leq b \bigr\}.
\]
L'operateur de Picard $T : B \to E$ est defini par~\eqref{eq:operateur-picard}.

\medskip
\textbf{Etape 3 : $T$ envoie $B$ dans $B$.}
Pour $\varphi \in B$ et $t \in I$, on a :
\[
  \norm{(T\varphi)(t) - y_0} = \norm{\int_{t_0}^{t} f(s, \varphi(s))\, \dd s} \leq M \abs{t - t_0} \leq M \alpha \leq b.
\]
Donc $T\varphi \in B$.

\medskip
\textbf{Etape 4 : Convergence des iterees de Picard.}
Montrons par recurrence que pour tout $k \geq 0$ et $t \in I$ :
\begin{equation}\label{eq:picard-estimation}
\norm{y_{k+1}(t) - y_k(t)} \leq \frac{M L^k}{(k+1)!} \abs{t - t_0}^{k+1}.
\end{equation}
\emph{Initialisation} ($k = 0$) : $\norm{y_1(t) - y_0} = \norm{\int_{t_0}^t f(s, y_0)\, \dd s} \leq M \abs{t - t_0}$. \checkmark

\emph{Heredite} : supposons~\eqref{eq:picard-estimation} vraie au rang $k$. Alors :
\begin{align*}
\norm{y_{k+2}(t) - y_{k+1}(t)}
  &= \norm{\int_{t_0}^{t} \bigl[f(s, y_{k+1}(s)) - f(s, y_k(s))\bigr]\, \dd s} \\
  &\leq \int_{t_0}^{t} L \, \norm{y_{k+1}(s) - y_k(s)} \, \dd s \\
  &\leq \int_{t_0}^{t} L \cdot \frac{M L^k}{(k+1)!} \abs{s - t_0}^{k+1} \, \dd s \\
  &= \frac{M L^{k+1}}{(k+1)!} \cdot \frac{\abs{t - t_0}^{k+2}}{k+2}
  = \frac{M L^{k+1}}{(k+2)!} \abs{t - t_0}^{k+2}.
\end{align*}

La serie telescopique $y_0 + \sum_{k=0}^{\infty} (y_{k+1} - y_k)$ converge donc normalement sur $I$, car :
\[
  \sum_{k=0}^{\infty} \norm{y_{k+1} - y_k}_\infty \leq \sum_{k=0}^{\infty} \frac{M L^k \alpha^{k+1}}{(k+1)!} \leq \frac{M}{L} \sum_{k=0}^{\infty} \frac{(L\alpha)^{k+1}}{(k+1)!} \leq \frac{M}{L}\bigl(e^{L\alpha} - 1\bigr) < \infty.
\]
Soit $y^* = \lim_{k \to \infty} y_k$. Cette fonction est continue sur $I$.

\medskip
\textbf{Etape 5 : $y^*$ est un point fixe de $T$.}
Par convergence uniforme, on peut passer a la limite dans~\eqref{eq:picard-recurrence} :
\[
  y^*(t) = \lim_{k \to \infty} y_{k+1}(t) = y_0 + \lim_{k \to \infty} \int_{t_0}^{t} f(s, y_k(s))\, \dd s = y_0 + \int_{t_0}^{t} f(s, y^*(s))\, \dd s.
\]
L'interversion limite-integrale est justifiee par la convergence uniforme de $f(s, y_k(s))$ vers $f(s, y^*(s))$ (car $f$ est continue et $y_k \to y^*$ uniformement).

Donc $y^*$ est solution du probleme de Cauchy.

\medskip
\textbf{Etape 6 : Unicite.}
Soient $y$ et $z$ deux solutions sur $I$. Posons $\varphi(t) = \norm{y(t) - z(t)}$. Alors :
\[
  \varphi(t) = \norm{\int_{t_0}^{t} \bigl[f(s, y(s)) - f(s, z(s))\bigr]\, \dd s} \leq L \int_{t_0}^{t} \varphi(s)\, \dd s.
\]
Par le lemme de Gronwall (theoreme~\ref{thm:gronwall} ci-dessous), $\varphi(t) = 0$ pour tout $t \in I$, d'ou $y = z$.
\end{proof}

\begin{remarque}\label{rem:norme-ponderee}
On peut aussi utiliser une \emph{norme ponderee}\index{norme ponderee} $\norm{\varphi}_\lambda = \sup_{t \in I} e^{-\lambda\abs{t-t_0}} \norm{\varphi(t)}$ avec $\lambda > L$ pour montrer que $T$ est une contraction stricte, et conclure par le theoreme du point fixe de Banach\index{Banach!theoreme du point fixe}.
\end{remarque}

%% -------------------------------------------------------------------
\section{Methode de Picard : exemples detailles}\label{sec:picard-exemples}

\begin{exemple}[Picard pour $y' = y$, $y(0) = 1$]\label{ex:picard-exp}
\index{Picard!exemple exponentielle}
Ici $f(t,y) = y$. Les iterees de Picard sont :
\begin{align*}
y_0(t) &= 1, \\
y_1(t) &= 1 + \int_0^t 1\, \dd s = 1 + t, \\
y_2(t) &= 1 + \int_0^t (1+s)\, \dd s = 1 + t + \frac{t^2}{2}, \\
y_3(t) &= 1 + \int_0^t \Bigl(1 + s + \frac{s^2}{2}\Bigr)\, \dd s = 1 + t + \frac{t^2}{2} + \frac{t^3}{6}.
\end{align*}
Par recurrence, $y_k(t) = \sum_{j=0}^{k} \frac{t^j}{j!}$, qui converge vers $e^t$.
\end{exemple}

\begin{exemple}[Picard pour $y' = 2ty$, $y(0) = 1$]\label{ex:picard-gauss}
Ici $f(t,y) = 2ty$. Les iterees :
\begin{align*}
y_0(t) &= 1, \\
y_1(t) &= 1 + \int_0^t 2s \cdot 1 \, \dd s = 1 + t^2, \\
y_2(t) &= 1 + \int_0^t 2s(1 + s^2)\, \dd s = 1 + t^2 + \frac{t^4}{2}, \\
y_3(t) &= 1 + t^2 + \frac{t^4}{2} + \frac{t^6}{6}.
\end{align*}
On reconnait les sommes partielles de $\sum_{k=0}^\infty \frac{t^{2k}}{k!} = e^{t^2}$.
\end{exemple}

\begin{exercice}[$\star\star$]\label{exo:picard-iteration}
Calculer les trois premieres iterees de Picard pour les problemes suivants et conjecturer la limite.
\begin{enumerate}[label=(\alph*)]
  \item $y' = -y + 1$, $y(0) = 0$.
  \item $y' = t + y^2$, $y(0) = 0$ (calculer $y_1$, $y_2$, $y_3$).
  \item $y' = 1 + y^2$, $y(0) = 0$.
\end{enumerate}
\end{exercice}

%% -------------------------------------------------------------------
\section{Theoreme de Cauchy--Peano}\label{sec:cauchy-peano}
\index{Cauchy--Peano!theoreme de}
\index{Peano!theoreme de}

\begin{theoreme}[Cauchy--Peano]\label{thm:cauchy-peano}
Soit $\Omega \subset \R \times \R^n$ un ouvert et $f : \Omega \to \R^n$ une fonction \emph{continue}. Alors, pour tout $(t_0, y_0) \in \Omega$, le probleme de Cauchy~\eqref{eq:cauchy-general} admet \emph{au moins une} solution locale.
\end{theoreme}

\begin{remarque}\label{rem:peano-pas-unicite}
La continuite seule ne suffit \emph{pas} pour l'unicite, comme le montre le contre-exemple classique suivant.
\end{remarque}

\begin{exemple}[Non-unicite : $y' = y^{2/3}$]\label{ex:non-unicite}
\index{non-unicite}
Considerons le probleme de Cauchy :
\[
  y' = y^{2/3}, \quad y(0) = 0.
\]
La fonction $f(t,y) = y^{2/3}$ est continue mais n'est \emph{pas} lipschitzienne en $y = 0$ (car $\frac{\partial f}{\partial y} = \frac{2}{3} y^{-1/3} \to \infty$ quand $y \to 0$).

Ce probleme admet la solution triviale $y(t) = 0$ et aussi, pour tout $c \geq 0$, la solution :
\[
  y_c(t) = \begin{cases} 0 & \text{si } t \leq c, \\ \left(\dfrac{t - c}{3}\right)^3 & \text{si } t > c. \end{cases}
\]
On verifie aisement que chaque $y_c$ est de classe $\mathcal{C}^1$ et satisfait l'equation. Il y a donc une \emph{infinite} de solutions.
\end{exemple}

% --- Figure TikZ : solutions non uniques ---
\begin{figure}[ht]
\centering
\begin{tikzpicture}[scale=1.0]
  \draw[->] (-1.5,0) -- (5,0) node[right] {$t$};
  \draw[->] (0,-0.3) -- (0,3.5) node[above] {$y$};
  % Solution triviale
  \draw[thick, blue] (-1.5,0) -- (5,0) node[below left] {$y \equiv 0$};
  % Solutions non triviales
  \foreach \c in {0, 0.8, 1.6, 2.4} {
    \draw[thick, red!70!black, domain=\c:4.5, samples=80]
      plot (\x, {((\x-\c)/3)^3});
  }
  \node[above right, red!70!black] at (3.2, 1.5) {$y_c(t) = \left(\frac{t-c}{3}\right)^{\!3}$};
  \node at (2.5, -0.8) {Infinite de solutions passant par $(0,0)$};
\end{tikzpicture}
\caption{Non-unicite pour $y' = y^{2/3}$, $y(0) = 0$.}\label{fig:non-unicite}
\end{figure}

%% -------------------------------------------------------------------
\section{Solutions maximales et explosion en temps fini}\label{sec:solutions-maximales}
\index{solution maximale}
\index{explosion en temps fini}

\begin{definition}[Solution maximale]\label{def:solution-maximale}
Une solution $y : J \to \R^n$ du probleme de Cauchy~\eqref{eq:cauchy-general} est dite \emph{maximale} si elle ne peut pas etre prolongee a un intervalle strictement plus grand. L'intervalle $J$ est appele \emph{intervalle maximal d'existence}\index{intervalle maximal}.
\end{definition}

\begin{theoreme}[Alternative d'explosion]\label{thm:explosion}
\index{alternative d'explosion}
Soit $y : {]}\omega^-, \omega^+{[} \to \R^n$ une solution maximale du probleme de Cauchy~\eqref{eq:cauchy-general}, ou $f$ est localement lipschitzienne. Si $\omega^+ < +\infty$, alors $y(t)$ \og{}explose\fg{} : pour tout compact $K \subset \Omega$, il existe $t_K < \omega^+$ tel que $(t, y(t)) \notin K$ pour $t > t_K$.

En particulier, $\norm{y(t)} \to +\infty$ quand $t \to \omega^+$ (si $\Omega = \R \times \R^n$).
\end{theoreme}

\begin{exemple}[Explosion en temps fini : $y' = y^2$]\label{ex:explosion-y2}
\index{explosion!exemple}
Le probleme $y' = y^2$, $y(0) = 1$ a pour solution :
\[
  y(t) = \frac{1}{1 - t}, \quad t \in \left]{-\infty}, 1\right[.
\]
L'intervalle maximal est $\left]{-\infty}, 1\right[$ et $y(t) \to +\infty$ quand $t \to 1^-$. La solution \og{}explose\fg{} en temps fini $\omega^+ = 1$.
\end{exemple}

\begin{exemple}[Solution globale vs locale]\label{ex:global-vs-local}
Comparons deux equations :
\begin{enumerate}[label=(\alph*)]
  \item $y' = y$ : solution $y(t) = e^t$, definie sur $\R$ tout entier (solution \emph{globale}).
  \item $y' = 1 + y^2$ : solution $y(t) = \tan(t)$ avec $y(0) = 0$, definie sur ${]}{-\pi/2}, \pi/2{[}$ seulement.
\end{enumerate}
\end{exemple}

%% -------------------------------------------------------------------
\section{Lemme de Gronwall et dependance continue}\label{sec:gronwall}
\index{Gronwall!lemme de}

\begin{lemme}[Gronwall, forme differentielle]\label{thm:gronwall}
Soit $\varphi : [t_0, T] \to \R_+$ une fonction continue verifiant :
\[
  \varphi(t) \leq \alpha + \beta \int_{t_0}^{t} \varphi(s)\, \dd s, \quad \forall\, t \in [t_0, T],
\]
ou $\alpha \geq 0$ et $\beta \geq 0$ sont des constantes. Alors :
\begin{equation}\label{eq:gronwall}
  \varphi(t) \leq \alpha \, e^{\beta(t - t_0)}, \quad \forall\, t \in [t_0, T].
\end{equation}
\end{lemme}

\begin{proof}
Posons $\Phi(t) = \int_{t_0}^{t} \varphi(s)\, \dd s$. Alors $\Phi'(t) = \varphi(t) \leq \alpha + \beta \Phi(t)$, donc :
\[
  \frac{\dd}{\dd t}\bigl[e^{-\beta(t-t_0)} \Phi(t)\bigr] = e^{-\beta(t-t_0)}\bigl[\Phi'(t) - \beta \Phi(t)\bigr] \leq \alpha\, e^{-\beta(t-t_0)}.
\]
En integrant de $t_0$ a $t$ et en utilisant $\Phi(t_0) = 0$ :
\[
  e^{-\beta(t-t_0)} \Phi(t) \leq \frac{\alpha}{\beta}\bigl(1 - e^{-\beta(t-t_0)}\bigr).
\]
Donc $\beta\Phi(t) \leq \alpha\bigl(e^{\beta(t-t_0)} - 1\bigr)$, et :
\[
  \varphi(t) \leq \alpha + \beta\Phi(t) \leq \alpha + \alpha\bigl(e^{\beta(t-t_0)} - 1\bigr) = \alpha\, e^{\beta(t-t_0)}. \qedhere
\]
\end{proof}

\begin{corollaire}[Gronwall, cas $\alpha = 0$]\label{cor:gronwall-zero}
Si $\varphi(t) \leq \beta \int_{t_0}^{t} \varphi(s)\, \dd s$ pour tout $t \in [t_0, T]$, alors $\varphi(t) = 0$ sur $[t_0, T]$.
\end{corollaire}

\begin{theoreme}[Dependance continue par rapport aux conditions initiales]\label{thm:dependance-continue}
\index{dependance continue}
Sous les hypotheses du theoreme de Cauchy--Lipschitz, soient $y$ et $z$ les solutions des problemes de Cauchy avec conditions initiales $y(t_0) = y_0$ et $z(t_0) = z_0$ respectivement. Alors, sur tout intervalle commun de definition $[t_0, T]$ :
\[
  \norm{y(t) - z(t)} \leq \norm{y_0 - z_0}\, e^{L(t - t_0)}, \quad \forall\, t \in [t_0, T].
\]
\end{theoreme}

\begin{proof}
Posons $\varphi(t) = \norm{y(t) - z(t)}$. Alors :
\begin{align*}
\varphi(t) &= \norm{y_0 - z_0 + \int_{t_0}^{t}\bigl[f(s, y(s)) - f(s, z(s))\bigr]\, \dd s} \\
  &\leq \norm{y_0 - z_0} + L \int_{t_0}^{t} \varphi(s)\, \dd s.
\end{align*}
Le lemme de Gronwall avec $\alpha = \norm{y_0 - z_0}$ et $\beta = L$ donne le resultat.
\end{proof}

\begin{remarque}\label{rem:sensibilite}
Ce resultat montre que l'erreur croit au plus exponentiellement avec le temps. Pour les systemes chaotiques, cette borne est essentiellement optimale.
\end{remarque}

%% -------------------------------------------------------------------
\section{Theoremes de comparaison}\label{sec:comparaison}
\index{comparaison!theoreme de}

\begin{theoreme}[Comparaison en dimension 1]\label{thm:comparaison}
Soient $f, g : [t_0, T] \times \R \to \R$ continues et localement lipschitziennes en $y$, avec $f(t, y) \leq g(t, y)$ pour tout $(t, y)$. Soient $y$ et $z$ les solutions de $y' = f(t, y)$ et $z' = g(t, z)$ respectivement, avec $y(t_0) \leq z(t_0)$. Alors :
\[
  y(t) \leq z(t), \quad \forall\, t \in [t_0, T].
\]
\end{theoreme}

\begin{proof}
Posons $w(t) = z(t) - y(t)$. Alors $w(t_0) \geq 0$ et :
\[
  w'(t) = g(t, z(t)) - f(t, y(t)) \geq f(t, z(t)) - f(t, y(t)).
\]
Si $w(t_1) < 0$ pour un certain $t_1 > t_0$, soit $\tau = \inf\{t > t_0 : w(t) < 0\}$. Alors $w(\tau) = 0$ et $w'(\tau) \leq 0$. Mais $w'(\tau) = g(\tau, z(\tau)) - f(\tau, y(\tau)) \geq g(\tau, z(\tau)) - f(\tau, z(\tau)) \geq 0$, car $y(\tau) = z(\tau)$. Contradiction avec $w'(\tau) \leq 0$ et $w$ strictement decroissante en $\tau$.

Une preuve rigoureuse utilise une perturbation $z_\varepsilon' = g(t, z_\varepsilon) + \varepsilon$ avec passage a la limite $\varepsilon \to 0$.
\end{proof}

\begin{exemple}\label{ex:comparaison}
Soit $y' = -y + \sin(t)$, $y(0) = 0$. Comme $-1 \leq \sin(t) \leq 1$, en comparant avec $z' = -z + 1$, $z(0) = 0$ (dont la solution est $z(t) = 1 - e^{-t}$), on obtient $y(t) \leq 1 - e^{-t} < 1$ pour tout $t > 0$.
\end{exemple}

%% -------------------------------------------------------------------
\section{Contre-exemples fondamentaux}\label{sec:contre-exemples}

\begin{exemple}[Recapitulatif : non-unicite]\label{ex:recap-non-unicite}
Pour $y' = 3 y^{2/3}$, $y(0) = 0$, la fonction $f(t,y) = 3y^{2/3}$ est continue mais pas lipschitzienne en $y = 0$. Les solutions $y \equiv 0$ et $y(t) = t^3$ coexistent. Plus generalement, pour tout $c \geq 0$, la fonction
\[
  y_c(t) = \begin{cases} 0 & \text{si } t \leq c, \\ (t - c)^3 & \text{si } t > c \end{cases}
\]
est solution.
\end{exemple}

\begin{exemple}[Recapitulatif : explosion en temps fini]\label{ex:recap-explosion}
Pour $y' = y^2$, $y(0) = y_0 > 0$, la solution $y(t) = \frac{y_0}{1 - y_0 t}$ explose en $t^* = 1/y_0$. Plus $y_0$ est grand, plus l'explosion est rapide.
\end{exemple}

% --- Figure TikZ : iterees de Picard convergeant ---
\begin{figure}[ht]
\centering
\begin{tikzpicture}[scale=1.2]
  \draw[->] (-0.3,0) -- (3.5,0) node[right] {$t$};
  \draw[->] (0,-0.3) -- (0,4) node[above] {$y$};
  % Solution exacte e^t
  \draw[thick, black, domain=0:1.3, samples=60] plot (\x, {exp(\x)}) node[right] {$e^t$};
  % y_0
  \draw[dashed, blue!60, domain=0:1.5] plot (\x, 1) node[right] {\tiny $y_0$};
  % y_1 = 1 + t
  \draw[dashed, red!60, domain=0:1.5] plot (\x, {1 + \x}) node[right] {\tiny $y_1$};
  % y_2 = 1 + t + t^2/2
  \draw[dashed, green!60!black, domain=0:1.4] plot (\x, {1 + \x + \x*\x/2}) node[right] {\tiny $y_2$};
  % y_3
  \draw[dashed, orange!80!black, domain=0:1.35] plot (\x, {1 + \x + \x*\x/2 + \x*\x*\x/6}) node[right] {\tiny $y_3$};
  \node at (1.8, -0.6) {Iterees de Picard pour $y' = y$};
\end{tikzpicture}
\caption{Convergence des iterees de Picard vers $e^t$.}\label{fig:picard-iterees}
\end{figure}

%% -------------------------------------------------------------------
\section{Exercices}\label{sec:exos-ch3}

\begin{exercice}[$\star$]\label{exo:lipschitz-verif}
Montrer que $f(t, y) = t^2 \cos(y)$ est globalement lipschitzienne en $y$ sur $[-T, T] \times \R$ et donner une constante de Lipschitz.
\end{exercice}

\begin{exercice}[$\star$]\label{exo:gronwall-application}
Soit $\varphi : [0, 1] \to \R_+$ continue verifiant $\varphi(t) \leq 2 + 3\int_0^t \varphi(s)\, \dd s$. Montrer que $\varphi(t) \leq 2e^{3t}$ pour tout $t \in [0,1]$.
\end{exercice}

\begin{exercice}[$\star\star$]\label{exo:picard-sin}
Considerer $y' = \sin(y)$, $y(0) = \pi/2$.
\begin{enumerate}[label=(\alph*)]
  \item Verifier que la condition de Lipschitz est satisfaite.
  \item Calculer les deux premieres iterees de Picard $y_1(t)$ et $y_2(t)$.
\end{enumerate}
\end{exercice}

\begin{exercice}[$\star\star$]\label{exo:explosion-temps}
Pour l'equation $y' = 1 + y^2$, $y(0) = 0$ :
\begin{enumerate}[label=(\alph*)]
  \item Resoudre explicitement.
  \item Determiner l'intervalle maximal d'existence.
  \item Calculer les trois premieres iterees de Picard et comparer avec la solution exacte.
\end{enumerate}
\end{exercice}

\begin{exercice}[$\star\star$]\label{exo:dependance-ci}
On considere $y' = -y + e^{-t}$, $y(0) = y_0$. Resoudre, puis montrer directement (sans utiliser le theoreme~\ref{thm:dependance-continue}) que si $\abs{y_0 - z_0} \leq \varepsilon$, alors les solutions correspondantes satisfont $\abs{y(t) - z(t)} \leq \varepsilon \, e^{-t}$ pour $t \geq 0$.
\end{exercice}

\begin{exercice}[$\star\star\star$]\label{exo:gronwall-generalise}
\textbf{(Gronwall generalise.)} Soit $\varphi, \psi : [t_0, T] \to \R_+$ des fonctions continues avec $\psi$ croissante, verifiant :
\[
  \varphi(t) \leq \psi(t) + \beta \int_{t_0}^{t} \varphi(s)\, \dd s.
\]
Montrer que $\varphi(t) \leq \psi(t) \, e^{\beta(t - t_0)}$.
\end{exercice}

\begin{exercice}[$\star\star\star$]\label{exo:cauchy-lipschitz-systeme}
Soit le systeme $\mathbf{y}' = A(t)\mathbf{y} + \mathbf{b}(t)$ ou $A : [0,T] \to \mathcal{M}_n(\R)$ et $\mathbf{b} : [0,T] \to \R^n$ sont continues. Montrer que la solution existe sur $[0, T]$ tout entier (pas d'explosion en temps fini).
\end{exercice}

%% -------------------------------------------------------------------
\section*{Resume du chapitre}

\begin{center}
\renewcommand{\arraystretch}{1.4}
\begin{tabular}{@{}p{4.5cm}p{9cm}@{}}
\toprule
\textbf{Resultat} & \textbf{Enonce succinct} \\
\midrule
Cauchy--Lipschitz & $f$ continue + Lipschitz en $y$ $\Rightarrow$ existence et unicite locale \\
Cauchy--Peano & $f$ continue $\Rightarrow$ existence (pas d'unicite en general) \\
Gronwall & $\varphi \leq \alpha + \beta \int \varphi$ $\Rightarrow$ $\varphi \leq \alpha e^{\beta t}$ \\
Dependance continue & Petite perturbation de $y_0$ $\Rightarrow$ petite perturbation de $y(t)$ \\
Alternative d'explosion & Solution maximale : ou bien globale, ou bien $\norm{y} \to \infty$ \\
Comparaison & $f \leq g$ et $y_0 \leq z_0$ $\Rightarrow$ $y(t) \leq z(t)$ \\
\bottomrule
\end{tabular}
\end{center}


%%%%%%%%%%%%%%%%%%%%%%%%%%%%%%%%%%%%%%%%%%%%%%%%%%%%%%%%%%%%%%%%%%%%
%  Chapitre 4 : Equations Lineaires du 2nd Ordre a Coefficients Constants
%%%%%%%%%%%%%%%%%%%%%%%%%%%%%%%%%%%%%%%%%%%%%%%%%%%%%%%%%%%%%%%%%%%%
\chapter{Equations Lineaires du Second Ordre a Coefficients Constants}\label{ch:edl2}
\index{equation lineaire!du second ordre}
\index{coefficients constants}

\section*{Introduction et motivation}

Les equations differentielles lineaires du second ordre a coefficients constants interviennent dans de nombreux domaines de la physique et de l'ingenierie.

\begin{exemple}[Systeme masse-ressort]\label{ex:masse-ressort}
\index{masse-ressort!systeme}
Considerons une masse $m$ attachee a un ressort de raideur $k$, avec un amortissement visqueux de coefficient $c$. Le principe fondamental de la dynamique donne :
\[
  m\, x''(t) + c\, x'(t) + k\, x(t) = F(t),
\]
ou $x(t)$ est le deplacement par rapport a la position d'equilibre et $F(t)$ une force exterieure.
\end{exemple}

\begin{exemple}[Circuit RLC]\label{ex:rlc}
\index{circuit RLC}
Pour un circuit RLC serie avec une tension $e(t)$, la charge $q(t)$ verifie :
\[
  L\, q''(t) + R\, q'(t) + \frac{1}{C}\, q(t) = e(t).
\]
\end{exemple}

Dans toute la suite, nous etudions l'equation generale :
\begin{equation}\label{eq:edl2-generale}
  a\, y''(t) + b\, y'(t) + c\, y(t) = f(t), \quad a \neq 0,
\end{equation}
que l'on divise par $a$ pour obtenir la forme standard :
\begin{equation}\label{eq:edl2-standard}
  y''(t) + p\, y'(t) + q\, y(t) = g(t),
\end{equation}
ou $p = b/a$, $q = c/a$ et $g = f/a$.

%% -------------------------------------------------------------------
\section{Equation homogene}\label{sec:edl2-homogene}
\index{equation homogene!du second ordre}

\subsection{Equation caracteristique}\label{subsec:eq-caract}
\index{equation caracteristique}

Considerons l'equation homogene :
\begin{equation}\label{eq:edl2-homogene}
  a\, y'' + b\, y' + c\, y = 0, \quad a \neq 0.
\end{equation}
Cherchons des solutions de la forme $y(t) = e^{rt}$. En substituant :
\[
  a\, r^2 e^{rt} + b\, r\, e^{rt} + c\, e^{rt} = 0 \implies e^{rt}(ar^2 + br + c) = 0.
\]
Comme $e^{rt} \neq 0$, on obtient l'\emph{equation caracteristique} :
\begin{equation}\label{eq:caract}
  ar^2 + br + c = 0.
\end{equation}

\begin{definition}[Discriminant]\label{def:discriminant}
\index{discriminant}
Le \emph{discriminant} de l'equation caracteristique est $\Delta = b^2 - 4ac$.
\end{definition}

Trois cas se presentent.

\subsection{Cas 1 : $\Delta > 0$ (deux racines reelles distinctes)}\label{subsec:delta-positif}
\index{racines!reelles distinctes}

\begin{theoreme}[Solution, cas $\Delta > 0$]\label{thm:delta-pos}
Si $\Delta > 0$, l'equation caracteristique a deux racines reelles distinctes :
\[
  r_1 = \frac{-b - \sqrt{\Delta}}{2a}, \quad r_2 = \frac{-b + \sqrt{\Delta}}{2a}.
\]
La solution generale de~\eqref{eq:edl2-homogene} est :
\begin{equation}\label{eq:sol-delta-pos}
  y(t) = C_1\, e^{r_1 t} + C_2\, e^{r_2 t}, \quad C_1, C_2 \in \R.
\end{equation}
\end{theoreme}

\begin{proof}
Les fonctions $y_1(t) = e^{r_1 t}$ et $y_2(t) = e^{r_2 t}$ sont solutions (par construction). Elles sont lineairement independantes car leur wronskien (voir section~\ref{subsec:wronskien}) est :
\[
  W(t) = (r_2 - r_1) e^{(r_1 + r_2)t} \neq 0
\]
puisque $r_1 \neq r_2$. Par le theoreme de structure des equations lineaires (l'espace des solutions est de dimension $2$), toute solution est combinaison lineaire de $y_1$ et $y_2$.
\end{proof}

\begin{exemple}\label{ex:delta-pos}
\index{amortissement!surcritique}
Resoudre $y'' - 5y' + 6y = 0$.

Equation caracteristique : $r^2 - 5r + 6 = 0$, donc $\Delta = 25 - 24 = 1 > 0$ et $r_1 = 2$, $r_2 = 3$.

Solution generale : $y(t) = C_1 e^{2t} + C_2 e^{3t}$.
\end{exemple}

\subsection{Cas 2 : $\Delta = 0$ (racine double)}\label{subsec:delta-nul}
\index{racine double}

\begin{theoreme}[Solution, cas $\Delta = 0$]\label{thm:delta-zero}
Si $\Delta = 0$, l'equation caracteristique a une racine double $r_0 = -b/(2a)$. La solution generale est :
\begin{equation}\label{eq:sol-delta-zero}
  y(t) = (C_1 + C_2\, t)\, e^{r_0 t}, \quad C_1, C_2 \in \R.
\end{equation}
\end{theoreme}

\begin{proof}
La fonction $y_1(t) = e^{r_0 t}$ est solution. Pour trouver une seconde solution independante, utilisons la \emph{reduction d'ordre} (section~\ref{sec:reduction-ordre}) : posons $y(t) = v(t) e^{r_0 t}$. En substituant dans~\eqref{eq:edl2-homogene} et en utilisant $2ar_0 + b = 0$ (car $r_0 = -b/(2a)$) :
\[
  a\, v''(t)\, e^{r_0 t} = 0 \implies v''(t) = 0 \implies v(t) = C_1 + C_2\, t.
\]
Les fonctions $e^{r_0 t}$ et $t\, e^{r_0 t}$ sont lineairement independantes.
\end{proof}

\begin{exemple}\label{ex:delta-zero}
\index{amortissement!critique}
Resoudre $y'' + 4y' + 4y = 0$.

Equation caracteristique : $r^2 + 4r + 4 = (r+2)^2 = 0$, donc $r_0 = -2$.

Solution generale : $y(t) = (C_1 + C_2\, t)\, e^{-2t}$.
\end{exemple}

\subsection{Cas 3 : $\Delta < 0$ (racines complexes conjuguees)}\label{subsec:delta-negatif}
\index{racines!complexes conjuguees}

\begin{theoreme}[Solution, cas $\Delta < 0$]\label{thm:delta-neg}
Si $\Delta < 0$, l'equation caracteristique a deux racines complexes conjuguees :
\[
  r = \alpha \pm i\beta, \quad \text{ou} \quad \alpha = -\frac{b}{2a}, \quad \beta = \frac{\sqrt{-\Delta}}{2a}.
\]
La solution generale \emph{reelle} est :
\begin{equation}\label{eq:sol-delta-neg}
  y(t) = e^{\alpha t}\bigl(C_1\, \cos(\beta t) + C_2\, \sin(\beta t)\bigr), \quad C_1, C_2 \in \R.
\end{equation}
\end{theoreme}

\begin{proof}
Formellement, les solutions complexes sont $e^{(\alpha + i\beta)t}$ et $e^{(\alpha - i\beta)t}$. Par la formule d'Euler :
\[
  e^{(\alpha + i\beta)t} = e^{\alpha t}(\cos(\beta t) + i\sin(\beta t)).
\]
En prenant la partie reelle et la partie imaginaire, on obtient deux solutions reelles lineairement independantes : $y_1(t) = e^{\alpha t}\cos(\beta t)$ et $y_2(t) = e^{\alpha t}\sin(\beta t)$.
\end{proof}

\begin{remarque}[Forme amplitude-phase]\label{rem:amplitude-phase}
\index{amplitude-phase}
La solution peut s'ecrire sous la forme :
\[
  y(t) = A\, e^{\alpha t} \cos(\beta t - \phi),
\]
ou $A = \sqrt{C_1^2 + C_2^2}$ est l'amplitude et $\phi = \arctan(C_2/C_1)$ la phase.
\end{remarque}

\begin{exemple}\label{ex:delta-neg}
\index{oscillations!amorties}
Resoudre $y'' + 2y' + 5y = 0$.

Equation caracteristique : $r^2 + 2r + 5 = 0$, donc $\Delta = 4 - 20 = -16 < 0$. Les racines sont $r = -1 \pm 2i$.

Solution generale : $y(t) = e^{-t}(C_1 \cos(2t) + C_2 \sin(2t))$.
\end{exemple}

% --- Figure TikZ : trois regimes ---
\begin{figure}[ht]
\centering
\begin{tikzpicture}[scale=0.9]
  % Surcritique : Delta > 0
  \begin{scope}[xshift=0cm]
    \draw[->] (-0.3,0) -- (4.2,0) node[right] {\small $t$};
    \draw[->] (0,-1.2) -- (0,2) node[above] {\small $y$};
    \draw[thick, blue, domain=0:4, samples=80] plot (\x, {1.5*exp(-0.5*\x) - 0.5*exp(-2*\x)});
    \node[below] at (2, -1.5) {\small Surcritique ($\Delta > 0$)};
  \end{scope}
  % Critique : Delta = 0
  \begin{scope}[xshift=5.5cm]
    \draw[->] (-0.3,0) -- (4.2,0) node[right] {\small $t$};
    \draw[->] (0,-1.2) -- (0,2) node[above] {\small $y$};
    \draw[thick, red!70!black, domain=0:4, samples=80] plot (\x, {(1 + 2*\x)*exp(-\x)});
    \node[below] at (2, -1.5) {\small Critique ($\Delta = 0$)};
  \end{scope}
  % Sous-critique : Delta < 0
  \begin{scope}[xshift=11cm]
    \draw[->] (-0.3,0) -- (4.2,0) node[right] {\small $t$};
    \draw[->] (0,-1.5) -- (0,2) node[above] {\small $y$};
    \draw[thick, green!50!black, domain=0:4, samples=100] plot (\x, {1.5*exp(-0.5*\x)*cos(2*\x r)});
    % Enveloppe
    \draw[dashed, gray, domain=0:4, samples=40] plot (\x, {1.5*exp(-0.5*\x)});
    \draw[dashed, gray, domain=0:4, samples=40] plot (\x, {-1.5*exp(-0.5*\x)});
    \node[below] at (2, -1.8) {\small Sous-critique ($\Delta < 0$)};
  \end{scope}
\end{tikzpicture}
\caption{Les trois regimes d'amortissement.}\label{fig:trois-regimes}
\end{figure}

%% -------------------------------------------------------------------
\subsection{Wronskien}\label{subsec:wronskien}
\index{wronskien}

\begin{definition}[Wronskien]\label{def:wronskien}
Soient $y_1, y_2$ deux solutions de~\eqref{eq:edl2-homogene}. Le \emph{wronskien} de $y_1$ et $y_2$ est :
\[
  W(t) = W[y_1, y_2](t) = \begin{vmatrix} y_1(t) & y_2(t) \\ y_1'(t) & y_2'(t) \end{vmatrix} = y_1(t)y_2'(t) - y_2(t)y_1'(t).
\]
\end{definition}

\begin{theoreme}[Formule d'Abel]\label{thm:abel}
\index{Abel!formule de}
Si $y_1, y_2$ sont solutions de $y'' + p\, y' + q\, y = 0$, alors :
\begin{equation}\label{eq:abel}
  W(t) = W(t_0)\, e^{-\int_{t_0}^{t} p\, \dd s} = W(t_0)\, e^{-p(t - t_0)}.
\end{equation}
En particulier, soit $W(t) = 0$ pour tout $t$, soit $W(t) \neq 0$ pour tout $t$.
\end{theoreme}

\begin{proof}
Calculons $W'(t)$ :
\begin{align*}
W'(t) &= y_1'(t)y_2'(t) + y_1(t)y_2''(t) - y_2'(t)y_1'(t) - y_2(t)y_1''(t) \\
  &= y_1(t)y_2''(t) - y_2(t)y_1''(t).
\end{align*}
Comme $y_i'' = -p\, y_i' - q\, y_i$ pour $i = 1, 2$ :
\begin{align*}
W'(t) &= y_1(-py_2' - qy_2) - y_2(-py_1' - qy_1) \\
  &= -p(y_1 y_2' - y_2 y_1') = -p\, W(t).
\end{align*}
Donc $W' = -pW$, d'ou $W(t) = W(t_0)e^{-p(t-t_0)}$.
\end{proof}

\begin{proposition}\label{prop:wronskien-indep}
Deux solutions $y_1, y_2$ de~\eqref{eq:edl2-homogene} forment un \emph{systeme fondamental de solutions}\index{systeme fondamental} si et seulement si $W(t_0) \neq 0$ pour un (et donc pour tout) $t_0$.
\end{proposition}

%% -------------------------------------------------------------------
\section{Problemes de Cauchy}\label{sec:problemes-cauchy-edl2}
\index{probleme de Cauchy!lineaire du second ordre}

\begin{theoreme}[Existence et unicite pour l'equation lineaire]\label{thm:existence-unicite-edl2}
Pour tous $t_0 \in \R$, $y_0, y_1 \in \R$, le probleme de Cauchy
\[
  y'' + p\, y' + q\, y = g(t), \quad y(t_0) = y_0, \quad y'(t_0) = y_1,
\]
ou $g$ est continue, admet une unique solution definie sur $\R$ tout entier.
\end{theoreme}

\begin{methode}[Determination des constantes]\label{meth:constantes}
\index{constantes!determination des}
Pour resoudre un probleme de Cauchy :
\begin{enumerate}
  \item Ecrire la solution generale $y(t) = C_1 y_1(t) + C_2 y_2(t)$.
  \item Imposer $y(t_0) = y_0$ : on obtient une premiere equation en $C_1, C_2$.
  \item Imposer $y'(t_0) = y_1$ : on obtient une seconde equation.
  \item Resoudre le systeme $2 \times 2$ (le determinant est $W(t_0) \neq 0$).
\end{enumerate}
\end{methode}

\begin{exemple}\label{ex:cauchy-edl2}
Resoudre $y'' - 3y' + 2y = 0$, $y(0) = 1$, $y'(0) = 0$.

Equation caracteristique : $r^2 - 3r + 2 = (r-1)(r-2) = 0$, donc $r_1 = 1$, $r_2 = 2$.

Solution generale : $y(t) = C_1 e^t + C_2 e^{2t}$.
\begin{align*}
  y(0) = 1 &\implies C_1 + C_2 = 1, \\
  y'(0) = 0 &\implies C_1 + 2C_2 = 0.
\end{align*}
D'ou $C_2 = -1$ et $C_1 = 2$. Solution : $y(t) = 2e^t - e^{2t}$.
\end{exemple}

\begin{exemple}\label{ex:cauchy-oscillation}
Resoudre $y'' + 4y = 0$, $y(0) = 3$, $y'(0) = -2$.

Equation caracteristique : $r^2 + 4 = 0$, donc $r = \pm 2i$ ($\alpha = 0$, $\beta = 2$).

Solution generale : $y(t) = C_1 \cos(2t) + C_2 \sin(2t)$.
\begin{align*}
  y(0) = 3 &\implies C_1 = 3, \\
  y'(0) = -2 &\implies 2C_2 = -2 \implies C_2 = -1.
\end{align*}
Solution : $y(t) = 3\cos(2t) - \sin(2t)$.
\end{exemple}

%% -------------------------------------------------------------------
\section{Equation non homogene : methode des coefficients indetermines}\label{sec:coefficients-indetermines}
\index{coefficients indetermines@coefficients indetermines!methode des}

\begin{theoreme}[Structure de la solution generale]\label{thm:structure-sol}
\index{solution!generale}
La solution generale de l'equation non homogene $ay'' + by' + cy = f(t)$ est :
\[
  y(t) = y_h(t) + y_p(t),
\]
ou $y_h$ est la solution generale de l'equation homogene associee et $y_p$ est une solution particuliere quelconque de l'equation non homogene.
\end{theoreme}

\begin{theoreme}[Principe de superposition]\label{thm:superposition}
\index{superposition!principe de}
Si $y_{p_1}$ est une solution particuliere de $ay'' + by' + cy = f_1(t)$ et $y_{p_2}$ est une solution particuliere de $ay'' + by' + cy = f_2(t)$, alors $\lambda_1 y_{p_1} + \lambda_2 y_{p_2}$ est une solution particuliere de $ay'' + by' + cy = \lambda_1 f_1(t) + \lambda_2 f_2(t)$.
\end{theoreme}

\subsection{Table des solutions particulieres}\label{subsec:table-sp}
\index{solution particuliere!table des}

Pour l'equation $ay'' + by' + cy = f(t)$, voici les formes a essayer :

\begin{center}
\renewcommand{\arraystretch}{1.6}
\begin{tabular}{@{}p{5.5cm}p{7.5cm}@{}}
\toprule
\textbf{Second membre $f(t)$} & \textbf{Forme de $y_p(t)$ a essayer} \\
\midrule
$P_n(t)$ (polynome degre $n$) & $t^s Q_n(t)$ \\
$e^{\gamma t}$ & $t^s A\, e^{\gamma t}$ \\
$P_n(t)\, e^{\gamma t}$ & $t^s Q_n(t)\, e^{\gamma t}$ \\
$e^{\alpha t}\cos(\beta t)$ ou $e^{\alpha t}\sin(\beta t)$ & $t^s e^{\alpha t}(A\cos(\beta t) + B\sin(\beta t))$ \\
$P_n(t)\, e^{\alpha t}\cos(\beta t)$ & $t^s e^{\alpha t}(Q_n(t)\cos(\beta t) + R_n(t)\sin(\beta t))$ \\
\bottomrule
\end{tabular}
\end{center}

\begin{remarque}[Valeur de $s$]\label{rem:valeur-s}
\index{multiplicite de la racine}
L'entier $s$ est la \emph{multiplicite} de $\gamma$ (ou $\alpha + i\beta$) comme racine de l'equation caracteristique :
\begin{itemize}
  \item $s = 0$ si $\gamma$ (ou $\alpha + i\beta$) n'est \emph{pas} racine de $ar^2 + br + c = 0$,
  \item $s = 1$ si c'est une racine simple,
  \item $s = 2$ si c'est une racine double.
\end{itemize}
\end{remarque}

\begin{exemple}[Second membre polynomiale]\label{ex:sp-polynome}
Resoudre $y'' - 3y' + 2y = 4t^2$.

\emph{Solution homogene} : $r^2 - 3r + 2 = 0$ donne $r_1 = 1$, $r_2 = 2$, donc $y_h = C_1 e^t + C_2 e^{2t}$.

\emph{Solution particuliere} : comme $\gamma = 0$ n'est pas racine de l'equation caracteristique, $s = 0$. On cherche $y_p = At^2 + Bt + C$. Alors $y_p' = 2At + B$, $y_p'' = 2A$.

En substituant :
\[
  2A - 3(2At + B) + 2(At^2 + Bt + C) = 4t^2.
\]
Identification : $t^2 : 2A = 4$, $t^1 : -6A + 2B = 0$, $t^0 : 2A - 3B + 2C = 0$.

D'ou $A = 2$, $B = 6$, $C = 7$. Solution particuliere : $y_p = 2t^2 + 6t + 7$.

Solution generale : $y(t) = C_1 e^t + C_2 e^{2t} + 2t^2 + 6t + 7$.
\end{exemple}

\begin{exemple}[Second membre exponentiel]\label{ex:sp-exp}
Resoudre $y'' - 3y' + 2y = e^{3t}$.

Equation caracteristique : $r_1 = 1$, $r_2 = 2$. Comme $\gamma = 3$ n'est pas racine, $s = 0$. On cherche $y_p = Ae^{3t}$.

$y_p'' - 3y_p' + 2y_p = 9Ae^{3t} - 9Ae^{3t} + 2Ae^{3t} = 2Ae^{3t} = e^{3t}$.

Donc $A = 1/2$ et $y_p = \frac{1}{2}e^{3t}$.
\end{exemple}

\begin{exemple}[Cas de resonance partielle ($s = 1$)]\label{ex:sp-resonance}
\index{resonance!partielle}
Resoudre $y'' - 3y' + 2y = e^{2t}$.

Ici $\gamma = 2$ est racine \emph{simple} de l'equation caracteristique, donc $s = 1$. On cherche $y_p = At\, e^{2t}$.

$y_p' = A(1 + 2t)e^{2t}$, $y_p'' = A(4 + 4t)e^{2t}$.

En substituant :
\[
  A(4 + 4t)e^{2t} - 3A(1 + 2t)e^{2t} + 2At\, e^{2t} = Ae^{2t}(4 + 4t - 3 - 6t + 2t) = Ae^{2t}.
\]
Donc $A = 1$ et $y_p = te^{2t}$.
\end{exemple}

\begin{exemple}[Cas de resonance double ($s = 2$)]\label{ex:sp-resonance-double}
Resoudre $y'' + 4y' + 4y = e^{-2t}$.

Equation caracteristique : $(r+2)^2 = 0$, racine double $r_0 = -2$. Comme $\gamma = -2$ est racine de multiplicite $2$, $s = 2$. On cherche $y_p = At^2 e^{-2t}$.

$y_p' = A(2t - 2t^2)e^{-2t}$, $y_p'' = A(2 - 8t + 4t^2)e^{-2t}$.

En substituant :
\[
  A(2 - 8t + 4t^2 + 8t - 8t^2 + 4t^2)e^{-2t} = 2Ae^{-2t} = e^{-2t}.
\]
Donc $A = 1/2$ et $y_p = \frac{t^2}{2}e^{-2t}$.
\end{exemple}

\begin{exemple}[Second membre trigonometrique]\label{ex:sp-trigo}
Resoudre $y'' + y = \cos(2t)$.

Equation caracteristique : $r^2 + 1 = 0$, racines $r = \pm i$. Comme $\alpha + i\beta = 2i$ n'est pas racine ($\beta = 2 \neq 1$), $s = 0$. On cherche $y_p = A\cos(2t) + B\sin(2t)$.

$y_p'' = -4A\cos(2t) - 4B\sin(2t)$.

$y_p'' + y_p = -3A\cos(2t) - 3B\sin(2t) = \cos(2t)$.

Donc $A = -1/3$, $B = 0$. Solution particuliere : $y_p = -\frac{1}{3}\cos(2t)$.

Solution generale : $y(t) = C_1\cos(t) + C_2\sin(t) - \frac{1}{3}\cos(2t)$.
\end{exemple}

%% -------------------------------------------------------------------
\section{Phenomene de resonance}\label{sec:resonance}
\index{resonance}

\begin{definition}[Resonance]\label{def:resonance}
Il y a \emph{resonance} lorsque la frequence de l'excitation coincide avec une frequence propre du systeme. Mathematiquement, cela correspond au cas ou $\alpha + i\beta$ est racine de l'equation caracteristique.
\end{definition}

\begin{exemple}[Resonance pure]\label{ex:resonance-pure}
\index{resonance!pure}
Considerons $y'' + \omega_0^2 y = \cos(\omega_0 t)$ (pas d'amortissement, excitation a la frequence propre $\omega_0$).

Comme $i\omega_0$ est racine simple, $s = 1$. On cherche $y_p = t(A\cos(\omega_0 t) + B\sin(\omega_0 t))$.

Apres calcul : $y_p = \frac{t}{2\omega_0}\sin(\omega_0 t)$.

L'amplitude croit \emph{lineairement} avec le temps. C'est la \textbf{resonance pure}\index{resonance!pure}.
\end{exemple}

% --- Figure TikZ : resonance ---
\begin{figure}[ht]
\centering
\begin{tikzpicture}[scale=0.8]
  \draw[->] (-0.3,0) -- (10,0) node[right] {$t$};
  \draw[->] (0,-3.5) -- (0,3.5) node[above] {$y$};
  % Enveloppe
  \draw[dashed, red, domain=0:9.5, samples=60] plot (\x, {0.35*\x});
  \draw[dashed, red, domain=0:9.5, samples=60] plot (\x, {-0.35*\x});
  % Signal
  \draw[thick, blue, domain=0:9.5, samples=200] plot (\x, {0.35*\x*sin(2*\x r)});
  \node[above right] at (5, 2.5) {$y_p = \frac{t}{2\omega_0}\sin(\omega_0 t)$};
  \node at (5, -4.2) {Resonance pure : amplitude croissante};
\end{tikzpicture}
\caption{Phenomene de resonance pure (sans amortissement).}\label{fig:resonance}
\end{figure}

\begin{exemple}[Resonance amortie]\label{ex:resonance-amortie}
\index{resonance!amortie}
Pour $y'' + 2\gamma y' + \omega_0^2 y = F_0 \cos(\omega t)$ avec $\gamma > 0$, l'amplitude de la reponse stationnaire est :
\[
  A(\omega) = \frac{F_0}{\sqrt{(\omega_0^2 - \omega^2)^2 + 4\gamma^2 \omega^2}}.
\]
Le maximum est atteint pour $\omega_{\mathrm{res}} = \sqrt{\omega_0^2 - 2\gamma^2}$ (si $\omega_0^2 > 2\gamma^2$).
\end{exemple}

%% -------------------------------------------------------------------
\section{Variation des constantes}\label{sec:variation-constantes}
\index{variation des constantes}

\begin{methode}[Variation des constantes (Lagrange)]\label{meth:variation-constantes}
\index{Lagrange!variation des constantes}
Pour trouver une solution particuliere de $y'' + py' + qy = g(t)$ quand la methode des coefficients indetermines ne s'applique pas, on pose :
\[
  y_p(t) = c_1(t)\, y_1(t) + c_2(t)\, y_2(t),
\]
ou $\{y_1, y_2\}$ est un systeme fondamental, et on impose :
\begin{equation}\label{eq:systeme-vc}
\begin{cases}
c_1'(t)\, y_1(t) + c_2'(t)\, y_2(t) = 0, \\
c_1'(t)\, y_1'(t) + c_2'(t)\, y_2'(t) = g(t).
\end{cases}
\end{equation}
Ce systeme admet une unique solution car le determinant est le wronskien $W(t) \neq 0$.
\end{methode}

\begin{proposition}\label{prop:formules-vc}
Les solutions du systeme~\eqref{eq:systeme-vc} sont :
\begin{equation}\label{eq:vc-formules}
  c_1'(t) = -\frac{y_2(t)\, g(t)}{W(t)}, \quad c_2'(t) = \frac{y_1(t)\, g(t)}{W(t)}.
\end{equation}
La solution particuliere est donc :
\[
  y_p(t) = -y_1(t)\int \frac{y_2(s)\, g(s)}{W(s)}\, \dd s + y_2(t)\int \frac{y_1(s)\, g(s)}{W(s)}\, \dd s.
\]
\end{proposition}

\begin{exemple}\label{ex:vc}
Resoudre $y'' + y = \frac{1}{\cos(t)}$.

\emph{Solution homogene} : $y_1 = \cos(t)$, $y_2 = \sin(t)$, $W = 1$.

Par variation des constantes :
\begin{align*}
c_1'(t) &= -\sin(t) \cdot \frac{1}{\cos(t)} = -\tan(t), &\quad c_1(t) &= \ln\abs{\cos(t)}, \\
c_2'(t) &= \cos(t) \cdot \frac{1}{\cos(t)} = 1, &\quad c_2(t) &= t.
\end{align*}
Solution particuliere : $y_p(t) = \cos(t)\ln\abs{\cos(t)} + t\sin(t)$.

Solution generale : $y(t) = C_1\cos(t) + C_2\sin(t) + \cos(t)\ln\abs{\cos(t)} + t\sin(t)$.
\end{exemple}

%% -------------------------------------------------------------------
\section{Equation d'Euler--Cauchy}\label{sec:euler-cauchy}
\index{Euler--Cauchy!equation d'}

\begin{definition}[Equation d'Euler--Cauchy]\label{def:euler-cauchy}
L'\emph{equation d'Euler--Cauchy} est une equation de la forme :
\begin{equation}\label{eq:euler-cauchy}
  t^2 y'' + \alpha\, t\, y' + \beta\, y = 0, \quad t > 0,
\end{equation}
ou $\alpha, \beta \in \R$.
\end{definition}

\begin{methode}[Resolution par le changement $t = e^u$]\label{meth:euler-cauchy}
\index{changement de variable!$t = e^u$}
Le changement de variable $t = e^u$ (soit $u = \ln t$) transforme l'equation d'Euler--Cauchy en une equation a coefficients constants. En posant $z(u) = y(e^u)$ :
\[
  ty' = \dot{z}, \quad t^2 y'' = \ddot{z} - \dot{z},
\]
ou le point designe $\frac{\dd}{\dd u}$. L'equation~\eqref{eq:euler-cauchy} devient :
\begin{equation}\label{eq:euler-cauchy-transformee}
  \ddot{z} + (\alpha - 1)\dot{z} + \beta z = 0.
\end{equation}
\end{methode}

\begin{remarque}\label{rem:euler-puissance}
Alternativement, on cherche directement des solutions de la forme $y(t) = t^r$. En substituant dans~\eqref{eq:euler-cauchy} :
\[
  t^r\bigl[r(r-1) + \alpha r + \beta\bigr] = 0,
\]
ce qui donne l'\emph{equation indicielle}\index{equation indicielle} :
\begin{equation}\label{eq:indicielle}
  r(r-1) + \alpha r + \beta = r^2 + (\alpha - 1)r + \beta = 0.
\end{equation}
\end{remarque}

Les trois cas du discriminant $\Delta' = (\alpha - 1)^2 - 4\beta$ donnent :

\begin{center}
\renewcommand{\arraystretch}{1.5}
\begin{tabular}{@{}lll@{}}
\toprule
\textbf{Cas} & \textbf{Racines} & \textbf{Solution generale ($t > 0$)} \\
\midrule
$\Delta' > 0$ & $r_1, r_2 \in \R$, $r_1 \neq r_2$ & $y(t) = C_1 t^{r_1} + C_2 t^{r_2}$ \\
$\Delta' = 0$ & $r_0$ double & $y(t) = (C_1 + C_2 \ln t)\, t^{r_0}$ \\
$\Delta' < 0$ & $r = \mu \pm i\nu$ & $y(t) = t^\mu\bigl(C_1 \cos(\nu\ln t) + C_2\sin(\nu\ln t)\bigr)$ \\
\bottomrule
\end{tabular}
\end{center}

\begin{exemple}\label{ex:euler-cauchy-1}
Resoudre $t^2 y'' + 3ty' + y = 0$ pour $t > 0$.

Equation indicielle : $r^2 + 2r + 1 = (r+1)^2 = 0$, racine double $r_0 = -1$.

Solution generale : $y(t) = (C_1 + C_2 \ln t)\, t^{-1} = \frac{C_1 + C_2 \ln t}{t}$.
\end{exemple}

\begin{exemple}\label{ex:euler-cauchy-2}
Resoudre $t^2 y'' + ty' + y = 0$ pour $t > 0$.

Equation indicielle : $r^2 + 1 = 0$, racines $r = \pm i$ ($\mu = 0$, $\nu = 1$).

Solution generale : $y(t) = C_1 \cos(\ln t) + C_2 \sin(\ln t)$.
\end{exemple}

%% -------------------------------------------------------------------
\section{Reduction d'ordre}\label{sec:reduction-ordre}
\index{reduction d'ordre}

\begin{methode}[Reduction d'ordre]\label{meth:reduction-ordre}
Si l'on connait \emph{une} solution $y_1 \neq 0$ de l'equation homogene $y'' + py' + qy = 0$, on peut trouver la solution generale en posant $y(t) = v(t)\, y_1(t)$. En substituant :
\[
  v'' y_1 + (2y_1' + py_1)v' = 0.
\]
C'est une equation du premier ordre en $v' = w$ :
\[
  w' + \left(\frac{2y_1'}{y_1} + p\right) w = 0,
\]
d'ou $w(t) = \frac{K}{y_1(t)^2} e^{-\int p\, \dd t}$ et $v(t) = \int w(s)\, \dd s$.
\end{methode}

\begin{exemple}\label{ex:reduction-ordre}
On sait que $y_1(t) = e^t$ est solution de $y'' - 2y' + y = 0$. Trouvons la solution generale.

Posons $y = ve^t$. Alors $y' = (v' + v)e^t$, $y'' = (v'' + 2v' + v)e^t$.

$(v'' + 2v' + v - 2v' - 2v + v)e^t = v''e^t = 0$.

Donc $v'' = 0$, $v(t) = C_1 + C_2 t$, et $y(t) = (C_1 + C_2 t)e^t$.
\end{exemple}

%% -------------------------------------------------------------------
\section{Exemples recapitulatifs completement resolus}\label{sec:exemples-recapitulatifs}

\begin{exemple}[Probleme complet avec conditions initiales]\label{ex:complet-1}
Resoudre $y'' + 2y' + 5y = 10\cos(t)$, $y(0) = 0$, $y'(0) = 0$.

\emph{Equation caracteristique} : $r^2 + 2r + 5 = 0$, $\Delta = 4 - 20 = -16$, $r = -1 \pm 2i$.

\emph{Solution homogene} : $y_h(t) = e^{-t}(C_1\cos(2t) + C_2\sin(2t))$.

\emph{Solution particuliere} : $\alpha + i\beta = i$ n'est pas racine. On cherche $y_p = A\cos(t) + B\sin(t)$.

$y_p'' + 2y_p' + 5y_p = (-A + 2B + 5A)\cos(t) + (-B - 2A + 5B)\sin(t)$
$= (4A + 2B)\cos(t) + (-2A + 4B)\sin(t) = 10\cos(t)$.

Systeme : $4A + 2B = 10$ et $-2A + 4B = 0$. D'ou $A = 2$, $B = 1$.

$y_p(t) = 2\cos(t) + \sin(t)$.

\emph{Conditions initiales} :
\begin{align*}
  y(0) = C_1 + 2 = 0 &\implies C_1 = -2, \\
  y'(0) = -C_1 + 2C_2 + 1 = 0 &\implies 2 + 2C_2 + 1 = 0 \implies C_2 = -\frac{3}{2}.
\end{align*}

\emph{Solution} : $y(t) = e^{-t}\bigl(-2\cos(2t) - \tfrac{3}{2}\sin(2t)\bigr) + 2\cos(t) + \sin(t)$.

Le terme transitoire $e^{-t}(\ldots)$ tend vers $0$ ; la reponse stationnaire est $2\cos(t) + \sin(t)$.
\end{exemple}

\begin{exemple}[Equation d'Euler--Cauchy non homogene]\label{ex:euler-cauchy-nh}
Resoudre $t^2 y'' - 2y = t^3$ pour $t > 0$.

\emph{Equation homogene} : $t^2y'' - 2y = 0$. Equation indicielle : $r(r-1) - 2 = r^2 - r - 2 = (r-2)(r+1) = 0$.

$y_h(t) = C_1 t^2 + C_2 t^{-1}$.

\emph{Solution particuliere} par variation des constantes. Avec $y_1 = t^2$, $y_2 = t^{-1}$ :
\[
  W = t^2 \cdot (-t^{-2}) - t^{-1} \cdot 2t = -1 - 2 = -3.
\]
On ecrit l'equation sous forme standard : $y'' - 2t^{-2}y = t$. Donc $g(t) = t$.
\begin{align*}
  c_1'(t) &= -\frac{t^{-1} \cdot t}{-3} = \frac{1}{3}, &\quad c_1(t) &= \frac{t}{3}, \\
  c_2'(t) &= \frac{t^2 \cdot t}{-3} = -\frac{t^3}{3}, &\quad c_2(t) &= -\frac{t^4}{12}.
\end{align*}
$y_p(t) = \frac{t}{3} \cdot t^2 + \left(-\frac{t^4}{12}\right) t^{-1} = \frac{t^3}{3} - \frac{t^3}{12} = \frac{t^3}{4}$.

\emph{Solution generale} : $y(t) = C_1 t^2 + C_2 t^{-1} + \frac{t^3}{4}$.
\end{exemple}

%% -------------------------------------------------------------------
\section{Exercices}\label{sec:exos-ch4}

\begin{exercice}[$\star$]\label{exo:edl2-homogene}
Resoudre les equations homogenes suivantes :
\begin{enumerate}[label=(\alph*)]
  \item $y'' - 4y' + 3y = 0$.
  \item $y'' + 6y' + 9y = 0$.
  \item $y'' + 2y' + 10y = 0$.
  \item $4y'' + 4y' + y = 0$.
\end{enumerate}
\end{exercice}

\begin{exercice}[$\star$]\label{exo:cauchy-lineaire}
Resoudre les problemes de Cauchy :
\begin{enumerate}[label=(\alph*)]
  \item $y'' - y = 0$, $y(0) = 2$, $y'(0) = -1$.
  \item $y'' + 9y = 0$, $y(0) = 0$, $y'(0) = 3$.
  \item $y'' + 4y' + 4y = 0$, $y(0) = 1$, $y'(0) = 1$.
\end{enumerate}
\end{exercice}

\begin{exercice}[$\star\star$]\label{exo:coefficients-indetermines}
Trouver la solution generale des equations :
\begin{enumerate}[label=(\alph*)]
  \item $y'' - 4y = e^{2t}$.
  \item $y'' + y = 3\sin(t)$.
  \item $y'' - 2y' + y = e^t$.
  \item $y'' + 4y = \cos(2t)$.
  \item $y'' + y' - 2y = t^2 + t$.
\end{enumerate}
\end{exercice}

\begin{exercice}[$\star\star$]\label{exo:variation-constantes}
Resoudre par variation des constantes :
\begin{enumerate}[label=(\alph*)]
  \item $y'' + y = \tan(t)$.
  \item $y'' - y = \frac{1}{\cosh(t)}$.
  \item $y'' + 4y = \frac{1}{\sin(2t)}$.
\end{enumerate}
\end{exercice}

\begin{exercice}[$\star\star$]\label{exo:euler-cauchy}
Resoudre pour $t > 0$ :
\begin{enumerate}[label=(\alph*)]
  \item $t^2 y'' - 2ty' + 2y = 0$.
  \item $t^2 y'' + 5ty' + 4y = 0$.
  \item $t^2 y'' + ty' + 4y = 0$.
\end{enumerate}
\end{exercice}

\begin{exercice}[$\star\star\star$]\label{exo:wronskien-calcul}
Soit $y'' + p(t)y' + q(t)y = 0$ ou $p$ et $q$ sont continues sur $\R$.
\begin{enumerate}[label=(\alph*)]
  \item Montrer que si $y_1$ et $y_2$ sont deux solutions, alors $W[y_1, y_2]$ verifie $W' = -p(t)W$.
  \item En deduire que si $W(t_0) = 0$ pour un $t_0$, alors $y_1$ et $y_2$ sont lineairement dependantes.
  \item Application : montrer que $\sin(t)$ et $\cos(t)$ ne peuvent pas etre simultanément solutions d'une equation $y'' + p(t)y' + q(t)y = 0$ avec $p$ non identiquement nul.
\end{enumerate}
\end{exercice}

\begin{exercice}[$\star\star\star$]\label{exo:resonance-amortie}
On considere l'oscillateur amorti force $y'' + 2\gamma y' + \omega_0^2 y = F_0\cos(\omega t)$ avec $\gamma > 0$.
\begin{enumerate}[label=(\alph*)]
  \item Trouver la solution particuliere stationnaire sous la forme $y_p = A\cos(\omega t - \phi)$.
  \item Exprimer $A$ et $\phi$ en fonction de $\omega$, $\omega_0$, $\gamma$, $F_0$.
  \item Montrer que $A(\omega)$ est maximale pour $\omega_{\mathrm{res}} = \sqrt{\omega_0^2 - 2\gamma^2}$ (si cette quantite est positive).
  \item Que se passe-t-il quand $\gamma \to 0^+$ ?
\end{enumerate}
\end{exercice}

\begin{exercice}[$\star\star\star$]\label{exo:reduction-ordre-general}
On sait que $y_1(t) = t$ est solution de $t^2 y'' - t(t+2)y' + (t+2)y = 0$ pour $t > 0$.
\begin{enumerate}[label=(\alph*)]
  \item Utiliser la reduction d'ordre pour trouver une seconde solution lineairement independante.
  \item Ecrire la solution generale.
\end{enumerate}
\end{exercice}

%% -------------------------------------------------------------------
\section*{Resume du chapitre}

\begin{center}
\renewcommand{\arraystretch}{1.5}
\begin{tabular}{@{}p{4.5cm}p{9cm}@{}}
\toprule
\textbf{Resultat/Methode} & \textbf{Enonce succinct} \\
\midrule
Equation caracteristique & $ar^2 + br + c = 0$ determine la nature des solutions \\
$\Delta > 0$ & $y_h = C_1 e^{r_1 t} + C_2 e^{r_2 t}$ \\
$\Delta = 0$ & $y_h = (C_1 + C_2 t)e^{r_0 t}$ \\
$\Delta < 0$ & $y_h = e^{\alpha t}(C_1\cos\beta t + C_2\sin\beta t)$ \\
Wronskien (Abel) & $W(t) = W(t_0)e^{-p(t-t_0)}$ \\
Coefficients indetermines & Forme de $y_p$ selon $f(t)$ et multiplicite $s$ \\
Variation des constantes & Methode universelle : systeme~\eqref{eq:systeme-vc} \\
Equation d'Euler--Cauchy & $t = e^u$ ou $y = t^r$ pour se ramener a coefficients constants \\
Resonance & $s \geq 1$ : amplification ; sans amortissement : $A \to \infty$ \\
\bottomrule
\end{tabular}
\end{center}

% === FIN DU CHUNK ch3-4 ===
% === DÉBUT DU CHUNK ch5-6 ===
% chapters5_6.tex — Chapitres 5 et 6 : Systèmes différentiels linéaires,
% Méthodes de résolution explicite (Wronskien, variation des constantes)

% ████████████████████████████████████████████████████████████████████████
\chapter{Systèmes différentiels linéaires}\label{ch:systemes-lineaires}
% ████████████████████████████████████████████████████████████████████████

\minitoc
\vspace{1cm}

\section*{Introduction}
\addcontentsline{toc}{section}{Introduction}

Jusqu'à présent, nous avons principalement étudié des équations
différentielles scalaires. De nombreux phénomènes physiques, biologiques
ou économiques font cependant intervenir \emph{plusieurs grandeurs couplées},
dont l'évolution est décrite simultanément par un \emph{système}
d'équations différentielles.\index{système différentiel}

\medskip

\noindent\textbf{Oscillateurs couplés.}\index{oscillateurs couplés}
Considérons deux masses $m_1$ et $m_2$ reliées entre elles et à deux
murs par des ressorts de raideurs $k_1$, $k_2$, $k_3$. En notant
$x_1(t)$ et $x_2(t)$ les déplacements par rapport aux positions
d'équilibre, le principe fondamental de la dynamique donne~:
\[
\begin{cases}
m_1 x_1'' = -k_1 x_1 + k_2(x_2 - x_1), \\
m_2 x_2'' = -k_3 x_2 - k_2(x_2 - x_1).
\end{cases}
\]
En posant $X = (x_1, x_1', x_2, x_2')^\top$, ce système d'ordre~2 se
ramène à un système d'ordre~1 de dimension~4~: $X' = AX$.

\medskip

\noindent\textbf{Modèle proie--prédateur linéarisé.}\index{Lotka--Volterra}
Le modèle de Lotka--Volterra linéarisé autour d'un point d'équilibre
s'écrit
\[
\begin{pmatrix} x' \\ y' \end{pmatrix}
= \begin{pmatrix} 0 & -\alpha \\ \beta & 0 \end{pmatrix}
  \begin{pmatrix} x \\ y \end{pmatrix},
\]
où $x(t)$ représente la variation de la population de proies et $y(t)$
celle de prédateurs. L'analyse des valeurs propres de la matrice
$A = \bigl(\begin{smallmatrix} 0 & -\alpha \\ \beta & 0 \end{smallmatrix}\bigr)$
détermine la nature du point d'équilibre.

\medskip

Ce chapitre développe la théorie des systèmes différentiels linéaires,
en se concentrant sur les systèmes à coefficients constants, et propose
une classification complète des portraits de phase en dimension~2.

% ============================================================
\section{Forme générale et premières propriétés}\label{sec:syst-forme-generale}
% ============================================================

\begin{definition}[Système différentiel linéaire]\label{def:syst-diff-lin}
  \index{système différentiel!linéaire}
  Soient $n \geq 1$, $I$ un intervalle ouvert de~$\R$,
  $A \colon I \to \mathcal{M}_n(\R)$ et
  $B \colon I \to \R^n$ deux fonctions continues.
  On appelle \emph{système différentiel linéaire} l'équation
  \begin{equation}\label{eq:syst-lin-general}
    X'(t) = A(t)\,X(t) + B(t),
  \end{equation}
  où $X \colon I \to \R^n$ est l'inconnue. Lorsque $B \equiv 0$,
  le système est dit \emph{homogène}\index{système homogène}.
\end{definition}

\begin{theoreme}[Cauchy--Lipschitz pour les systèmes linéaires]\label{thm:CL-systemes}
  \index{Cauchy--Lipschitz!systèmes linéaires}
  Soit le système~\eqref{eq:syst-lin-general} avec $A$ et $B$
  continues sur~$I$. Pour tout $(t_0, X_0) \in I \times \R^n$,
  le problème de Cauchy
  \[
    X'(t) = A(t)\,X(t) + B(t), \qquad X(t_0) = X_0,
  \]
  admet une \emph{unique} solution définie sur~$I$ tout entier.
\end{theoreme}

\begin{proof}
  La fonction $F(t,X) = A(t)X + B(t)$ est continue en $(t,X)$ et
  globalement lipschitzienne en~$X$ sur tout compact $[a,b] \subset I$~:
  \[
    \norm{F(t,X) - F(t,Y)} = \norm{A(t)(X-Y)}
      \leq \norm{A(t)}\,\norm{X-Y}
      \leq M\,\norm{X-Y},
  \]
  où $M = \sup_{t \in [a,b]} \norm{A(t)} < +\infty$ par continuité
  sur un compact. Le théorème de Cauchy--Lipschitz global
  s'applique et donne l'existence et l'unicité sur~$I$ tout entier.
\end{proof}

\begin{corollaire}[Structure de l'espace des solutions]\label{cor:structure-solutions-syst}
  \index{espace des solutions!système linéaire}
  L'ensemble $\mathcal{S}_H$ des solutions du système homogène
  $X' = A(t)X$ est un sous-espace vectoriel de $\mathcal{C}^1(I,\R^n)$
  de dimension~$n$. L'ensemble des solutions du système complet
  \eqref{eq:syst-lin-general} est un sous-espace affine de direction
  $\mathcal{S}_H$.
\end{corollaire}

% ============================================================
\section{Exponentielle de matrice}\label{sec:exp-matrice}
% ============================================================

Pour résoudre le système homogène $X' = AX$ à coefficients constants,
l'outil central est l'\emph{exponentielle de matrice}.

\begin{definition}[Exponentielle de matrice]\label{def:exp-matrice}
  \index{exponentielle de matrice}
  Pour toute matrice $A \in \mathcal{M}_n(\R)$ (ou $\mathcal{M}_n(\C)$),
  on définit
  \begin{equation}\label{eq:exp-matrice}
    e^A = \exp(A) = \sum_{k=0}^{+\infty} \frac{A^k}{k!}.
  \end{equation}
  Cette série converge absolument pour toute norme matricielle
  sous-multiplicative.
\end{definition}

\begin{proposition}[Propriétés de l'exponentielle]\label{prop:proprietes-exp}
  \index{exponentielle de matrice!propriétés}
  Soient $A, B \in \mathcal{M}_n(\R)$ et $t, s \in \R$.
  \begin{enumerate}[label=(\roman*)]
    \item $e^{0_n} = I_n$.
    \item Si $AB = BA$, alors $e^{A+B} = e^A e^B$.
    \item $e^A$ est inversible et $(e^A)^{-1} = e^{-A}$.
    \item $\dfrac{\dd}{\dd t} e^{tA} = A\,e^{tA} = e^{tA}\,A$.
    \item Si $P$ est inversible, alors $e^{P^{-1}AP} = P^{-1} e^A P$.
    \item $\det(e^A) = e^{\tr(A)}$.\index{trace!et exponentielle}
  \end{enumerate}
\end{proposition}

\begin{proof}
  \emph{(iv)} On dérive terme à terme la série convergente~:
  \[
    \frac{\dd}{\dd t} \sum_{k=0}^{\infty} \frac{(tA)^k}{k!}
    = \sum_{k=1}^{\infty} \frac{k\,t^{k-1}A^k}{k!}
    = A \sum_{k=1}^{\infty} \frac{t^{k-1}A^{k-1}}{(k-1)!}
    = A\,e^{tA}.
  \]
  La convergence uniforme sur tout compact de $\R$ justifie la dérivation.

  \emph{(vi)} On utilise la formule de Jacobi
  $\det(e^A) = \lim_{m\to\infty}\det\bigl((I+A/m)^m\bigr)
  = \lim_{m\to\infty}(1+\tr(A)/m + O(1/m^2))^m = e^{\tr(A)}$.
  Une preuve plus rigoureuse passe par la forme de Jordan.
\end{proof}

\begin{theoreme}[Solution du système homogène à coefficients constants]\label{thm:sol-syst-homogene}
  \index{système homogène!solution}
  La solution du problème de Cauchy $X' = AX$, $X(t_0) = X_0$
  ($A \in \mathcal{M}_n(\R)$ constante) est
  \begin{equation}\label{eq:sol-syst-homogene}
    X(t) = e^{(t-t_0)A}\,X_0.
  \end{equation}
  En particulier, si $t_0 = 0$, alors $X(t) = e^{tA}\,X_0$.
\end{theoreme}

\begin{proof}
  Posons $\varphi(t) = e^{(t-t_0)A}X_0$. Alors
  $\varphi(t_0) = e^0 X_0 = X_0$ et
  $\varphi'(t) = A\,e^{(t-t_0)A}X_0 = A\,\varphi(t)$
  d'après la proposition~\ref{prop:proprietes-exp}~(iv).
  L'unicité de Cauchy--Lipschitz conclut.
\end{proof}

% ============================================================
\section{Systèmes homogènes à coefficients constants~: méthode des
  valeurs propres}\label{sec:methode-vp}
% ============================================================

\subsection{Cas diagonalisable}\label{subsec:cas-diag}

\begin{proposition}[Exponentielle d'une matrice diagonalisable]\label{prop:exp-diag}
  \index{exponentielle de matrice!cas diagonalisable}
  Si $A = P\,D\,P^{-1}$ avec $D = \diag(\lambda_1,\ldots,\lambda_n)$,
  alors
  \[
    e^{tA} = P\,\diag(e^{\lambda_1 t},\ldots,e^{\lambda_n t})\,P^{-1}.
  \]
\end{proposition}

\begin{methode}[Résolution par diagonalisation]\label{meth:diag}
  \index{méthode!diagonalisation}
  Pour résoudre $X' = AX$~:
  \begin{enumerate}
    \item Calculer les valeurs propres $\lambda_1,\ldots,\lambda_n$ de~$A$.
    \item Trouver une base de vecteurs propres $(V_1,\ldots,V_n)$.
    \item La solution générale est
      $X(t) = c_1 e^{\lambda_1 t} V_1 + \cdots + c_n e^{\lambda_n t} V_n$,
      $c_i \in \R$.
  \end{enumerate}
\end{methode}

\begin{exemple}[Système $2\times 2$ diagonalisable]\label{ex:syst-2x2-diag}
  Résolvons $X' = AX$ avec $A = \begin{pmatrix} 1 & 2 \\ 0 & 3 \end{pmatrix}$.

  \emph{Valeurs propres.}
  $\det(A - \lambda I) = (1-\lambda)(3-\lambda) = 0$,
  d'où $\lambda_1 = 1$ et $\lambda_2 = 3$.

  \emph{Vecteurs propres.}
  Pour $\lambda_1 = 1$~: $(A - I)V = 0 \Rightarrow
  \begin{pmatrix} 0 & 2 \\ 0 & 2 \end{pmatrix} V = 0$,
  d'où $V_1 = \begin{pmatrix} 1 \\ 0 \end{pmatrix}$.
  Pour $\lambda_2 = 3$~: $(A - 3I)V = 0 \Rightarrow
  \begin{pmatrix} -2 & 2 \\ 0 & 0 \end{pmatrix} V = 0$,
  d'où $V_2 = \begin{pmatrix} 1 \\ 1 \end{pmatrix}$.

  \emph{Solution générale~:}
  \[
    X(t) = c_1 e^t \begin{pmatrix} 1 \\ 0 \end{pmatrix}
         + c_2 e^{3t} \begin{pmatrix} 1 \\ 1 \end{pmatrix}
    = \begin{pmatrix} c_1 e^t + c_2 e^{3t} \\ c_2 e^{3t} \end{pmatrix}.
  \]
\end{exemple}

\subsection{Valeurs propres complexes}\label{subsec:vp-complexes}

\begin{proposition}[Forme réelle des solutions complexes]\label{prop:vp-complexes}
  \index{valeurs propres complexes!forme réelle}
  Soit $A \in \mathcal{M}_n(\R)$ avec une valeur propre complexe
  $\lambda = \alpha + i\beta$ ($\beta \neq 0$) et un vecteur propre
  associé $V = U + iW \in \C^n$. Alors $\bar\lambda = \alpha - i\beta$
  est aussi valeur propre, de vecteur propre $\bar V = U - iW$, et
  les deux solutions réelles indépendantes associées sont
  \begin{align}\label{eq:sol-reelle-complexe}
    X_1(t) &= e^{\alpha t}\bigl(\cos(\beta t)\,U - \sin(\beta t)\,W\bigr), \\
    X_2(t) &= e^{\alpha t}\bigl(\sin(\beta t)\,U + \cos(\beta t)\,W\bigr).
  \end{align}
\end{proposition}

\begin{proof}
  La solution complexe $e^{\lambda t}V = e^{(\alpha+i\beta)t}(U+iW)$
  se développe en
  \[
    e^{\alpha t}\bigl[\cos(\beta t) + i\sin(\beta t)\bigr](U + iW)
    = e^{\alpha t}\bigl[(\cos(\beta t)\,U - \sin(\beta t)\,W)
      + i(\sin(\beta t)\,U + \cos(\beta t)\,W)\bigr].
  \]
  Les parties réelle et imaginaire sont des solutions réelles du système.
\end{proof}

\begin{exemple}[Valeurs propres complexes conjuguées]\label{ex:vp-complexes}
  Soit $A = \begin{pmatrix} 0 & -2 \\ 2 & 0 \end{pmatrix}$.
  On a $\det(A - \lambda I) = \lambda^2 + 4 = 0$, d'où
  $\lambda = \pm 2i$.

  Pour $\lambda = 2i$~: $(A - 2iI)V = 0 \Rightarrow
  \begin{pmatrix} -2i & -2 \\ 2 & -2i \end{pmatrix}V = 0$,
  d'où $V = \begin{pmatrix} 1 \\ -i \end{pmatrix}
  = \begin{pmatrix} 1 \\ 0 \end{pmatrix}
    + i\begin{pmatrix} 0 \\ -1 \end{pmatrix}$.

  Ici $\alpha = 0$, $\beta = 2$, $U = \bigl(\begin{smallmatrix}1\\0\end{smallmatrix}\bigr)$,
  $W = \bigl(\begin{smallmatrix}0\\-1\end{smallmatrix}\bigr)$.
  La solution générale est~:
  \[
    X(t) = c_1 \begin{pmatrix} \cos 2t \\ \sin 2t \end{pmatrix}
         + c_2 \begin{pmatrix} \sin 2t \\ -\cos 2t \end{pmatrix}.
  \]
  Les trajectoires sont des \emph{cercles} centrés à l'origine~: c'est un \emph{centre}.
\end{exemple}

\subsection{Valeurs propres répétées~: vecteurs propres généralisés}\label{subsec:vp-repetees}

\begin{definition}[Vecteur propre généralisé]\label{def:vecteur-propre-gen}
  \index{vecteur propre généralisé}
  Soit $\lambda$ une valeur propre de~$A$ de multiplicité algébrique~$m$.
  Un vecteur $V \neq 0$ est un \emph{vecteur propre généralisé} de rang~$k$
  si $(A - \lambda I)^k V = 0$ mais $(A - \lambda I)^{k-1} V \neq 0$.
\end{definition}

\begin{proposition}[Solution associée à un bloc de Jordan]\label{prop:sol-jordan}
  \index{bloc de Jordan!solution}
  Soit $J_m(\lambda) = \lambda I_m + N$ le bloc de Jordan de taille~$m$
  associé à~$\lambda$, où $N$ est nilpotente ($N^m = 0$). Alors
  \[
    e^{tJ_m(\lambda)} = e^{\lambda t}
    \begin{pmatrix}
      1 & t & \frac{t^2}{2!} & \cdots & \frac{t^{m-1}}{(m-1)!} \\
      0 & 1 & t & \cdots & \frac{t^{m-2}}{(m-2)!} \\
      \vdots & & \ddots & \ddots & \vdots \\
      0 & 0 & \cdots & 1 & t \\
      0 & 0 & \cdots & 0 & 1
    \end{pmatrix}.
  \]
\end{proposition}

\begin{proof}
  Puisque $\lambda I$ et $N$ commutent,
  $e^{tJ_m(\lambda)} = e^{t\lambda I} e^{tN} = e^{\lambda t} e^{tN}$.
  Or $N^m = 0$, donc $e^{tN} = \sum_{k=0}^{m-1} \frac{t^k N^k}{k!}$,
  et les coefficients $(i,j)$ de $N^k$ valent $\delta_{j,i+k}$.
\end{proof}

\begin{methode}[Résolution avec valeurs propres répétées ($2 \times 2$)]\label{meth:vp-repetees}
  \index{méthode!valeurs propres répétées}
  Si $A$ a une valeur propre double $\lambda$ avec un seul vecteur
  propre $V_1$ (à proportionnalité près), on cherche un vecteur propre
  généralisé $V_2$ tel que $(A - \lambda I)V_2 = V_1$. La solution
  générale est alors~:
  \[
    X(t) = c_1 e^{\lambda t} V_1
         + c_2 e^{\lambda t}(V_2 + tV_1).
  \]
\end{methode}

\begin{exemple}[Valeur propre double, un seul vecteur propre]\label{ex:vp-double}
  Soit $A = \begin{pmatrix} 3 & 1 \\ 0 & 3 \end{pmatrix}$.
  La seule valeur propre est $\lambda = 3$ (double). On a
  $A - 3I = \begin{pmatrix} 0 & 1 \\ 0 & 0 \end{pmatrix}$,
  d'où $V_1 = \begin{pmatrix} 1 \\ 0 \end{pmatrix}$.

  Cherchons $V_2$ tel que $(A - 3I)V_2 = V_1$~:
  $\begin{pmatrix} 0 & 1 \\ 0 & 0 \end{pmatrix}
  \begin{pmatrix} a \\ b \end{pmatrix}
  = \begin{pmatrix} 1 \\ 0 \end{pmatrix}$,
  d'où $b = 1$ et $a$ libre, on prend $V_2 = \begin{pmatrix} 0 \\ 1 \end{pmatrix}$.

  Solution générale~:
  \[
    X(t) = c_1 e^{3t} \begin{pmatrix} 1 \\ 0 \end{pmatrix}
         + c_2 e^{3t} \left[\begin{pmatrix} 0 \\ 1 \end{pmatrix}
           + t\begin{pmatrix} 1 \\ 0 \end{pmatrix}\right]
    = e^{3t} \begin{pmatrix} c_1 + c_2 t \\ c_2 \end{pmatrix}.
  \]
\end{exemple}

% ============================================================
\section{Classification des systèmes $2\times 2$}\label{sec:classification-2x2}
% ============================================================

\begin{theoreme}[Classification des points d'équilibre en dimension~2]\label{thm:classif-2d}
  \index{point d'équilibre!classification}
  Soit $A \in \mathcal{M}_2(\R)$ avec $\det(A) \neq 0$. Notons
  $\tau = \tr(A)$, $\delta = \det(A)$ et $\Delta = \tau^2 - 4\delta$
  le discriminant du polynôme caractéristique $\lambda^2 - \tau\lambda + \delta = 0$.

  \begin{center}
  \renewcommand{\arraystretch}{1.4}
  \begin{tabular}{llll}
    \toprule
    \textbf{Condition} & \textbf{Valeurs propres} & \textbf{Type} & \textbf{Stabilité} \\
    \midrule
    $\Delta > 0$, $\delta > 0$, $\tau < 0$
      & $\lambda_1, \lambda_2 < 0$ réels
      & N\oe ud stable\index{noeud stable} & Asympt.\ stable \\
    $\Delta > 0$, $\delta > 0$, $\tau > 0$
      & $\lambda_1, \lambda_2 > 0$ réels
      & N\oe ud instable\index{noeud instable} & Instable \\
    $\Delta > 0$, $\delta < 0$
      & $\lambda_1 < 0 < \lambda_2$
      & Point-selle\index{point-selle} & Instable \\
    $\Delta < 0$, $\tau < 0$
      & $\alpha \pm i\beta$, $\alpha < 0$
      & Foyer stable\index{foyer stable} & Asympt.\ stable \\
    $\Delta < 0$, $\tau > 0$
      & $\alpha \pm i\beta$, $\alpha > 0$
      & Foyer instable\index{foyer instable} & Instable \\
    $\Delta < 0$, $\tau = 0$
      & $\pm i\beta$ purs imaginaires
      & Centre\index{centre} & Stable (non asympt.) \\
    $\Delta = 0$, $\tau < 0$
      & $\lambda < 0$ double
      & N\oe ud dégénéré stable & Asympt.\ stable \\
    $\Delta = 0$, $\tau > 0$
      & $\lambda > 0$ double
      & N\oe ud dégénéré instable & Instable \\
    \bottomrule
  \end{tabular}
  \end{center}
\end{theoreme}

\begin{remarque}\label{rem:trace-det}
  Le plan $(\tau, \delta)$ permet une visualisation synthétique~:
  la parabole $\delta = \tau^2/4$ sépare les valeurs propres réelles
  ($\Delta > 0$, au-dessous) des valeurs propres complexes
  ($\Delta < 0$, au-dessus). L'axe $\delta = 0$ correspond au cas
  dégénéré $\det(A)=0$.
\end{remarque}

% ============================================================
\section{Portraits de phase en dimension~2}\label{sec:portraits-phase}
% ============================================================

\index{portrait de phase}

Nous illustrons chaque type de point d'équilibre par un portrait de phase.

\subsection{N\oe ud stable}\label{subsec:noeud-stable}

\begin{exemple}[N\oe ud stable]\label{ex:noeud-stable}
  Considérons $A = \begin{pmatrix} -2 & 0 \\ 0 & -1 \end{pmatrix}$,
  de valeurs propres $\lambda_1 = -2$ et $\lambda_2 = -1$, toutes deux
  négatives.
\end{exemple}

\begin{center}
\begin{tikzpicture}[>=Stealth, scale=1.1]
  \draw[->] (-2.5,0) -- (2.5,0) node[right]{$x_1$};
  \draw[->] (0,-2.5) -- (0,2.5) node[above]{$x_2$};
  \node at (0,-3){N\oe ud stable ($\lambda_1=-2,\;\lambda_2=-1$)};
  \foreach \angle in {0,30,...,330} {
    \draw[->, blue!70!black, thick]
      ({2*cos(\angle)},{2*sin(\angle)})
      .. controls ({0.8*cos(\angle)},{1.2*sin(\angle)})
      .. (0,0);
  }
  % Axes propres
  \draw[red, very thick, ->] (2.2,0) -- (0.1,0);
  \draw[red, very thick, ->] (-2.2,0) -- (-0.1,0);
  \draw[orange, very thick, ->] (0,2.2) -- (0,0.1);
  \draw[orange, very thick, ->] (0,-2.2) -- (0,-0.1);
  \node[red] at (2.5,-0.4) {$V_1$};
  \node[orange] at (0.5,2.3) {$V_2$};
  \filldraw (0,0) circle (2pt);
\end{tikzpicture}
\end{center}

\subsection{N\oe ud instable}\label{subsec:noeud-instable}

\begin{exemple}[N\oe ud instable]\label{ex:noeud-instable}
  $A = \begin{pmatrix} 2 & 0 \\ 0 & 1 \end{pmatrix}$~:
  $\lambda_1 = 2 > 0$, $\lambda_2 = 1 > 0$. Les trajectoires
  s'éloignent de l'origine.
\end{exemple}

\begin{center}
\begin{tikzpicture}[>=Stealth, scale=1.1]
  \draw[->] (-2.5,0) -- (2.5,0) node[right]{$x_1$};
  \draw[->] (0,-2.5) -- (0,2.5) node[above]{$x_2$};
  \node at (0,-3){N\oe ud instable ($\lambda_1=2,\;\lambda_2=1$)};
  \foreach \angle in {0,30,...,330} {
    \draw[->, blue!70!black, thick]
      (0,0)
      .. controls ({0.8*cos(\angle)},{1.2*sin(\angle)})
      .. ({2*cos(\angle)},{2*sin(\angle)});
  }
  \draw[red, very thick, ->] (0.1,0) -- (2.2,0);
  \draw[red, very thick, ->] (-0.1,0) -- (-2.2,0);
  \draw[orange, very thick, ->] (0,0.1) -- (0,2.2);
  \draw[orange, very thick, ->] (0,-0.1) -- (0,-2.2);
  \filldraw (0,0) circle (2pt);
\end{tikzpicture}
\end{center}

\subsection{Point-selle}\label{subsec:point-selle}

\begin{exemple}[Point-selle]\label{ex:point-selle}
  $A = \begin{pmatrix} 1 & 0 \\ 0 & -2 \end{pmatrix}$~:
  $\lambda_1 = 1 > 0$, $\lambda_2 = -2 < 0$. L'axe $x_2$ est la
  variété stable, l'axe $x_1$ est la variété instable.
\end{exemple}

\begin{center}
\begin{tikzpicture}[>=Stealth, scale=1.1]
  \draw[->] (-2.5,0) -- (2.5,0) node[right]{$x_1$};
  \draw[->] (0,-2.5) -- (0,2.5) node[above]{$x_2$};
  \node at (0,-3){Point-selle ($\lambda_1=1,\;\lambda_2=-2$)};
  % Variété instable (axe x1)
  \draw[red, very thick, ->] (0.1,0) -- (2.3,0);
  \draw[red, very thick, ->] (-0.1,0) -- (-2.3,0);
  % Variété stable (axe x2)
  \draw[green!50!black, very thick, ->] (0,2.3) -- (0,0.1);
  \draw[green!50!black, very thick, ->] (0,-2.3) -- (0,-0.1);
  % Trajectoires hyperboliques
  \foreach \s in {1,-1} {
    \foreach \r in {1,-1} {
      \draw[->, blue!70!black, thick]
        ({\s*0.3},{\r*2}) .. controls ({\s*0.8},{\r*0.8}) .. ({\s*2},{\r*0.3});
    }
  }
  \filldraw (0,0) circle (2pt);
  \node[red] at (2.3,0.4){instable};
  \node[green!50!black] at (0.9,2.3){stable};
\end{tikzpicture}
\end{center}

\subsection{Foyer stable}\label{subsec:foyer-stable}

\begin{exemple}[Foyer stable]\label{ex:foyer-stable}
  $A = \begin{pmatrix} -1 & -2 \\ 2 & -1 \end{pmatrix}$,
  de valeurs propres $\lambda = -1 \pm 2i$.
  Les trajectoires sont des spirales convergeant vers l'origine.
\end{exemple}

\begin{center}
\begin{tikzpicture}[>=Stealth, scale=1.1]
  \draw[->] (-2.5,0) -- (2.5,0) node[right]{$x_1$};
  \draw[->] (0,-2.5) -- (0,2.5) node[above]{$x_2$};
  \node at (0,-3){Foyer stable ($\lambda = -1 \pm 2i$)};
  % Spirale paramétrée
  \draw[->, blue!70!black, thick, samples=50, smooth, domain=0:6.28]
    plot ({2*exp(-0.25*\x)*cos(2*\x r)},{2*exp(-0.25*\x)*sin(2*\x r)});
  \draw[->, blue!70!black, thick, samples=50, smooth, domain=0:6.28]
    plot ({-2*exp(-0.25*\x)*cos(2*\x r)},{-2*exp(-0.25*\x)*sin(2*\x r)});
  \draw[->, red!70!black, thick, samples=50, smooth, domain=0:6.28]
    plot ({1.5*exp(-0.25*\x)*cos(2*\x r + 90)},{1.5*exp(-0.25*\x)*sin(2*\x r + 90)});
  \filldraw (0,0) circle (2pt);
\end{tikzpicture}
\end{center}

\subsection{Foyer instable}\label{subsec:foyer-instable}

\begin{exemple}[Foyer instable]\label{ex:foyer-instable}
  $A = \begin{pmatrix} 1 & -2 \\ 2 & 1 \end{pmatrix}$,
  de valeurs propres $\lambda = 1 \pm 2i$.
  Les spirales s'éloignent de l'origine.
\end{exemple}

\begin{center}
\begin{tikzpicture}[>=Stealth, scale=1.1]
  \draw[->] (-2.5,0) -- (2.5,0) node[right]{$x_1$};
  \draw[->] (0,-2.5) -- (0,2.5) node[above]{$x_2$};
  \node at (0,-3){Foyer instable ($\lambda = 1 \pm 2i$)};
  % Spirale inversée (dessiner en remontant le temps)
  \draw[->, blue!70!black, thick, samples=50, smooth, domain=6.28:0]
    plot ({2*exp(-0.25*\x)*cos(2*\x r)},{2*exp(-0.25*\x)*sin(2*\x r)});
  \draw[->, blue!70!black, thick, samples=50, smooth, domain=6.28:0]
    plot ({-2*exp(-0.25*\x)*cos(2*\x r)},{-2*exp(-0.25*\x)*sin(2*\x r)});
  \draw[->, red!70!black, thick, samples=50, smooth, domain=6.28:0]
    plot ({1.5*exp(-0.25*\x)*cos(2*\x r+90)},{1.5*exp(-0.25*\x)*sin(2*\x r+90)});
  \filldraw (0,0) circle (2pt);
\end{tikzpicture}
\end{center}

\subsection{Centre}\label{subsec:centre}

\begin{exemple}[Centre]\label{ex:centre}
  $A = \begin{pmatrix} 0 & -2 \\ 2 & 0 \end{pmatrix}$,
  de valeurs propres $\lambda = \pm 2i$.
  Les trajectoires sont des ellipses (ici des cercles).
\end{exemple}

\begin{center}
\begin{tikzpicture}[>=Stealth, scale=1.1]
  \draw[->] (-2.5,0) -- (2.5,0) node[right]{$x_1$};
  \draw[->] (0,-2.5) -- (0,2.5) node[above]{$x_2$};
  \node at (0,-3){Centre ($\lambda = \pm 2i$)};
  % Orbites circulaires
  \foreach \r in {0.5, 1.0, 1.5, 2.0} {
    \draw[blue!70!black, thick] (0,0) circle (\r);
  }
  % Flèches de direction
  \draw[->, blue!70!black, very thick] (0,1) -- (-0.15,1);
  \draw[->, blue!70!black, very thick] (0,1.5) -- (-0.15,1.5);
  \draw[->, blue!70!black, very thick] (0,2) -- (-0.15,2);
  \draw[->, blue!70!black, very thick] (0,-1) -- (0.15,-1);
  \draw[->, blue!70!black, very thick] (0,-1.5) -- (0.15,-1.5);
  \filldraw (0,0) circle (2pt);
\end{tikzpicture}
\end{center}

\subsection{N\oe ud dégénéré (étoile et non-étoile)}\label{subsec:noeud-degenere}

\begin{exemple}[N\oe ud dégénéré étoile]\label{ex:noeud-etoile}
  $A = \begin{pmatrix} -1 & 0 \\ 0 & -1 \end{pmatrix} = -I_2$.
  La valeur propre $\lambda = -1$ est double avec deux vecteurs propres
  indépendants. Toutes les droites passant par l'origine sont
  des trajectoires rectilignes convergentes~: c'est un
  \emph{n\oe ud étoile}\index{noeud étoile}.
\end{exemple}

\begin{center}
\begin{tikzpicture}[>=Stealth, scale=1.1]
  \draw[->] (-2.5,0) -- (2.5,0) node[right]{$x_1$};
  \draw[->] (0,-2.5) -- (0,2.5) node[above]{$x_2$};
  \node at (0,-3){N\oe ud étoile stable ($\lambda = -1$ double)};
  \foreach \angle in {0,22.5,...,337.5} {
    \draw[->, blue!70!black, thick]
      ({2*cos(\angle)},{2*sin(\angle)}) -- ({0.1*cos(\angle)},{0.1*sin(\angle)});
  }
  \filldraw (0,0) circle (2pt);
\end{tikzpicture}
\end{center}

\begin{exemple}[N\oe ud dégénéré non-étoile]\label{ex:noeud-non-etoile}
  $A = \begin{pmatrix} -1 & 1 \\ 0 & -1 \end{pmatrix}$.
  Valeur propre $\lambda = -1$ double, un seul vecteur propre
  $V_1 = \bigl(\begin{smallmatrix}1\\0\end{smallmatrix}\bigr)$.
  Les trajectoires sont tangentes à la direction propre en l'origine.
\end{exemple}

\begin{center}
\begin{tikzpicture}[>=Stealth, scale=1.1]
  \draw[->] (-2.5,0) -- (2.5,0) node[right]{$x_1$};
  \draw[->] (0,-2.5) -- (0,2.5) node[above]{$x_2$};
  \node at (0,-3){N\oe ud dégénéré stable ($\lambda=-1$ double, Jordan)};
  % Direction propre
  \draw[red, very thick, ->] (2.2,0) -- (0.1,0);
  \draw[red, very thick, ->] (-2.2,0) -- (-0.1,0);
  % Trajectoires courbes tangentes à l'axe x1
  \foreach \s in {0.5, 1.0, 1.5} {
    \draw[->, blue!70!black, thick]
      ({-2},{\s}) .. controls ({-0.5},{\s*0.3}) .. (0,0);
    \draw[->, blue!70!black, thick]
      ({2},{-\s}) .. controls ({0.5},{-\s*0.3}) .. (0,0);
    \draw[->, blue!70!black, thick]
      ({-2},{-\s}) .. controls ({-0.5},{-\s*0.3}) .. (0,0);
    \draw[->, blue!70!black, thick]
      ({2},{\s}) .. controls ({0.5},{\s*0.3}) .. (0,0);
  }
  \filldraw (0,0) circle (2pt);
\end{tikzpicture}
\end{center}

% ============================================================
\section{Matrice fondamentale et wronskien matriciel}\label{sec:matrice-fondamentale}
% ============================================================

\begin{definition}[Système fondamental, matrice fondamentale]\label{def:matrice-fondamentale}
  \index{matrice fondamentale}\index{système fondamental de solutions}
  Soit $X' = A(t)X$ un système linéaire homogène sur un intervalle~$I$.
  \begin{enumerate}[label=(\roman*)]
    \item Un \emph{système fondamental de solutions} est une famille
      $(X_1,\ldots,X_n)$ de $n$ solutions formant, pour tout
      $t \in I$, une base de~$\R^n$.
    \item La \emph{matrice fondamentale} (ou résolvante) associée est
      la matrice $\Phi(t) = \bigl(X_1(t) \mid \cdots \mid X_n(t)\bigr)
      \in \mathcal{M}_n(\R)$, dont les colonnes sont les solutions
      du système fondamental.
  \end{enumerate}
\end{definition}

\begin{remarque}\label{rem:resolvante-exp}
  Pour le système à coefficients constants $X' = AX$, la matrice
  résolvante vérifiant $\Phi(0) = I_n$ est $\Phi(t) = e^{tA}$.
\end{remarque}

\begin{definition}[Wronskien matriciel]\label{def:wronskien-matriciel}
  \index{wronskien!matriciel}
  Le \emph{wronskien} du système fondamental $(X_1,\ldots,X_n)$ est
  la fonction
  \[
    W(t) = \det\bigl(\Phi(t)\bigr) = \det\bigl(X_1(t),\ldots,X_n(t)\bigr).
  \]
\end{definition}

\begin{theoreme}[Formule d'Abel--Liouville]\label{thm:abel-liouville}
  \index{formule d'Abel--Liouville}\index{Abel}\index{Liouville}
  Soient $(X_1,\ldots,X_n)$ des solutions de $X' = A(t)X$ sur~$I$ et
  $W(t) = \det(X_1(t),\ldots,X_n(t))$. Alors, pour tout $t_0 \in I$~:
  \begin{equation}\label{eq:abel-liouville}
    W(t) = W(t_0)\,\exp\!\left(\int_{t_0}^{t} \tr\bigl(A(s)\bigr)\,\dd s\right).
  \end{equation}
  En particulier, $W$ ne s'annule jamais ou est identiquement nul.
\end{theoreme}

\begin{proof}
  Posons $\Phi = (X_1 \mid \cdots \mid X_n)$. On a
  $\Phi' = A(t)\Phi$. Dérivons le déterminant~:
  \[
    W'(t) = \frac{\dd}{\dd t}\det(\Phi(t)).
  \]
  En utilisant la multilinéarité du déterminant par rapport aux colonnes~:
  \[
    W'(t) = \sum_{j=1}^{n} \det\bigl(X_1,\ldots,X_j',\ldots,X_n\bigr).
  \]
  Or $X_j' = A(t)X_j = \sum_{i=1}^n a_{ij}(t) X_i$, donc le $j$-ème
  terme vaut
  \[
    \det\bigl(X_1,\ldots,\textstyle\sum_i a_{ij} X_i,\ldots,X_n\bigr)
    = a_{jj}(t)\,\det(X_1,\ldots,X_n) = a_{jj}(t)\,W(t),
  \]
  car les termes $i \neq j$ donnent un déterminant avec deux colonnes
  égales, donc nul. On obtient $W'(t) = \tr(A(t))\,W(t)$, qui est
  une EDO linéaire scalaire dont la solution est~\eqref{eq:abel-liouville}.
\end{proof}

\begin{corollaire}\label{cor:wronskien-non-nul}
  Les solutions $X_1,\ldots,X_n$ forment un système fondamental si et
  seulement si $W(t_0) \neq 0$ pour un (et donc tout) $t_0 \in I$.
\end{corollaire}

% ============================================================
\section{Systèmes non homogènes~: variation des constantes}\label{sec:syst-variation-cst}
% ============================================================

\begin{theoreme}[Variation des constantes pour les systèmes]\label{thm:var-cst-syst}
  \index{variation des constantes!systèmes}
  La solution du système complet $X' = A(t)X + B(t)$ vérifiant
  $X(t_0) = X_0$ est
  \begin{equation}\label{eq:var-cst-syst}
    X(t) = \Phi(t)\,\Phi(t_0)^{-1}\,X_0
         + \Phi(t)\int_{t_0}^{t} \Phi(s)^{-1}\,B(s)\,\dd s,
  \end{equation}
  où $\Phi$ est une matrice fondamentale du système homogène associé.
\end{theoreme}

\begin{proof}
  On cherche la solution sous la forme $X(t) = \Phi(t)\,C(t)$ avec
  $C \colon I \to \R^n$ (méthode de variation des constantes).
  En dérivant~:
  \[
    X' = \Phi'\,C + \Phi\,C' = A\Phi\,C + \Phi\,C'.
  \]
  Or $X' = AX + B = A\Phi C + B$, d'où $\Phi\,C' = B$,
  soit $C'(t) = \Phi(t)^{-1}B(t)$. En intégrant~:
  \[
    C(t) = C(t_0) + \int_{t_0}^t \Phi(s)^{-1}B(s)\,\dd s.
  \]
  La condition $X(t_0) = X_0$ donne $\Phi(t_0)C(t_0) = X_0$,
  donc $C(t_0) = \Phi(t_0)^{-1}X_0$, et on obtient~\eqref{eq:var-cst-syst}.
\end{proof}

\begin{corollaire}[Cas des coefficients constants]\label{cor:var-cst-cst}
  Si $A$ est constante et $t_0 = 0$~:
  \begin{equation}\label{eq:duhamel}
    X(t) = e^{tA}X_0 + \int_0^t e^{(t-s)A}B(s)\,\dd s.
  \end{equation}
  Cette formule est appelée \emph{formule de Duhamel}\index{formule de Duhamel}.
\end{corollaire}

\begin{exemple}[Système non homogène $2\times 2$]\label{ex:syst-non-homogene}
  Résolvons $X' = AX + B(t)$ avec
  $A = \begin{pmatrix} 1 & 0 \\ 0 & 2 \end{pmatrix}$ et
  $B(t) = \begin{pmatrix} e^t \\ 0 \end{pmatrix}$, $X(0) = 0$.

  La matrice fondamentale est $\Phi(t) = \begin{pmatrix} e^t & 0 \\ 0 & e^{2t} \end{pmatrix}$, d'où
  $\Phi(t)^{-1} = \begin{pmatrix} e^{-t} & 0 \\ 0 & e^{-2t} \end{pmatrix}$.

  On calcule~:
  \[
    \int_0^t \Phi(s)^{-1}B(s)\,\dd s
    = \int_0^t \begin{pmatrix} e^{-s} \cdot e^s \\ 0 \end{pmatrix}\dd s
    = \int_0^t \begin{pmatrix} 1 \\ 0 \end{pmatrix}\dd s
    = \begin{pmatrix} t \\ 0 \end{pmatrix}.
  \]
  Donc
  $X(t) = \Phi(t)\begin{pmatrix}t\\0\end{pmatrix}
  = \begin{pmatrix}te^t\\0\end{pmatrix}$.
\end{exemple}

% ============================================================
\section{Exercices}\label{sec:exo-ch5}
% ============================================================

\begin{exercice}[Résolution par diagonalisation]\label{exo:syst-diag}
  Résoudre le système $X' = AX$ avec $X(0) = \begin{pmatrix}1\\2\end{pmatrix}$ et~:
  \begin{enumerate}[label=\alph*)]
    \item $A = \begin{pmatrix} 3 & -1 \\ -1 & 3 \end{pmatrix}$
    \item $A = \begin{pmatrix} 1 & 4 \\ 1 & 1 \end{pmatrix}$
  \end{enumerate}
\end{exercice}

\begin{exercice}[Valeurs propres complexes]\label{exo:syst-complexe}
  Résoudre $X' = AX$ pour $A = \begin{pmatrix} 1 & -5 \\ 1 & -1 \end{pmatrix}$.
  Identifier le type du point d'équilibre et esquisser le portrait de phase.
\end{exercice}

\begin{exercice}[Valeur propre double]\label{exo:syst-jordan}
  Résoudre $X' = AX$ avec $A = \begin{pmatrix} 5 & -4 \\ 1 & 1 \end{pmatrix}$.
  \emph{Indication~:} la valeur propre est $\lambda = 3$ (double).
\end{exercice}

\begin{exercice}[Exponentielle de matrice]\label{exo:exp-matrice}
  Calculer $e^{tA}$ dans les cas suivants~:
  \begin{enumerate}[label=\alph*)]
    \item $A = \begin{pmatrix} 0 & 1 \\ -1 & 0 \end{pmatrix}$
    \item $A = \begin{pmatrix} 2 & 1 \\ 0 & 2 \end{pmatrix}$
    \item $A = \begin{pmatrix} 1 & 1 & 0 \\ 0 & 1 & 1 \\ 0 & 0 & 1 \end{pmatrix}$
  \end{enumerate}
\end{exercice}

\begin{exercice}[Formule d'Abel--Liouville]\label{exo:abel-liouville}
  Soit $X' = A(t)X$ avec $A(t) = \begin{pmatrix} 0 & 1 \\ -t & 0 \end{pmatrix}$.
  \begin{enumerate}[label=\alph*)]
    \item Montrer que $\tr(A(t)) = 0$ et en déduire que le wronskien
      de tout système fondamental est constant.
    \item Si $X_1(t) = \begin{pmatrix}\varphi(t)\\\varphi'(t)\end{pmatrix}$
      est une solution, exprimer une seconde solution indépendante à l'aide
      de~$\varphi$.
  \end{enumerate}
\end{exercice}

\begin{exercice}[Système non homogène $\bigstar\bigstar$]\label{exo:syst-non-homog}
  Résoudre $X' = AX + B(t)$ avec
  $A = \begin{pmatrix} 0 & 1 \\ -1 & 0 \end{pmatrix}$ et
  $B(t) = \begin{pmatrix} 0 \\ \sin t \end{pmatrix}$, $X(0) = 0$.
\end{exercice}

\begin{exercice}[Classification $\bigstar\bigstar$]\label{exo:classification}
  Pour chacune des matrices suivantes, déterminer le type du point
  d'équilibre de $X' = AX$ et esquisser le portrait de phase~:
  \begin{enumerate}[label=\alph*)]
    \item $A = \begin{pmatrix} -3 & 2 \\ -2 & 1 \end{pmatrix}$
    \item $A = \begin{pmatrix} 1 & -1 \\ 2 & -1 \end{pmatrix}$
    \item $A = \begin{pmatrix} -1 & -1 \\ 4 & -1 \end{pmatrix}$
  \end{enumerate}
\end{exercice}

\begin{exercice}[Système $3\times 3$ $\bigstar\bigstar\bigstar$]\label{exo:syst-3x3}
  Résoudre le système $X' = AX$ pour
  $A = \begin{pmatrix} 2 & 1 & 0 \\ 0 & 2 & 1 \\ 0 & 0 & 2 \end{pmatrix}$.
  En déduire l'exponentielle $e^{tA}$.
\end{exercice}

% ============================================================
\section*{Résumé du chapitre~5}
\addcontentsline{toc}{section}{Résumé du chapitre~5}
% ============================================================

\begin{center}
\fbox{\parbox{0.92\textwidth}{
\textbf{Points clés du chapitre~5~:}
\begin{itemize}
  \item Un système $X' = A(t)X + B(t)$ admet une unique solution
    globale pour toute condition initiale (Cauchy--Lipschitz).
  \item L'espace des solutions du système homogène est un espace
    vectoriel de dimension~$n$.
  \item Pour $A$ constante, la solution est $X(t) = e^{tA}X_0$.
    Le calcul de $e^{tA}$ repose sur la réduction de $A$
    (diagonalisation ou forme de Jordan).
  \item La formule d'Abel--Liouville relie le wronskien à la trace
    de~$A(t)$~: $W(t) = W(t_0)\exp\int_{t_0}^t \tr A(s)\,\dd s$.
  \item La classification des points d'équilibre en dimension~2
    dépend de $\tau = \tr(A)$, $\delta = \det(A)$ et du signe de
    $\Delta = \tau^2 - 4\delta$.
  \item La méthode de variation des constantes donne la solution
    du système complet~:
    $X(t) = \Phi(t)\Phi(t_0)^{-1}X_0
           + \Phi(t)\int_{t_0}^t \Phi(s)^{-1}B(s)\,\dd s$.
\end{itemize}
}}
\end{center}

\newpage

% ████████████████████████████████████████████████████████████████████████
\chapter{Méthodes de résolution explicite~: wronskien, variation des
  constantes et séries entières}\label{ch:methodes-resolution}
% ████████████████████████████████████████████████████████████████████████

\minitoc
\vspace{1cm}

\section*{Introduction}
\addcontentsline{toc}{section}{Introduction}

Ce chapitre rassemble les principales méthodes de résolution
\emph{explicite} des équations et systèmes différentiels linéaires.
Si le chapitre précédent a développé la théorie des systèmes, nous
adoptons ici un point de vue plus calculatoire, en détaillant~:
le \emph{wronskien}\index{wronskien} et ses applications,
la méthode de \emph{variation des constantes}\index{variation des constantes}
(pour les équations d'ordre~$n$ et les systèmes),
une introduction aux \emph{fonctions de Green}\index{fonction de Green},
la méthode des \emph{séries entières}\index{séries entières!méthode},
et la \emph{réduction d'ordre}\index{réduction d'ordre}.

% ============================================================
\section{Wronskien~: définition, propriétés, théorème d'Abel}\label{sec:wronskien}
% ============================================================

\subsection{Wronskien d'une équation scalaire d'ordre~$n$}\label{subsec:wronskien-scalaire}

\begin{definition}[Wronskien]\label{def:wronskien}
  \index{wronskien!définition}
  Soient $y_1,\ldots,y_n$ des fonctions $(n-1)$ fois dérivables sur un
  intervalle~$I$. Le \emph{wronskien} de $(y_1,\ldots,y_n)$ est la
  fonction
  \begin{equation}\label{eq:wronskien-def}
    W(t) = W[y_1,\ldots,y_n](t)
    = \begin{vmatrix}
        y_1(t) & y_2(t) & \cdots & y_n(t) \\
        y_1'(t) & y_2'(t) & \cdots & y_n'(t) \\
        \vdots & \vdots & \ddots & \vdots \\
        y_1^{(n-1)}(t) & y_2^{(n-1)}(t) & \cdots & y_n^{(n-1)}(t)
      \end{vmatrix}.
  \end{equation}
\end{definition}

\begin{proposition}[Wronskien et indépendance linéaire]\label{prop:wronskien-indep}
  \index{wronskien!et indépendance linéaire}
  Soient $y_1,\ldots,y_n$ des solutions de l'équation linéaire homogène
  d'ordre~$n$
  \begin{equation}\label{eq:edo-lin-n}
    y^{(n)} + a_{n-1}(t)\,y^{(n-1)} + \cdots + a_1(t)\,y' + a_0(t)\,y = 0,
  \end{equation}
  où les $a_i$ sont continues sur~$I$. Les propriétés suivantes sont
  équivalentes~:
  \begin{enumerate}[label=(\roman*)]
    \item $(y_1,\ldots,y_n)$ est un système fondamental de solutions.
    \item $W(t_0) \neq 0$ pour un $t_0 \in I$.
    \item $W(t) \neq 0$ pour tout $t \in I$.
  \end{enumerate}
\end{proposition}

\begin{theoreme}[Théorème d'Abel pour les EDO scalaires]\label{thm:abel-scalaire}
  \index{Abel!théorème (EDO scalaire)}\index{wronskien!formule d'Abel}
  Soient $y_1,\ldots,y_n$ des solutions de~\eqref{eq:edo-lin-n}. Alors
  \begin{equation}\label{eq:abel-scalaire}
    W(t) = W(t_0)\,\exp\!\left(-\int_{t_0}^t a_{n-1}(s)\,\dd s\right).
  \end{equation}
\end{theoreme}

\begin{proof}
  On ramène l'équation~\eqref{eq:edo-lin-n} au système $X' = A(t)X$
  par le changement de variables $X = (y, y', \ldots, y^{(n-1)})^\top$,
  avec la matrice compagnon\index{matrice compagnon}
  \[
    A(t) = \begin{pmatrix}
      0 & 1 & 0 & \cdots & 0 \\
      0 & 0 & 1 & \cdots & 0 \\
      \vdots & & & \ddots & \vdots \\
      0 & 0 & 0 & \cdots & 1 \\
      -a_0(t) & -a_1(t) & -a_2(t) & \cdots & -a_{n-1}(t)
    \end{pmatrix}.
  \]
  On a $\tr(A(t)) = -a_{n-1}(t)$. Le théorème d'Abel--Liouville
  (théorème~\ref{thm:abel-liouville}) appliqué au système donne
  immédiatement~\eqref{eq:abel-scalaire}, car le wronskien du système
  fondamental coïncide avec le wronskien scalaire.
\end{proof}

\begin{exemple}[Application du théorème d'Abel]\label{ex:abel-application}
  Considérons l'équation $y'' + p(t)y' + q(t)y = 0$ (ordre~2).
  Si $y_1$ est une solution connue non nulle, le wronskien est
  $W(t) = W_0\,e^{-\int_{t_0}^t p(s)\,\dd s}$.
  Or $W = y_1 y_2' - y_1' y_2$, d'où
  \[
    \frac{\dd}{\dd t}\!\left(\frac{y_2}{y_1}\right)
    = \frac{y_1 y_2' - y_1' y_2}{y_1^2}
    = \frac{W(t)}{y_1(t)^2}.
  \]
  En intégrant~:
  \begin{equation}\label{eq:2e-solution-abel}
    y_2(t) = y_1(t)\int_{t_0}^t \frac{W_0\,e^{-\int_{t_0}^s p(u)\,\dd u}}
    {y_1(s)^2}\,\dd s.
  \end{equation}
  Ceci fournit une seconde solution indépendante à partir d'une seule
  solution connue.
\end{exemple}

% ============================================================
\section{Variation des constantes pour les EDO d'ordre~$n$}\label{sec:var-cst-scalaire}
% ============================================================

\begin{theoreme}[Variation des constantes — ordre $n$]\label{thm:var-cst-ordre-n}
  \index{variation des constantes!ordre $n$}
  Considérons l'équation linéaire complète
  \begin{equation}\label{eq:edo-complete-n}
    y^{(n)} + a_{n-1}(t)\,y^{(n-1)} + \cdots + a_0(t)\,y = f(t),
  \end{equation}
  et soit $(y_1,\ldots,y_n)$ un système fondamental de l'équation
  homogène associée. Une solution particulière de~\eqref{eq:edo-complete-n}
  est
  \begin{equation}\label{eq:var-cst-n}
    y_p(t) = \sum_{k=1}^{n} y_k(t)\int_{t_0}^t
      \frac{W_k(s)}{W(s)}\,f(s)\,\dd s,
  \end{equation}
  où $W(t)$ est le wronskien de $(y_1,\ldots,y_n)$ et $W_k(t)$ est
  le déterminant obtenu en remplaçant la $k$-ème colonne de la matrice
  wronskienne par $(0,0,\ldots,0,1)^\top$.
\end{theoreme}

\begin{proof}
  On cherche une solution particulière sous la forme
  $y_p(t) = \sum_{k=1}^n c_k(t)\,y_k(t)$ avec des fonctions $c_k$
  à déterminer. En imposant les $n-1$ conditions\footnote{Ces conditions
  de Lagrange\index{conditions de Lagrange} garantissent que
  $y_p^{(j)}$ a la même forme que si les $c_k$ étaient constantes,
  pour $j = 0,\ldots,n-2$.}
  \[
    \sum_{k=1}^n c_k'(t)\,y_k^{(j)}(t) = 0, \qquad j = 0,1,\ldots,n-2,
  \]
  et en substituant dans~\eqref{eq:edo-complete-n}, on obtient le système
  de Cramer~:
  \[
    \sum_{k=1}^n c_k'(t)\,y_k^{(n-1)}(t) = f(t).
  \]
  Le système complet pour $(c_1',\ldots,c_n')$ s'écrit sous forme
  matricielle~:
  \[
    \begin{pmatrix}
      y_1 & \cdots & y_n \\
      y_1' & \cdots & y_n' \\
      \vdots & & \vdots \\
      y_1^{(n-1)} & \cdots & y_n^{(n-1)}
    \end{pmatrix}
    \begin{pmatrix} c_1' \\ \vdots \\ c_n' \end{pmatrix}
    = \begin{pmatrix} 0 \\ \vdots \\ 0 \\ f(t) \end{pmatrix}.
  \]
  Le déterminant de ce système est $W(t) \neq 0$. Par la règle de
  Cramer, $c_k'(t) = W_k(t)\,f(t)/W(t)$, et on conclut en intégrant.
\end{proof}

\begin{exemple}[Variation des constantes pour une EDO d'ordre~2]\label{ex:var-cst-ordre2}
  Résolvons $y'' + y = \frac{1}{\cos t}$ sur $]-\pi/2, \pi/2[$.

  \emph{Système fondamental~:} $y_1 = \cos t$, $y_2 = \sin t$.

  \emph{Wronskien~:}
  $W = \begin{vmatrix} \cos t & \sin t \\ -\sin t & \cos t \end{vmatrix}
  = \cos^2 t + \sin^2 t = 1$.

  \emph{Calcul des $W_k$~:}
  \[
    W_1 = \begin{vmatrix} 0 & \sin t \\ 1 & \cos t \end{vmatrix} = -\sin t,
    \qquad
    W_2 = \begin{vmatrix} \cos t & 0 \\ -\sin t & 1 \end{vmatrix} = \cos t.
  \]

  \emph{Solution particulière~:}
  \begin{align*}
    y_p &= \cos t \int \frac{-\sin t}{\cos t}\,\dd t
         + \sin t \int \frac{\cos t}{\cos t}\,\dd t \\
        &= \cos t \cdot \ln\abs{\cos t} + t\sin t.
  \end{align*}

  \emph{Solution générale~:}
  $y(t) = c_1 \cos t + c_2 \sin t + \cos t \cdot \ln\abs{\cos t} + t\sin t$.
\end{exemple}

% ============================================================
\section{Variation des constantes pour les systèmes}\label{sec:var-cst-systemes}
% ============================================================

La méthode a déjà été établie au théorème~\ref{thm:var-cst-syst}.
Nous la reformulons ici dans un cadre plus général et donnons un
exemple détaillé supplémentaire.

\begin{methode}[Variation des constantes pour $X' = A(t)X + B(t)$]\label{meth:var-cst-syst-detail}
  \index{variation des constantes!systèmes (méthode)}
  \begin{enumerate}
    \item Résoudre le système homogène $X' = A(t)X$ pour obtenir une
      matrice fondamentale~$\Phi(t)$.
    \item Poser $X(t) = \Phi(t)\,C(t)$ et dériver~:
      $\Phi'\,C + \Phi\,C' = A\Phi\,C + B$,
      d'où $\Phi\,C' = B$.
    \item Résoudre $C'(t) = \Phi(t)^{-1}B(t)$ par intégration.
    \item La solution est $X(t) = \Phi(t)\bigl[C(t_0) + \int_{t_0}^t \Phi(s)^{-1}B(s)\,\dd s\bigr]$.
  \end{enumerate}
\end{methode}

\begin{exemple}[Variation des constantes~: système $2\times 2$]\label{ex:var-cst-syst-2x2}
  Résolvons $X' = AX + B(t)$ avec
  $A = \begin{pmatrix} 0 & 1 \\ -1 & 0 \end{pmatrix}$,
  $B(t) = \begin{pmatrix} 0 \\ t \end{pmatrix}$, $X(0) = 0$.

  \emph{Étape 1.} Les valeurs propres de $A$ sont $\pm i$.
  La matrice fondamentale avec $\Phi(0) = I_2$ est
  \[
    \Phi(t) = e^{tA} = \begin{pmatrix}
      \cos t & \sin t \\ -\sin t & \cos t
    \end{pmatrix}.
  \]

  \emph{Étape 2.}
  $\Phi(t)^{-1} = \begin{pmatrix}
    \cos t & -\sin t \\ \sin t & \cos t
  \end{pmatrix}$.

  \emph{Étape 3.}
  \begin{align*}
    \Phi(s)^{-1}B(s) &= \begin{pmatrix}
      \cos s & -\sin s \\ \sin s & \cos s
    \end{pmatrix}\begin{pmatrix}0\\s\end{pmatrix}
    = \begin{pmatrix} -s\sin s \\ s\cos s \end{pmatrix}.
  \end{align*}

  Intégrons composante par composante (intégration par parties)~:
  \begin{align*}
    \int_0^t (-s\sin s)\,\dd s &= s\cos s\big|_0^t - \int_0^t \cos s\,\dd s
      = t\cos t - \sin t, \\
    \int_0^t s\cos s\,\dd s &= s\sin s\big|_0^t - \int_0^t \sin s\,\dd s
      = t\sin t + \cos t - 1.
  \end{align*}

  \emph{Étape 4.}
  \begin{align*}
    X(t) &= \Phi(t)\begin{pmatrix}t\cos t - \sin t \\ t\sin t + \cos t - 1\end{pmatrix} \\
    &= \begin{pmatrix}
      \cos t & \sin t \\ -\sin t & \cos t
    \end{pmatrix}
    \begin{pmatrix}t\cos t - \sin t \\ t\sin t + \cos t - 1\end{pmatrix}.
  \end{align*}

  Développons la première composante~:
  \begin{align*}
    x_1(t) &= \cos t(t\cos t - \sin t) + \sin t(t\sin t + \cos t - 1) \\
    &= t\cos^2 t - \sin t\cos t + t\sin^2 t + \sin t\cos t - \sin t \\
    &= t - \sin t.
  \end{align*}
  Et la seconde~:
  \begin{align*}
    x_2(t) &= -\sin t(t\cos t - \sin t) + \cos t(t\sin t + \cos t - 1) \\
    &= -t\sin t\cos t + \sin^2 t + t\sin t\cos t + \cos^2 t - \cos t \\
    &= 1 - \cos t.
  \end{align*}

  \emph{Résultat~:} $X(t) = \begin{pmatrix} t - \sin t \\ 1 - \cos t \end{pmatrix}$.
\end{exemple}

% ============================================================
\section{Fonction de Green}\label{sec:green}
% ============================================================

\begin{definition}[Fonction de Green]\label{def:green}
  \index{fonction de Green}
  Considérons l'équation $y^{(n)} + a_{n-1}(t)y^{(n-1)} + \cdots + a_0(t)y = f(t)$
  sur~$I$, avec un système fondamental $(y_1,\ldots,y_n)$. La
  \emph{fonction de Green} associée est
  \begin{equation}\label{eq:green-def}
    G(t,s) = \sum_{k=1}^n y_k(t)\,\frac{W_k(s)}{W(s)},
  \end{equation}
  où $W$ et $W_k$ sont définis comme au théorème~\ref{thm:var-cst-ordre-n}.
  La solution particulière vérifiant $y_p(t_0) = y_p'(t_0) = \cdots
  = y_p^{(n-1)}(t_0) = 0$ s'écrit alors~:
  \begin{equation}\label{eq:sol-green}
    y_p(t) = \int_{t_0}^t G(t,s)\,f(s)\,\dd s.
  \end{equation}
\end{definition}

\begin{proposition}[Propriétés de la fonction de Green]\label{prop:green-proprietes}
  \index{fonction de Green!propriétés}
  La fonction $G$ vérifie~:
  \begin{enumerate}[label=(\roman*)]
    \item Pour $s$ fixé, $t \mapsto G(t,s)$ est solution de l'équation
      homogène pour $t \neq s$.
    \item $G(s,s) = \frac{\partial G}{\partial t}(s,s)
      = \cdots = \frac{\partial^{n-2} G}{\partial t^{n-2}}(s,s) = 0$
      et $\frac{\partial^{n-1} G}{\partial t^{n-1}}(s,s) = 1$.
    \item La fonction de Green «~encode~» l'opérateur de résolution
      du problème non homogène comme un opérateur intégral.
  \end{enumerate}
\end{proposition}

\begin{exemple}[Fonction de Green pour $y'' + y = f(t)$]\label{ex:green-ordre2}
  Le système fondamental est $y_1 = \cos t$, $y_2 = \sin t$,
  $W = 1$. On a~:
  \[
    G(t,s) = \cos t\cdot(-\sin s) + \sin t \cdot \cos s = \sin(t - s).
  \]
  La solution particulière nulle en $t_0 = 0$ est~:
  \[
    y_p(t) = \int_0^t \sin(t-s)\,f(s)\,\dd s.
  \]
  On reconnaît un \emph{produit de convolution}\index{convolution}
  avec le noyau $\sin$.
\end{exemple}

% ============================================================
\section{Solutions en séries entières}\label{sec:series-entieres}
% ============================================================

\index{séries entières!résolution d'EDO}

Lorsque les coefficients d'une EDO ne sont pas constants mais
analytiques, on peut chercher des solutions sous forme de séries
entières.

\subsection{Points ordinaires}\label{subsec:points-ordinaires}

\begin{definition}[Point ordinaire]\label{def:point-ordinaire}
  \index{point ordinaire}
  Considérons l'équation
  \begin{equation}\label{eq:edo-series}
    y'' + p(t)\,y' + q(t)\,y = 0.
  \end{equation}
  Un point $t_0$ est dit \emph{ordinaire} si $p$ et $q$ sont
  analytiques (développables en séries entières) au voisinage de~$t_0$.
\end{definition}

\begin{theoreme}[Solutions en séries entières aux points ordinaires]\label{thm:series-point-ord}
  \index{séries entières!point ordinaire}
  Si $t_0$ est un point ordinaire de~\eqref{eq:edo-series}, alors
  l'équation admet deux solutions linéairement indépendantes de la forme
  \[
    y(t) = \sum_{n=0}^{\infty} a_n (t - t_0)^n,
  \]
  dont le rayon de convergence est au moins égal au minimum des rayons
  de convergence des séries de $p$ et $q$ en~$t_0$.
\end{theoreme}

\begin{methode}[Recherche de solutions en séries entières]\label{meth:series}
  \index{méthode!séries entières}
  \begin{enumerate}
    \item Poser $y(t) = \sum_{n=0}^\infty a_n t^n$ (en prenant $t_0 = 0$
      pour simplifier).
    \item Calculer $y'(t) = \sum_{n=1}^\infty n\,a_n t^{n-1}$ et
      $y''(t) = \sum_{n=2}^\infty n(n-1)\,a_n t^{n-2}$.
    \item Substituer dans l'équation et regrouper les termes par
      puissances de~$t$.
    \item Identifier les coefficients~: on obtient une \emph{relation
      de récurrence}\index{relation de récurrence} pour les~$a_n$.
    \item Déterminer $a_0$ et $a_1$ à l'aide des conditions initiales.
  \end{enumerate}
\end{methode}

\begin{exemple}[Équation d'Airy]\label{ex:airy}
  \index{équation d'Airy}
  Résolvons $y'' - ty = 0$ en série entière autour de $t_0 = 0$.

  Posons $y = \sum_{n \geq 0} a_n t^n$, d'où
  $y'' = \sum_{n \geq 2} n(n-1)a_n t^{n-2} = \sum_{m \geq 0}(m+2)(m+1)a_{m+2}t^m$.

  Le terme $ty = \sum_{n \geq 0} a_n t^{n+1} = \sum_{m \geq 1} a_{m-1}t^m$.

  L'équation $y'' - ty = 0$ donne~:
  \begin{itemize}
    \item $m = 0$~: $2 \cdot 1 \cdot a_2 = 0$, donc $a_2 = 0$.
    \item $m \geq 1$~: $(m+2)(m+1)a_{m+2} = a_{m-1}$, soit
      \begin{equation}\label{eq:recurrence-airy}
        a_{m+2} = \frac{a_{m-1}}{(m+2)(m+1)}, \qquad m \geq 1.
      \end{equation}
  \end{itemize}

  Les deux suites indépendantes (en posant $a_0$ et $a_1$ libres)~:

  \emph{Suite $a_0$~:} $a_0, 0, 0, \frac{a_0}{3\cdot 2},
  0, 0, \frac{a_0}{6\cdot 5\cdot 3\cdot 2}, \ldots$

  \emph{Suite $a_1$~:} $0, a_1, 0, 0, \frac{a_1}{4\cdot 3},
  0, 0, \frac{a_1}{7\cdot 6\cdot 4\cdot 3}, \ldots$

  Ces séries convergent pour tout $t \in \R$ (rayon de convergence
  infini). Les deux solutions indépendantes sont les \emph{fonctions
  d'Airy}\index{Airy!fonctions} $\mathrm{Ai}(t)$ et $\mathrm{Bi}(t)$.
\end{exemple}

\begin{exemple}[Équation de Hermite]\label{ex:hermite}
  \index{équation de Hermite}\index{Hermite!polynômes de}
  L'équation $y'' - 2ty' + 2\alpha y = 0$ ($\alpha \in \R$) se résout
  en posant $y = \sum a_n t^n$. On obtient la récurrence~:
  \[
    a_{n+2} = \frac{2(n - \alpha)}{(n+2)(n+1)}\,a_n, \qquad n \geq 0.
  \]
  Pour $\alpha = N \in \N$, la série s'interrompt et fournit un
  \emph{polynôme de Hermite}~$H_N(t)$.
\end{exemple}

\subsection{Points singuliers réguliers~: méthode de Frobenius}\label{subsec:frobenius}

\begin{definition}[Point singulier régulier]\label{def:point-sing-reg}
  \index{point singulier régulier}\index{Frobenius}
  Le point $t_0$ est un \emph{point singulier régulier}
  de~\eqref{eq:edo-series} si les fonctions $(t-t_0)\,p(t)$ et
  $(t-t_0)^2\,q(t)$ sont analytiques en~$t_0$, c'est-à-dire si $p$ a
  au plus un pôle simple et $q$ un pôle d'ordre au plus~2 en~$t_0$.
\end{definition}

\begin{theoreme}[Méthode de Frobenius]\label{thm:frobenius}
  \index{Frobenius!méthode de}
  Si $t_0$ est un point singulier régulier de~\eqref{eq:edo-series},
  il existe au moins une solution de la forme
  \begin{equation}\label{eq:frobenius}
    y(t) = (t-t_0)^r \sum_{n=0}^{\infty} a_n (t-t_0)^n, \qquad a_0 \neq 0,
  \end{equation}
  où $r$ est une racine de l'\emph{équation indicielle}\index{équation indicielle}
  \[
    r(r-1) + p_0\,r + q_0 = 0, \qquad
    p_0 = \lim_{t\to t_0}(t-t_0)p(t), \quad
    q_0 = \lim_{t\to t_0}(t-t_0)^2 q(t).
  \]
  Si les deux racines $r_1, r_2$ (avec $\operatorname{Re}(r_1) \geq
  \operatorname{Re}(r_2)$) de l'équation indicielle ne diffèrent pas
  d'un entier, on obtient deux solutions indépendantes de la
  forme~\eqref{eq:frobenius}. Sinon, la seconde solution peut contenir
  un terme logarithmique.
\end{theoreme}

\begin{exemple}[Équation de Bessel d'ordre~$\nu$]\label{ex:bessel}
  \index{Bessel!équation de}\index{Bessel!fonctions de}
  L'équation de Bessel $t^2 y'' + ty' + (t^2 - \nu^2)y = 0$ a un
  point singulier régulier en $t_0 = 0$. Sous forme normalisée~:
  $y'' + \frac{1}{t}y' + \frac{t^2 - \nu^2}{t^2}y = 0$, d'où
  $p_0 = 1$ et $q_0 = -\nu^2$.

  L'équation indicielle est $r(r-1) + r - \nu^2 = r^2 - \nu^2 = 0$,
  de racines $r_1 = \nu$ et $r_2 = -\nu$.

  Pour $r = \nu$, la méthode de Frobenius donne la \emph{fonction de
  Bessel de première espèce}~:
  \[
    J_\nu(t) = \sum_{m=0}^{\infty}
      \frac{(-1)^m}{m!\,\Gamma(m+\nu+1)}
      \left(\frac{t}{2}\right)^{2m+\nu}.
  \]
\end{exemple}

% ============================================================
\section{Réduction d'ordre}\label{sec:reduction-ordre}
% ============================================================

\begin{theoreme}[Réduction d'ordre]\label{thm:reduction-ordre}
  \index{réduction d'ordre}
  Soit $y_1$ une solution non nulle de $y'' + p(t)y' + q(t)y = 0$.
  En posant $y = y_1\,v$, on obtient pour~$v$ l'équation
  \begin{equation}\label{eq:reduction-ordre}
    v'' + \left(2\frac{y_1'}{y_1} + p\right)v' = 0,
  \end{equation}
  qui est du premier ordre en $w = v'$. Sa résolution donne
  une seconde solution indépendante~:
  \begin{equation}\label{eq:2e-solution}
    y_2(t) = y_1(t) \int \frac{e^{-\int p(t)\,\dd t}}{y_1(t)^2}\,\dd t.
  \end{equation}
\end{theoreme}

\begin{proof}
  Posons $y = y_1 v$. Alors $y' = y_1' v + y_1 v'$ et
  $y'' = y_1'' v + 2y_1' v' + y_1 v''$. En substituant~:
  \begin{align*}
    y_1'' v + 2y_1' v' + y_1 v'' + p(y_1' v + y_1 v') + q y_1 v &= 0 \\
    \underbrace{(y_1'' + py_1' + qy_1)}_{=\,0} v + y_1 v'' + (2y_1' + py_1)v' &= 0.
  \end{align*}
  Il reste $y_1 v'' + (2y_1' + py_1)v' = 0$, soit
  $v'' + (2y_1'/y_1 + p)v' = 0$.
  C'est une EDO du premier ordre en $w = v'$~:
  $w' + (2y_1'/y_1 + p)w = 0$, d'où
  $w(t) = \frac{C}{y_1(t)^2}\,e^{-\int p(t)\,\dd t}$.
  En intégrant~$w$, on obtient~$v$, puis $y_2 = y_1 v$.
\end{proof}

\begin{exemple}[Réduction d'ordre]\label{ex:reduction-ordre}
  Considérons $y'' - \frac{2}{t}y' + \frac{2}{t^2}y = 0$ sur
  $]0,+\infty[$, qui admet $y_1(t) = t$ comme solution.

  Posons $y = tv$, alors~:
  \[
    v'' + \left(\frac{2}{t} - \frac{2}{t}\right)v' = 0,
    \qquad\text{soit } v'' = 0.
  \]
  Donc $v(t) = c_1 + c_2 t$, et la seconde solution est
  $y_2(t) = t \cdot t = t^2$.

  Solution générale~: $y(t) = c_1 t + c_2 t^2$.
\end{exemple}

% ============================================================
\section{Tableau récapitulatif des méthodes}\label{sec:tableau-methodes}
% ============================================================

\begin{center}
\renewcommand{\arraystretch}{1.5}
\begin{tabular}{p{4.5cm}p{4.5cm}p{4.5cm}}
  \toprule
  \textbf{Type d'équation} & \textbf{Méthode} & \textbf{Référence} \\
  \midrule
  $y' = f(t,y)$ séparable
    & Séparation des variables
    & Chap.~précédents \\
  $y' + p(t)y = q(t)$ (linéaire ordre~1)
    & Facteur intégrant $e^{\int p}$
    & Chap.~précédents \\
  $y'' + ay' + by = 0$ (coeff.\ constants)
    & Polynôme caractéristique
    & Chap.~précédents \\
  $y'' + ay' + by = f(t)$ (coeff.\ constants)
    & Variation des constantes\newline ou second membre particulier
    & \S\ref{sec:var-cst-scalaire} \\
  $y^{(n)} + a_{n-1}y^{(n-1)} + \cdots = f$
    & Variation des constantes\newline (formule de Lagrange--Cramer)
    & Thm~\ref{thm:var-cst-ordre-n} \\
  $X' = AX + B$ (système)
    & Variation des constantes\newline matricielle
    & \S\ref{sec:var-cst-systemes} \\
  $y'' + p(t)y' + q(t)y = 0$,\newline $y_1$ connue
    & Réduction d'ordre
    & \S\ref{sec:reduction-ordre} \\
  $y'' + p(t)y' + q(t)y = 0$,\newline $p, q$ analytiques
    & Séries entières
    & \S\ref{sec:series-entieres} \\
  $y'' + p(t)y' + q(t)y = 0$,\newline point singulier régulier
    & Frobenius
    & \S\ref{subsec:frobenius} \\
  Problème non homogène,\newline cadre intégral
    & Fonction de Green
    & \S\ref{sec:green} \\
  \bottomrule
\end{tabular}
\end{center}

% ============================================================
\section{Exemples supplémentaires}\label{sec:exemples-supp}
% ============================================================

\begin{exemple}[Variation des constantes — ordre~3]\label{ex:var-cst-ordre3}
  Résolvons $y''' - y' = e^t$.

  \emph{Équation caractéristique~:} $r^3 - r = r(r-1)(r+1) = 0$,
  racines $r = 0, 1, -1$.

  \emph{Système fondamental~:} $y_1 = 1$, $y_2 = e^t$, $y_3 = e^{-t}$.

  \emph{Wronskien~:}
  \[
    W = \begin{vmatrix}
      1 & e^t & e^{-t} \\
      0 & e^t & -e^{-t} \\
      0 & e^t & e^{-t}
    \end{vmatrix}
    = 1 \cdot (e^t \cdot e^{-t} - (-e^{-t})\cdot e^t) = 2.
  \]

  \emph{Calcul de $W_k$~:}
  \begin{align*}
    W_1 &= \begin{vmatrix}
      0 & e^t & e^{-t} \\
      0 & e^t & -e^{-t} \\
      1 & e^t & e^{-t}
    \end{vmatrix} = 1 \cdot (e^t \cdot (-e^{-t}) - e^{-t}\cdot e^t) = -2, \\[3pt]
    W_2 &= \begin{vmatrix}
      1 & 0 & e^{-t} \\
      0 & 0 & -e^{-t} \\
      0 & 1 & e^{-t}
    \end{vmatrix} = 1\cdot(0\cdot e^{-t} - (-e^{-t})\cdot 1) = e^{-t}, \\[3pt]
    W_3 &= \begin{vmatrix}
      1 & e^t & 0 \\
      0 & e^t & 0 \\
      0 & e^t & 1
    \end{vmatrix} = 1\cdot(e^t \cdot 1 - 0) = e^t.
  \end{align*}

  Les $c_k'$~:
  \[
    c_1' = \frac{-2\cdot e^t}{2} = -e^t, \quad
    c_2' = \frac{e^{-t}\cdot e^t}{2} = \frac{1}{2}, \quad
    c_3' = \frac{e^t \cdot e^t}{2} = \frac{e^{2t}}{2}.
  \]

  Intégration~: $c_1 = -e^t$, $c_2 = t/2$, $c_3 = e^{2t}/4$.

  Solution particulière~:
  \[
    y_p = 1\cdot(-e^t) + e^t\cdot \frac{t}{2} + e^{-t}\cdot\frac{e^{2t}}{4}
    = -e^t + \frac{t}{2}e^t + \frac{e^t}{4}
    = e^t\!\left(\frac{t}{2} - \frac{3}{4}\right).
  \]

  \emph{Solution générale~:}
  $y(t) = C_1 + C_2 e^t + C_3 e^{-t} + e^t\!\left(\frac{t}{2} - \frac{3}{4}\right)$.
\end{exemple}

\begin{exemple}[Équation d'Euler]\label{ex:euler}
  \index{équation d'Euler}
  L'équation $t^2 y'' + \alpha t y' + \beta y = 0$ ($t > 0$) est dite
  \emph{d'Euler} (ou de Cauchy--Euler). Le changement de variable
  $t = e^s$ la transforme en une EDO à coefficients constants~:
  \[
    \frac{\dd^2 y}{\dd s^2} + (\alpha - 1)\frac{\dd y}{\dd s} + \beta y = 0.
  \]
  L'équation indicielle est $r(r-1) + \alpha r + \beta = 0$,
  soit $r^2 + (\alpha-1)r + \beta = 0$.

  \emph{Exemple concret~:} $t^2 y'' + 3ty' + y = 0$. Ici $\alpha=3$,
  $\beta=1$, donc $r^2+2r+1=(r+1)^2=0$, racine double $r=-1$.
  Solution~: $y(t) = (C_1 + C_2 \ln t)/t$.
\end{exemple}

\begin{exemple}[Série entière~: équation de Legendre]\label{ex:legendre}
  \index{Legendre!équation de}\index{Legendre!polynômes de}
  L'équation de Legendre $(1-t^2)y'' - 2ty' + n(n+1)y = 0$ ($n \in \N$)
  a $t_0 = 0$ comme point ordinaire. La récurrence est~:
  \[
    a_{k+2} = \frac{k(k+1) - n(n+1)}{(k+2)(k+1)}\,a_k.
  \]
  Pour $k = n$, le numérateur s'annule, et la série s'interrompt~:
  on obtient le \emph{polynôme de Legendre} $P_n(t)$.

  Par exemple, pour $n = 2$~: $a_4 = 0$ et, en normalisant $P_2(1)=1$~:
  $P_2(t) = \frac{1}{2}(3t^2 - 1)$.
\end{exemple}

% ============================================================
\section{Exercices}\label{sec:exo-ch6}
% ============================================================

\begin{exercice}[$\bigstar$ Wronskien]\label{exo:wronskien-calcul}
  Calculer le wronskien de~:
  \begin{enumerate}[label=\alph*)]
    \item $y_1 = e^t$, $y_2 = e^{2t}$, $y_3 = e^{3t}$.
    \item $y_1 = t$, $y_2 = t^2$, $y_3 = t^3$ sur $]0,+\infty[$.
  \end{enumerate}
  Vérifier la formule d'Abel dans chaque cas.
\end{exercice}

\begin{exercice}[$\bigstar$ Variation des constantes]\label{exo:var-cst-1}
  Résoudre par variation des constantes~:
  \begin{enumerate}[label=\alph*)]
    \item $y'' + 4y = \tan(2t)$ sur $]-\pi/4, \pi/4[$.
    \item $y'' - 3y' + 2y = \dfrac{e^{3t}}{1+e^t}$.
  \end{enumerate}
\end{exercice}

\begin{exercice}[$\bigstar$ Réduction d'ordre]\label{exo:reduction}
  L'équation $t^2 y'' - 2ty' + 2y = 0$ admet $y_1(t) = t$ comme
  solution. Trouver la solution générale par réduction d'ordre.
\end{exercice}

\begin{exercice}[$\bigstar$ Séries entières]\label{exo:series-1}
  Résoudre en séries entières autour de $t_0 = 0$~:
  \begin{enumerate}[label=\alph*)]
    \item $y'' + ty = 0$ (équation d'Airy).
    \item $y'' + t^2 y = 0$.
  \end{enumerate}
  Déterminer le rayon de convergence des solutions.
\end{exercice}

\begin{exercice}[$\bigstar\bigstar$ Formule d'Abel et reconstruction]\label{exo:abel-reconstr}
  Soit $y'' + p(t)y' + q(t)y = 0$ sur $]0,+\infty[$ avec
  $y_1(t) = \sin t / t$.
  \begin{enumerate}[label=\alph*)]
    \item Sachant que $p(t) = 2/t$, utiliser la formule~\eqref{eq:2e-solution-abel}
      pour trouver une seconde solution~$y_2$.
    \item En déduire que $y_2(t) = -\cos t / t$ (à une constante
      multiplicative près).
  \end{enumerate}
\end{exercice}

\begin{exercice}[$\bigstar\bigstar$ Fonction de Green]\label{exo:green}
  \begin{enumerate}[label=\alph*)]
    \item Calculer la fonction de Green $G(t,s)$ de l'équation
      $y'' + \omega^2 y = f(t)$ pour $\omega > 0$.
    \item En déduire la solution du problème
      $y'' + \omega^2 y = \cos(\omega t)$, $y(0) = y'(0) = 0$.
  \end{enumerate}
\end{exercice}

\begin{exercice}[$\bigstar\bigstar$ Méthode de Frobenius]\label{exo:frobenius}
  Appliquer la méthode de Frobenius à l'équation de Bessel d'ordre
  $\nu = 0$~: $ty'' + y' + ty = 0$.
  \begin{enumerate}[label=\alph*)]
    \item Écrire l'équation indicielle et trouver ses racines.
    \item Déterminer les premiers termes de la solution $J_0(t)$.
    \item Montrer que $J_0(t) = \sum_{m=0}^\infty \frac{(-1)^m}{(m!)^2}
      \left(\frac{t}{2}\right)^{2m}$.
  \end{enumerate}
\end{exercice}

\begin{exercice}[$\bigstar\bigstar$ Identité du wronskien]\label{exo:identite-wronsk}
  Soient $y_1, y_2$ deux solutions de $y'' + q(t)y = 0$ ($p \equiv 0$).
  \begin{enumerate}[label=\alph*)]
    \item Montrer que $W[y_1,y_2]$ est constant.
    \item En déduire que si $y_1$ a deux zéros consécutifs
      $\alpha < \beta$, alors $y_2$ a au moins un zéro dans
      $]\alpha, \beta[$. \emph{(Théorème de séparation de Sturm.)}
      \index{Sturm!théorème de séparation}
  \end{enumerate}
\end{exercice}

\begin{exercice}[$\bigstar\bigstar\bigstar$ Équation de Tchebychev]\label{exo:tchebychev}
  \index{Tchebychev!polynômes de}
  L'équation de Tchebychev est $(1-t^2)y'' - ty' + n^2 y = 0$,
  $n \in \N$.
  \begin{enumerate}[label=\alph*)]
    \item Montrer que $t_0 = 0$ est un point ordinaire.
    \item Trouver la récurrence sur les coefficients de la série entière.
    \item Montrer que pour $n \in \N$, une des deux séries s'interrompt
      et donne le polynôme de Tchebychev~$T_n(t)$.
    \item Vérifier que $T_n(\cos\theta) = \cos(n\theta)$.
  \end{enumerate}
\end{exercice}

\begin{exercice}[$\bigstar\bigstar\bigstar$ Opérateur de résolution]\label{exo:operateur-resol}
  Soit $L[y] = y'' + p(t)y' + q(t)y$ un opérateur différentiel linéaire
  d'ordre~2, et $G(t,s)$ sa fonction de Green.
  \begin{enumerate}[label=\alph*)]
    \item Montrer que $y_p(t) = \int_{t_0}^t G(t,s)f(s)\,\dd s$
      vérifie $L[y_p] = f$ et $y_p(t_0) = y_p'(t_0) = 0$.
    \item En déduire que l'opérateur $T\colon f \mapsto y_p$ est un
      opérateur intégral borné de $L^2(I)$ dans $\mathcal{C}^1(I)$
      (sous des hypothèses de régularité convenables).
    \item Montrer que si $q(t) > 0$ et $p \equiv 0$, la fonction de
      Green vérifie $G(t,s) \geq 0$ pour $t \geq s$.
  \end{enumerate}
\end{exercice}

% ============================================================
\section*{Résumé du chapitre~6}
\addcontentsline{toc}{section}{Résumé du chapitre~6}
% ============================================================

\begin{center}
\fbox{\parbox{0.92\textwidth}{
\textbf{Points clés du chapitre~6~:}
\begin{itemize}
  \item Le \textbf{wronskien} $W[y_1,\ldots,y_n]$ est non nul si et
    seulement si les $y_i$ sont linéairement indépendantes (pour des
    solutions d'une même EDO linéaire). Le \textbf{théorème d'Abel}
    donne $W(t) = W(t_0)\exp\bigl(-\int_{t_0}^t a_{n-1}\bigr)$.
  \item La \textbf{variation des constantes} fournit une solution
    particulière via le système de Cramer~:
    $c_k'(t) = W_k(t)\,f(t)/W(t)$.
  \item La \textbf{fonction de Green} reformule la solution
    particulière comme opérateur intégral~:
    $y_p(t) = \int_{t_0}^t G(t,s)\,f(s)\,\dd s$.
  \item Aux \textbf{points ordinaires}, on cherche des solutions en
    séries entières $\sum a_n(t-t_0)^n$, qui convergent au moins
    aussi loin que les coefficients.
  \item Aux \textbf{points singuliers réguliers}, la
    \textbf{méthode de Frobenius} cherche $y = (t-t_0)^r \sum a_n(t-t_0)^n$,
    où $r$ est racine de l'équation indicielle.
  \item La \textbf{réduction d'ordre}~: connaissant $y_1$, on obtient
    $y_2 = y_1\int\frac{e^{-\int p}}{y_1^2}$.
  \item Consulter le \textbf{tableau récapitulatif} (\S\ref{sec:tableau-methodes})
    pour choisir la méthode adaptée à chaque type d'équation.
\end{itemize}
}}
\end{center}

% === FIN DU CHUNK ch5-6 ===
% === DÉBUT DU CHUNK ch7-8 ===

%%%%%%%%%%%%%%%%%%%%%%%%%%%%%%%%%%%%%%%%%%%%%%%%%%%%%%%%%%%%%%%%%%%%%
%  Chapitre 7 — Transformées de Laplace et Applications
%%%%%%%%%%%%%%%%%%%%%%%%%%%%%%%%%%%%%%%%%%%%%%%%%%%%%%%%%%%%%%%%%%%%%

\chapter{Transformées de Laplace et Applications}\label{ch:laplace}

\section*{Introduction}
\addcontentsline{toc}{section}{Introduction}

Considérons le circuit \emph{RLC}\index{circuit RLC} série suivant~:
une résistance $R$, une inductance $L$ et un condensateur $C$ sont
alimentés par une source de tension $E(t)$ qui vaut $E_0$ pour
$0\leq t < T$ et $0$ pour $t\geq T$. La loi de Kirchhoff donne
\[
  L\,q''(t) + R\,q'(t) + \frac{1}{C}\,q(t) = E(t),
  \qquad q(0)=0,\quad q'(0)=0,
\]
où $q(t)$ est la charge du condensateur. La discontinuité de $E$ en
$t=T$ empêche toute approche directe par les méthodes classiques des
chapitres précédents. La \emph{transformée de Laplace}\index{transformée de Laplace}
transforme cette équation intégro-différentielle en une équation
\emph{algébrique} en $s$, que l'on résout aisément ; on revient
ensuite à $q(t)$ par transformation inverse. C'est cette stratégie
puissante que nous développons dans ce chapitre.

\medskip

\noindent\textbf{Conventions.} Dans tout ce chapitre, les fonctions
considérées sont définies sur $[0,+\infty[$ et à valeurs dans $\R$
(ou $\C$). On note $s$ la variable complexe de la transformée.

% ============================================================
\section{Définition et existence}\label{sec:laplace-def}
% ============================================================

\begin{definition}[Transformée de Laplace]\label{def:laplace}
  \index{transformée de Laplace!définition}
  Soit $f:[0,+\infty[\to\R$ une fonction. La \emph{transformée de
  Laplace} de $f$ est la fonction $F$ définie par
  \[
    F(s) = \mathcal{L}\{f\}(s)
    \;=\; \int_{0}^{+\infty} e^{-st}\,f(t)\,\dd t,
  \]
  pour tout $s\in\C$ tel que l'intégrale converge.
\end{definition}

\begin{definition}[Fonction d'ordre exponentiel]\label{def:ordre-exp}
  \index{ordre exponentiel}
  Une fonction $f:[0,+\infty[\to\R$ est dite \emph{d'ordre exponentiel}
  $\alpha\in\R$ s'il existe des constantes $M>0$ et $T\geq 0$ telles
  que
  \[
    \abs{f(t)} \leq M\,e^{\alpha t}, \qquad \forall\, t\geq T.
  \]
\end{definition}

\begin{definition}[Continuité par morceaux]\label{def:cont-morceaux}
  \index{continuité par morceaux}
  Une fonction $f:[0,+\infty[\to\R$ est \emph{continue par morceaux}
  sur tout intervalle borné $[0,A]$ si elle possède au plus un nombre
  fini de discontinuités sur $[0,A]$, et si les limites à gauche et à
  droite existent en chaque point de discontinuité.
\end{definition}

\begin{theoreme}[Existence de la transformée de Laplace]\label{thm:existence-laplace}
  \index{transformée de Laplace!existence}
  Si $f:[0,+\infty[\to\R$ est continue par morceaux sur tout
  intervalle borné et d'ordre exponentiel~$\alpha$, alors
  $\mathcal{L}\{f\}(s)$ existe pour tout $s>\alpha$ et
  \[
    \abs{\mathcal{L}\{f\}(s)} \leq \frac{M}{s-\alpha},
    \qquad s>\alpha.
  \]
  En particulier, $\lim_{s\to+\infty}\mathcal{L}\{f\}(s)=0$.
\end{theoreme}

\begin{proof}
  Pour $t\geq T$ et $s>\alpha$, on a
  $\abs{e^{-st}f(t)}\leq M\,e^{-(s-\alpha)t}$.
  La fonction $t\mapsto M\,e^{-(s-\alpha)t}$ est intégrable sur
  $[T,+\infty[$ car $s-\alpha>0$. Par ailleurs, $f$ est bornée sur
  $[0,T]$ car continue par morceaux sur un compact, donc
  $\int_0^{T}\abs{e^{-st}f(t)}\,\dd t<+\infty$. Par comparaison,
  l'intégrale $\int_0^{+\infty}e^{-st}f(t)\,\dd t$ converge
  absolument. De plus,
  \[
    \abs{\mathcal{L}\{f\}(s)}
    \leq \int_0^{T}\abs{f(t)}\,e^{-st}\,\dd t
         + M\int_{T}^{+\infty}e^{-(s-\alpha)t}\,\dd t
    \leq C_T + \frac{M\,e^{-(s-\alpha)T}}{s-\alpha}
    \leq \frac{M'}{s-\alpha}
  \]
  pour une constante $M'$ appropriée. La limite en $+\infty$ en
  découle immédiatement.
\end{proof}

\begin{remarque}[Unicité]\label{rem:unicite-laplace}
  \index{transformée de Laplace!unicité}
  Si $f$ et $g$ sont continues par morceaux et si
  $\mathcal{L}\{f\}(s)=\mathcal{L}\{g\}(s)$ pour $s$ assez grand,
  alors $f(t)=g(t)$ en tout point de continuité commun. C'est le
  \emph{théorème de Lerch}\index{Lerch (théorème de)}.
\end{remarque}

% ============================================================
\section{Table des transformées usuelles}\label{sec:table-laplace}
% ============================================================

Nous démontrons les transformées les plus importantes ; les autres
se déduisent des propriétés de la section~\ref{sec:proprietes-laplace}.

\begin{proposition}[Transformée de la constante]\label{prop:laplace-1}
  \index{transformée de Laplace!de $1$}
  Pour $s>0$,
  $\displaystyle\mathcal{L}\{1\}(s)=\frac{1}{s}$.
\end{proposition}

\begin{proof}
  $\displaystyle\mathcal{L}\{1\}(s)
  =\int_0^{+\infty}e^{-st}\,\dd t
  =\left[-\frac{e^{-st}}{s}\right]_0^{+\infty}
  =\frac{1}{s}$.
\end{proof}

\begin{proposition}[Transformée de $t^n$]\label{prop:laplace-tn}
  \index{transformée de Laplace!de $t^n$}
  Pour $n\in\N$ et $s>0$,
  $\displaystyle\mathcal{L}\{t^n\}(s)=\frac{n!}{s^{n+1}}$.
\end{proposition}

\begin{proof}
  Par récurrence. Le cas $n=0$ est la proposition~\ref{prop:laplace-1}.
  En intégrant par parties,
  \[
    \mathcal{L}\{t^n\}(s)
    = \int_0^{+\infty}t^n\,e^{-st}\,\dd t
    = \left[-\frac{t^n\,e^{-st}}{s}\right]_0^{+\infty}
      + \frac{n}{s}\int_0^{+\infty}t^{n-1}\,e^{-st}\,\dd t
    = \frac{n}{s}\,\mathcal{L}\{t^{n-1}\}(s).
  \]
  L'hypothèse de récurrence donne
  $\mathcal{L}\{t^n\}(s)=\frac{n}{s}\cdot\frac{(n-1)!}{s^n}=\frac{n!}{s^{n+1}}$.
\end{proof}

\begin{proposition}[Transformée de $e^{at}$]\label{prop:laplace-eat}
  \index{transformée de Laplace!de $e^{at}$}
  Pour $a\in\R$ et $s>a$,
  $\displaystyle\mathcal{L}\{e^{at}\}(s)=\frac{1}{s-a}$.
\end{proposition}

\begin{proof}
  $\displaystyle\mathcal{L}\{e^{at}\}(s)
  =\int_0^{+\infty}e^{-(s-a)t}\,\dd t
  =\frac{1}{s-a}$ pour $s>a$.
\end{proof}

\begin{proposition}[Transformées de $\sin(bt)$ et $\cos(bt)$]\label{prop:laplace-sincos}
  \index{transformée de Laplace!de $\sin(bt)$}
  \index{transformée de Laplace!de $\cos(bt)$}
  Pour $b\in\R$ et $s>0$,
  \[
    \mathcal{L}\{\sin(bt)\}(s)=\frac{b}{s^2+b^2},
    \qquad
    \mathcal{L}\{\cos(bt)\}(s)=\frac{s}{s^2+b^2}.
  \]
\end{proposition}

\begin{proof}
  On utilise la formule d'Euler $e^{ibt}=\cos(bt)+i\sin(bt)$ et la
  linéarité. On a $\mathcal{L}\{e^{ibt}\}(s)=\frac{1}{s-ib}$ pour
  $s>0$. En rationalisant~:
  \[
    \frac{1}{s-ib}=\frac{s+ib}{s^2+b^2}
    = \frac{s}{s^2+b^2} + i\,\frac{b}{s^2+b^2}.
  \]
  En identifiant parties réelle et imaginaire, on obtient le résultat.
\end{proof}

\begin{proposition}[Transformée de $t^n e^{at}$]\label{prop:laplace-tneat}
  \index{transformée de Laplace!de $t^n e^{at}$}
  Pour $n\in\N$, $a\in\R$ et $s>a$,
  $\displaystyle\mathcal{L}\{t^n e^{at}\}(s)=\frac{n!}{(s-a)^{n+1}}$.
\end{proposition}

\begin{proof}
  Ceci résulte directement de la propriété de translation en $s$
  (théorème~\ref{thm:s-shift}) appliquée à $\mathcal{L}\{t^n\}$.
\end{proof}

% ============================================================
\section{Propriétés de la transformée de Laplace}\label{sec:proprietes-laplace}
% ============================================================

\begin{theoreme}[Linéarité]\label{thm:linearite-laplace}
  \index{transformée de Laplace!linéarité}
  Soient $f,g$ deux fonctions admettant des transformées de Laplace
  et $\alpha,\beta\in\R$. Alors
  \[
    \mathcal{L}\{\alpha f+\beta g\}(s)
    = \alpha\,\mathcal{L}\{f\}(s)+\beta\,\mathcal{L}\{g\}(s).
  \]
\end{theoreme}

\begin{proof}
  Immédiat par linéarité de l'intégrale.
\end{proof}

\begin{theoreme}[Translation en $s$ ($s$-shift)]\label{thm:s-shift}
  \index{transformée de Laplace!translation en $s$}
  \index{$s$-shift}
  Si $\mathcal{L}\{f\}(s)=F(s)$, alors
  \[
    \mathcal{L}\{e^{at}f(t)\}(s) = F(s-a), \qquad s>s_0+a.
  \]
\end{theoreme}

\begin{proof}
  $\displaystyle
  \mathcal{L}\{e^{at}f(t)\}(s)
  = \int_0^{+\infty}e^{-st}\,e^{at}\,f(t)\,\dd t
  = \int_0^{+\infty}e^{-(s-a)t}\,f(t)\,\dd t
  = F(s-a)$.
\end{proof}

\begin{theoreme}[Translation en $t$ ($t$-shift)]\label{thm:t-shift}
  \index{transformée de Laplace!translation en $t$}
  \index{$t$-shift}
  Si $\mathcal{L}\{f\}(s)=F(s)$ et $a\geq 0$, alors
  \[
    \mathcal{L}\{u_a(t)\,f(t-a)\}(s) = e^{-as}\,F(s),
  \]
  où $u_a$ désigne la fonction de Heaviside\index{Heaviside (fonction de)}
  (définition~\ref{def:heaviside}).
\end{theoreme}

\begin{proof}
  On effectue le changement de variable $\tau=t-a$~:
  \[
    \mathcal{L}\{u_a(t)\,f(t-a)\}(s)
    = \int_a^{+\infty}e^{-st}\,f(t-a)\,\dd t
    = \int_0^{+\infty}e^{-s(\tau+a)}\,f(\tau)\,\dd\tau
    = e^{-as}\,F(s).
    \qedhere
  \]
\end{proof}

\begin{theoreme}[Transformée des dérivées]\label{thm:laplace-derivee}
  \index{transformée de Laplace!des dérivées}
  Si $f$ est continue, $f'$ est continue par morceaux, et les deux
  sont d'ordre exponentiel, alors
  \[
    \mathcal{L}\{f'\}(s) = s\,\mathcal{L}\{f\}(s) - f(0).
  \]
  Plus généralement, si $f,f',\ldots,f^{(n-1)}$ sont continues et
  $f^{(n)}$ est continue par morceaux, toutes d'ordre exponentiel,
  \[
    \mathcal{L}\{f^{(n)}\}(s)
    = s^n F(s) - s^{n-1}f(0) - s^{n-2}f'(0) - \cdots - f^{(n-1)}(0).
  \]
\end{theoreme}

\begin{proof}
  Pour $n=1$, on intègre par parties~:
  \begin{align*}
    \mathcal{L}\{f'\}(s)
    &= \int_0^{+\infty}e^{-st}\,f'(t)\,\dd t
    = \bigl[e^{-st}f(t)\bigr]_0^{+\infty}
      + s\int_0^{+\infty}e^{-st}f(t)\,\dd t \\
    &= -f(0) + s\,F(s).
  \end{align*}
  Le terme en $+\infty$ est nul car $f$ est d'ordre exponentiel et
  $s$ est suffisamment grand. Le cas général s'obtient par récurrence~:
  on pose $g=f'$ et on applique le résultat à $g^{(n-1)}$.
\end{proof}

\begin{corollaire}[Dérivée seconde]\label{cor:laplace-f2}
  $\mathcal{L}\{f''\}(s)=s^2F(s)-sf(0)-f'(0)$.
\end{corollaire}

\begin{theoreme}[Transformée d'une intégrale]\label{thm:laplace-integrale}
  \index{transformée de Laplace!d'une intégrale}
  Si $\mathcal{L}\{f\}(s)=F(s)$, alors
  \[
    \mathcal{L}\!\left\{\int_0^t f(\tau)\,\dd\tau\right\}(s)
    = \frac{F(s)}{s}.
  \]
\end{theoreme}

\begin{proof}
  Posons $g(t)=\int_0^t f(\tau)\,\dd\tau$. Alors $g'(t)=f(t)$ et
  $g(0)=0$. Par le théorème~\ref{thm:laplace-derivee},
  $\mathcal{L}\{g'\}(s)=s\,G(s)-g(0)=s\,G(s)$, d'où
  $G(s)=F(s)/s$.
\end{proof}

\begin{theoreme}[Dérivation de la transformée]\label{thm:deriv-F}
  \index{transformée de Laplace!dérivation en $s$}
  Si $\mathcal{L}\{f\}(s)=F(s)$, alors
  \[
    \mathcal{L}\{t\,f(t)\}(s) = -F'(s).
  \]
  Plus généralement, $\mathcal{L}\{t^n f(t)\}(s)=(-1)^n F^{(n)}(s)$.
\end{theoreme}

\begin{proof}
  On dérive sous le signe intégral (justifié par convergence dominée)~:
  \[
    F'(s) = \frac{\dd}{\dd s}\int_0^{+\infty}e^{-st}f(t)\,\dd t
    = \int_0^{+\infty}(-t)\,e^{-st}f(t)\,\dd t
    = -\mathcal{L}\{t\,f(t)\}(s).
    \qedhere
  \]
\end{proof}

\begin{theoreme}[Convolution]\label{thm:convolution}
  \index{convolution}\index{transformée de Laplace!convolution}
  Soient $f$ et $g$ deux fonctions admettant des transformées de
  Laplace. Le \emph{produit de convolution} est défini par
  \[
    (f*g)(t) = \int_0^t f(\tau)\,g(t-\tau)\,\dd\tau.
  \]
  Alors $\mathcal{L}\{f*g\}(s)=F(s)\cdot G(s)$.
\end{theoreme}

\begin{proof}
  On calcule~:
  \begin{align*}
    \mathcal{L}\{f*g\}(s)
    &= \int_0^{+\infty}e^{-st}\left(\int_0^t f(\tau)\,g(t-\tau)\,\dd\tau\right)\dd t.
  \end{align*}
  On échange l'ordre d'intégration (Fubini)~: le domaine est
  $\{(\tau,t):0\leq\tau\leq t<+\infty\}=\{(\tau,t):0\leq\tau<+\infty,\;t\geq\tau\}$.
  Ainsi,
  \begin{align*}
    \mathcal{L}\{f*g\}(s)
    &= \int_0^{+\infty}f(\tau)\left(\int_\tau^{+\infty}e^{-st}\,g(t-\tau)\,\dd t\right)\dd\tau.
  \end{align*}
  Dans l'intégrale intérieure, on pose $u=t-\tau$~:
  \[
    \int_\tau^{+\infty}e^{-st}\,g(t-\tau)\,\dd t
    = \int_0^{+\infty}e^{-s(u+\tau)}\,g(u)\,\dd u
    = e^{-s\tau}\,G(s).
  \]
  Par conséquent,
  $\mathcal{L}\{f*g\}(s)
  = \int_0^{+\infty}f(\tau)\,e^{-s\tau}\,\dd\tau\cdot G(s)
  = F(s)\,G(s)$.
\end{proof}

\begin{theoreme}[Théorème de la valeur initiale]\label{thm:val-initiale}
  \index{théorème de la valeur initiale}
  Si $f$ et $f'$ sont d'ordre exponentiel et si $f$ est continue en $0$, alors
  \[
    f(0^+) = \lim_{s\to+\infty} s\,F(s).
  \]
\end{theoreme}

\begin{proof}
  D'après le théorème~\ref{thm:laplace-derivee},
  $\mathcal{L}\{f'\}(s)=sF(s)-f(0)$. Quand $s\to+\infty$,
  $\mathcal{L}\{f'\}(s)\to 0$ (théorème~\ref{thm:existence-laplace}),
  d'où $\lim_{s\to+\infty}sF(s)=f(0)$.
\end{proof}

\begin{theoreme}[Théorème de la valeur finale]\label{thm:val-finale}
  \index{théorème de la valeur finale}
  Si $\lim_{t\to+\infty}f(t)$ existe (et est finie), alors
  \[
    \lim_{t\to+\infty}f(t) = \lim_{s\to 0^+} s\,F(s).
  \]
\end{theoreme}

\begin{proof}
  On a $\mathcal{L}\{f'\}(s)=sF(s)-f(0)$. En faisant $s\to 0^+$~:
  \[
    \lim_{s\to 0^+}\bigl(sF(s)-f(0)\bigr)
    = \lim_{s\to 0^+}\int_0^{+\infty}e^{-st}f'(t)\,\dd t
    = \int_0^{+\infty}f'(t)\,\dd t
    = \lim_{T\to+\infty}f(T)-f(0),
  \]
  d'où le résultat.
\end{proof}

% ============================================================
\section{Transformée inverse et fractions partielles}\label{sec:laplace-inverse}
% ============================================================

\begin{definition}[Transformée inverse]\label{def:laplace-inv}
  \index{transformée de Laplace!inverse}
  Si $F(s)=\mathcal{L}\{f\}(s)$, on écrit
  $f(t)=\mathcal{L}^{-1}\{F\}(t)$. Grâce à l'unicité
  (remarque~\ref{rem:unicite-laplace}), cette inverse est
  essentiellement unique.
\end{definition}

\begin{remarque}[Linéarité de l'inverse]\label{rem:inv-lin}
  $\mathcal{L}^{-1}\{\alpha F+\beta G\}=\alpha\mathcal{L}^{-1}\{F\}
  +\beta\mathcal{L}^{-1}\{G\}$.
\end{remarque}

La méthode des \emph{fractions partielles}\index{fractions partielles}
est l'outil principal pour calculer des transformées inverses. Soit
$F(s)=P(s)/Q(s)$ une fraction rationnelle avec $\deg P<\deg Q$.

\begin{methode}[Décomposition en fractions partielles]\label{meth:frac-part}
  \index{fractions partielles!méthode}
  \begin{enumerate}[label=\textbf{(\arabic*)}]
    \item Factoriser $Q(s)$ en produit de facteurs linéaires et
      quadratiques irréductibles.
    \item Pour chaque facteur linéaire simple $(s-a)$~: terme
      $\dfrac{A}{s-a}$.
    \item Pour chaque facteur linéaire répété $(s-a)^k$~: termes
      $\dfrac{A_1}{s-a}+\dfrac{A_2}{(s-a)^2}+\cdots+\dfrac{A_k}{(s-a)^k}$.
    \item Pour chaque facteur quadratique $(s^2+ps+q)$~: terme
      $\dfrac{As+B}{s^2+ps+q}$.
    \item Identifier les coefficients et appliquer les transformées
      inverses connues.
  \end{enumerate}
\end{methode}

\begin{exemple}[Fractions partielles]\label{ex:frac-part-1}
  Calculons $\mathcal{L}^{-1}\!\left\{\dfrac{5s+3}{(s-1)(s+2)}\right\}$.

  On écrit $\dfrac{5s+3}{(s-1)(s+2)}=\dfrac{A}{s-1}+\dfrac{B}{s+2}$.
  En multipliant par $(s-1)(s+2)$~: $5s+3=A(s+2)+B(s-1)$.

  Pour $s=1$~: $8=3A$, soit $A=8/3$.
  Pour $s=-2$~: $-7=-3B$, soit $B=7/3$.

  Donc
  $\mathcal{L}^{-1}\!\left\{\dfrac{5s+3}{(s-1)(s+2)}\right\}
  = \dfrac{8}{3}\,e^{t}+\dfrac{7}{3}\,e^{-2t}$.
\end{exemple}

% ============================================================
\section{Fonctions de Heaviside et distribution de Dirac}\label{sec:heaviside-dirac}
% ============================================================

\begin{definition}[Fonction échelon de Heaviside]\label{def:heaviside}
  \index{Heaviside (fonction de)}\index{fonction échelon}
  Pour $a\geq 0$, la \emph{fonction de Heaviside} (ou fonction
  échelon unité) est
  \[
    u_a(t) = u(t-a) =
    \begin{cases}
      0 & \text{si } t<a, \\
      1 & \text{si } t\geq a.
    \end{cases}
  \]
  On note parfois $\theta(t-a)$ ou $H(t-a)$.
\end{definition}

\begin{proposition}[Transformée de Heaviside]\label{prop:laplace-heaviside}
  \index{transformée de Laplace!de Heaviside}
  Pour $a\geq 0$ et $s>0$,
  $\displaystyle\mathcal{L}\{u_a(t)\}(s)=\frac{e^{-as}}{s}$.
\end{proposition}

\begin{proof}
  $\displaystyle\mathcal{L}\{u_a(t)\}(s)
  =\int_a^{+\infty}e^{-st}\,\dd t
  =\frac{e^{-as}}{s}$.
\end{proof}

\begin{remarque}[Fenêtre rectangulaire]\label{rem:fenetre}
  La fonction porte $\Pi_{a,b}(t)=u_a(t)-u_b(t)$ ($0\leq a<b$)
  vaut $1$ sur $[a,b[$ et $0$ ailleurs. Sa transformée est
  $\frac{e^{-as}-e^{-bs}}{s}$.
\end{remarque}

\begin{definition}[Distribution de Dirac]\label{def:dirac}
  \index{Dirac (distribution de)}\index{distribution de Dirac}
  La \emph{distribution de Dirac} $\delta_a$ (ou $\delta(t-a)$) est
  définie comme la «~limite~» de fonctions concentrées en $t=a$~:
  formellement, elle vérifie
  \[
    \int_{-\infty}^{+\infty}\delta(t-a)\,\varphi(t)\,\dd t = \varphi(a)
  \]
  pour toute fonction test $\varphi$ continue. En particulier,
  $\delta_0=\delta$ vérifie $\int\delta(t)\,\dd t=1$.
\end{definition}

\begin{proposition}[Transformée de Dirac]\label{prop:laplace-dirac}
  \index{transformée de Laplace!de Dirac}
  Pour $a\geq 0$,
  $\mathcal{L}\{\delta(t-a)\}(s)=e^{-as}$.
  En particulier, $\mathcal{L}\{\delta(t)\}(s)=1$.
\end{proposition}

\begin{proof}
  Formellement,
  $\mathcal{L}\{\delta(t-a)\}(s)
  =\int_0^{+\infty}e^{-st}\,\delta(t-a)\,\dd t=e^{-as}$
  par la propriété caractéristique de~$\delta$.
\end{proof}

% ============================================================
\section{Résolution d'EDO par la transformée de Laplace}\label{sec:resolution-laplace}
% ============================================================

\begin{methode}[Résolution d'un PVI par Laplace]\label{meth:resolution-laplace}
  \index{transformée de Laplace!résolution d'EDO}
  \begin{enumerate}[label=\textbf{Étape \arabic*.}]
    \item Appliquer $\mathcal{L}$ aux deux membres de l'EDO.
    \item Utiliser le théorème~\ref{thm:laplace-derivee} pour exprimer
      $\mathcal{L}\{y'\}$, $\mathcal{L}\{y''\}$, etc., en fonction de
      $Y(s)=\mathcal{L}\{y\}(s)$ et des conditions initiales.
    \item Résoudre l'équation algébrique en $Y(s)$.
    \item Calculer $y(t)=\mathcal{L}^{-1}\{Y\}(t)$ par fractions
      partielles.
  \end{enumerate}
\end{methode}

\begin{exemple}[EDO du second ordre à coefficients constants]\label{ex:laplace-edo2}
  Résolvons le problème de Cauchy
  \[
    y''-3y'+2y=0, \qquad y(0)=1,\quad y'(0)=0.
  \]

  \textbf{Étape 1--2.} On applique $\mathcal{L}$~:
  \[
    s^2Y-sy(0)-y'(0)-3\bigl(sY-y(0)\bigr)+2Y=0,
  \]
  soit $(s^2-3s+2)Y=s-3$.

  \textbf{Étape 3.}
  $\displaystyle Y(s)=\frac{s-3}{s^2-3s+2}=\frac{s-3}{(s-1)(s-2)}$.

  \textbf{Étape 4.} Décomposition~:
  $\frac{s-3}{(s-1)(s-2)}=\frac{A}{s-1}+\frac{B}{s-2}$.

  $s=1$~: $-2=-A$, soit $A=2$. \quad $s=2$~: $-1=B$, soit $B=-1$.
  \[
    y(t)=2e^{t}-e^{2t}.
  \]
\end{exemple}

\begin{exemple}[Forçage discontinu]\label{ex:laplace-discontinu}
  \index{forçage discontinu}
  Considérons
  \[
    y''+y = g(t), \quad y(0)=0,\; y'(0)=0,
    \quad\text{où}\quad
    g(t)=\begin{cases}1&\text{si } 0\leq t<\pi,\\0&\text{si } t\geq\pi.\end{cases}
  \]
  On écrit $g(t)=1-u_\pi(t)$, de sorte que
  $G(s)=\frac{1}{s}-\frac{e^{-\pi s}}{s}$.

  En appliquant $\mathcal{L}$~:
  $(s^2+1)Y(s)=\frac{1-e^{-\pi s}}{s}$, d'où
  \[
    Y(s)=\frac{1}{s(s^2+1)}-\frac{e^{-\pi s}}{s(s^2+1)}.
  \]
  Or $\frac{1}{s(s^2+1)}=\frac{1}{s}-\frac{s}{s^2+1}$, dont
  l'inverse est $1-\cos t$. Par le théorème~\ref{thm:t-shift},
  \[
    y(t) = (1-\cos t) - u_\pi(t)\,\bigl(1-\cos(t-\pi)\bigr)
         = (1-\cos t) - u_\pi(t)\,(1+\cos t).
  \]
  Pour $t\geq\pi$, on obtient $y(t)=-2\cos t$.
\end{exemple}

\begin{exemple}[Système d'EDO]\label{ex:laplace-systeme}
  \index{transformée de Laplace!système}
  Résolvons
  \[
    \begin{cases}
      x'=2x+y, \\
      y'=3x,
    \end{cases}
    \qquad x(0)=1,\quad y(0)=0.
  \]
  En appliquant $\mathcal{L}$ et notant $X=\mathcal{L}\{x\}$,
  $\mathscr{Y}=\mathcal{L}\{y\}$~:
  \[
    sX-1=2X+\mathscr{Y}, \qquad s\mathscr{Y}=3X.
  \]
  De la seconde~: $\mathscr{Y}=3X/s$. En substituant~:
  $sX-1=2X+3X/s$, soit $X(s^2-2s-3)=s$, d'où
  \[
    X(s) = \frac{s}{(s-3)(s+1)}
    = \frac{3/4}{s-3}+\frac{1/4}{s+1}.
  \]
  Ainsi $x(t)=\frac{3}{4}e^{3t}+\frac{1}{4}e^{-t}$.
  Puis $\mathscr{Y}=3X/s$ et par fractions partielles~:
  \[
    \mathscr{Y}(s) = \frac{3}{(s-3)(s+1)}
    = \frac{3/4}{s-3}-\frac{3/4}{s+1},
  \]
  d'où $y(t)=\frac{3}{4}e^{3t}-\frac{3}{4}e^{-t}$.
\end{exemple}

\begin{exemple}[Forçage périodique]\label{ex:laplace-periodique}
  \index{forçage périodique}
  Soit $f$ une fonction périodique de période $T>0$. Alors
  \[
    \mathcal{L}\{f\}(s) = \frac{1}{1-e^{-sT}}\int_0^{T}e^{-st}f(t)\,\dd t.
  \]
  En effet, en découpant l'intégrale sur les intervalles $[nT,(n+1)T]$
  et en effectuant le changement $t\mapsto t-nT$~:
  \[
    \mathcal{L}\{f\}(s)
    = \sum_{n=0}^{+\infty}e^{-nsT}\int_0^{T}e^{-st}f(t)\,\dd t
    = \frac{1}{1-e^{-sT}}\int_0^{T}e^{-st}f(t)\,\dd t.
  \]

  \textbf{Application.} Considérons l'onde carrée
  $f(t)=(-1)^{\lfloor t/\pi\rfloor}$ de période $2\pi$.
  Sur $[0,2\pi[$, $f(t)=1$ pour $t\in[0,\pi[$ et $f(t)=-1$
  pour $t\in[\pi,2\pi[$. Alors
  \[
    \int_0^{2\pi}e^{-st}f(t)\,\dd t
    = \frac{1-e^{-\pi s}}{s}-\frac{e^{-\pi s}-e^{-2\pi s}}{s}
    = \frac{(1-e^{-\pi s})^2}{s},
  \]
  d'où $\mathcal{L}\{f\}(s)
  =\frac{(1-e^{-\pi s})^2}{s(1-e^{-2\pi s})}
  =\frac{1-e^{-\pi s}}{s(1+e^{-\pi s})}
  =\frac{1}{s}\tanh\!\left(\frac{\pi s}{2}\right)$.
\end{exemple}

\begin{exemple}[Impulsion de Dirac]\label{ex:laplace-dirac}
  \index{Dirac (distribution de)!dans une EDO}
  Résolvons $y''+4y=\delta(t-\pi)$, $y(0)=0$, $y'(0)=0$.

  On applique $\mathcal{L}$~: $(s^2+4)Y(s)=e^{-\pi s}$, d'où
  $Y(s)=\frac{e^{-\pi s}}{s^2+4}$.

  Or $\mathcal{L}^{-1}\left\{\frac{1}{s^2+4}\right\}=\frac{1}{2}\sin(2t)$.
  Par le $t$-shift~:
  \[
    y(t)=\frac{1}{2}\,u_\pi(t)\,\sin\bigl(2(t-\pi)\bigr)
    = \frac{1}{2}\,u_\pi(t)\,\sin(2t).
  \]
\end{exemple}

% ============================================================
\section{Retour à l'exemple introductif}\label{sec:retour-circuit}
% ============================================================

Reprenons le circuit RLC de l'introduction. On a l'EDO
$Lq''+Rq'+\frac{1}{C}q = E_0\bigl(1-u_T(t)\bigr)$ avec
$q(0)=q'(0)=0$. Après application de $\mathcal{L}$~:
\[
  \left(Ls^2+Rs+\frac{1}{C}\right)Q(s)
  = E_0\,\frac{1-e^{-Ts}}{s}.
\]
Posons $Z(s)=Ls^2+Rs+1/C$ (l'\emph{impédance}\index{impédance}).
Alors $Q(s)=E_0\,\frac{1-e^{-Ts}}{s\,Z(s)}$. L'inversion par
fractions partielles dépend des racines de $Z(s)$~: cas
sur-amorti ($R^2>4L/C$), critique ($R^2=4L/C$), ou
sous-amorti ($R^2<4L/C$), qui donnent respectivement des
exponentielles réelles, un terme $te^{-\alpha t}$, ou des
oscillations amorties $e^{-\alpha t}\sin(\omega t+\varphi)$.

% ============================================================
\section{Table complète des transformées de Laplace}\label{sec:table-complete}
% ============================================================

\index{transformée de Laplace!table}
\begin{center}
\renewcommand{\arraystretch}{1.5}
\begin{tabular}{lll}
  \toprule
  $f(t)$ & $F(s)=\mathcal{L}\{f\}(s)$ & Domaine \\
  \midrule
  $1$                     & $\dfrac{1}{s}$                   & $s>0$ \\[4pt]
  $t^n$                   & $\dfrac{n!}{s^{n+1}}$            & $s>0$ \\[4pt]
  $e^{at}$                & $\dfrac{1}{s-a}$                 & $s>a$ \\[4pt]
  $t^n e^{at}$            & $\dfrac{n!}{(s-a)^{n+1}}$       & $s>a$ \\[4pt]
  $\sin(bt)$              & $\dfrac{b}{s^2+b^2}$            & $s>0$ \\[4pt]
  $\cos(bt)$              & $\dfrac{s}{s^2+b^2}$            & $s>0$ \\[4pt]
  $e^{at}\sin(bt)$        & $\dfrac{b}{(s-a)^2+b^2}$        & $s>a$ \\[4pt]
  $e^{at}\cos(bt)$        & $\dfrac{s-a}{(s-a)^2+b^2}$      & $s>a$ \\[4pt]
  $t\sin(bt)$             & $\dfrac{2bs}{(s^2+b^2)^2}$      & $s>0$ \\[4pt]
  $t\cos(bt)$             & $\dfrac{s^2-b^2}{(s^2+b^2)^2}$  & $s>0$ \\[4pt]
  $\sinh(bt)$             & $\dfrac{b}{s^2-b^2}$            & $s>\abs{b}$ \\[4pt]
  $\cosh(bt)$             & $\dfrac{s}{s^2-b^2}$            & $s>\abs{b}$ \\[4pt]
  $u_a(t)$                & $\dfrac{e^{-as}}{s}$            & $s>0$ \\[4pt]
  $\delta(t-a)$           & $e^{-as}$                       & tout $s$ \\[4pt]
  $u_a(t)\,f(t-a)$        & $e^{-as}F(s)$                   & \\[4pt]
  $e^{at}f(t)$            & $F(s-a)$                        & \\[4pt]
  $t^n f(t)$              & $(-1)^n F^{(n)}(s)$             & \\[4pt]
  $f'(t)$                 & $sF(s)-f(0)$                    & \\[4pt]
  $f''(t)$                & $s^2F(s)-sf(0)-f'(0)$           & \\[4pt]
  $(f*g)(t)$              & $F(s)\,G(s)$                    & \\
  \bottomrule
\end{tabular}
\end{center}

% ============================================================
\section{Exercices}\label{sec:exo-laplace}
% ============================================================

\begin{exercice}[Calculs directs]\label{exo:laplace-direct}
  Calculer les transformées de Laplace suivantes~:
  \begin{enumerate}[label=\alph*)]
    \item $\mathcal{L}\{3t^2-2e^{-t}+5\sin(3t)\}$.
    \item $\mathcal{L}\{t^2 e^{3t}\}$.
    \item $\mathcal{L}\{e^{-2t}\cos(4t)\}$.
    \item $\mathcal{L}\{t\,\sin(2t)\}$.
  \end{enumerate}
\end{exercice}

\begin{exercice}[Transformées inverses]\label{exo:laplace-inv}
  Calculer les transformées inverses~:
  \begin{enumerate}[label=\alph*)]
    \item $\mathcal{L}^{-1}\!\left\{\dfrac{3s+2}{s^2+4s+13}\right\}$.
    \item $\mathcal{L}^{-1}\!\left\{\dfrac{s}{(s+1)^3}\right\}$.
    \item $\mathcal{L}^{-1}\!\left\{\dfrac{2s+1}{(s-1)(s^2+1)}\right\}$.
    \item $\mathcal{L}^{-1}\!\left\{\dfrac{e^{-2s}}{s^2+4}\right\}$.
  \end{enumerate}
\end{exercice}

\begin{exercice}[Résolution de PVI]\label{exo:laplace-pvi}
  Résoudre par la transformée de Laplace~:
  \begin{enumerate}[label=\alph*)]
    \item $y''+4y'+4y=e^{-2t}$, $y(0)=0$, $y'(0)=1$.
    \item $y''-2y'+5y=0$, $y(0)=1$, $y'(0)=3$.
    \item $y''+y=u_2(t)\sin(t-2)$, $y(0)=0$, $y'(0)=0$.
  \end{enumerate}
\end{exercice}

\begin{exercice}[Convolution]\label{exo:convolution}
  \begin{enumerate}[label=\alph*)]
    \item Montrer que $\sin t * \sin t = \frac{1}{2}(\sin t - t\cos t)$.
    \item Utiliser la convolution pour résoudre
      $y''+y=f(t)$, $y(0)=y'(0)=0$.
  \end{enumerate}
\end{exercice}

\begin{exercice}[Système]\label{exo:laplace-syst}
  Résoudre par la transformée de Laplace~:
  \[
    \begin{cases}
      x'+y'=-x+2y, \\
      x'-y'=4x-y,
    \end{cases}
    \qquad x(0)=2,\; y(0)=-1.
  \]
\end{exercice}

\begin{exercice}[Forçage périodique]\label{exo:laplace-period}
  Soit $f(t)=t$ pour $t\in[0,1[$ prolongée par périodicité de
  période $1$ (onde en dents de scie).
  \begin{enumerate}[label=\alph*)]
    \item Calculer $\mathcal{L}\{f\}(s)$.
    \item Résoudre $y'+y=f(t)$, $y(0)=0$.
  \end{enumerate}
\end{exercice}

% ============================================================
\section*{Résumé du chapitre}
\addcontentsline{toc}{section}{Résumé du chapitre}
% ============================================================

\begin{itemize}
  \item La transformée de Laplace $\mathcal{L}\{f\}(s)=\int_0^{+\infty}
    e^{-st}f(t)\,\dd t$ convertit une EDO en équation algébrique.
  \item Les conditions d'existence reposent sur la continuité par
    morceaux et l'ordre exponentiel.
  \item Les propriétés clés sont~: linéarité, translation en $s$ et
    en $t$, transformée des dérivées ($\mathcal{L}\{f'\}=sF-f(0)$),
    convolution ($\mathcal{L}\{f*g\}=FG$).
  \item La fonction de Heaviside $u_a$ et la distribution de Dirac
    $\delta$ permettent de traiter les forçages discontinus et
    impulsionnels.
  \item La méthode de résolution~: appliquer $\mathcal{L}$, résoudre
    en $Y(s)$, inverser par fractions partielles.
  \item Les théorèmes de la valeur initiale et finale relient le
    comportement asymptotique de $f$ à celui de $sF(s)$.
\end{itemize}


%%%%%%%%%%%%%%%%%%%%%%%%%%%%%%%%%%%%%%%%%%%%%%%%%%%%%%%%%%%%%%%%%%%%%
%  Chapitre 8 — Stabilité des Solutions et Équilibres
%%%%%%%%%%%%%%%%%%%%%%%%%%%%%%%%%%%%%%%%%%%%%%%%%%%%%%%%%%%%%%%%%%%%%

\chapter{Stabilité des Solutions et Équilibres}\label{ch:stabilite}

\section*{Introduction}
\addcontentsline{toc}{section}{Introduction}

Considérons un pendule pesant de longueur $\ell$. L'angle $\theta$
vérifie $\theta''+\frac{g}{\ell}\sin\theta=0$. Cette équation admet
deux positions d'équilibre~: $\theta=0$ (pendule en bas) et
$\theta=\pi$ (pendule en haut). L'intuition physique dit que la
première est \emph{stable} et la seconde \emph{instable}. Mais
comment formaliser mathématiquement cette notion~? Et que se passe-t-il
pour des systèmes plus complexes où l'intuition fait défaut~?

\index{stabilité}\index{équilibre}

Ce chapitre développe la théorie de la stabilité au sens de
\emph{Lyapunov}\index{Lyapunov, A.\,M.}, qui fournit des outils
rigoureux pour analyser le comportement des solutions au voisinage
des points d'équilibre. Nous étudierons deux approches
complémentaires~: la \emph{méthode directe} de Lyapunov (construction
de fonctions d'énergie) et la \emph{linéarisation} (étude de la
matrice jacobienne).

\medskip

\noindent\textbf{Cadre.} On considère le système autonome
\begin{equation}\label{eq:syst-autonome}
  \mathbf{x}'=\mathbf{f}(\mathbf{x}),
  \qquad \mathbf{x}\in\R^n,
\end{equation}
où $\mathbf{f}:\Omega\to\R^n$ est de classe $\mathcal{C}^1$ sur un
ouvert $\Omega\subset\R^n$.

% ============================================================
\section{Points d'équilibre~: définition et classification}\label{sec:pts-equilibre}
% ============================================================

\begin{definition}[Point d'équilibre]\label{def:pt-equilibre}
  \index{point d'équilibre}\index{équilibre!point d'}
  Un point $\mathbf{x}_0\in\Omega$ est un \emph{point d'équilibre}
  (ou \emph{point stationnaire}, ou \emph{point critique}) du système
  \eqref{eq:syst-autonome} si $\mathbf{f}(\mathbf{x}_0)=\mathbf{0}$.
  La solution constante $\mathbf{x}(t)\equiv\mathbf{x}_0$ est alors
  une solution du système.
\end{definition}

\begin{remarque}[Translation]\label{rem:translation-equilibre}
  Par le changement de variable $\mathbf{y}=\mathbf{x}-\mathbf{x}_0$,
  on peut toujours supposer que le point d'équilibre est l'origine
  $\mathbf{0}$. On suppose cette normalisation dans la suite.
\end{remarque}

\begin{exemple}[Points d'équilibre du pendule]\label{ex:equil-pendule}
  \index{pendule!points d'équilibre}
  Le pendule non linéaire $\theta''+\frac{g}{\ell}\sin\theta=0$ se
  réécrit comme le système
  \[
    \begin{cases}
      x_1'=x_2, \\
      x_2'=-\frac{g}{\ell}\sin x_1,
    \end{cases}
  \]
  avec $x_1=\theta$, $x_2=\theta'$. Les points d'équilibre sont les
  $(k\pi,0)$ pour $k\in\Z$~: positions verticales haute et basse.
\end{exemple}

\begin{exemple}[Lotka-Volterra]\label{ex:equil-lotka}
  \index{Lotka-Volterra}
  Le système prédateur-proie
  \[
    \begin{cases}
      x'=\alpha x-\beta xy, \\
      y'=-\gamma y+\delta xy,
    \end{cases}
    \qquad \alpha,\beta,\gamma,\delta>0,
  \]
  admet deux équilibres~: $(0,0)$ (extinction) et
  $(\gamma/\delta,\alpha/\beta)$ (coexistence).
\end{exemple}

% ============================================================
\section{Définitions de stabilité}\label{sec:def-stabilite}
% ============================================================

\begin{definition}[Stabilité au sens de Lyapunov]\label{def:stab-lyapunov}
  \index{stabilité!au sens de Lyapunov}
  Le point d'équilibre $\mathbf{0}$ est \emph{stable au sens de
  Lyapunov} si, pour tout $\varepsilon>0$, il existe $\delta>0$ tel que
  \[
    \norm{\mathbf{x}(0)}<\delta
    \implies
    \norm{\mathbf{x}(t)}<\varepsilon,
    \qquad\forall\,t\geq 0.
  \]
\end{definition}

\begin{definition}[Stabilité asymptotique]\label{def:stab-asymptotique}
  \index{stabilité!asymptotique}
  Le point d'équilibre $\mathbf{0}$ est \emph{asymptotiquement stable}
  s'il est stable au sens de Lyapunov et s'il existe $\delta_0>0$ tel
  que
  \[
    \norm{\mathbf{x}(0)}<\delta_0
    \implies
    \lim_{t\to+\infty}\mathbf{x}(t)=\mathbf{0}.
  \]
\end{definition}

\begin{definition}[Instabilité]\label{def:instabilite}
  \index{instabilité}\index{stabilité!instabilité}
  Le point d'équilibre $\mathbf{0}$ est \emph{instable} s'il n'est
  pas stable au sens de Lyapunov~: il existe $\varepsilon_0>0$ tel que,
  pour tout $\delta>0$, il existe une solution $\mathbf{x}(t)$ avec
  $\norm{\mathbf{x}(0)}<\delta$ et un temps $t_0>0$ tel que
  $\norm{\mathbf{x}(t_0)}\geq\varepsilon_0$.
\end{definition}

\begin{definition}[Stabilité asymptotique globale]\label{def:stab-globale}
  \index{stabilité!asymptotique globale}
  Le point d'équilibre $\mathbf{0}$ est \emph{globalement
  asymptotiquement stable} s'il est stable et si
  $\lim_{t\to+\infty}\mathbf{x}(t)=\mathbf{0}$ pour toute condition
  initiale $\mathbf{x}(0)\in\R^n$.
\end{definition}

\begin{remarque}[Hiérarchie]\label{rem:hierarchie-stab}
  On a les implications~:
  \[
    \text{glob. asympt. stable}
    \implies
    \text{asympt. stable}
    \implies
    \text{stable (Lyapunov)}.
  \]
  Les réciproques sont fausses en général.
\end{remarque}

\begin{center}
\begin{tikzpicture}[>=Stealth, scale=0.9]
  % Stable
  \begin{scope}[xshift=-4.5cm]
    \draw[->] (-1.8,0)--(1.8,0) node[right]{$x_1$};
    \draw[->] (0,-1.8)--(0,1.8) node[above]{$x_2$};
    \draw[blue, thick] (0,0) circle (0.5);
    \draw[red, dashed] (0,0) circle (1.2);
    \draw[->, thick, blue!60!black, domain=0:5, samples=50]
      plot ({0.8*cos(80*\x)*exp(-0.01*\x)},{0.8*sin(80*\x)*exp(-0.01*\x)});
    \node[below] at (0,-2) {Stable};
    \node[blue] at (0.5,0.7) {\small$\delta$};
    \node[red] at (1.2,1.0) {\small$\varepsilon$};
  \end{scope}
  % Asymptotiquement stable
  \begin{scope}
    \draw[->] (-1.8,0)--(1.8,0) node[right]{$x_1$};
    \draw[->] (0,-1.8)--(0,1.8) node[above]{$x_2$};
    \draw[->, thick, blue!60!black, domain=0:8, samples=50]
      plot ({1.2*cos(60*\x)*exp(-0.15*\x)},{1.2*sin(60*\x)*exp(-0.15*\x)});
    \fill (0,0) circle (2pt);
    \node[below] at (0,-2) {Asympt. stable};
  \end{scope}
  % Instable
  \begin{scope}[xshift=4.5cm]
    \draw[->] (-1.8,0)--(1.8,0) node[right]{$x_1$};
    \draw[->] (0,-1.8)--(0,1.8) node[above]{$x_2$};
    \draw[->, thick, red!70!black, domain=0:4, samples=50]
      plot ({0.2*cos(70*\x)*exp(0.15*\x)},{0.2*sin(70*\x)*exp(0.15*\x)});
    \fill (0,0) circle (2pt);
    \node[below] at (0,-2) {Instable};
  \end{scope}
\end{tikzpicture}
\end{center}

% ============================================================
\section{Méthode directe de Lyapunov}\label{sec:lyapunov-directe}
% ============================================================

\begin{definition}[Fonction de Lyapunov]\label{def:fct-lyapunov}
  \index{Lyapunov!fonction de}\index{fonction de Lyapunov}
  Soit $V:\mathcal{U}\to\R$ de classe $\mathcal{C}^1$ définie sur un
  voisinage ouvert $\mathcal{U}$ de $\mathbf{0}$.
  \begin{enumerate}[label=(\roman*)]
    \item $V$ est \emph{définie positive} si $V(\mathbf{0})=0$ et
      $V(\mathbf{x})>0$ pour tout $\mathbf{x}\in\mathcal{U}\setminus\{\mathbf{0}\}$.
    \item $V$ est \emph{semi-définie positive} si $V(\mathbf{0})=0$ et
      $V(\mathbf{x})\geq 0$ pour tout $\mathbf{x}\in\mathcal{U}$.
    \item La \emph{dérivée orbitale}\index{dérivée orbitale} de $V$
      le long des trajectoires du système~\eqref{eq:syst-autonome} est
      \[
        \dot{V}(\mathbf{x})
        = \nabla V(\mathbf{x})\cdot\mathbf{f}(\mathbf{x})
        = \sum_{i=1}^n \frac{\partial V}{\partial x_i}(\mathbf{x})\,f_i(\mathbf{x}).
      \]
  \end{enumerate}
\end{definition}

\begin{remarque}[Interprétation]\label{rem:interpret-lyapunov}
  La dérivée orbitale $\dot{V}$ mesure la variation de $V$ le long
  des solutions sans qu'il soit nécessaire de résoudre le système.
  C'est la puissance de la méthode~: on raisonne sur $V$ et $\dot{V}$
  sans connaître explicitement les trajectoires.
\end{remarque}

\begin{theoreme}[Stabilité de Lyapunov]\label{thm:lyapunov-stab}
  \index{Lyapunov!théorème de stabilité}
  Si $V$ est définie positive et $\dot{V}\leq 0$ sur
  $\mathcal{U}\setminus\{\mathbf{0}\}$, alors l'origine est
  \emph{stable} au sens de Lyapunov.
\end{theoreme}

\begin{proof}
  Soit $\varepsilon>0$ assez petit pour que la boule fermée
  $\bar{B}_\varepsilon\subset\mathcal{U}$. Sur la sphère
  $S_\varepsilon=\{\mathbf{x}:\norm{\mathbf{x}}=\varepsilon\}$,
  la fonction $V$ est continue et strictement positive, donc
  elle atteint son minimum $m=\min_{S_\varepsilon}V>0$.

  Par continuité de $V$ en $\mathbf{0}$ (avec $V(\mathbf{0})=0$),
  il existe $\delta\in\,]0,\varepsilon[$ tel que
  $\norm{\mathbf{x}}<\delta\implies V(\mathbf{x})<m$.

  Soit $\mathbf{x}(t)$ une solution avec $\norm{\mathbf{x}(0)}<\delta$.
  Comme $\dot{V}\leq 0$, la fonction $t\mapsto V(\mathbf{x}(t))$ est
  décroissante, donc
  $V(\mathbf{x}(t))\leq V(\mathbf{x}(0))<m$ pour tout $t\geq 0$.

  Or si $\mathbf{x}(t)$ atteignait $S_\varepsilon$ pour un certain
  $t$, on aurait $V(\mathbf{x}(t))\geq m$, contradiction. Donc
  $\norm{\mathbf{x}(t)}<\varepsilon$ pour tout $t\geq 0$.
\end{proof}

\begin{theoreme}[Stabilité asymptotique de Lyapunov]\label{thm:lyapunov-asympt}
  \index{Lyapunov!stabilité asymptotique}
  Si $V$ est définie positive et $\dot{V}$ est \emph{définie
  négative} (c'est-à-dire $\dot{V}(\mathbf{x})<0$ pour tout
  $\mathbf{x}\neq\mathbf{0}$), alors l'origine est
  \emph{asymptotiquement stable}.
\end{theoreme}

\begin{proof}
  La stabilité découle du théorème~\ref{thm:lyapunov-stab}. Il reste
  à montrer la convergence. Prenons $\delta$ comme dans la preuve
  précédente et $\norm{\mathbf{x}(0)}<\delta$. Alors
  $\norm{\mathbf{x}(t)}<\varepsilon$ et $t\mapsto V(\mathbf{x}(t))$
  est strictement décroissante, bornée inférieurement par $0$, donc
  converge vers une limite $c\geq 0$.

  Montrons que $c=0$. Par l'absurde, si $c>0$, alors
  $\mathbf{x}(t)$ reste dans l'anneau compact
  $K=\{\mathbf{x}:c\leq V(\mathbf{x})\leq V(\mathbf{x}(0))\}\subset
  \bar{B}_\varepsilon\setminus\{\mathbf{0}\}$.
  Sur $K$, $\dot{V}$ atteint son maximum $-\mu<0$ ($\dot{V}$ est
  continue et strictement négative sur un compact ne contenant pas
  l'origine). Donc
  $V(\mathbf{x}(t))\leq V(\mathbf{x}(0))-\mu t\to-\infty$, ce qui
  contredit $V\geq 0$. Ainsi $c=0$, et comme $V$ est définie positive,
  $\mathbf{x}(t)\to\mathbf{0}$.
\end{proof}

\begin{theoreme}[Instabilité de Chetaev]\label{thm:chetaev}
  \index{Chetaev (théorème d'instabilité)}\index{instabilité!Chetaev}
  Supposons qu'il existe une fonction $V$ de classe $\mathcal{C}^1$
  sur un voisinage $\mathcal{U}$ de $\mathbf{0}$ et un ouvert
  $\mathcal{D}\subset\mathcal{U}$ tels que~:
  \begin{enumerate}[label=(\roman*)]
    \item $\mathbf{0}\in\partial\mathcal{D}$ ;
    \item $V(\mathbf{x})>0$ et $\dot{V}(\mathbf{x})>0$ pour tout
      $\mathbf{x}\in\mathcal{D}\cap\mathcal{U}$ ;
    \item $V(\mathbf{x})=0$ pour $\mathbf{x}\in\partial\mathcal{D}\cap\mathcal{U}$.
  \end{enumerate}
  Alors l'origine est \emph{instable}.
\end{theoreme}

\begin{proof}
  Soit $B_\varepsilon\subset\mathcal{U}$. Comme
  $\mathbf{0}\in\partial\mathcal{D}$, il existe
  $\mathbf{x}_0\in\mathcal{D}\cap B_\varepsilon$ arbitrairement
  proche de l'origine. Considérons la solution $\mathbf{x}(t)$ avec
  $\mathbf{x}(0)=\mathbf{x}_0$.

  Tant que $\mathbf{x}(t)\in\mathcal{D}\cap\mathcal{U}$,
  $\dot{V}>0$ donc $V(\mathbf{x}(t))>V(\mathbf{x}_0)>0$, ce qui
  empêche $\mathbf{x}(t)$ de toucher $\partial\mathcal{D}$ (où $V=0$).
  Donc $\mathbf{x}(t)$ reste dans $\mathcal{D}$ tant qu'il reste
  dans $\mathcal{U}$. Comme $V$ croît strictement et est bornée dans
  tout compact inclus dans $B_\varepsilon\cap\mathcal{D}$, la
  trajectoire doit finir par quitter $B_\varepsilon$. Ceci montre
  l'instabilité.
\end{proof}

\begin{theoreme}[Principe d'invariance de La Salle]\label{thm:lasalle}
  \index{La Salle (principe d'invariance)}\index{principe d'invariance}
  Soit $V$ une fonction de Lyapunov définie positive avec $\dot{V}\leq 0$
  sur un ensemble compact positivement invariant $\mathcal{K}$. Posons
  \[
    E = \{\mathbf{x}\in\mathcal{K}:\dot{V}(\mathbf{x})=0\}.
  \]
  Soit $M$ le plus grand sous-ensemble invariant contenu dans $E$.
  Alors toute trajectoire issue de $\mathcal{K}$ converge vers $M$
  quand $t\to+\infty$.
\end{theoreme}

\begin{proof}[Idée de la preuve]
  Soit $\mathbf{x}(t)$ une trajectoire dans $\mathcal{K}$.
  Comme $V$ est décroissante et bornée, elle converge vers $c\geq 0$.
  L'ensemble $\omega$-limite\index{ensemble $\omega$-limite}
  $\omega(\mathbf{x}_0)$ est non vide (compacité de $\mathcal{K}$),
  invariant, et $V\equiv c$ sur $\omega(\mathbf{x}_0)$. Donc
  $\dot{V}=0$ sur $\omega(\mathbf{x}_0)$, ce qui implique
  $\omega(\mathbf{x}_0)\subset E$. Par invariance de
  $\omega(\mathbf{x}_0)$, on a $\omega(\mathbf{x}_0)\subset M$.
\end{proof}

\begin{remarque}[Utilité]\label{rem:lasalle-utilite}
  Le principe de La Salle est crucial quand $\dot{V}$ est seulement
  semi-définie négative~: on ne peut pas conclure directement par le
  théorème~\ref{thm:lyapunov-asympt}, mais le principe d'invariance
  permet souvent de prouver la stabilité asymptotique malgré tout.
\end{remarque}

\begin{exemple}[Pendule amorti et La Salle]\label{ex:pendule-amorti}
  \index{pendule!amorti}
  Considérons le pendule avec frottement~:
  \[
    \theta''+b\theta'+\sin\theta=0, \qquad b>0.
  \]
  Posons $x_1=\theta$, $x_2=\theta'$ et l'énergie
  $V(x_1,x_2)=\frac{1}{2}x_2^2+(1-\cos x_1)$.
  Alors $V$ est définie positive au voisinage de l'origine et
  \[
    \dot{V}=x_2(-\sin x_1-bx_2)+x_2\sin x_1=-bx_2^2\leq 0.
  \]
  Ici $\dot{V}$ n'est que semi-définie négative ($\dot{V}=0$ si
  $x_2=0$). L'ensemble $E=\{(x_1,x_2):\dot{V}=0\}=\{x_2=0\}$.
  Sur $E$, $x_2=0$ implique $x_2'=-\sin x_1-bx_2=-\sin x_1$, et
  pour rester dans $E$ (c'est-à-dire $x_2'\equiv 0$), il faut
  $\sin x_1=0$. Localement autour de l'origine, $M=\{(0,0)\}$.
  Par le principe de La Salle, l'origine est asymptotiquement stable.
\end{exemple}

% ============================================================
\section{Méthode de linéarisation}\label{sec:linearisation}
% ============================================================

\begin{definition}[Linéarisation]\label{def:linearisation}
  \index{linéarisation}\index{matrice jacobienne}
  La \emph{linéarisation} du système~\eqref{eq:syst-autonome} en
  l'équilibre $\mathbf{0}$ est le système linéaire
  \[
    \mathbf{y}'=A\mathbf{y},
    \qquad A=D\mathbf{f}(\mathbf{0})
    = \left(\frac{\partial f_i}{\partial x_j}(\mathbf{0})\right)_{1\leq i,j\leq n},
  \]
  où $A$ est la matrice jacobienne de $\mathbf{f}$ en $\mathbf{0}$.
\end{definition}

\begin{definition}[Équilibre hyperbolique]\label{def:hyperbolique}
  \index{équilibre hyperbolique}\index{hyperbolique (équilibre)}
  Le point d'équilibre $\mathbf{0}$ est dit \emph{hyperbolique} si
  toutes les valeurs propres de $A=D\mathbf{f}(\mathbf{0})$ ont une
  partie réelle non nulle.
\end{definition}

\begin{theoreme}[Hartman--Grobman]\label{thm:hartman-grobman}
  \index{Hartman--Grobman (théorème de)}
  Si $\mathbf{0}$ est un point d'équilibre \emph{hyperbolique} du
  système~\eqref{eq:syst-autonome}, alors au voisinage de l'origine,
  le portrait de phase du système non linéaire est
  \emph{topologiquement conjugué} à celui du système linéarisé
  $\mathbf{y}'=A\mathbf{y}$~: il existe un homéomorphisme local qui
  envoie les trajectoires de l'un sur celles de l'autre en préservant
  le sens du temps.
\end{theoreme}

\begin{remarque}[Portée et limites]\label{rem:hartman-limites}
  Ce théorème garantit que la linéarisation donne la bonne
  \emph{topologie} du portrait de phase pour les équilibres
  hyperboliques. L'homéomorphisme n'est pas nécessairement un
  difféomorphisme~: les angles et les distances ne sont pas préservés.
\end{remarque}

\begin{corollaire}[Stabilité par linéarisation]\label{cor:stab-linearisation}
  \index{stabilité!par linéarisation}
  Soit $\mathbf{0}$ un point d'équilibre du
  système~\eqref{eq:syst-autonome} et $A=D\mathbf{f}(\mathbf{0})$.
  \begin{enumerate}[label=(\roman*)]
    \item Si toutes les valeurs propres de $A$ ont une partie réelle
      strictement négative ($\operatorname{Re}(\lambda_i)<0$ pour
      tout $i$), alors $\mathbf{0}$ est \emph{asymptotiquement stable}.
    \item Si au moins une valeur propre a une partie réelle
      strictement positive, alors $\mathbf{0}$ est \emph{instable}.
  \end{enumerate}
\end{corollaire}

\begin{remarque}[Cas non hyperbolique]\label{rem:cas-centre}
  \index{centre (cas non concluant)}
  Lorsque certaines valeurs propres sont purement imaginaires
  (partie réelle nulle) et les autres à partie réelle négative, la
  linéarisation ne permet pas de conclure. Le système linéarisé
  présente un \emph{centre}, mais le système non linéaire peut être
  stable, asymptotiquement stable, ou instable selon les termes
  d'ordre supérieur.
\end{remarque}

% ============================================================
\section{Stabilité des systèmes linéaires}\label{sec:stab-lineaire}
% ============================================================

\begin{theoreme}[Stabilité et valeurs propres]\label{thm:stab-vp}
  \index{stabilité!système linéaire}
  Pour le système linéaire $\mathbf{x}'=A\mathbf{x}$ avec
  $A\in\R^{n\times n}$~:
  \begin{enumerate}[label=(\roman*)]
    \item L'origine est \emph{asymptotiquement stable} si et
      seulement si toutes les valeurs propres de $A$ satisfont
      $\operatorname{Re}(\lambda)<0$.
    \item L'origine est \emph{stable} (mais pas asymptotiquement) si
      et seulement si $\operatorname{Re}(\lambda)\leq 0$ pour toute
      valeur propre, et les valeurs propres à partie réelle nulle
      sont semi-simples (blocs de Jordan de taille $1$).
    \item L'origine est \emph{instable} dans tous les autres cas.
  \end{enumerate}
\end{theoreme}

\begin{proof}
  La solution générale est $\mathbf{x}(t)=e^{tA}\mathbf{x}(0)$. Le
  comportement de $e^{tA}$ est déterminé par la forme de Jordan de
  $A$. Un bloc de Jordan $J_k(\lambda)$ contribue des termes de la
  forme $t^j e^{\lambda t}$ avec $0\leq j<k$.
  \begin{itemize}
    \item Si $\operatorname{Re}(\lambda)<0$, alors
      $t^j e^{\operatorname{Re}(\lambda)t}\to 0$ pour tout $j$.
    \item Si $\operatorname{Re}(\lambda)=0$ et $k=1$, les termes
      $e^{i\operatorname{Im}(\lambda)t}$ restent bornés.
    \item Si $\operatorname{Re}(\lambda)=0$ et $k>1$, les termes
      $t^j e^{i\operatorname{Im}(\lambda)t}$ sont non bornés.
    \item Si $\operatorname{Re}(\lambda)>0$, les termes explosent.
      \qedhere
  \end{itemize}
\end{proof}

% ============================================================
\section{Analyse du plan de phase pour les systèmes $2\times 2$}\label{sec:plan-phase}
% ============================================================

\index{plan de phase}\index{portrait de phase}

Considérons le système linéaire plan $\mathbf{x}'=A\mathbf{x}$ avec
$A\in\R^{2\times 2}$. Soient $\lambda_1,\lambda_2$ les valeurs propres
de $A$. On note $\tau=\tr(A)=\lambda_1+\lambda_2$ et
$\Delta=\det(A)=\lambda_1\lambda_2$. Les valeurs propres sont les
racines de $\lambda^2-\tau\lambda+\Delta=0$.

\begin{proposition}[Classification des portraits de phase plans]\label{prop:classif-phase}
  \index{classification des équilibres}
  \begin{enumerate}[label=\textbf{(\arabic*)}]
    \item \textbf{N\oe ud stable}\index{noeud stable} ($\Delta>0$,
      $\tau<0$, $\tau^2>4\Delta$)~: deux valeurs propres réelles
      négatives. Toutes les trajectoires convergent vers l'origine.
    \item \textbf{N\oe ud instable}\index{noeud instable} ($\Delta>0$,
      $\tau>0$, $\tau^2>4\Delta$)~: deux valeurs propres réelles
      positives. Divergence.
    \item \textbf{Point selle}\index{point selle} ($\Delta<0$)~:
      une valeur propre positive, une négative. Instable.
    \item \textbf{Foyer stable}\index{foyer stable} ($\Delta>0$,
      $\tau<0$, $\tau^2<4\Delta$)~: valeurs propres complexes à
      partie réelle négative. Spirales convergentes.
    \item \textbf{Foyer instable}\index{foyer instable} ($\Delta>0$,
      $\tau>0$, $\tau^2<4\Delta$)~: valeurs propres complexes à
      partie réelle positive. Spirales divergentes.
    \item \textbf{Centre}\index{centre} ($\Delta>0$, $\tau=0$)~:
      valeurs propres purement imaginaires. Orbites périodiques
      (stable mais pas asymptotiquement, pour le système linéaire).
  \end{enumerate}
\end{proposition}

\begin{center}
\begin{tikzpicture}[>=Stealth, scale=0.75]
  % Noeud stable
  \begin{scope}[xshift=-5cm, yshift=3cm]
    \draw[->] (-1.6,0)--(1.6,0); \draw[->] (0,-1.6)--(0,1.6);
    \foreach \a in {0,45,...,315}{
      \draw[->, blue!70!black, thick, domain=0:1.5, samples=30]
        plot ({1.4*cos(\a)*exp(-1.2*\x)},{1.4*sin(\a)*exp(-0.8*\x)});
    }
    \node[below] at (0,-1.9) {\small N\oe ud stable};
  \end{scope}
  % Point selle
  \begin{scope}[xshift=0cm, yshift=3cm]
    \draw[->] (-1.6,0)--(1.6,0); \draw[->] (0,-1.6)--(0,1.6);
    \draw[->, red!70!black, thick, domain=-1.3:0, samples=20]
      plot ({exp(\x)},{exp(-\x)});
    \draw[->, red!70!black, thick, domain=-1.3:0, samples=20]
      plot ({-exp(\x)},{-exp(-\x)});
    \draw[->, red!70!black, thick, domain=0:1.3, samples=20]
      plot ({exp(-\x)},{-exp(\x)});
    \draw[->, red!70!black, thick, domain=0:1.3, samples=20]
      plot ({-exp(-\x)},{exp(\x)});
    \node[below] at (0,-1.9) {\small Point selle};
  \end{scope}
  % Foyer stable
  \begin{scope}[xshift=5cm, yshift=3cm]
    \draw[->] (-1.6,0)--(1.6,0); \draw[->] (0,-1.6)--(0,1.6);
    \draw[->, blue!60!black, thick, domain=0:8, samples=50]
      plot ({1.3*exp(-0.15*\x)*cos(50*\x)},{1.3*exp(-0.15*\x)*sin(50*\x)});
    \node[below] at (0,-1.9) {\small Foyer stable};
  \end{scope}
  % Noeud instable
  \begin{scope}[xshift=-5cm, yshift=-2.5cm]
    \draw[->] (-1.6,0)--(1.6,0); \draw[->] (0,-1.6)--(0,1.6);
    \foreach \a in {0,45,...,315}{
      \draw[->, red!70!black, thick, domain=0:1.5, samples=30]
        plot ({0.15*cos(\a)*exp(1.2*\x)},{0.15*sin(\a)*exp(0.8*\x)});
    }
    \node[below] at (0,-1.9) {\small N\oe ud instable};
  \end{scope}
  % Centre
  \begin{scope}[xshift=0cm, yshift=-2.5cm]
    \draw[->] (-1.6,0)--(1.6,0); \draw[->] (0,-1.6)--(0,1.6);
    \foreach \r in {0.4,0.7,1.1}{
      \draw[green!50!black, thick] (0,0) circle (\r);
    }
    \draw[->, green!50!black, thick] (0.4,0) arc (0:30:0.4);
    \draw[->, green!50!black, thick] (0.7,0) arc (0:30:0.7);
    \draw[->, green!50!black, thick] (1.1,0) arc (0:30:1.1);
    \node[below] at (0,-1.9) {\small Centre};
  \end{scope}
  % Foyer instable
  \begin{scope}[xshift=5cm, yshift=-2.5cm]
    \draw[->] (-1.6,0)--(1.6,0); \draw[->] (0,-1.6)--(0,1.6);
    \draw[->, red!60!black, thick, domain=0:8, samples=50]
      plot ({0.1*exp(0.15*\x)*cos(50*\x)},{0.1*exp(0.15*\x)*sin(50*\x)});
    \node[below] at (0,-1.9) {\small Foyer instable};
  \end{scope}
\end{tikzpicture}
\end{center}

\begin{center}
\begin{tikzpicture}[>=Stealth, scale=0.85]
  \draw[->] (-3.5,0)--(3.5,0) node[right]{$\tau$};
  \draw[->] (0,-1.5)--(0,3.5) node[above]{$\Delta$};
  \draw[thick, blue!70!black, domain=-3.2:3.2, samples=50]
    plot (\x, {0.25*\x*\x}) node[right]{\small$\tau^2=4\Delta$};
  \node[blue!60!black] at (-1.8,2.5) {\small Foyer stable};
  \node[red!60!black] at (1.8,2.5) {\small Foyer instable};
  \node[blue!60!black] at (-2.8,0.7) {\small N\oe ud};
  \node[blue!60!black] at (-2.8,0.3) {\small stable};
  \node[red!60!black] at (2.8,0.7) {\small N\oe ud};
  \node[red!60!black] at (2.8,0.3) {\small instable};
  \node[red!60!black] at (0,-0.9) {\small Point selle};
  \node[green!50!black] at (0,1.5) {\small Centre};
  \draw[green!50!black, thick, dashed] (0,0.05)--(0,3.2);
\end{tikzpicture}
\end{center}

% ============================================================
\section{Exemples détaillés}\label{sec:exemples-stab}
% ============================================================

\begin{exemple}[Pendule non linéaire]\label{ex:stab-pendule}
  \index{pendule!stabilité}
  Reprenons le pendule $\theta''+\omega_0^2\sin\theta=0$ avec
  $\omega_0^2=g/\ell$, soit le système
  \[
    x_1'=x_2, \qquad x_2'=-\omega_0^2\sin x_1.
  \]

  \textbf{Équilibre $(0,0)$} (pendule en bas). La jacobienne est
  \[
    A=\begin{pmatrix}0&1\\-\omega_0^2&0\end{pmatrix},
    \qquad \lambda=\pm i\omega_0.
  \]
  Les valeurs propres sont purement imaginaires~: le linéarisé donne
  un \emph{centre}. La linéarisation seule ne conclut pas.
  Utilisons une fonction de Lyapunov~: l'énergie mécanique
  \[
    V(x_1,x_2)=\frac{1}{2}x_2^2+\omega_0^2(1-\cos x_1).
  \]
  On vérifie~: $V(0,0)=0$, $V>0$ au voisinage de l'origine, et
  \[
    \dot{V}=x_2\cdot(-\omega_0^2\sin x_1)+\omega_0^2\sin x_1\cdot x_2=0.
  \]
  Donc $\dot{V}\equiv 0$~: l'énergie est conservée. Par le
  théorème~\ref{thm:lyapunov-stab}, l'origine est stable. Les
  trajectoires sont les courbes de niveau de $V$~: orbites
  périodiques fermées.

  \textbf{Équilibre $(\pi,0)$} (pendule en haut). La jacobienne est
  \[
    A=\begin{pmatrix}0&1\\\omega_0^2&0\end{pmatrix},
    \qquad \lambda=\pm\omega_0.
  \]
  Une valeur propre réelle positive~: l'équilibre est un
  \emph{point selle}, donc \emph{instable}.

  \begin{center}
  \begin{tikzpicture}[>=Stealth, scale=0.7]
    \draw[->] (-4,0)--(4,0) node[right]{$\theta$};
    \draw[->] (0,-2.5)--(0,2.5) node[above]{$\dot\theta$};
    % Orbites fermées autour de l'origine
    \foreach \E in {0.3,0.7,1.2}{
      \draw[blue!70!black, thick, domain=-180:180, samples=50]
        plot ({(\E)*sin(\x)}, {(\E)*cos(\x)});
    }
    % Séparatrice
    \draw[red, thick, domain=-170:170, samples=50]
      plot ({3.14*\x/180},{1.8*cos(\x/2)});
    \draw[red, thick, domain=-170:170, samples=50]
      plot ({3.14*\x/180},{-1.8*cos(\x/2)});
    \fill (0,0) circle (2.5pt) node[below right]{\small$(0,0)$};
    \fill (3.14,0) circle (2.5pt) node[below right]{\small$(\pi,0)$};
    \fill (-3.14,0) circle (2.5pt) node[below left]{\small$(-\pi,0)$};
    \node at (-3,-2.2) {\small\textcolor{red}{Séparatrices}};
  \end{tikzpicture}
  \end{center}
\end{exemple}

\begin{exemple}[Système de Lotka--Volterra]\label{ex:stab-lotka}
  \index{Lotka-Volterra!stabilité}
  Reprenons le système prédateur-proie avec $\alpha=\beta=\gamma=\delta=1$~:
  \[
    x'=x(1-y), \qquad y'=y(-1+x).
  \]
  L'équilibre de coexistence est $(1,1)$. Posons $u=x-1$, $v=y-1$~:
  \[
    u'=(1+u)(-v)=-v-uv, \qquad v'=(1+v)\,u=u+uv.
  \]
  La jacobienne en $(0,0)$ est
  $A=\begin{pmatrix}0&-1\\1&0\end{pmatrix}$, de valeurs propres
  $\pm i$~: c'est un \emph{centre}. La linéarisation ne conclut pas.

  Construisons une fonction de Lyapunov. Posons
  \[
    V(x,y)=x-\ln x+y-\ln y-2 \qquad (x,y>0).
  \]
  On vérifie $V(1,1)=0$ et $V>0$ pour $(x,y)\neq(1,1)$ (par convexité
  de $t\mapsto t-\ln t$, dont le minimum est $1$ en $t=1$). Puis
  \begin{align*}
    \dot{V}&=\left(1-\frac{1}{x}\right)x(1-y)
            +\left(1-\frac{1}{y}\right)y(-1+x) \\
           &=(x-1)(1-y)+(y-1)(x-1) \\
           &=(x-1)(1-y)+(x-1)(y-1)=0.
  \end{align*}
  L'énergie $V$ est conservée~: l'équilibre $(1,1)$ est stable et les
  trajectoires sont des courbes fermées.
\end{exemple}

\begin{exemple}[Oscillateur de van der Pol]\label{ex:van-der-pol}
  \index{van der Pol (oscillateur de)}
  L'équation de van der Pol est
  \[
    x''-\mu(1-x^2)x'+x=0, \qquad \mu>0.
  \]
  Le système associé est
  $x_1'=x_2$,
  $x_2'=-x_1+\mu(1-x_1^2)x_2$.

  L'unique équilibre est l'origine. La jacobienne est
  \[
    A=\begin{pmatrix}0&1\\-1&\mu\end{pmatrix},
    \qquad \lambda=\frac{\mu\pm\sqrt{\mu^2-4}}{2}.
  \]
  Pour $0<\mu<2$, les valeurs propres sont complexes de partie réelle
  $\mu/2>0$~: l'origine est un \emph{foyer instable}. Les trajectoires
  spiralent vers l'extérieur et convergent vers un \emph{cycle
  limite}\index{cycle limite} stable (existence garantie par le
  théorème de Poincaré--Bendixson\index{Poincaré--Bendixson}).

  \begin{center}
  \begin{tikzpicture}[>=Stealth, scale=0.65]
    \draw[->] (-3.5,0)--(3.5,0) node[right]{$x_1$};
    \draw[->] (0,-3.5)--(0,3.5) node[above]{$x_2$};
    % Spirale sortante depuis l'intérieur
    \draw[->, blue!60!black, thick, domain=0:15, samples=50]
      plot ({(0.3+0.12*\x)*cos(40*\x)},{(0.3+0.12*\x)*sin(40*\x)});
    % Cycle limite (approximation)
    \draw[red, very thick] (0,0) circle (2.2);
    % Spirale entrante depuis l'extérieur
    \draw[->, green!50!black, thick, domain=0:10, samples=50]
      plot ({(3.0-0.08*\x)*cos(40*\x)},{(3.0-0.08*\x)*sin(40*\x)});
    \fill (0,0) circle (2pt);
    \node[red] at (2.8,1.5) {\small Cycle limite};
  \end{tikzpicture}
  \end{center}
\end{exemple}

\begin{exemple}[Construction d'une fonction de Lyapunov]\label{ex:construction-lyapunov}
  \index{Lyapunov!construction}
  Considérons le système
  \[
    x'=-x+y^2, \qquad y'=-x\,y-y^3.
  \]
  L'origine est un équilibre. Essayons la forme quadratique
  $V(x,y)=\frac{1}{2}(x^2+y^2)$. Alors
  \begin{align*}
    \dot{V} &= x(-x+y^2)+y(-xy-y^3) \\
            &= -x^2+xy^2-xy^2-y^4 \\
            &= -x^2-y^4.
  \end{align*}
  On a $\dot{V}\leq 0$ et $\dot{V}=0$ seulement si $x=0$ et $y=0$.
  Donc $\dot{V}$ est définie négative, et par le
  théorème~\ref{thm:lyapunov-asympt}, l'origine est
  \emph{asymptotiquement stable}.
\end{exemple}

\begin{exemple}[Instabilité par Chetaev]\label{ex:chetaev}
  \index{Chetaev (théorème d'instabilité)!exemple}
  Considérons $x'=x+y^2$, $y'=-y+x^2$. L'origine est un équilibre.
  Posons $V(x,y)=\frac{1}{2}(x^2-y^2)$. Alors
  \[
    \dot{V}=x(x+y^2)-y(-y+x^2)=x^2+xy^2+y^2-x^2y=x^2+y^2+xy^2-x^2y.
  \]
  Au voisinage de l'origine, les termes dominants sont $x^2+y^2>0$.
  Prenons $\mathcal{D}=\{(x,y):x^2>y^2,\;x>0\}$. Sur
  $\partial\mathcal{D}\cap\mathcal{U}$, on a $x^2=y^2$ donc $V=0$.
  Dans $\mathcal{D}$, $V>0$ et $\dot{V}>0$ pour $\norm{(x,y)}$ assez
  petit. Par le théorème de Chetaev~\ref{thm:chetaev}, l'origine est
  \emph{instable}.
\end{exemple}

% ============================================================
\section{Courbes de niveau de Lyapunov}\label{sec:courbes-lyapunov}
% ============================================================

\index{Lyapunov!courbes de niveau}

Les courbes de niveau d'une fonction de Lyapunov $V$ fournissent des
informations géométriques précieuses~: lorsque $\dot{V}<0$, les
trajectoires traversent les courbes de niveau vers l'intérieur (vers
les valeurs plus faibles de $V$), ce qui illustre la convergence.

\begin{center}
\begin{tikzpicture}[>=Stealth, scale=0.85]
  \draw[->] (-2.8,0)--(2.8,0) node[right]{$x_1$};
  \draw[->] (0,-2.8)--(0,2.8) node[above]{$x_2$};
  % Courbes de niveau elliptiques
  \foreach \r/\c in {0.6/blue!20, 1.1/blue!35, 1.6/blue!50, 2.1/blue!65}{
    \draw[\c, thick] (0,0) ellipse ({\r} and {0.7*\r});
  }
  % Flèches montrant la décroissance
  \foreach \a in {30,75,120,165,210,255,300,345}{
    \draw[->, red!70!black, thick]
      ({1.8*cos(\a)},{1.26*sin(\a)}) -- ({1.4*cos(\a)},{0.98*sin(\a)});
  }
  \fill (0,0) circle (2pt);
  \node at (2.5,1.8) {\small$V=c_4$};
  \node at (1.8,1.2) {\small$V=c_3$};
  \node at (1.3,0.6) {\small$V=c_2$};
  \node at (0.8,0.2) {\small$V=c_1$};
  \node[below right] at (0,0) {\small$\mathbf{0}$};
\end{tikzpicture}
\end{center}

% ============================================================
\section{Exercices}\label{sec:exo-stabilite}
% ============================================================

\begin{exercice}[Points d'équilibre et linéarisation]\label{exo:equil-lin}
  Pour chacun des systèmes suivants, trouver les points d'équilibre,
  calculer la matrice jacobienne, et déterminer la stabilité par
  linéarisation quand c'est possible.
  \begin{enumerate}[label=\alph*)]
    \item $x'=x(1-x-y)$, $y'=y(2-y-3x)$.
    \item $x'=\sin y$, $y'=-x-y$.
    \item $x'=y$, $y'=-\sin x-y$.
  \end{enumerate}
\end{exercice}

\begin{exercice}[Fonctions de Lyapunov]\label{exo:lyapunov}
  \begin{enumerate}[label=\alph*)]
    \item Montrer que $V(x,y)=x^2+2y^2$ est une fonction de Lyapunov
      pour le système $x'=-x+2xy$, $y'=-y-x^2$ au voisinage de
      l'origine. Conclure sur la stabilité.
    \item Pour $x'=-x^3$, $y'=-y$, montrer que l'origine est
      asymptotiquement stable en utilisant $V(x,y)=x^4+y^2$.
    \item Soit $x''+h(x')x'+g(x)=0$ où $g(0)=0$, $xg(x)>0$ pour
      $x\neq 0$, et $h(v)\geq h_0>0$. Montrer que l'origine est
      asymptotiquement stable en utilisant
      $V(x,v)=\int_0^x g(\xi)\,\dd\xi+\frac{1}{2}v^2$.
  \end{enumerate}
\end{exercice}

\begin{exercice}[Classification des portraits de phase]\label{exo:classif-phase}
  Pour chaque matrice $A$, déterminer le type de portrait de phase et
  esquisser quelques trajectoires~:
  \begin{enumerate}[label=\alph*)]
    \item $A=\begin{pmatrix}-1&0\\0&-3\end{pmatrix}$.
    \item $A=\begin{pmatrix}0&1\\-4&0\end{pmatrix}$.
    \item $A=\begin{pmatrix}1&-1\\1&1\end{pmatrix}$.
    \item $A=\begin{pmatrix}2&1\\0&2\end{pmatrix}$.
  \end{enumerate}
\end{exercice}

\begin{exercice}[Principe de La Salle]\label{exo:lasalle}
  Considérons le système
  \[
    x''+x'+x+x^3=0.
  \]
  \begin{enumerate}[label=\alph*)]
    \item Montrer que $V(x,y)=\frac{1}{2}(x^2+y^2)+\frac{1}{4}x^4$
      (avec $y=x'$) est une fonction de Lyapunov.
    \item Calculer $\dot{V}$ et appliquer le principe de La Salle pour
      prouver la stabilité asymptotique de l'origine.
  \end{enumerate}
\end{exercice}

\begin{exercice}[Instabilité]\label{exo:instabilite}
  Montrer que l'origine du système $x'=y+x(x^2+y^2)$,
  $y'=-x+y(x^2+y^2)$ est instable. \textit{Indication~: passer en
  coordonnées polaires et étudier $r'$.}
\end{exercice}

\begin{exercice}[Système de Lorenz simplifié]\label{exo:lorenz}
  \index{Lorenz (système de)}
  Considérons le système simplifié
  \[
    x'=-\sigma x+\sigma y,\qquad y'=rx-y-xz,\qquad z'=xy-bz,
  \]
  avec $\sigma,b>0$ et $r>0$.
  \begin{enumerate}[label=\alph*)]
    \item Trouver les points d'équilibre.
    \item Pour $r<1$, montrer que l'origine est globalement
      asymptotiquement stable en utilisant
      $V(x,y,z)=\frac{1}{2\sigma}x^2+\frac{1}{2}y^2+\frac{1}{2}z^2$.
    \item Pour $r>1$, étudier la stabilité des trois équilibres par
      linéarisation.
  \end{enumerate}
\end{exercice}

\begin{exercice}[Compétition entre deux espèces]\label{exo:competition}
  \index{compétition entre espèces}
  Deux espèces en compétition sont modélisées par
  \[
    x'=x(1-x-\alpha y),\qquad y'=y(1-y-\beta x),
    \qquad \alpha,\beta>0.
  \]
  \begin{enumerate}[label=\alph*)]
    \item Trouver tous les points d'équilibre dans le quadrant
      $\{x\geq 0,y\geq 0\}$.
    \item Analyser la stabilité de chaque équilibre en fonction des
      paramètres $\alpha$ et $\beta$.
    \item Interpréter biologiquement les cas $\alpha\beta<1$ et
      $\alpha\beta>1$.
  \end{enumerate}
\end{exercice}

% ============================================================
\section*{Résumé du chapitre}
\addcontentsline{toc}{section}{Résumé du chapitre}
% ============================================================

\begin{itemize}
  \item Un point d'équilibre $\mathbf{x}_0$ vérifie
    $\mathbf{f}(\mathbf{x}_0)=\mathbf{0}$.
  \item La \emph{stabilité au sens de Lyapunov}~: les trajectoires
    proches restent proches. La \emph{stabilité asymptotique}~:
    de plus, elles convergent vers l'équilibre.
  \item \emph{Méthode de Lyapunov}~: si $V$ est définie positive
    et $\dot{V}\leq 0$, alors stabilité ; si $\dot{V}<0$, alors
    stabilité asymptotique.
  \item Le \emph{principe de La Salle} étend la méthode au cas
    $\dot{V}\leq 0$ semi-définie~: les trajectoires convergent
    vers le plus grand invariant dans $\{\dot{V}=0\}$.
  \item Le théorème de \emph{Chetaev} fournit un critère
    d'instabilité.
  \item \emph{Linéarisation}~: si $A=D\mathbf{f}(\mathbf{0})$ a
    toutes ses valeurs propres à partie réelle $<0$, alors
    stabilité asymptotique. Si une valeur propre a
    $\operatorname{Re}>0$, instabilité.
  \item Le théorème de \emph{Hartman--Grobman}~: pour un équilibre
    hyperbolique, le portrait de phase non linéaire est
    topologiquement conjugué au linéarisé.
  \item En dimension $2$, la classification des portraits de phase
    (n\oe ud, selle, foyer, centre) dépend du discriminant et de la
    trace de la jacobienne.
\end{itemize}

% === FIN DU CHUNK ch7-8 ===
% === DÉBUT DU CHUNK ch9-10+end ===
%%%%%%%%%%%%%%%%%%%%%%%%%%%%%%%%%%%%%%%%%%%%%%%%%%%%%%%%%%%%%%%%%%%%
%  Chapitres 9–10 + Formulaire + Bibliographie + Index
%%%%%%%%%%%%%%%%%%%%%%%%%%%%%%%%%%%%%%%%%%%%%%%%%%%%%%%%%%%%%%%%%%%%

%% ═══════════════════════════════════════════════════════════════════
%%  CHAPITRE 9 — Systèmes Non-Linéaires et Bifurcations
%% ═══════════════════════════════════════════════════════════════════
\chapter{Introduction aux Systèmes Non-Linéaires et Bifurcations}
\label{chap:nonlineaire}
\index{systèmes non-linéaires}

\section*{Introduction}

Les chapitres précédents ont traité des équations différentielles
linéaires, pour lesquelles la théorie est essentiellement complète.
Cependant, la grande majorité des modèles issus de la physique, de
la biologie ou de la chimie sont \emph{non-linéaires}. Le comportement
qualitatif de ces systèmes --- cycles limites, bifurcations, chaos ---
constitue l'un des domaines les plus riches des mathématiques appliquées
contemporaines.
\index{cycle limite}
\index{bifurcation}

\begin{exemple}[Dynamique des populations]
\label{ex:motivation-nonlin}
Le modèle de Lotka--Volterra pour un système proie--prédateur s'écrit
\[
\begin{cases}
\dot{x} = \alpha x - \beta xy, \\
\dot{y} = \delta xy - \gamma y,
\end{cases}
\]
où $x(t)$ et $y(t)$ désignent les populations de proies et de
prédateurs respectivement. Le terme $xy$ rend ce système intrinsèquement
non-linéaire.
\end{exemple}

\begin{exemple}[Réactions chimiques]
\label{ex:motivation-chimie}
La cinétique de Michaelis--Menten, modélisant l'action d'une enzyme
sur un substrat, conduit au système
\[
\begin{cases}
\dot{s} = -k_1 e\, s + k_{-1} c, \\
\dot{c} = k_1 e\, s - (k_{-1} + k_2) c,
\end{cases}
\]
où $s$ est la concentration du substrat, $c$ celle du complexe
enzyme-substrat, et $e = e_0 - c$ est la concentration en enzyme libre.
\end{exemple}

% ─── 9.1 Systèmes non-linéaires dans R² ─────────────────────────
\section{Systèmes non-linéaires dans \texorpdfstring{$\R^2$}{R²}}
\label{sec:syst-nonlin-R2}
\index{système autonome}

On considère un système autonome dans $\R^2$ :
\begin{equation}
\label{eq:syst-nonlin}
\dot{\mathbf{x}} = \mathbf{f}(\mathbf{x}),
\qquad \mathbf{x} = \begin{pmatrix} x_1 \\ x_2 \end{pmatrix},
\quad \mathbf{f} : \Omega \subset \R^2 \to \R^2.
\end{equation}

\begin{theoreme}[Existence et unicité --- Cauchy--Lipschitz]
\label{thm:cauchy-lip-nonlin}
\index{Cauchy--Lipschitz (théorème)}
Soit $\Omega \subset \R^2$ un ouvert et
$\mathbf{f} : \Omega \to \R^2$ localement lipschitzienne.
Pour tout $\mathbf{x}_0 \in \Omega$, il existe un intervalle maximal
$I = (\alpha, \omega)$ contenant $0$ et une unique solution maximale
$\mathbf{x} : I \to \Omega$ du problème de Cauchy
\[
\dot{\mathbf{x}} = \mathbf{f}(\mathbf{x}), \qquad
\mathbf{x}(0) = \mathbf{x}_0.
\]
\end{theoreme}

\begin{definition}[Flot]
\label{def:flot}
\index{flot}
Le \emph{flot} du système \eqref{eq:syst-nonlin} est l'application
$\varphi : D \subset \R \times \Omega \to \Omega$ définie par
$\varphi(t, \mathbf{x}_0) = \mathbf{x}(t)$, où $\mathbf{x}$ est
la solution maximale issue de $\mathbf{x}_0$. Le flot vérifie :
\begin{enumerate}[label=(\roman*)]
  \item $\varphi(0, \mathbf{x}_0) = \mathbf{x}_0$ ;
  \item $\varphi(t+s, \mathbf{x}_0) = \varphi(t, \varphi(s, \mathbf{x}_0))$
        (propriété de semi-groupe).
\end{enumerate}
\end{definition}

\begin{definition}[Points d'équilibre]
\label{def:pts-equil}
\index{point d'équilibre}
Un point $\mathbf{x}^* \in \Omega$ est un \emph{point d'équilibre}
(ou point singulier, point critique) du système \eqref{eq:syst-nonlin}
si $\mathbf{f}(\mathbf{x}^*) = \mathbf{0}$.
\end{definition}

% ─── 9.2 Analyse du plan de phase ───────────────────────────────
\section{Analyse du plan de phase}
\label{sec:plan-phase}
\index{plan de phase}

\subsection{Isoclines nulles}
\label{subsec:isoclines}
\index{isocline nulle}

\begin{definition}[Isoclines nulles]
\label{def:isoclines}
Pour le système $\dot{x}_1 = f_1(x_1, x_2)$,
$\dot{x}_2 = f_2(x_1, x_2)$, on appelle :
\begin{itemize}
  \item \emph{isocline nulle de $x_1$} : l'ensemble
        $\{(x_1, x_2) : f_1(x_1, x_2) = 0\}$ ;
  \item \emph{isocline nulle de $x_2$} : l'ensemble
        $\{(x_1, x_2) : f_2(x_1, x_2) = 0\}$.
\end{itemize}
Les points d'équilibre sont les intersections de ces deux ensembles.
\end{definition}

\begin{remarque}
Les isoclines nulles décomposent le plan en régions où le signe de
chaque composante du champ de vecteurs est constant. Cela permet de
déterminer la direction qualitative du flot dans chaque région.
\end{remarque}

\begin{exemple}[Isoclines du système de Lotka--Volterra]
\label{ex:isoclines-LV}
Pour le système $\dot{x} = x(\alpha - \beta y)$,
$\dot{y} = y(\delta x - \gamma)$, les isoclines nulles sont :
\begin{itemize}
  \item $\dot{x} = 0$ : $x = 0$ ou $y = \alpha/\beta$ ;
  \item $\dot{y} = 0$ : $y = 0$ ou $x = \gamma/\delta$.
\end{itemize}
Les points d'équilibre sont $(0,0)$ et
$(\gamma/\delta,\, \alpha/\beta)$.
\end{exemple}

\begin{center}
\begin{tikzpicture}[scale=1.1]
  \draw[->] (-0.3,0) -- (4.5,0) node[right] {$x$};
  \draw[->] (0,-0.3) -- (0,3.5) node[above] {$y$};
  % Isoclines
  \draw[blue, thick, dashed] (0,0) -- (0,3.3)
    node[left, black, font=\small] {$x=0$};
  \draw[blue, thick, dashed] (-0.2,2) -- (4.3,2)
    node[right, black, font=\small] {$y=\alpha/\beta$};
  \draw[red, thick, dashed] (0,0) -- (4.3,0);
  \draw[red, thick, dashed] (2.5,-0.2) -- (2.5,3.3)
    node[above, black, font=\small] {$x=\gamma/\delta$};
  % Point d'équilibre
  \fill (2.5,2) circle (2pt) node[above right] {$E^*$};
  \fill (0,0) circle (2pt) node[below left] {$O$};
  % Flèches directionnelles
  \draw[->, thick, green!50!black] (1.2,1) -- (1.6,1.3);
  \draw[->, thick, green!50!black] (3.5,1) -- (3.8,1.4);
  \draw[->, thick, green!50!black] (3.5,2.8) -- (3.1,2.9);
  \draw[->, thick, green!50!black] (1.2,2.8) -- (0.9,2.5);
\end{tikzpicture}
\end{center}

\subsection{Champ de directions}
\label{subsec:champ-dir}
\index{champ de directions}

L'analyse du champ de directions consiste à tracer en chaque point
$(x_1, x_2)$ un petit vecteur proportionnel à
$\mathbf{f}(x_1, x_2)$. Combinée aux isoclines nulles, cette
technique fournit une description qualitative complète du portrait
de phase.

\subsection{Théorème de Poincaré--Bendixson}
\label{subsec:poincare-bendixson}
\index{Poincaré--Bendixson (théorème)}

Le théorème suivant est l'un des résultats les plus puissants de la
théorie qualitative des EDO en dimension deux.

\begin{theoreme}[Poincaré--Bendixson]
\label{thm:poincare-bendixson}
Soit $\mathbf{f} \in \mathcal{C}^1(\Omega, \R^2)$ et soit
$\mathbf{x}(t)$ une solution définie pour tout $t \geq 0$ dont
l'ensemble $\omega$-limite\index{ensemble oméga-limite}
$\omega(\mathbf{x})$ est un compact non vide contenu dans $\Omega$
et ne contenant qu'un nombre fini de points d'équilibre.
Alors $\omega(\mathbf{x})$ est l'un des ensembles suivants :
\begin{enumerate}[label=(\roman*)]
  \item un point d'équilibre ;
  \item une orbite périodique (cycle) ;
  \item un ensemble formé de points d'équilibre et d'orbites
        homoclines ou hétéroclines les reliant.
\end{enumerate}
\end{theoreme}

\begin{corollaire}
\label{cor:PB-piege}
Si le flot d'un système plan $\mathcal{C}^1$ est piégé dans un
compact $K$ (c'est-à-dire $\varphi(t, \mathbf{x}_0) \in K$ pour
tout $t \geq 0$) et si $K$ ne contient aucun point d'équilibre,
alors $K$ contient au moins une orbite périodique.
\end{corollaire}

\begin{remarque}
Le théorème de Poincaré--Bendixson est \emph{spécifique à la
dimension 2}. En dimension $3$ ou plus, des comportements beaucoup
plus complexes (attracteurs étranges, chaos) sont possibles.
\end{remarque}

\subsection{Cycles limites}
\label{subsec:cycles-limites}
\index{cycle limite}

\begin{definition}[Cycle limite]
\label{def:cycle-limite}
Un \emph{cycle limite}\index{cycle limite!définition} est une orbite
périodique isolée, c'est-à-dire une orbite fermée $\Gamma$ telle qu'il
existe un voisinage de $\Gamma$ ne contenant aucune autre orbite
fermée.

Un cycle limite est dit :
\begin{itemize}
  \item \emph{stable} (attractif) si les trajectoires voisines
        spiralent vers $\Gamma$ quand $t \to +\infty$ ;
  \item \emph{instable} (répulsif) si elles spiralent vers $\Gamma$
        quand $t \to -\infty$ ;
  \item \emph{semi-stable} sinon.
\end{itemize}
\end{definition}

\begin{exemple}[Oscillateur de van der Pol]
\label{ex:van-der-pol-cycle}
\index{van der Pol (oscillateur)}
L'équation de van der Pol
\begin{equation}
\label{eq:van-der-pol}
\ddot{x} - \mu(1 - x^2)\dot{x} + x = 0, \qquad \mu > 0,
\end{equation}
se réécrit comme le système
\[
\begin{cases}
\dot{x} = y, \\
\dot{y} = \mu(1 - x^2)y - x.
\end{cases}
\]
L'unique point d'équilibre est l'origine, qui est un foyer instable
(pour $\mu > 0$). Le théorème de Poincaré--Bendixson permet de
montrer l'existence d'un cycle limite stable entourant l'origine.
\end{exemple}

\begin{center}
\begin{tikzpicture}[scale=0.85]
  \draw[->] (-3.5,0) -- (3.5,0) node[right] {$x$};
  \draw[->] (0,-3.5) -- (0,3.5) node[above] {$y$};
  % Cycle limite stylisé (approximation elliptique)
  \draw[blue, very thick] (0,0) ellipse (2.2 and 2.8);
  \node[blue] at (2.8,2.2) {cycle limite};
  % Spirale intérieure (sortante)
  \draw[red, thick, ->] (0.3,0.1) arc[start angle=10, end angle=350,
    x radius=0.6, y radius=0.8];
  \draw[red, thick, ->] (0.8,0.1) arc[start angle=5, end angle=350,
    x radius=1.2, y radius=1.5];
  % Spirale extérieure (entrante)
  \draw[green!50!black, thick, ->]
    (2.8,0.1) arc[start angle=5, end angle=340,
    x radius=2.8, y radius=3.2];
  \fill (0,0) circle (2pt) node[below right] {$O$};
  \node[red, font=\small] at (0,-1.2) {instable};
  \node[green!50!black, font=\small] at (0,3.3) {entrante};
\end{tikzpicture}
\end{center}

\subsection{Critère de Dulac}
\label{subsec:dulac}
\index{Dulac (critère)}

\begin{theoreme}[Critère de Dulac--Bendixson]
\label{thm:dulac}
Soit $D \subset \R^2$ un domaine simplement connexe et soit
$B : D \to \R$ une fonction de classe $\mathcal{C}^1$ telle que
\[
\frac{\partial(Bf_1)}{\partial x_1}
+ \frac{\partial(Bf_2)}{\partial x_2}
\]
ne change pas de signe dans $D$ et ne soit pas identiquement nulle
sur un ouvert. Alors le système $\dot{\mathbf{x}} = \mathbf{f}(\mathbf{x})$
n'a pas d'orbite périodique entièrement contenue dans $D$.
\end{theoreme}

\begin{proof}
Par l'absurde, supposons l'existence d'une orbite périodique
$\Gamma$ contenue dans $D$ et soit $\Sigma$ le domaine qu'elle
délimite. Par le théorème de Green :
\[
\oint_\Gamma (Bf_1\, \dd x_2 - Bf_2\, \dd x_1)
= \iint_\Sigma \left(
  \frac{\partial(Bf_1)}{\partial x_1}
  + \frac{\partial(Bf_2)}{\partial x_2}
\right) \dd x_1\, \dd x_2.
\]
Le membre de droite est de signe constant et non nul. Or, sur
$\Gamma$, on a $\dd x_i = f_i\, \dd t$, d'où le membre de gauche
vaut $\oint_\Gamma B(f_1 f_2 - f_2 f_1)\, \dd t = 0$. Contradiction.
\end{proof}

\begin{exemple}
\label{ex:dulac-application}
Pour le système $\dot{x} = y + x^3$, $\dot{y} = -x + y^3$,
choisissons $B \equiv 1$. Alors
\[
\frac{\partial f_1}{\partial x} + \frac{\partial f_2}{\partial y}
= 3x^2 + 3y^2 > 0
\]
pour $(x,y) \neq (0,0)$. Par le critère de Dulac, il n'y a aucune
orbite périodique.
\end{exemple}

% ─── 9.3 Théorie des bifurcations ───────────────────────────────
\section{Théorie des bifurcations}
\label{sec:bifurcations}
\index{bifurcation!théorie}

On considère un système dépendant d'un paramètre $\mu \in \R$ :
\begin{equation}
\label{eq:syst-param}
\dot{\mathbf{x}} = \mathbf{f}(\mathbf{x}, \mu).
\end{equation}

\begin{definition}[Stabilité structurelle]
\label{def:stab-struct}
\index{stabilité structurelle}
Un système est dit \emph{structurellement stable} si sa structure
qualitative (nombre et type de points d'équilibre, orbites périodiques)
ne change pas sous de petites perturbations du champ de vecteurs.
\end{definition}

\begin{definition}[Bifurcation]
\label{def:bifurcation}
\index{bifurcation!définition}
Une \emph{bifurcation} se produit en une valeur $\mu = \mu_0$ du
paramètre si le portrait de phase du système change qualitativement
lorsque $\mu$ traverse $\mu_0$. La valeur $\mu_0$ est appelée
\emph{valeur de bifurcation}.
\end{definition}

\subsection{Bifurcation n\oe ud-col}
\label{subsec:bif-noeud-col}
\index{bifurcation!noeud-col}

La forme normale unidimensionnelle est :
\begin{equation}
\label{eq:bif-noeud-col}
\dot{x} = \mu + x^2.
\end{equation}

\begin{itemize}
  \item Pour $\mu < 0$ : deux points d'équilibre $x^* = \pm\sqrt{-\mu}$
        (un stable, un instable).
  \item Pour $\mu = 0$ : un unique point d'équilibre $x^* = 0$
        (semi-stable).
  \item Pour $\mu > 0$ : aucun point d'équilibre.
\end{itemize}

\begin{center}
\begin{tikzpicture}[scale=1.2]
  \draw[->] (-2.5,0) -- (1.5,0) node[right] {$\mu$};
  \draw[->] (0,-1.8) -- (0,1.8) node[above] {$x^*$};
  % Branche stable (trait plein)
  \draw[blue, very thick, domain=-2.2:0, samples=40]
    plot (\x, {-sqrt(-\x)});
  % Branche instable (tirets)
  \draw[red, very thick, dashed, domain=-2.2:0, samples=40]
    plot (\x, {sqrt(-\x)});
  \fill (0,0) circle (2pt);
  \node[below right] at (0.2,-0.2) {$\mu_0 = 0$};
  \node[blue, font=\small] at (-1.5,-1.6) {stable};
  \node[red, font=\small] at (-1.5,1.6) {instable};
  \node[font=\bfseries] at (0,2.2) {Bifurcation n\oe ud-col};
\end{tikzpicture}
\end{center}

\subsection{Bifurcation transcritique}
\label{subsec:bif-transcritique}
\index{bifurcation!transcritique}

La forme normale est :
\begin{equation}
\label{eq:bif-transcritique}
\dot{x} = \mu x - x^2.
\end{equation}

Les points d'équilibre sont $x^* = 0$ et $x^* = \mu$. Pour $\mu < 0$,
$x^* = 0$ est stable et $x^* = \mu$ instable ; pour $\mu > 0$,
les rôles s'échangent. Il y a échange de stabilité en $\mu = 0$.

\begin{center}
\begin{tikzpicture}[scale=1.2]
  \draw[->] (-2.2,0) -- (2.2,0) node[right] {$\mu$};
  \draw[->] (0,-1.8) -- (0,1.8) node[above] {$x^*$};
  % Branche x*=0
  \draw[blue, very thick] (-2,0) -- (0,0);
  \draw[red, very thick, dashed] (0,0) -- (2,0);
  % Branche x*=mu
  \draw[red, very thick, dashed] (-2,-2) -- (0,0);
  \draw[blue, very thick] (0,0) -- (1.8,1.8);
  \fill (0,0) circle (2pt);
  \node[font=\bfseries] at (0,2.2) {Bifurcation transcritique};
\end{tikzpicture}
\end{center}

\subsection{Bifurcation fourche}
\label{subsec:bif-fourche}
\index{bifurcation!fourche (pitchfork)}

\paragraph{Fourche supercritique.}
Forme normale : $\dot{x} = \mu x - x^3$.

\begin{center}
\begin{tikzpicture}[scale=1.2]
  \draw[->] (-1.5,0) -- (2.5,0) node[right] {$\mu$};
  \draw[->] (0,-1.8) -- (0,1.8) node[above] {$x^*$};
  % x*=0 stable puis instable
  \draw[blue, very thick] (-1.3,0) -- (0,0);
  \draw[red, very thick, dashed] (0,0) -- (2.2,0);
  % Branches x*=±sqrt(mu) stables
  \draw[blue, very thick, domain=0.01:2.2, samples=40]
    plot (\x, {sqrt(\x)});
  \draw[blue, very thick, domain=0.01:2.2, samples=40]
    plot (\x, {-sqrt(\x)});
  \fill (0,0) circle (2pt);
  \node[font=\bfseries] at (0.5,2.2) {Fourche supercritique};
\end{tikzpicture}
\end{center}

\paragraph{Fourche sous-critique.}
Forme normale : $\dot{x} = \mu x + x^3$.

\begin{center}
\begin{tikzpicture}[scale=1.2]
  \draw[->] (-2.5,0) -- (1.5,0) node[right] {$\mu$};
  \draw[->] (0,-1.8) -- (0,1.8) node[above] {$x^*$};
  % x*=0 instable puis stable
  \draw[red, very thick, dashed] (-2.2,0) -- (0,0);
  \draw[blue, very thick] (0,0) -- (1.3,0);
  % Branches instables x*=±sqrt(-mu)
  \draw[red, very thick, dashed, domain=-2.2:-0.01, samples=40]
    plot (\x, {sqrt(-\x)});
  \draw[red, very thick, dashed, domain=-2.2:-0.01, samples=40]
    plot (\x, {-sqrt(-\x)});
  \fill (0,0) circle (2pt);
  \node[font=\bfseries] at (-0.5,2.2) {Fourche sous-critique};
\end{tikzpicture}
\end{center}

\subsection{Bifurcation de Hopf}
\label{subsec:bif-hopf}
\index{bifurcation!Hopf}

\begin{theoreme}[Bifurcation de Hopf]
\label{thm:hopf}
Soit $\mathbf{x}^*(\mu)$ une famille de points d'équilibre du
système $\dot{\mathbf{x}} = \mathbf{f}(\mathbf{x}, \mu)$ dans
$\R^2$. Supposons que la matrice jacobienne
$D_\mathbf{x}\mathbf{f}(\mathbf{x}^*(\mu), \mu)$ possède une paire
de valeurs propres conjuguées $\lambda(\mu) = \alpha(\mu) \pm i\beta(\mu)$
vérifiant :
\begin{enumerate}[label=(\roman*)]
  \item $\alpha(\mu_0) = 0$ et $\beta(\mu_0) \neq 0$ (paire
        purement imaginaire en $\mu = \mu_0$) ;
  \item $\alpha'(\mu_0) \neq 0$ (condition de transversalité).
\end{enumerate}
Alors il existe un voisinage de $(\mathbf{x}^*(\mu_0), \mu_0)$ dans
lequel naît (ou disparaît) une famille d'orbites périodiques dont la
période tend vers $2\pi/\abs{\beta(\mu_0)}$ quand $\mu \to \mu_0$.

La bifurcation est dite :
\begin{itemize}
  \item \emph{supercritique} si un cycle limite \emph{stable} naît
        du côté où l'équilibre devient instable ;
  \item \emph{sous-critique} si un cycle limite \emph{instable} naît
        du côté où l'équilibre est encore stable.
\end{itemize}
\end{theoreme}

\begin{exemple}[Bifurcation de Hopf --- forme normale]
\label{ex:hopf-normale}
Le système en coordonnées polaires
\[
\dot{r} = \mu r - r^3, \qquad \dot{\theta} = 1,
\]
possède l'origine comme équilibre pour tout $\mu$. Pour $\mu > 0$,
un cycle limite stable $r = \sqrt{\mu}$ apparaît. C'est une bifurcation
de Hopf supercritique.
\end{exemple}

\begin{center}
\begin{tikzpicture}[scale=1.0]
  % Diagramme de bifurcation de Hopf
  \draw[->] (-2.5,0) -- (2.5,0) node[right] {$\mu$};
  \draw[->] (0,-0.3) -- (0,2.5) node[above] {$r$};
  % Equilibre stable
  \draw[blue, very thick] (-2.2,0) -- (0,0);
  % Equilibre instable
  \draw[red, very thick, dashed] (0,0) -- (2.2,0);
  % Cycle limite (amplitude = sqrt(mu))
  \draw[blue, very thick, domain=0.01:2.2, samples=40]
    plot (\x, {sqrt(\x)});
  \fill (0,0) circle (2pt);
  \node[font=\bfseries] at (0,2.9) {Bifurcation de Hopf supercritique};
  \node[blue, font=\small] at (2,1.8) {$r = \sqrt{\mu}$};
\end{tikzpicture}
\end{center}

% ─── 9.4 Théorie de l'indice ────────────────────────────────────
\section{Théorie de l'indice (introduction)}
\label{sec:indice}
\index{indice d'un point singulier}

\begin{definition}[Indice d'un point singulier]
\label{def:indice}
Soit $\mathbf{x}^*$ un point d'équilibre isolé et $\Gamma$ une
courbe fermée simple orientée positivement entourant $\mathbf{x}^*$
et ne passant par aucun autre point d'équilibre. L'\emph{indice} de
$\mathbf{x}^*$ est le nombre de tours que fait le vecteur
$\mathbf{f}(\mathbf{x})$ lorsque $\mathbf{x}$ parcourt $\Gamma$
une fois :
\[
I_{\mathbf{x}^*} = \frac{1}{2\pi}
\oint_\Gamma \dd\left[\arctan\frac{f_2}{f_1}\right].
\]
\end{definition}

\begin{proposition}[Propriétés de l'indice]
\label{prop:indice-prop}
\begin{enumerate}[label=(\roman*)]
  \item L'indice d'un n\oe ud, d'un foyer ou d'un centre est $+1$.
  \item L'indice d'un col est $-1$.
  \item Si une courbe fermée simple ne contient aucun point
        d'équilibre, son indice est $0$.
  \item Si une orbite périodique $\Gamma$ existe, alors la somme des
        indices des points d'équilibre à l'intérieur de $\Gamma$
        vaut $+1$.
\end{enumerate}
\end{proposition}

\begin{corollaire}
\label{cor:indice-cycle}
Toute orbite périodique entoure au moins un point d'équilibre.
Si elle n'entoure qu'un seul point d'équilibre, celui-ci est
nécessairement un n\oe ud, un foyer ou un centre (indice $+1$).
\end{corollaire}

% ─── 9.5 Exemples développés ────────────────────────────────────
\section{Exemples développés}
\label{sec:exemples-nonlin}

\subsection{Système de Lotka--Volterra}
\label{subsec:lotka-volterra-detail}
\index{Lotka--Volterra (système)}

Considérons le système classique avec $\alpha = \beta = \delta = \gamma = 1$ :
\begin{equation}
\label{eq:LV-classique}
\dot{x} = x(1 - y), \qquad \dot{y} = y(x - 1).
\end{equation}

\paragraph{Points d'équilibre.}
$(0,0)$ et $(1,1)$.

\paragraph{Linéarisation en $(1,1)$.}
La matrice jacobienne est
\[
J = \begin{pmatrix}
1 - y & -x \\ y & x - 1
\end{pmatrix}\bigg|_{(1,1)}
= \begin{pmatrix} 0 & -1 \\ 1 & 0 \end{pmatrix}.
\]
Les valeurs propres sont $\lambda = \pm i$ : le point $(1,1)$ est un
centre pour le système linéarisé. On montre que c'est aussi un centre
pour le système non-linéaire grâce à l'intégrale première
\begin{equation}
\label{eq:LV-integrale}
H(x,y) = x - \ln x + y - \ln y,
\end{equation}
qui est constante le long des trajectoires.

\paragraph{Portrait de phase.}
Les trajectoires sont des courbes fermées autour de $(1,1)$
dans le quadrant $x > 0$, $y > 0$.

\begin{center}
\begin{tikzpicture}[scale=1.3]
  \draw[->] (-0.3,0) -- (4,0) node[right] {$x$ (proies)};
  \draw[->] (0,-0.3) -- (0,3.5) node[above] {$y$ (prédateurs)};
  % Courbes de niveau de H (ellipses déformées, schématiques)
  \draw[blue, thick] (1,1) ellipse (0.4 and 0.35);
  \draw[blue, thick] (1,1) ellipse (0.8 and 0.65);
  \draw[blue, thick] (1,1) ellipse (1.4 and 1.05);
  \draw[blue, thick] (1,1) ellipse (2.1 and 1.5);
  % Flèches de direction
  \draw[->, thick, red] (1,0.3) -- (1.3,0.3);
  \draw[->, thick, red] (2.2,1) -- (2.2,1.3);
  \draw[->, thick, red] (1,2.1) -- (0.7,2.1);
  \draw[->, thick, red] (0.3,1) -- (0.3,0.7);
  \fill (1,1) circle (1.5pt) node[below right, font=\small] {$(1,1)$};
  \node at (1,1) [above left, font=\small] {centre};
\end{tikzpicture}
\end{center}

\subsection{Oscillateur de van der Pol}
\label{subsec:van-der-pol-detail}
\index{van der Pol (oscillateur)!analyse détaillée}

Reprenons l'équation \eqref{eq:van-der-pol} avec $\mu = 1$ :
\[
\dot{x} = y, \qquad \dot{y} = (1 - x^2)y - x.
\]

\paragraph{Linéarisation en $(0,0)$.}
\[
J(0,0) = \begin{pmatrix} 0 & 1 \\ -1 & 1 \end{pmatrix},
\qquad \lambda = \frac{1 \pm i\sqrt{3}}{2}.
\]
Les parties réelles sont positives : l'origine est un foyer instable.

\paragraph{Existence du cycle limite.}
On construit un anneau de piégeage\index{anneau de piégeage}
$K = \{\mathbf{x} : r_1 \leq \norm{\mathbf{x}} \leq r_2\}$
tel que le champ pointe vers l'intérieur sur les deux cercles
frontières. Par le corollaire~\ref{cor:PB-piege} du théorème de
Poincaré--Bendixson, $K$ contient un cycle limite.

\subsection{Modèle SIR}
\label{subsec:SIR}
\index{SIR (modèle)}

Le modèle SIR de Kermack--McKendrick\index{Kermack--McKendrick} pour
la propagation d'une épidémie s'écrit :
\begin{equation}
\label{eq:SIR}
\dot{S} = -\beta SI, \qquad
\dot{I} = \beta SI - \gamma I, \qquad
\dot{R} = \gamma I,
\end{equation}
où $S$, $I$, $R$ désignent les proportions de susceptibles, infectés
et rétablis, avec $S + I + R = 1$.

Puisque $R = 1 - S - I$, on étudie le système plan :
\[
\dot{S} = -\beta SI, \qquad
\dot{I} = (\beta S - \gamma) I.
\]

\paragraph{Nombre de reproduction de base.}
\index{nombre de reproduction de base}
Le \emph{nombre de reproduction de base} est
$\mathcal{R}_0 = \beta / \gamma$.
L'épidémie se propage si et seulement si $\mathcal{R}_0 S(0) > 1$.

\paragraph{Analyse qualitative.}
Les isoclines nulles sont $I = 0$ (axe des $S$) et
$S = \gamma/\beta = 1/\mathcal{R}_0$. Les trajectoires dans le
plan $(S, I)$ satisfont
\[
\frac{\dd I}{\dd S} = -1 + \frac{1}{\mathcal{R}_0 S},
\]
d'où $I + S - \frac{1}{\mathcal{R}_0}\ln S = \text{const}$.
Les courbes de phase montent (l'épidémie croît) tant que
$S > 1/\mathcal{R}_0$, puis décroissent.

% ─── 9.6 Exercices ──────────────────────────────────────────────
\section{Exercices}
\label{sec:exo-chap9}

\begin{exercice}
\label{exo:isoclines}
Pour le système $\dot{x} = x(3 - x - 2y)$,
$\dot{y} = y(2 - x - y)$, tracer les isoclines nulles, trouver tous
les points d'équilibre et déterminer leur nature par linéarisation.
\end{exercice}

\begin{exercice}
\label{exo:dulac-exo}
Montrer, à l'aide du critère de Dulac avec $B(x,y) = e^{-2x}$,
que le système $\dot{x} = y$, $\dot{y} = -x - y + x^2 + y^2$
n'admet pas d'orbite périodique dans $\R^2$.
\end{exercice}

\begin{exercice}
\label{exo:hopf-exo}
Considérer le système
\[
\dot{x} = \mu x - y - x(x^2 + y^2), \qquad
\dot{y} = x + \mu y - y(x^2 + y^2).
\]
\begin{enumerate}[label=\alph*)]
  \item Passer en coordonnées polaires et montrer que l'équation
        en $r$ admet une bifurcation de Hopf en $\mu = 0$.
  \item Déterminer si la bifurcation est supercritique ou
        sous-critique.
  \item Tracer le diagramme de bifurcation.
\end{enumerate}
\end{exercice}

\begin{exercice}
\label{exo:LV-generique}
Pour le système de Lotka--Volterra généralisé
\[
\dot{x} = x(\alpha - \beta y), \qquad \dot{y} = y(-\gamma + \delta x),
\]
avec $\alpha, \beta, \gamma, \delta > 0$ :
\begin{enumerate}[label=\alph*)]
  \item Montrer que $H(x,y) = \delta x - \gamma \ln x
        + \beta y - \alpha \ln y$ est une intégrale première dans le
        quadrant $x > 0$, $y > 0$.
  \item En déduire que toutes les trajectoires sont fermées autour
        du point d'équilibre intérieur.
\end{enumerate}
\end{exercice}

\begin{exercice}
\label{exo:bif-noeud-col-exo}
Étudier les bifurcations du système
$\dot{x} = \mu - x^2$, $\dot{y} = -y$,
en fonction du paramètre $\mu \in \R$. Tracer les portraits de phase
pour $\mu < 0$, $\mu = 0$ et $\mu > 0$.
\end{exercice}

\begin{exercice}
\label{exo:indice-exo}
Soit $\dot{x} = -y + x(x^2 + y^2 - 1)$,
$\dot{y} = x + y(x^2 + y^2 - 1)$.
\begin{enumerate}[label=\alph*)]
  \item Montrer que le cercle unité est une orbite périodique.
  \item Calculer l'indice de l'origine.
  \item Déterminer la stabilité du cycle.
\end{enumerate}
\end{exercice}

% ─── 9.7 Résumé ─────────────────────────────────────────────────
\section*{Résumé du chapitre}

\begin{itemize}
  \item Les \textbf{isoclines nulles} permettent de localiser les
        points d'équilibre et de déterminer la direction qualitative
        du flot.
  \item Le \textbf{théorème de Poincaré--Bendixson} garantit, en
        dimension 2, qu'un flot piégé dans un compact sans point
        d'équilibre engendre un cycle limite.
  \item Le \textbf{critère de Dulac} permet d'exclure l'existence
        d'orbites périodiques.
  \item Les \textbf{bifurcations} classiques (n\oe ud-col,
        transcritique, fourche, Hopf) décrivent les changements
        qualitatifs du portrait de phase en fonction d'un paramètre.
  \item La \textbf{théorie de l'indice} fournit des contraintes
        topologiques sur la disposition des points d'équilibre et
        des orbites périodiques.
\end{itemize}


%% ═══════════════════════════════════════════════════════════════════
%%  CHAPITRE 10 — Applications
%% ═══════════════════════════════════════════════════════════════════
\chapter{Applications --- Mécanique, Électricité, Biologie, Démographie}
\label{chap:applications}
\index{applications des EDO}

\section*{Introduction}

Ce chapitre illustre la puissance des méthodes développées tout au
long de ce cours en les appliquant à des problèmes concrets issus
de la mécanique, de l'électrotechnique, de la biologie et de la
démographie. L'accent est mis sur la \emph{modélisation} ---
c'est-à-dire la traduction d'un problème physique ou biologique
en une équation différentielle --- et sur l'\emph{interprétation}
des solutions obtenues.

% ─── 10.1 Oscillations mécaniques ───────────────────────────────
\section{Oscillations mécaniques}
\label{sec:oscillations}
\index{oscillations mécaniques}

\subsection{Le pendule simple}
\label{subsec:pendule-simple}
\index{pendule!simple}

L'équation du pendule simple de longueur $\ell$ dans un champ de
gravité $g$ est :
\begin{equation}
\label{eq:pendule-simple}
\ddot{\theta} + \frac{g}{\ell}\sin\theta = 0.
\end{equation}

\begin{remarque}
Pour de petites oscillations ($\abs{\theta} \ll 1$), l'approximation
$\sin\theta \approx \theta$ donne l'oscillateur harmonique
$\ddot{\theta} + \omega_0^2 \theta = 0$ avec
$\omega_0 = \sqrt{g/\ell}$.
\end{remarque}

\paragraph{Analyse dans le plan de phase.}
En posant $x = \theta$, $y = \dot{\theta}$, on obtient le système
\[
\dot{x} = y, \qquad \dot{y} = -\omega_0^2 \sin x.
\]
L'énergie $E = \frac{1}{2}y^2 - \omega_0^2 \cos x$ est une intégrale
première. Le portrait de phase présente :
\begin{itemize}
  \item des centres en $(2k\pi, 0)$, $k \in \Z$ (oscillations) ;
  \item des cols en $((2k+1)\pi, 0)$ (positions instables) ;
  \item des \emph{séparatrices}\index{séparatrice} reliant les cols
        (mouvement apériodique limite).
\end{itemize}

\begin{center}
\begin{tikzpicture}[scale=0.75]
  \draw[->] (-4.5,0) -- (4.5,0) node[right] {$\theta$};
  \draw[->] (0,-2.8) -- (0,2.8) node[above] {$\dot{\theta}$};
  % Orbites fermées autour de l'origine
  \foreach \a in {0.3,0.7,1.1} {
    \draw[blue, thick] (0,0) ellipse ({\a*2.5} and {\a*1.8});
  }
  % Séparatrices (schématiques)
  \draw[red, very thick, domain=-3.14:3.14, samples=50]
    plot (\x, {1.8*cos(\x/2 r)});
  \draw[red, very thick, domain=-3.14:3.14, samples=50]
    plot (\x, {-1.8*cos(\x/2 r)});
  % Points
  \fill (0,0) circle (2pt) node[below right, font=\small] {centre};
  \fill (3.14,0) circle (2pt) node[below, font=\small] {col};
  \fill (-3.14,0) circle (2pt) node[below, font=\small] {col};
  \node at (0,-3.3) {Portrait de phase du pendule};
\end{tikzpicture}
\end{center}

\subsection{Pendule amorti et forcé}
\label{subsec:pendule-amorti}
\index{pendule!amorti}
\index{pendule!forcé}

Avec amortissement visqueux de coefficient $c$ et forçage harmonique :
\begin{equation}
\label{eq:pendule-force}
\ddot{\theta} + \frac{c}{m\ell}\dot{\theta}
+ \frac{g}{\ell}\sin\theta = \frac{F_0}{m\ell}\cos(\omega t).
\end{equation}

Dans le régime linéarisé, on retrouve l'oscillateur harmonique amorti
forcé, dont la solution en régime permanent est
\[
\theta_p(t) = A(\omega)\cos(\omega t - \varphi(\omega)),
\]
avec l'amplitude de résonance\index{résonance}
\[
A(\omega) = \frac{F_0/(m\ell)}
{\sqrt{(\omega_0^2 - \omega^2)^2 + (c\omega/(m\ell))^2}}.
\]

\subsection{Oscillateurs couplés}
\label{subsec:oscillateurs-couples}
\index{oscillateurs couplés}

Deux masses $m$ reliées par des ressorts de raideur $k$ (aux murs)
et $k_c$ (couplage) conduisent au système :
\begin{equation}
\label{eq:oscillateurs-couples}
\begin{cases}
m\ddot{x}_1 = -kx_1 - k_c(x_1 - x_2), \\
m\ddot{x}_2 = -kx_2 - k_c(x_2 - x_1).
\end{cases}
\end{equation}

Les \emph{modes normaux}\index{modes normaux} s'obtiennent par le
changement de variables $q_1 = x_1 + x_2$, $q_2 = x_1 - x_2$ :
\[
\ddot{q}_1 + \omega_1^2 q_1 = 0, \qquad
\ddot{q}_2 + \omega_2^2 q_2 = 0,
\]
avec $\omega_1 = \sqrt{k/m}$ et
$\omega_2 = \sqrt{(k + 2k_c)/m}$.

\subsection{Résonance et battements}
\label{subsec:resonance-battements}
\index{résonance}
\index{battements}

Lorsque la fréquence d'excitation $\omega$ est proche de la fréquence
propre $\omega_0$, la superposition des deux modes donne le phénomène
de \emph{battements} :
\[
x(t) = \frac{2F_0}{m(\omega_0^2 - \omega^2)}
\sin\!\left(\frac{\omega_0 - \omega}{2}t\right)
\cos\!\left(\frac{\omega_0 + \omega}{2}t\right).
\]
L'enveloppe oscille à la fréquence lente
$\abs{\omega_0 - \omega}/2$.

\begin{center}
\begin{tikzpicture}[scale=0.8]
  \draw[->] (-0.5,0) -- (10,0) node[right] {$t$};
  \draw[->] (0,-2.2) -- (0,2.2) node[above] {$x(t)$};
  % Battements
  \draw[blue, thick, domain=0:9.5, samples=50]
    plot (\x, {1.8*sin(30*\x)*cos(4*\x r)});
  % Enveloppe
  \draw[red, dashed, thick, domain=0:9.5, samples=50]
    plot (\x, {1.8*sin(30*\x)});
  \draw[red, dashed, thick, domain=0:9.5, samples=50]
    plot (\x, {-1.8*sin(30*\x)});
  \node[red, font=\small] at (7,2) {enveloppe};
\end{tikzpicture}
\end{center}

% ─── 10.2 Circuits électriques ──────────────────────────────────
\section{Circuits électriques}
\label{sec:circuits}
\index{circuits électriques}

\subsection{Circuit RC}
\label{subsec:circuit-RC}
\index{circuit RC}

Un condensateur $C$ en série avec une résistance $R$ soumis à une
tension $E(t)$ obéit à :
\begin{equation}
\label{eq:RC}
RC\,\dot{v}_C + v_C = E(t),
\end{equation}
où $v_C$ est la tension aux bornes du condensateur. C'est une EDO
linéaire du premier ordre de constante de temps $\tau = RC$.

\begin{exemple}[Charge d'un condensateur]
\label{ex:charge-RC}
Pour $E(t) = E_0$ (échelon) et $v_C(0) = 0$ :
\[
v_C(t) = E_0(1 - e^{-t/\tau}).
\]
Le condensateur se charge exponentiellement avec la constante de
temps $\tau = RC$.
\end{exemple}

\subsection{Circuit RL}
\label{subsec:circuit-RL}
\index{circuit RL}

Une inductance $L$ en série avec une résistance $R$ :
\begin{equation}
\label{eq:RL}
L\,\dot{i} + Ri = E(t).
\end{equation}
Constante de temps : $\tau = L/R$.

\subsection{Circuit RLC série}
\label{subsec:circuit-RLC}
\index{circuit RLC}

\begin{equation}
\label{eq:RLC}
L\ddot{q} + R\dot{q} + \frac{q}{C} = E(t),
\end{equation}
où $q$ est la charge du condensateur. C'est l'analogue électrique
exact de l'oscillateur mécanique amorti forcé avec les
correspondances :
\[
m \leftrightarrow L, \quad c \leftrightarrow R, \quad
k \leftrightarrow 1/C, \quad F(t) \leftrightarrow E(t).
\]

\paragraph{Régimes de fonctionnement.}
\index{régime!apériodique}
\index{régime!critique}
\index{régime!pseudo-périodique}
Posons $\omega_0 = 1/\sqrt{LC}$ et $\xi = R/(2\sqrt{L/C})$
(facteur d'amortissement). L'équation caractéristique
$Lr^2 + Rr + 1/C = 0$ donne :
\begin{itemize}
  \item $\xi < 1$ : régime pseudo-périodique
        ($q(t) = e^{-\xi\omega_0 t}[A\cos\omega_d t + B\sin\omega_d t]$
        avec $\omega_d = \omega_0\sqrt{1-\xi^2}$) ;
  \item $\xi = 1$ : régime critique
        ($q(t) = (A + Bt)e^{-\omega_0 t}$) ;
  \item $\xi > 1$ : régime apériodique.
\end{itemize}

\subsection{Fonctions de transfert et lien avec Laplace}
\label{subsec:transfert}
\index{fonction de transfert}

En appliquant la transformée de Laplace à \eqref{eq:RLC} avec
conditions initiales nulles :
\[
H(s) = \frac{Q(s)}{E(s)} = \frac{1}{Ls^2 + Rs + 1/C}
= \frac{1/L}{s^2 + (R/L)s + 1/(LC)}.
\]
Les pôles de $H(s)$ correspondent aux valeurs propres du système
autonome associé. La réponse fréquentielle s'obtient en évaluant
$H(i\omega)$.

\begin{center}
\begin{tikzpicture}[scale=0.9]
  % Circuit RLC schématique
  \draw (0,0) -- (1,0);
  % Résistance
  \draw (1,0) -- (1.3,0.3) -- (1.6,-0.3) -- (1.9,0.3)
        -- (2.2,-0.3) -- (2.5,0.3) -- (2.8,0) -- (3,0);
  \node[above] at (1.9,0.4) {$R$};
  % Inductance
  \draw (3,0) -- (3.2,0);
  \foreach \x in {3.4,3.8,4.2,4.6} {
    \draw (\x-0.2,0) arc (180:0:0.2);
  }
  \draw (4.8,0) -- (5,0);
  \node[above] at (4.1,0.3) {$L$};
  % Condensateur
  \draw (5,0) -- (5.5,0);
  \draw (5.5,-0.5) -- (5.5,0.5);
  \draw (5.8,-0.5) -- (5.8,0.5);
  \draw (5.8,0) -- (6.3,0);
  \node[above] at (5.65,0.5) {$C$};
  % Retour
  \draw (6.3,0) -- (6.3,-1.5) -- (0,-1.5) -- (0,0);
  % Source
  \draw (0,-0.3) -- (0,-1.2);
  \node[left] at (0,-0.75) {$E(t)$};
  \draw[thick] (-0.15,-0.3) -- (0.15,-0.3);
  \draw[thick] (-0.08,-1.2) -- (0.08,-1.2);
\end{tikzpicture}
\end{center}

% ─── 10.3 Dynamique des populations ─────────────────────────────
\section{Dynamique des populations}
\label{sec:populations}
\index{dynamique des populations}

\subsection{Modèle de Malthus}
\label{subsec:malthus}
\index{Malthus (modèle)}

Le modèle le plus simple de croissance d'une population est :
\begin{equation}
\label{eq:malthus}
\dot{N} = rN, \qquad N(0) = N_0,
\end{equation}
où $r$ est le taux de croissance intrinsèque. La solution est
$N(t) = N_0 e^{rt}$ : croissance exponentielle si $r > 0$,
décroissance si $r < 0$.

\begin{remarque}
Le modèle de Malthus est irréaliste à long terme car il prédit une
croissance illimitée. Il est néanmoins valable sur de courtes
périodes pour des populations en phase de croissance initiale.
\end{remarque}

\subsection{Équation logistique de Verhulst}
\label{subsec:logistique}
\index{Verhulst (équation logistique)}
\index{équation logistique}

Pour tenir compte de la capacité limite $K$ du milieu :
\begin{equation}
\label{eq:logistique}
\dot{N} = rN\!\left(1 - \frac{N}{K}\right).
\end{equation}

\begin{theoreme}[Solution de l'équation logistique]
\label{thm:sol-logistique}
La solution du problème de Cauchy \eqref{eq:logistique} avec
$N(0) = N_0 > 0$ est
\[
N(t) = \frac{KN_0}{N_0 + (K - N_0)e^{-rt}}.
\]
Elle vérifie $\lim_{t \to +\infty} N(t) = K$.
\end{theoreme}

\begin{proof}
L'équation est à variables séparables. On écrit
\[
\frac{\dd N}{N(1 - N/K)} = r\, \dd t.
\]
Décomposition en éléments simples :
$\frac{1}{N(1-N/K)} = \frac{1}{N} + \frac{1}{K-N}$.
L'intégration donne $\ln\frac{N}{K-N} = rt + C$, d'où le résultat.
\end{proof}

\begin{center}
\begin{tikzpicture}[scale=0.85]
  \draw[->] (-0.5,0) -- (8,0) node[right] {$t$};
  \draw[->] (0,-0.3) -- (0,3.5) node[above] {$N$};
  \draw[dashed, gray] (0,3) -- (7.5,3) node[right] {$K$};
  % Courbe logistique (N0 petit)
  \draw[blue, very thick, domain=0:7.5, samples=50]
    plot (\x, {3*0.2/(0.2 + 0.8*exp(-0.8*\x))});
  % Courbe logistique (N0 > K)
  \draw[red, thick, domain=0:7.5, samples=50]
    plot (\x, {3*1.5/(1.5 + (-0.5)*exp(-0.8*\x))});
  \node[blue, font=\small] at (2,0.5) {$N_0 < K$};
  \node[red, font=\small] at (2,3.4) {$N_0 > K$};
  \node[left] at (0,3) {$K$};
\end{tikzpicture}
\end{center}

\subsection{Modèle de Lotka--Volterra proie--prédateur}
\label{subsec:LV-application}
\index{Lotka--Volterra (système)!application}

Reprenons le modèle de l'exemple~\ref{ex:motivation-nonlin} avec
des valeurs numériques : $\alpha = 1$, $\beta = 0{,}5$,
$\gamma = 0{,}75$, $\delta = 0{,}25$.

\begin{methode}[Analyse complète d'un système proie--prédateur]
\label{meth:LV}
\begin{enumerate}
  \item \textbf{Points d'équilibre.} $E_0 = (0,0)$ et
        $E^* = (\gamma/\delta, \alpha/\beta) = (3, 2)$.
  \item \textbf{Linéarisation en $E_0$.}
        $J(0,0) = \begin{psmallmatrix} 1 & 0 \\ 0 & -0{,}75
        \end{psmallmatrix}$ : col (instable).
  \item \textbf{Linéarisation en $E^*$.}
        $J(3,2) = \begin{psmallmatrix} 0 & -1{,}5 \\ 0{,}5 & 0
        \end{psmallmatrix}$ : valeurs propres $\pm i\sqrt{0{,}75}$,
        centre.
  \item \textbf{Intégrale première.}
        $H(x,y) = 0{,}25 x - 0{,}75\ln x + 0{,}5 y - \ln y = C$.
  \item \textbf{Conclusion.} Les trajectoires dans $\{x>0, y>0\}$
        sont des courbes fermées : les populations oscillent
        périodiquement.
\end{enumerate}
\end{methode}

\subsection{Espèces en compétition}
\label{subsec:competition}
\index{espèces en compétition}

Deux espèces exploitant la même ressource obéissent au modèle de
Lotka--Volterra compétitif :
\begin{equation}
\label{eq:competition}
\begin{cases}
\dot{x}_1 = r_1 x_1\!\left(1 - \dfrac{x_1 + \alpha_{12}x_2}{K_1}\right), \\[8pt]
\dot{x}_2 = r_2 x_2\!\left(1 - \dfrac{x_2 + \alpha_{21}x_1}{K_2}\right),
\end{cases}
\end{equation}
où $\alpha_{12}$ et $\alpha_{21}$ mesurent l'intensité de la
compétition interspécifique.

\begin{proposition}[Issue de la compétition]
\label{prop:competition}
\index{principe d'exclusion compétitive}
L'issue dépend des positions relatives des isoclines nulles :
\begin{enumerate}[label=(\roman*)]
  \item $K_1 > K_2/\alpha_{21}$ et $K_2 < K_1/\alpha_{12}$ :
        l'espèce 1 exclut l'espèce 2.
  \item $K_2 > K_1/\alpha_{12}$ et $K_1 < K_2/\alpha_{21}$ :
        l'espèce 2 exclut l'espèce 1.
  \item $K_1 > K_2/\alpha_{21}$ et $K_2 > K_1/\alpha_{12}$ :
        bistabilité (le résultat dépend des conditions initiales).
  \item $K_1 < K_2/\alpha_{21}$ et $K_2 < K_1/\alpha_{12}$ :
        coexistence stable.
\end{enumerate}
\end{proposition}

\subsection{Modèle SIR épidémiologique}
\label{subsec:SIR-app}
\index{SIR (modèle)!application}

Reprenons le modèle \eqref{eq:SIR} de la section~\ref{subsec:SIR}
avec des valeurs numériques pour illustrer la dynamique d'une
épidémie.

\begin{exemple}[Épidémie de grippe]
\label{ex:grippe-SIR}
Avec $\beta = 0{,}4$ (jour$^{-1}$), $\gamma = 0{,}1$ (jour$^{-1}$),
on a $\mathcal{R}_0 = 4$. Pour $S(0) = 0{,}99$, $I(0) = 0{,}01$,
$R(0) = 0$ :
\begin{itemize}
  \item La condition $\mathcal{R}_0 S(0) = 3{,}96 > 1$ est vérifiée :
        l'épidémie se propage.
  \item Le pic épidémique se produit quand
        $S = 1/\mathcal{R}_0 = 0{,}25$.
  \item La proportion finale de susceptibles $S_\infty$ satisfait
        l'équation transcendante
        $S_\infty = S_0 e^{-\mathcal{R}_0(1-S_\infty)}$.
\end{itemize}
\end{exemple}

\begin{center}
\begin{tikzpicture}[scale=0.85]
  \draw[->] (-0.3,0) -- (1.2,0) node[right] {$S$};
  \draw[->] (0,-0.3) -- (0,1.2) node[above] {$I$};
  % Trajectoire schématique dans le plan (S,I)
  \draw[blue, very thick, domain=0.99:0.08, samples=50]
    plot (\x, {0.01 + \x - 0.99 + 0.25*ln(\x/0.99)});
  % Isocline dI/dt = 0
  \draw[red, dashed, thick] (0.25,-0.2) -- (0.25,1.1)
    node[above, font=\small] {$S = 1/\mathcal{R}_0$};
  \draw[->, thick, green!50!black] (0.7,0.15) -- (0.5,0.3);
  \node[font=\small] at (0.6,-0.4)
    {Trajectoire dans le plan $(S,I)$};
\end{tikzpicture}
\end{center}

% ─── 10.4 Cinétique chimique ────────────────────────────────────
\section{Cinétique chimique}
\label{sec:cinetique}
\index{cinétique chimique}

\subsection{Réactions du premier et du second ordre}
\label{subsec:reactions-ordre}

\paragraph{Réaction du premier ordre.}
\index{réaction!premier ordre}
$A \to \text{produits}$, loi cinétique :
\begin{equation}
\label{eq:reaction-ordre1}
\frac{\dd [A]}{\dd t} = -k[A], \qquad
[A](t) = [A]_0 e^{-kt}.
\end{equation}
Temps de demi-vie : $t_{1/2} = \ln 2 / k$.

\paragraph{Réaction du second ordre.}
\index{réaction!second ordre}
$2A \to \text{produits}$, loi cinétique :
\begin{equation}
\label{eq:reaction-ordre2}
\frac{\dd [A]}{\dd t} = -k[A]^2, \qquad
\frac{1}{[A](t)} = \frac{1}{[A]_0} + kt.
\end{equation}

\subsection{Cinétique enzymatique de Michaelis--Menten}
\label{subsec:michaelis-menten}
\index{Michaelis--Menten (cinétique)}

Le mécanisme $E + S \underset{k_{-1}}{\overset{k_1}{\rightleftharpoons}}
ES \overset{k_2}{\to} E + P$ conduit, sous l'hypothèse de
quasi-état stationnaire\index{quasi-état stationnaire}, à la loi de
Michaelis--Menten :
\begin{equation}
\label{eq:michaelis-menten}
v = \frac{\dd [P]}{\dd t}
= \frac{V_{\max}[S]}{K_M + [S]},
\end{equation}
où $V_{\max} = k_2[E]_0$ et
$K_M = (k_{-1} + k_2)/k_1$ est la constante de Michaelis.

\begin{remarque}
Pour $[S] \ll K_M$, on retrouve une cinétique de pseudo-premier
ordre : $v \approx (V_{\max}/K_M)[S]$. Pour $[S] \gg K_M$, la
vitesse sature : $v \approx V_{\max}$.
\end{remarque}

\begin{exemple}[Équation différentielle du substrat]
\label{ex:edo-substrat}
L'évolution du substrat satisfait l'EDO
\[
\frac{\dd [S]}{\dd t} = -\frac{V_{\max}[S]}{K_M + [S]},
\qquad [S](0) = [S]_0.
\]
Cette équation est à variables séparables :
\[
\left(1 + \frac{K_M}{[S]}\right)\dd[S] = -V_{\max}\, \dd t,
\]
d'où la solution implicite
$[S] + K_M \ln [S] = [S]_0 + K_M \ln [S]_0 - V_{\max}\, t$.
\end{exemple}

% ─── 10.5 Autres applications ───────────────────────────────────
\section{Autres applications}
\label{sec:autres-applications}

\subsection{Courbes de poursuite}
\label{subsec:poursuite}
\index{courbe de poursuite}

Un objet $B$ se déplace à vitesse constante $v_B$ le long de l'axe
des ordonnées. Un poursuivant $A$, de vitesse $v_A$, se dirige à
chaque instant vers $B$. Si $y = y(x)$ décrit la trajectoire de $A$,
on obtient l'EDO du second ordre :
\begin{equation}
\label{eq:poursuite}
x\, y'' = \frac{v_B}{v_A}\sqrt{1 + (y')^2}.
\end{equation}

\begin{exemple}[Cas $v_A = v_B$]
\label{ex:poursuite-egal}
Pour $v_A = v_B$ et les conditions $y(a) = 0$, $y'(a) = 0$,
on pose $p = y'$ et l'on obtient l'équation de Bernoulli en $p$.
La solution est
\[
y(x) = \frac{a}{2}\left[\frac{1}{2}\left(\frac{x}{a}\right)^2
- \frac{1}{2} - \ln\frac{x}{a}\right],
\]
et le poursuivant n'atteint jamais sa cible.
\end{exemple}

\subsection{La chaînette}
\label{subsec:chainette}
\index{chaînette}

Un câble homogène de masse linéique $\rho$ suspendu entre deux
points sous son propre poids adopte la forme $y = y(x)$ satisfaisant
\begin{equation}
\label{eq:chainette}
y'' = \frac{\rho g}{T_0}\sqrt{1 + (y')^2},
\end{equation}
où $T_0$ est la tension horizontale. La solution est la
\emph{chaînette} :
\[
y(x) = \frac{T_0}{\rho g}\cosh\!\left(\frac{\rho g}{T_0}x\right)
= a\cosh\!\left(\frac{x}{a}\right),
\qquad a = \frac{T_0}{\rho g}.
\]

\begin{proof}[Résolution]
On pose $p = y'$, de sorte que $p' = (1/a)\sqrt{1+p^2}$. En séparant
les variables : $\frac{\dd p}{\sqrt{1+p^2}} = \frac{\dd x}{a}$, d'où
$\sinh^{-1}(p) = x/a + C$. Avec $p(0) = 0$ (symétrie), $C = 0$, soit
$p = \sinh(x/a)$ et $y = a\cosh(x/a) + C'$.
\end{proof}

\subsection{La brachistochrone}
\label{subsec:brachistochrone}
\index{brachistochrone}

Le problème de la brachistochrone consiste à trouver la courbe reliant
deux points $A$ et $B$ le long de laquelle une particule glisse sans
frottement sous l'effet de la gravité en un temps minimal. Ce
problème, historiquement posé par Jean Bernoulli en 1696, conduit par
le calcul des variations à l'EDO
\[
y(1 + (y')^2) = C,
\]
dont la solution est une \emph{cycloïde}\index{cycloïde} :
\[
x = R(\theta - \sin\theta), \qquad y = R(1 - \cos\theta).
\]

% ─── 10.6 Exemples développés supplémentaires ───────────────────
\section{Exemples développés}
\label{sec:exemples-app}

\begin{exemple}[Chute avec résistance de l'air]
\label{ex:chute-air}
\index{chute avec résistance}
Un objet de masse $m$ tombe dans l'air avec une force de résistance
proportionnelle au carré de la vitesse. L'équation du mouvement est
\[
m\dot{v} = mg - kv^2.
\]
C'est une EDO à variables séparables. La vitesse limite est
$v_\infty = \sqrt{mg/k}$ et la solution est
\[
v(t) = v_\infty \tanh\!\left(\frac{g}{v_\infty}t\right).
\]
\end{exemple}

\begin{exemple}[Datation au carbone 14]
\label{ex:carbone14}
\index{datation radioactive}
La désintégration radioactive suit une loi de premier ordre :
$\dot{N} = -\lambda N$, d'où $N(t) = N_0 e^{-\lambda t}$.
Pour le carbone 14, $t_{1/2} = 5730$ ans, soit
$\lambda = \ln 2 / 5730 \approx 1{,}21 \times 10^{-4}$ an$^{-1}$.
Si un échantillon contient une fraction $f$ de la quantité originale,
son âge est
\[
t = -\frac{1}{\lambda}\ln f = -\frac{t_{1/2}}{\ln 2}\ln f.
\]
\end{exemple}

\begin{exemple}[Refroidissement de Newton]
\label{ex:newton-refroid}
\index{refroidissement de Newton}
La loi de refroidissement de Newton s'écrit
$\dot{T} = -k(T - T_{\text{ext}})$, d'où
\[
T(t) = T_{\text{ext}} + (T_0 - T_{\text{ext}})e^{-kt}.
\]
\end{exemple}

\begin{exemple}[Mélange dans un réservoir]
\label{ex:reservoir}
\index{problème de mélange}
Un réservoir de volume $V$ contient initialement de l'eau pure.
On y verse une solution de concentration $c_{\text{in}}$ au débit
$Q$ (même débit sortant). Si $m(t)$ est la masse de soluté :
\[
\dot{m} = Qc_{\text{in}} - \frac{Q}{V}m, \qquad
m(t) = c_{\text{in}}V(1 - e^{-Qt/V}).
\]
La concentration tend vers $c_{\text{in}}$ avec la constante de
temps $V/Q$.
\end{exemple}

% ─── 10.7 Exercices ─────────────────────────────────────────────
\section{Exercices}
\label{sec:exo-chap10}

\begin{exercice}[Pendule amorti]
\label{exo:pendule-amorti}
Un pendule de longueur $\ell = 1$ m est soumis à un amortissement
de coefficient $c/(m\ell) = 0{,}5$ s$^{-1}$.
\begin{enumerate}[label=\alph*)]
  \item Linéariser l'équation autour de $\theta = 0$ et résoudre.
  \item Déterminer le régime (pseudo-périodique, critique ou
        apériodique).
  \item Calculer le temps nécessaire pour que l'amplitude soit
        réduite à $1\%$ de sa valeur initiale.
\end{enumerate}
\end{exercice}

\begin{exercice}[Circuit RLC]
\label{exo:RLC}
Un circuit RLC série a $R = 100\;\Omega$, $L = 0{,}1$ H,
$C = 10\;\mu$F.
\begin{enumerate}[label=\alph*)]
  \item Calculer $\omega_0$ et $\xi$.
  \item Déterminer le régime de fonctionnement.
  \item Résoudre pour une tension échelon $E(t) = 10$ V avec
        $q(0) = 0$, $\dot{q}(0) = 0$.
  \item Calculer la fonction de transfert $H(s)$ et tracer
        qualitativement le diagramme de Bode en amplitude.
\end{enumerate}
\end{exercice}

\begin{exercice}[Modèle logistique avec récolte]
\label{exo:logistique-recolte}
\index{logistique avec récolte}
On modifie l'équation logistique en ajoutant une récolte constante :
$\dot{N} = rN(1 - N/K) - h$.
\begin{enumerate}[label=\alph*)]
  \item Trouver les points d'équilibre en fonction de $h$.
  \item Montrer qu'il existe une valeur critique
        $h_c = rK/4$ au-delà de laquelle la population
        s'éteint.
  \item Identifier le type de bifurcation en $h = h_c$.
  \item Tracer le diagramme de bifurcation.
\end{enumerate}
\end{exercice}

\begin{exercice}[Espèces en compétition]
\label{exo:competition}
Soit le système
\[
\dot{x} = x(3 - x - 2y), \qquad \dot{y} = y(2 - x - y).
\]
\begin{enumerate}[label=\alph*)]
  \item Trouver tous les points d'équilibre.
  \item Linéariser en chaque point et déterminer leur stabilité.
  \item Tracer les isoclines nulles et le portrait de phase qualitatif.
  \item Interpréter biologiquement le résultat.
\end{enumerate}
\end{exercice}

\begin{exercice}[Cinétique de Michaelis--Menten]
\label{exo:michaelis}
Pour une réaction enzymatique avec $V_{\max} = 2$ mol/L/s et
$K_M = 0{,}5$ mol/L :
\begin{enumerate}[label=\alph*)]
  \item Écrire et résoudre (implicitement) l'EDO pour $[S](t)$
        avec $[S]_0 = 5$ mol/L.
  \item Approximer la solution pour les temps courts ($[S] \gg K_M$)
        et les temps longs ($[S] \ll K_M$).
  \item Estimer le temps nécessaire pour consommer $90\%$ du
        substrat.
\end{enumerate}
\end{exercice}

\begin{exercice}[Modèle SIR avec démographie]
\label{exo:SIR-demo}
\index{SIR (modèle)!avec démographie}
On ajoute un flux de naissances et de morts au modèle SIR :
\[
\dot{S} = \mu - \beta SI - \mu S, \quad
\dot{I} = \beta SI - \gamma I - \mu I, \quad
\dot{R} = \gamma I - \mu R.
\]
\begin{enumerate}[label=\alph*)]
  \item Montrer que $S + I + R = 1$ est invariant.
  \item Trouver l'équilibre sans maladie (DFE) et l'équilibre
        endémique (EE).
  \item Calculer le nombre de reproduction de base $\mathcal{R}_0$.
  \item Montrer que le DFE est stable si $\mathcal{R}_0 < 1$ et
        que l'EE est stable si $\mathcal{R}_0 > 1$.
\end{enumerate}
\end{exercice}

\begin{exercice}[Courbe de poursuite]
\label{exo:poursuite}
Un lapin court le long de l'axe des $y$ à vitesse $v$. Un chien,
partant de $(a, 0)$, court à vitesse $kv$ ($k > 1$) en se dirigeant
toujours vers le lapin.
\begin{enumerate}[label=\alph*)]
  \item Établir l'EDO vérifiée par la trajectoire $y(x)$ du chien.
  \item Résoudre par le changement de variable $p = y'$.
  \item Déterminer le point de capture et le temps de capture.
\end{enumerate}
\end{exercice}

\begin{exercice}[Modélisation complète]
\label{exo:modelisation}
Un lac de volume $V = 10^6$ m$^3$ est pollué. De l'eau propre
arrive au débit $Q_{\text{in}} = 500$ m$^3$/jour et de l'eau sort au
même débit. La concentration initiale de polluant est $c_0 = 200$
mg/m$^3$ et le polluant se dégrade chimiquement au taux
$\lambda = 0{,}001$ jour$^{-1}$.
\begin{enumerate}[label=\alph*)]
  \item Écrire l'EDO pour la concentration $c(t)$.
  \item Résoudre et tracer la courbe de dépollution.
  \item Calculer le temps nécessaire pour atteindre $c = 10$ mg/m$^3$.
\end{enumerate}
\end{exercice}

% ─── 10.8 Résumé ────────────────────────────────────────────────
\section*{Résumé du chapitre}

\begin{itemize}
  \item Les \textbf{oscillations mécaniques} (pendule, oscillateurs
        couplés) sont décrites par des EDO du second ordre ; le
        portrait de phase permet une analyse qualitative complète.
  \item Les \textbf{circuits électriques} (RC, RL, RLC) conduisent
        aux mêmes types d'EDO avec l'analogie
        $m \leftrightarrow L$, $c \leftrightarrow R$,
        $k \leftrightarrow 1/C$.
  \item En \textbf{dynamique des populations}, le modèle logistique
        de Verhulst, le système de Lotka--Volterra et le modèle SIR
        illustrent la richesse des comportements non-linéaires.
  \item La \textbf{cinétique chimique} produit des EDO à variables
        séparables (ordres 1 et 2) ou des systèmes non-linéaires
        (Michaelis--Menten).
  \item Les \textbf{courbes de poursuite}, la \textbf{chaînette} et
        la \textbf{brachistochrone} sont des applications géométriques
        classiques des EDO.
  \item La \emph{modélisation} --- c'est-à-dire le passage du
        problème physique à l'EDO, puis l'interprétation de la
        solution --- est la compétence centrale de ce chapitre.
\end{itemize}


%% ═══════════════════════════════════════════════════════════════════
%%  FORMULAIRE RÉCAPITULATIF
%% ═══════════════════════════════════════════════════════════════════
\chapter*{Formulaire Récapitulatif}
\addcontentsline{toc}{chapter}{Formulaire Récapitulatif}
\label{chap:formulaire}
\index{formulaire récapitulatif}

\section*{I. Équations du premier ordre}

\renewcommand{\arraystretch}{1.6}
\begin{center}
\begin{tabular}{p{4.2cm} p{4cm} p{5.5cm}}
\toprule
\textbf{Type} & \textbf{Forme} & \textbf{Méthode} \\
\midrule
Variables séparables &
$y' = f(x)g(y)$ &
$\displaystyle\int \frac{\dd y}{g(y)} = \int f(x)\,\dd x$ \\
Linéaire &
$y' + p(x)y = q(x)$ &
Facteur intégrant $e^{\int p\,\dd x}$ \\
Bernoulli &
$y' + p(x)y = q(x)y^n$ &
Substitution $z = y^{1-n}$ \\
Exacte &
$M\,\dd x + N\,\dd y = 0$ &
$\partial M/\partial y = \partial N/\partial x$ \\
Homogène &
$y' = F(y/x)$ &
Substitution $v = y/x$ \\
Ricatti &
$y' = P + Qy + Ry^2$ &
Solution part.\ $y_1$, puis $y = y_1 + 1/z$ \\
\bottomrule
\end{tabular}
\end{center}

\section*{II. Équations linéaires d'ordre $n$ à coefficients constants}

Équation : $a_n y^{(n)} + \cdots + a_1 y' + a_0 y = f(x)$.

\medskip
\noindent\textbf{Équation caractéristique :}
$a_n r^n + \cdots + a_1 r + a_0 = 0$.

\medskip
\renewcommand{\arraystretch}{1.5}
\begin{center}
\begin{tabular}{p{4.5cm} p{8cm}}
\toprule
\textbf{Racine} & \textbf{Solutions fondamentales} \\
\midrule
$r$ réelle simple & $e^{rx}$ \\
$r$ réelle de multiplicité $m$ &
$e^{rx}, xe^{rx}, \ldots, x^{m-1}e^{rx}$ \\
$\alpha \pm i\beta$ simples &
$e^{\alpha x}\cos\beta x,\; e^{\alpha x}\sin\beta x$ \\
$\alpha \pm i\beta$ de mult.\ $m$ &
$x^k e^{\alpha x}\cos\beta x,\;
 x^k e^{\alpha x}\sin\beta x,\; k=0,\ldots,m-1$ \\
\bottomrule
\end{tabular}
\end{center}

\section*{III. Table de transformées de Laplace}

\renewcommand{\arraystretch}{1.5}
\begin{center}
\begin{tabular}{cc|cc}
\toprule
$f(t)$ & $\mathcal{L}\{f\}(s)$ &
$f(t)$ & $\mathcal{L}\{f\}(s)$ \\
\midrule
$1$ & $\dfrac{1}{s}$ &
$t^n$ & $\dfrac{n!}{s^{n+1}}$ \\[6pt]
$e^{at}$ & $\dfrac{1}{s-a}$ &
$t^n e^{at}$ & $\dfrac{n!}{(s-a)^{n+1}}$ \\[6pt]
$\sin(bt)$ & $\dfrac{b}{s^2+b^2}$ &
$\cos(bt)$ & $\dfrac{s}{s^2+b^2}$ \\[6pt]
$e^{at}\sin(bt)$ & $\dfrac{b}{(s-a)^2+b^2}$ &
$e^{at}\cos(bt)$ & $\dfrac{s-a}{(s-a)^2+b^2}$ \\[6pt]
$t\sin(bt)$ & $\dfrac{2bs}{(s^2+b^2)^2}$ &
$t\cos(bt)$ & $\dfrac{s^2-b^2}{(s^2+b^2)^2}$ \\[6pt]
$\sinh(bt)$ & $\dfrac{b}{s^2-b^2}$ &
$\cosh(bt)$ & $\dfrac{s}{s^2-b^2}$ \\[6pt]
$u(t-a)$ & $\dfrac{e^{-as}}{s}$ &
$\delta(t-a)$ & $e^{-as}$ \\[6pt]
$f'(t)$ & $sF(s) - f(0)$ &
$f''(t)$ & $s^2F(s) - sf(0) - f'(0)$ \\[6pt]
$(f*g)(t)$ & $F(s)\cdot G(s)$ &
$t^n f(t)$ & $(-1)^n F^{(n)}(s)$ \\
\bottomrule
\end{tabular}
\end{center}

\section*{IV. Systèmes linéaires}

Système $\dot{\mathbf{x}} = A\mathbf{x}$, $A \in \R^{n \times n}$.

\medskip
\noindent\textbf{Solution :} $\mathbf{x}(t) = e^{At}\mathbf{x}_0$.

\medskip
\renewcommand{\arraystretch}{1.5}
\begin{center}
\begin{tabular}{p{5cm} p{4cm} p{4cm}}
\toprule
\textbf{Valeurs propres ($n=2$)} & \textbf{Type} & \textbf{Stabilité} \\
\midrule
$\lambda_1, \lambda_2 < 0$ réelles & N\oe ud stable & Stable \\
$\lambda_1, \lambda_2 > 0$ réelles & N\oe ud instable & Instable \\
$\lambda_1 < 0 < \lambda_2$ & Col & Instable \\
$\alpha \pm i\beta$, $\alpha < 0$ & Foyer stable & Stable \\
$\alpha \pm i\beta$, $\alpha > 0$ & Foyer instable & Instable \\
$\pm i\beta$ & Centre & Stable (non asymp.) \\
\bottomrule
\end{tabular}
\end{center}

\section*{V. Stabilité de Lyapunov}

\noindent\textbf{Théorème de Lyapunov :} Si $V(\mathbf{x})$ est définie
positive et $\dot{V}(\mathbf{x}) \leq 0$, alors l'équilibre est stable.
Si $\dot{V} < 0$ (définie négative), il est asymptotiquement stable.

\section*{VI. Bifurcations en dimension 1}

\renewcommand{\arraystretch}{1.5}
\begin{center}
\begin{tabular}{p{3.5cm} p{3.5cm} p{6cm}}
\toprule
\textbf{Bifurcation} & \textbf{Forme normale} & \textbf{Caractéristique} \\
\midrule
N\oe ud-col & $\dot{x} = \mu + x^2$ &
Création/destruction de 2 équilibres \\
Transcritique & $\dot{x} = \mu x - x^2$ &
Échange de stabilité \\
Fourche (sup.) & $\dot{x} = \mu x - x^3$ &
1 équilibre $\to$ 3 équilibres \\
Fourche (sous.) & $\dot{x} = \mu x + x^3$ &
Instabilité sous-critique \\
\bottomrule
\end{tabular}
\end{center}

\renewcommand{\arraystretch}{1.0}

%% ═══════════════════════════════════════════════════════════════════
%%  BIBLIOGRAPHIE
%% ═══════════════════════════════════════════════════════════════════
\begin{thebibliography}{99}

\bibitem{tenenbaum}
\textsc{Tenenbaum}, M. et \textsc{Pollard}, H.,
\textit{Ordinary Differential Equations},
Dover Publications, 1985.
\index{Tenenbaum}

\bibitem{arnold}
\textsc{Arnold}, V.\,I.,
\textit{Ordinary Differential Equations},
Springer-Verlag, 3\up{e} édition, 2006.
\index{Arnold}

\bibitem{chicone}
\textsc{Chicone}, C.,
\textit{Ordinary Differential Equations with Applications},
Springer, Texts in Applied Mathematics, vol.~34, 2\up{e} édition, 2006.
\index{Chicone}

\bibitem{hirsch-smale}
\textsc{Hirsch}, M.\,W., \textsc{Smale}, S. et \textsc{Devaney}, R.\,L.,
\textit{Differential Equations, Dynamical Systems, and an Introduction
to Chaos},
Academic Press, 3\up{e} édition, 2013.
\index{Hirsch}\index{Smale}

\bibitem{verhulst}
\textsc{Verhulst}, F.,
\textit{Nonlinear Differential Equations and Dynamical Systems},
Springer-Verlag, Universitext, 2\up{e} édition, 1996.
\index{Verhulst}

\bibitem{braun}
\textsc{Braun}, M.,
\textit{Differential Equations and Their Applications:
An Introduction to Applied Mathematics},
Springer-Verlag, 4\up{e} édition, 1993.
\index{Braun}

\bibitem{perko}
\textsc{Perko}, L.,
\textit{Differential Equations and Dynamical Systems},
Springer-Verlag, Texts in Applied Mathematics, vol.~7,
3\up{e} édition, 2001.
\index{Perko}

\end{thebibliography}

%% ═══════════════════════════════════════════════════════════════════
%%  INDEX
%% ═══════════════════════════════════════════════════════════════════
\printindex
\end{document}
